\documentclass[11pt,oneside]{article}
\usepackage[a4paper,margin=27mm,headheight=15pt]{geometry}
\usepackage[T1]{fontenc}
\usepackage[utf8]{inputenc}
\IfFileExists{libertinus.sty}{\usepackage{libertinus}}{\usepackage{lmodern}}
\usepackage{microtype}
\usepackage{amsmath,amssymb,amsthm,mathtools,mathrsfs}
\usepackage{bm}
\usepackage{booktabs,longtable,array,tabularx}
\usepackage{graphicx}
\usepackage{xcolor}
\usepackage{enumitem}
\usepackage{fancyhdr}
\usepackage{titlesec}
\usepackage{xurl}
\usepackage{needspace}
\usepackage{ragged2e}
\usepackage{hyperref}
\usepackage{bookmark}
\usepackage[nameinlink,noabbrev]{cleveref}

\definecolor{linkblue}{RGB}{25,75,135}
\definecolor{shadegray}{RGB}{246,247,249}
\definecolor{rulegray}{RGB}{195,201,208}
\definecolor{leanfill}{RGB}{219,238,249}
\definecolor{classicalfill}{RGB}{232,235,239}
\definecolor{mixedfill}{RGB}{231,224,244}
\definecolor{frontierfill}{RGB}{251,231,204}
\definecolor{conventionfill}{RGB}{237,237,225}
\definecolor{historicalfill}{RGB}{239,226,214}
\definecolor{computationalfill}{RGB}{224,239,225}
\definecolor{refutedfill}{RGB}{247,217,217}
\definecolor{leantext}{RGB}{18,76,112}
\definecolor{classicaltext}{RGB}{65,70,78}
\definecolor{mixedtext}{RGB}{78,50,111}
\definecolor{frontiertext}{RGB}{124,67,15}
\definecolor{conventiontext}{RGB}{83,83,51}
\definecolor{historicaltext}{RGB}{102,67,43}
\definecolor{computationaltext}{RGB}{35,88,38}
\definecolor{refutedtext}{RGB}{125,34,34}

\hypersetup{
  colorlinks=true,
  linkcolor=linkblue,
  citecolor=linkblue,
  urlcolor=linkblue,
  pdftitle={Mathematical Glossary for Analysis/FabiusFunction -- Complete Updated Edition},
  pdfsubject={Concept-based glossary for the current Lean development and active documentation corpus},
  pdfauthor={Terminology synthesized from Vladimir Reshetnikov's ProveIt repository},
  pdfkeywords={Fabius function, Rvachev up-function, Thue-Morse, q-series, glossary, Lean, Lambert W, transseries, spectral theory}
}

\pagestyle{fancy}
\fancyhf{}
\fancyhead[L]{\small Mathematical glossary for \texttt{Analysis/FabiusFunction}}
\fancyhead[R]{\small\nouppercase{\leftmark}}
\fancyfoot[C]{\thepage}
\renewcommand{\headrulewidth}{0.4pt}
\titleformat{\section}{\Large\bfseries}{\thesection}{0.7em}{}
\titleformat{\subsection}{\large\bfseries}{\thesubsection}{0.7em}{}
\titleformat{\subsubsection}{\normalsize\bfseries}{\thesubsubsection}{0.7em}{}
\setcounter{secnumdepth}{2}
\setcounter{tocdepth}{2}
\setlist[itemize]{leftmargin=2em,itemsep=0.25em,topsep=0.4em}
\setlist[enumerate]{leftmargin=2.3em,itemsep=0.25em,topsep=0.4em}
\setlength{\emergencystretch}{3em}
\allowdisplaybreaks
\raggedbottom

\newtheorem{definition}{Definition}[section]
\newcommand{\R}{\mathbb R}
\newcommand{\C}{\mathbb C}
\newcommand{\Q}{\mathbb Q}
\newcommand{\N}{\mathbb N}
\newcommand{\Z}{\mathbb Z}
\newcommand{\Up}{\operatorname{up}}
\newcommand{\sinc}{\operatorname{sinc}}
\newcommand{\wt}{\operatorname{wt}_2}
\newcommand{\e}{\mathrm e}
\newcommand{\ii}{\mathrm i}
\newcommand{\dd}{\,\mathrm d}
\newcommand{\Var}{\operatorname{Var}}
\newcommand{\E}{\mathbb E}
\newcommand{\Prob}{\mathbb P}
\newcommand{\bigO}{\mathcal O}
\newcommand{\smallo}{o}
\newcommand{\supp}{\operatorname{supp}}
\newcommand{\sgn}{\operatorname{sgn}}
\newcommand{\lcm}{\operatorname{lcm}}
\newcommand{\vtwo}{v_2}
\newcommand{\qbinom}[3]{\genfrac{[}{]}{0pt}{}{#1}{#2}_{#3}}
\newcommand{\Sha}{\mathop{\mathrm{III}}}
\newcommand{\repo}{\nolinkurl{Analysis/FabiusFunction}}
\newcommand{\crossref}[1]{\textit{See also:} #1.}

\newcommand{\statusbox}[3]{%
  \begingroup\setlength{\fboxsep}{2.1pt}%
  \colorbox{#1}{\textcolor{#2}{\sffamily\bfseries\scriptsize\strut #3}}%
  \endgroup}
\newcommand{\statuslean}{\statusbox{leanfill}{leantext}{LEAN-PROVED}}
\newcommand{\statusclassical}{\statusbox{classicalfill}{classicaltext}{CLASSICAL}}
\newcommand{\statusmixed}{\statusbox{mixedfill}{mixedtext}{PROJECT / MIXED}}
\newcommand{\statusfrontier}{\statusbox{frontierfill}{frontiertext}{FRONTIER}}
\newcommand{\statusconvention}{\statusbox{conventionfill}{conventiontext}{CONVENTION}}
\newcommand{\statushistorical}{\statusbox{historicalfill}{historicaltext}{HISTORICAL}}
\newcommand{\statuscomputational}{\statusbox{computationalfill}{computationaltext}{COMPUTATIONAL}}
\newcommand{\statusrefuted}{\statusbox{refutedfill}{refutedtext}{CORRECTED / REFUTED}}

\newcommand{\glossaryentry}[4]{%
  \par\addvspace{0.92\baselineskip}%
  \Needspace{5\baselineskip}%
  \phantomsection\hypertarget{gls:#1}{}%
  \noindent\begin{tabularx}{\linewidth}{@{}>{\RaggedRight\arraybackslash}Xr@{}}
    \textbf{\large #2.} & #3
  \end{tabularx}\par\nobreak
  \smallskip\noindent #4\par
  \nopagebreak[2]\vspace{0.12\baselineskip}%
  \noindent\color{rulegray}\rule{\linewidth}{0.25pt}\color{black}%
}

\newenvironment{scopebox}{%
  \par\smallskip\noindent\begin{center}\begin{minipage}{0.94\textwidth}%
  \color{black}\hrule\vspace{0.6em}\small
}{%
  \vspace{0.6em}\hrule\end{minipage}\end{center}\smallskip
}

\newcolumntype{L}[1]{>{\RaggedRight\arraybackslash}p{#1}}
\hyphenation{Fa-bi-us Rva-chev Thue--Morse}

\title{\bfseries Mathematical Glossary for \texttt{Analysis/FabiusFunction}\\[0.35em]
\large Complete updated edition for the current formal and documentation corpus}
\author{Terminology synthesized from Vladimir Reshetnikov's \texttt{ProveIt} repository}
\date{Repository state inspected: 3 September 2026\\
Formal census checkpoint: 1 September 2026}

\begin{document}
\maketitle

\begin{abstract}
This edition is a ground-up corpus-level update of the mathematical glossary for
\repo.  It contains \textbf{976 alphabetized headwords}, replacing the
prior 437-headword edition.  The audit covers the bounded and global Fabius
functions, Rvachev's up-function, inverse and computability theory, exact dyadic
arithmetic, Thue--Morse analysis, probability and operator models,
\(q\)-Pochhammer and Gaussian-binomial calculus, basic hypergeometric and theta
functions, Fourier and spectral theory, splines and sampling, orthogonal
polynomials, combinatorial coefficient calculus, Lambert-\(W\) and Mellin
endpoint analysis, polynomial--logarithmic transseries, and the active frontier
programs.  Definitions give a characteristic formula, normalization or branch
choice when material, and the role of the notion in the project.  The typography
and running-head design preserve the style of the repository's principal
Fabius--Rvachev documents.
\end{abstract}

\begin{scopebox}
\textbf{Repository snapshot and coverage.}
The live formal census used by this edition records 675 source modules and
8,909 public declarations at its 1 September 2026 checkpoint.  The active
documentation/notation inventory inspected on 3 September 2026 records 172
active \texttt{.tex} sources including the notation catalogue, of which 73 are
standalone documents; it also records 146 shared semantic notation commands and
30 validated document-local non-Fourier hat names.  Historical sources below
\texttt{docs/archive/} were used for provenance and alias control, not counted
as current theorem-bearing duplicates.
\end{scopebox}

\begin{scopebox}
\textbf{Editorial unit of completeness.}
Completeness is concept-based rather than token-based.  The glossary includes
every mathematically substantive concept, named object, normalization,
proof-status distinction, transform convention, asymptotic coordinate, and
recurring construction found in the active corpus.  It consolidates transparent
spelling and inflection variants; it does not create separate headwords for
elementary words such as ``number'' or ``function,'' local dummy variables,
one-use theorem labels, or mechanically predictable specializations.  A local
symbol not promoted to the shared notation API has no project-wide meaning
unless its defining document supplies one.
\end{scopebox}

\section*{How to read the status badges}
\addcontentsline{toc}{section}{How to read the status badges}
\begin{description}[leftmargin=3.5cm,style=nextline]
\item[\statuslean] The named object or core theorem has an active Lean source
and belongs to the formal facade.  The badge does not claim that every
expository sentence in the entry is itself a theorem statement.
\item[\statusclassical] Standard background mathematics used by the project.
\item[\statusmixed] A classical notion with repository-specific normalization,
a bridge between formal and analytic presentations, or a topic whose relevant
subresults have mixed proof status.
\item[\statusfrontier] Active research material, conjectural structure, or a
semi-formalized program.  No kernel-checked status is implied.
\item[\statusconvention] A notation, normalization, branch, equality mode, or
editorial governance rule.
\item[\statushistorical] A source-faithful or retired usage retained solely for
interpretation and provenance.
\end{description}

\tableofcontents
\newpage
\section*{Headword index}\addcontentsline{toc}{section}{Headword index}
\small This index is hyperlinked to the alphabetized definitions.  Status badges appear at each entry, not in the index.\normalsize\par\medskip
\begingroup\footnotesize\setlength{\tabcolsep}{5pt}\renewcommand{\arraystretch}{1.10}
\begin{longtable}{L{0.31\textwidth}L{0.31\textwidth}L{0.31\textwidth}}
\hyperlink{gls:abel-equation}{Abel equation} & \hyperlink{gls:absolute-continuity}{Absolute continuity} & \hyperlink{gls:absolute-convergence}{Absolute convergence} \\
\hyperlink{gls:absolute-legendre-summability}{Absolute Legendre summability} & \hyperlink{gls:activation-deficit}{Activation deficit} & \hyperlink{gls:activation-profile}{Activation profile} \\
\hyperlink{gls:active-documentation-corpus}{Active documentation corpus} & \hyperlink{gls:gaussian-binomial-adjacent-ratio}{Adjacent Gaussian-binomial ratio} & \hyperlink{gls:admissible-class}{Admissible class} \\
\hyperlink{gls:fractional-affine-covariance}{Affine covariance of fractional integrals} & \hyperlink{gls:affine-mixed-difference}{Affine mixed difference} & \hyperlink{gls:affine-transformation}{Affine transformation} \\
\hyperlink{gls:algebraic-function-branch}{Algebraic function branch} & \hyperlink{gls:aliasing}{Aliasing} & \hyperlink{gls:all-orders-asymptotic-expansion}{All-orders asymptotic expansion} \\
\hyperlink{gls:almost-everywhere}{Almost everywhere} & \hyperlink{gls:almost-sure-convergence}{Almost-sure convergence} & \hyperlink{gls:amplitude}{Amplitude} \\
\hyperlink{gls:analysis-operator}{Analysis operator} & \hyperlink{gls:analytic-continuation}{Analytic continuation} & \hyperlink{gls:analytic-germ}{Analytic germ} \\
\hyperlink{gls:analytic-locus}{Analytic locus} & \hyperlink{gls:analytic-moment}{Analytic moment} & \hyperlink{gls:analytic-zero-multiplicity}{Analytic zero multiplicity} \\
\hyperlink{gls:analyticity-at-a-point}{Analyticity at a point} & \hyperlink{gls:angular-frequency-fourier-transform}{Angular-frequency Fourier transform} & \hyperlink{gls:antidiagonal-polynomial}{Antidiagonal polynomial} \\
\hyperlink{gls:appell-deconvolution}{Appell deconvolution} & \hyperlink{gls:appell-sequence}{Appell sequence} & \hyperlink{gls:approximate-identity}{Approximate identity} \\
\hyperlink{gls:approximation-theory}{Approximation theory} & \hyperlink{gls:arbitrary-phase-synthesis}{Arbitrary-phase synthesis} & \hyperlink{gls:arbitrary-phase-up-synthesis}{Arbitrary-phase up synthesis} \\
\hyperlink{gls:arithmetic-ray}{Arithmetic ray} & \hyperlink{gls:associahedron}{Associahedron} & \hyperlink{gls:asymptotic-coefficient-jet}{Asymptotic coefficient jet} \\
\hyperlink{gls:asymptotic-equivalence}{Asymptotic equivalence} & \hyperlink{gls:asymptotic-expansion}{Asymptotic expansion} & \hyperlink{gls:asymptotic-scale}{Asymptotic scale} \\
\hyperlink{gls:atom-of-a-measure}{Atom of a measure} & \hyperlink{gls:atomic-function}{Atomic function} & \hyperlink{gls:atomless-law}{Atomless law} \\
\hyperlink{gls:automatic-sequence}{Automatic sequence} & \hyperlink{gls:averaging-operator}{Averaging operator} & \hyperlink{gls:b-spline}{B-spline} \\
\hyperlink{gls:bailey-lemma}{Bailey lemma} & \hyperlink{gls:bailey-pair}{Bailey pair} & \hyperlink{gls:balanced-basic-hypergeometric-series}{Balanced basic hypergeometric series} \\
\hyperlink{gls:banach-fixed-point-theorem}{Banach fixed-point theorem} & \hyperlink{gls:barnes-g-function}{Barnes \(G\)-function} & \hyperlink{gls:base-q-digit}{Base-\(q\) digit} \\
\hyperlink{gls:base-q-digit-sum}{Base-\(q\) digit sum} & \hyperlink{gls:basic-hypergeometric-series}{Basic hypergeometric series} & \hyperlink{gls:basic-hypergeometric-r-s-phi}{Basic hypergeometric series \({}_r\phi_s\)} \\
\hyperlink{gls:bell-hessenberg-determinant}{Bell--Hessenberg determinant} & \hyperlink{gls:bell-number}{Bell number} & \hyperlink{gls:bell-polynomial}{Bell polynomial} \\
\hyperlink{gls:bell-polynomial-asymptotic-coefficient}{Bell-polynomial asymptotic coefficient} & \hyperlink{gls:bernoulli-number}{Bernoulli number} & \hyperlink{gls:bernoulli-recurrence}{Bernoulli recurrence} \\
\hyperlink{gls:berry-esseen-bound}{Berry--Esseen bound} & \hyperlink{gls:bessel-inequality}{Bessel inequality} & \hyperlink{gls:big-o-and-little-o}{Big-O and little-o notation} \\
\hyperlink{gls:bijection}{Bijection} & \hyperlink{gls:bilateral-basic-hypergeometric-series}{Bilateral basic hypergeometric series} & \hyperlink{gls:bilateral-laplace-transform}{Bilateral Laplace transform} \\
\hyperlink{gls:binary-block-factorization}{Binary block factorization} & \hyperlink{gls:binary-digit-sum}{Binary digit sum} & \hyperlink{gls:binary-expansion}{Binary expansion} \\
\hyperlink{gls:binary-floor-digit-formula}{Binary floor-digit formula} & \hyperlink{gls:binary-reduction-series}{Binary reduction series} & \hyperlink{gls:binary-reflection-recursion}{Binary reflection recursion} \\
\hyperlink{gls:binomial-coefficient}{Binomial coefficient} & \hyperlink{gls:binomial-basis}{Binomial polynomial basis} & \hyperlink{gls:binomial-transform}{Binomial transform} \\
\hyperlink{gls:biorthogonal-system}{Biorthogonal system} & \hyperlink{gls:bit-extraction}{Bit extraction} & \hyperlink{gls:boolean-cube-difference}{Boolean-cube mixed difference} \\
\hyperlink{gls:boolean-cumulant}{Boolean cumulant} & \hyperlink{gls:borel-probability-measure}{Borel probability measure} & \hyperlink{gls:borel-summability}{Borel summability} \\
\hyperlink{gls:borel-transform}{Borel transform} & \hyperlink{gls:bose-einstein-kernel}{Bose-Einstein kernel} & \hyperlink{gls:bose-finite-part-integral}{Bose finite-part integral} \\
\hyperlink{gls:boundary-value-problem}{Boundary value problem} & \hyperlink{gls:bounded-fabius-candidate}{Bounded Fabius candidate} & \hyperlink{gls:bounded-fabius-function}{Bounded Fabius function} \\
\hyperlink{gls:bounded-variation}{Bounded variation} & \hyperlink{gls:box-spline}{Box spline} & \hyperlink{gls:bromwich-inversion-integral}{Bromwich inversion integral} \\
\hyperlink{gls:bromwich-saddle-method}{Bromwich saddle method} & \hyperlink{gls:bump-function}{Bump function} & \hyperlink{gls:c-infinity-smoothness}{C-infinity smoothness} \\
\hyperlink{gls:canonical-clamped-inverse}{Canonical clamped inverse} & \hyperlink{gls:canonical-correlation}{Canonical correlation} & \hyperlink{gls:canonical-fabius-fixed-point}{Canonical Fabius fixed point} \\
\hyperlink{gls:fabius-specification}{Canonical Fabius specification} & \hyperlink{gls:canonical-fabius-inverse}{Canonical inverse Fabius function} & \hyperlink{gls:canonical-solution}{Canonical solution} \\
\hyperlink{gls:caputo-fractional-derivative}{Caputo fractional derivative} & \hyperlink{gls:cardinal-interpolation}{Cardinal interpolation} & \hyperlink{gls:carleman-condition}{Carleman condition} \\
\hyperlink{gls:carleman-linearization}{Carleman linearization} & \hyperlink{gls:carleman-matrix}{Carleman matrix} & \hyperlink{gls:carry-count}{Carry count} \\
\hyperlink{gls:catalan-number}{Catalan number} & \hyperlink{gls:cauchy-repeated-integration-formula}{Cauchy formula for repeated integration} & \hyperlink{gls:cauchy-integral-formula}{Cauchy integral formula} \\
\hyperlink{gls:cauchy-product}{Cauchy product} & \hyperlink{gls:cauchy-theorem}{Cauchy theorem} & \hyperlink{gls:cauchy-transform}{Cauchy transform of a measure} \\
\hyperlink{gls:cdf-layer-cake}{CDF layer-cake formula} & \hyperlink{gls:centered-appell-sequence}{Centered Appell sequence} & \hyperlink{gls:centered-finite-spline}{Centered finite spline} \\
\hyperlink{gls:centered-geometric-uniform-law}{Centered geometric uniform law} & \hyperlink{gls:centered-periodic-correction}{Centered periodic correction} & \hyperlink{gls:centered-sinc-partial-product}{Centered sinc partial product} \\
\hyperlink{gls:centered-thue-morse-shell}{Centered Thue--Morse shell} & \hyperlink{gls:central-binomial-valuation}{Central binomial valuation} & \hyperlink{gls:central-gaussian-binomial-coefficient}{Central Gaussian binomial coefficient} \\
\hyperlink{gls:central-limit-theorem}{Central limit theorem} & \hyperlink{gls:chaos-multi-index}{Chaos multi-index} & \hyperlink{gls:characteristic-function}{Characteristic function} \\
\hyperlink{gls:christoffel-darboux-kernel}{Christoffel--Darboux kernel} & \hyperlink{gls:christoffel-function}{Christoffel function} & \hyperlink{gls:closed-form}{Closed form} \\
\hyperlink{gls:coefficient-extraction}{Coefficient extraction} & \hyperlink{gls:coercion}{Coercion and scalar cast} & \hyperlink{gls:compact-support}{Compact support} \\
\hyperlink{gls:pascal-companion-row}{Companion Pascal row} & \hyperlink{gls:compiled-tree}{Compiled tree} & \hyperlink{gls:complementary-bell-number}{Complementary Bell number} \\
\hyperlink{gls:complete-binary-block}{Complete binary block} & \hyperlink{gls:complete-exponential-bell-polynomial}{Complete exponential Bell polynomial} & \hyperlink{gls:complete-homogeneous-symmetric-polynomial}{Complete homogeneous symmetric polynomial} \\
\hyperlink{gls:complex-exponential-generating-function}{Complex exponential generating function} & \hyperlink{gls:complex-moment-generating-function}{Complex moment-generating function} & \hyperlink{gls:complex-shift}{Complex shift} \\
\hyperlink{gls:composite-mesh-sharpness}{Composite-mesh sharpness} & \hyperlink{gls:computable-fabius-inverse}{Computable inverse Fabius function} & \hyperlink{gls:computable-real}{Computable real} \\
\hyperlink{gls:computable-real-function}{Computable real function} & \hyperlink{gls:computable-sequence}{Computable sequence} & \hyperlink{gls:conditional-asymptotic}{Conditional asymptotic theorem} \\
\hyperlink{gls:conditional-law}{Conditional law} & \hyperlink{gls:gaussian-binomial-continuity}{Continuity of Gaussian coefficients in the base} & \hyperlink{gls:continuous-gaussian-coefficient}{Continuous Gaussian coefficient} \\
\hyperlink{gls:continuous-linear-naturality}{Continuous-linear naturality} & \hyperlink{gls:contour-deformation}{Contour deformation} & \hyperlink{gls:contour-integral}{Contour integral} \\
\hyperlink{gls:contraction-constant}{Contraction constant} & \hyperlink{gls:contraction-mapping}{Contraction mapping} & \hyperlink{gls:convexity-and-concavity}{Convexity and concavity} \\
\hyperlink{gls:convolution}{Convolution} & \hyperlink{gls:convolution-measure}{Convolution of measures} & \hyperlink{gls:convolution-power}{Convolution power} \\
\hyperlink{gls:cornish-fisher-expansion}{Cornish--Fisher expansion} & \hyperlink{gls:correction-term}{Correction term} & \hyperlink{gls:cosine-series}{Cosine series} \\
\hyperlink{gls:coupling}{Coupling} & \hyperlink{gls:covariance-kernel}{Covariance kernel} & \hyperlink{gls:critical-value}{Critical value} \\
\hyperlink{gls:cumulant}{Cumulant} & \hyperlink{gls:cumulant-generating-function}{Cumulant-generating function} & \hyperlink{gls:cumulant-scaling-law}{Cumulant scaling law} \\
\hyperlink{gls:cumulative-distribution-function}{Cumulative distribution function} & \hyperlink{gls:cyclotomic-inverse-branch}{Cyclotomic inverse branch} & \hyperlink{gls:cyclotomic-polynomial}{Cyclotomic polynomial} \\
\hyperlink{gls:cyclotomic-valuation}{Cyclotomic valuation} & \hyperlink{gls:d-finite-function}{\(D\)-finite function} & \hyperlink{gls:decay-estimate}{Decay estimate} \\
\hyperlink{gls:declaration-arity}{Declaration arity} & \hyperlink{gls:declaration-inventory}{Declaration inventory} & \hyperlink{gls:delay-differential-equation}{Delay differential equation} \\
\hyperlink{gls:delta-series}{Delta series} & \hyperlink{gls:denominator-bound}{Denominator bound} & \hyperlink{gls:denominator-clearing}{Denominator clearing} \\
\hyperlink{gls:dense-analytic-locus}{Dense analytic locus} & \hyperlink{gls:derivative-scaling}{Derivative scaling} & \hyperlink{gls:derivative-value-distribution}{Derivative-value distribution} \\
\hyperlink{gls:determinant-bell-identity}{Determinant--Bell identity} & \hyperlink{gls:determinate-moment-problem}{Determinate moment problem} & \hyperlink{gls:differential-entropy}{Differential entropy} \\
\hyperlink{gls:differential-equation}{Differential equation} & \hyperlink{gls:differential-algebraic-function}{Differentially algebraic function} & \hyperlink{gls:differentially-transcendental-function}{Differentially transcendental function} \\
\hyperlink{gls:differentiation-under-the-integral-sign}{Differentiation under the integral sign} & \hyperlink{gls:digit-character}{Digit character} & \hyperlink{gls:digit-multiplicity-profile}{Digit multiplicity profile} \\
\hyperlink{gls:digit-power-sum}{Digit power sum} & \hyperlink{gls:digitwise-factorization}{Digitwise factorization} & \hyperlink{gls:dirac-comb}{Dirac comb} \\
\hyperlink{gls:dirac-measure}{Dirac measure} & \hyperlink{gls:dirichlet-eta-function}{Dirichlet eta function} & \hyperlink{gls:discrete-antiderivative}{Discrete antiderivative} \\
\hyperlink{gls:discrete-b-spline}{Discrete B-spline} & \hyperlink{gls:discrete-convolution}{Discrete convolution} & \hyperlink{gls:discrete-convolution-power-kernel}{Discrete convolution-power kernel} \\
\hyperlink{gls:discrete-fourier-transform}{Discrete Fourier transform} & \hyperlink{gls:discrete-limit}{Discrete limit} & \hyperlink{gls:distribution-law}{Distribution law} \\
\hyperlink{gls:division-free-identity}{Division-free identity} & \hyperlink{gls:division-free-recurrence}{Division-free recurrence} & \hyperlink{gls:dominated-convergence-theorem}{Dominated convergence theorem} \\
\hyperlink{gls:double-factorial}{Double factorial} & \hyperlink{gls:dyadic-argument}{Dyadic argument} & \hyperlink{gls:dyadic-common-denominator}{Dyadic common denominator} \\
\hyperlink{gls:dyadic-derivative-cell}{Dyadic derivative cell} & \hyperlink{gls:dyadic-derivative-sign-law}{Dyadic derivative sign law} & \hyperlink{gls:direction-profile}{Dyadic direction profile} \\
\hyperlink{gls:dyadic-factorized-product-extrapolant}{Dyadic factorized-product extrapolant} & \hyperlink{gls:dyadic-grid}{Dyadic grid} & \hyperlink{gls:dyadic-harmonic-sum}{Dyadic harmonic sum} \\
\hyperlink{gls:dyadic-name}{Dyadic name} & \hyperlink{gls:dyadic-rational}{Dyadic rational} & \hyperlink{gls:dyadic-refinement-invariance}{Dyadic refinement invariance} \\
\hyperlink{gls:dyadic-spine}{Dyadic spine} & \hyperlink{gls:dyadic-taylor-evaluator}{Dyadic Taylor evaluator} & \hyperlink{gls:dyadic-valuation}{Dyadic valuation} \\
\hyperlink{gls:dyadic-denominator-lcm}{Dyadic-value denominator LCM} & \hyperlink{gls:dyadic-zero-classification}{Dyadic zero classification} & \hyperlink{gls:edgeworth-expansion}{Edgeworth expansion} \\
\hyperlink{gls:effective-flatness}{Effective flatness} & \hyperlink{gls:effective-modulus}{Effective modulus} & \hyperlink{gls:effective-uniform-continuity}{Effective uniform continuity} \\
\hyperlink{gls:elementary-function-class}{Elementary-function class} & \hyperlink{gls:elementary-logarithmic-asymptotic}{Elementary logarithmic asymptotic} & \hyperlink{gls:elementary-symmetric-polynomial}{Elementary symmetric polynomial} \\
\hyperlink{gls:empty-sum-product-conventions}{Empty sum, empty product, and zero-power conventions} & \hyperlink{gls:endpoint-asymptotic}{Endpoint asymptotic} & \hyperlink{gls:endpoint-convention}{Endpoint convention} \\
\hyperlink{gls:endpoint-derivative-jet}{Endpoint derivative jet} & \hyperlink{gls:endpoint-flatness}{Endpoint flatness} & \hyperlink{gls:entire-function}{Entire function} \\
\hyperlink{gls:entire-function-exponential-type}{Entire function of exponential type} & \hyperlink{gls:entire-product-zero-ledger}{Entire-product zero ledger} & \hyperlink{gls:entropy-power}{Entropy power} \\
\hyperlink{gls:equality-in-distribution}{Equality in distribution} & \hyperlink{gls:equality-of-measures}{Equality of measures} & \hyperlink{gls:equicontinuity}{Equicontinuity} \\
\hyperlink{gls:euler-maclaurin-formula}{Euler-Maclaurin formula} & \hyperlink{gls:euler-mascheroni-constant}{Euler-Mascheroni constant} & \hyperlink{gls:euler-pentagonal-theorem}{Euler pentagonal theorem} \\
\hyperlink{gls:euler-first-q-exponential}{Euler's first \(q\)-exponential} & \hyperlink{gls:euler-second-q-exponential}{Euler's second \(q\)-exponential} & \hyperlink{gls:eulerian-number}{Eulerian number} \\
\hyperlink{gls:eulerian-polynomial}{Eulerian polynomial} & \hyperlink{gls:exact-arithmetic}{Exact arithmetic} & \hyperlink{gls:exact-dyadic-derivative}{Exact dyadic derivative} \\
\hyperlink{gls:exact-rational-moment-table}{Exact rational moment table} & \hyperlink{gls:exact-topological-support}{Exact topological support} & \hyperlink{gls:expansive-matrix}{Expansive matrix} \\
\hyperlink{gls:expectation}{Expectation} & \hyperlink{gls:nonclaim}{Explicit non-claim} & \hyperlink{gls:explicit-uniform-error}{Explicit uniform error} \\
\hyperlink{gls:exponential-formula}{Exponential formula} & \hyperlink{gls:exponential-generating-function}{Exponential generating function} & \hyperlink{gls:exponential-riordan-array}{Exponential Riordan array} \\
\hyperlink{gls:tilted-measure}{Exponentially tilted measure} & \hyperlink{gls:extended-nonnegative-reals}{Extended nonnegative reals} & \hyperlink{gls:extension-by-zero}{Extension by zero} \\
\hyperlink{gls:exterior-constant-tails}{Exterior constant tails} & \hyperlink{gls:faa-di-bruno-formula}{Faa di Bruno formula} & \hyperlink{gls:faa-di-bruno-partition-propagation}{Fa\`a di Bruno partition propagation} \\
\hyperlink{gls:cdf-density-bridge}{Fabius CDF--density bridge} & \hyperlink{gls:fabius-distribution}{Fabius distribution} & \hyperlink{gls:fabius-fractional-shift}{Fabius fractional shift} \\
\hyperlink{gls:fabius-function}{Fabius function} & \hyperlink{gls:fabius-recurrence-sequence}{Fabius recurrence sequence} & \hyperlink{gls:fabius-fold-map}{Fabius--Rvachev fold map} \\
\hyperlink{gls:fabius-self-map-dynamics}{Fabius self-map dynamics} & \hyperlink{gls:factorial-normalization}{Factorial normalization} & \hyperlink{gls:falling-factorial}{Falling factorial} \\
\hyperlink{gls:fast-signed-dyadic-name}{Fast signed-dyadic name} & \hyperlink{gls:filter-convergence}{Filter convergence} & \hyperlink{gls:finite-difference}{Finite difference} \\
\hyperlink{gls:finite-difference-table}{Finite-difference table} & \hyperlink{gls:finite-field}{Finite field} & \hyperlink{gls:finite-part-integral}{Finite-part integral} \\
\hyperlink{gls:finite-q-binomial-theorem}{Finite \(q\)-binomial theorem} & \hyperlink{gls:finite-q-pochhammer-product}{Finite \(q\)-Pochhammer product} & \hyperlink{gls:finite-shifted-up-synthesis}{Finite shifted-up synthesis} \\
\hyperlink{gls:finite-sinc-product}{Finite sinc product} & \hyperlink{gls:finite-spline-approximation}{Finite spline approximation} & \hyperlink{gls:first-defect-monomial}{First-defect monomial} \\
\hyperlink{gls:first-nonzero-shell}{First nonzero shell} & \hyperlink{gls:first-saddle-correction}{First saddle correction} & \hyperlink{gls:fisher-information}{Fisher information} \\
\hyperlink{gls:fixed-point}{Fixed point} & \hyperlink{gls:fixed-scale-up-synthesis}{Fixed-scale up synthesis} & \hyperlink{gls:flat-function}{Flat function} \\
\hyperlink{gls:flat-jet}{Flat jet} & \hyperlink{gls:focused-import}{Focused import} & \hyperlink{gls:formal-asymptotic-series}{Formal asymptotic series} \\
\hyperlink{gls:formal-germ}{Formal germ} & \hyperlink{gls:formal-logarithm}{Formal logarithm} & \hyperlink{gls:formal-normal-form}{Formal normal form} \\
\hyperlink{gls:formal-power-series}{Formal power series} & \hyperlink{gls:formal-versus-analytic-identity}{Formal versus analytic identity} & \hyperlink{gls:formal-versus-realized-transseries}{Formal versus realized transseries} \\
\hyperlink{gls:formalization-crosswalk}{Formalization crosswalk} & \hyperlink{gls:fourier-coefficient}{Fourier coefficient} & \hyperlink{gls:fourier-energy}{Fourier energy} \\
\hyperlink{gls:fourier-inversion}{Fourier inversion} & \hyperlink{gls:fourier-legendre-coefficient}{Fourier-Legendre coefficient} & \hyperlink{gls:fourier-legendre-series}{Fourier-Legendre series} \\
\hyperlink{gls:fourier-mode}{Fourier mode} & \hyperlink{gls:fourier-series}{Fourier series} & \hyperlink{gls:fourier-shell}{Fourier shell} \\
\hyperlink{gls:fourier-transform}{Fourier transform} & \hyperlink{gls:fourier-convention}{Fourier-transform convention} & \hyperlink{gls:transform-of-a-distribution}{Fourier transform of a distribution} \\
\hyperlink{gls:fractional-cdf-layer-cake}{Fractional CDF layer-cake formula} & \hyperlink{gls:fractional-semigroup-law}{Fractional integration semigroup law} & \hyperlink{gls:fractional-iterate}{Fractional iterate} \\
\hyperlink{gls:fractional-order-lowering}{Fractional order lowering} & \hyperlink{gls:fractional-order-raising}{Fractional order raising} & \hyperlink{gls:fractional-part}{Fractional part} \\
\hyperlink{gls:frame-bound}{Frame bound} & \hyperlink{gls:free-cumulant}{Free cumulant} & \hyperlink{gls:free-infinite-divisibility}{Free infinite divisibility} \\
\hyperlink{gls:fubini-theorem}{Fubini theorem} & \hyperlink{gls:full-logarithmic-expansion}{Full logarithmic expansion} & \hyperlink{gls:full-moment-sequence-c}{Full moment sequence $c_n$} \\
\hyperlink{gls:functional-differential-equation}{Functional-differential equation} & \hyperlink{gls:functional-equation}{Functional equation} & \hyperlink{gls:fundamental-theorem-of-calculus}{Fundamental theorem of calculus} \\
\hyperlink{gls:gamma-derivative-tower}{Gamma-derivative tower} & \hyperlink{gls:gamma-function}{Gamma function} & \hyperlink{gls:gamma-ratio}{Gamma ratio} \\
\hyperlink{gls:gamma-tower-ratio}{Gamma-tower ratio} & \hyperlink{gls:gamma-zeta-coefficient}{Gamma-zeta coefficient} & \hyperlink{gls:gartner-ellis-theorem}{Gärtner--Ellis theorem} \\
\hyperlink{gls:gauss-lobatto-quadrature}{Gauss--Lobatto quadrature} & \hyperlink{gls:gauss-radau-quadrature}{Gauss--Radau quadrature} & \hyperlink{gls:gaussian-binomial-bounds}{Gaussian-binomial bounds} \\
\hyperlink{gls:gaussian-binomial-coefficient}{Gaussian binomial coefficient} & \hyperlink{gls:gaussian-binomial-palindromy}{Gaussian binomial palindromy} & \hyperlink{gls:gaussian-binomial-polynomial}{Gaussian binomial polynomial} \\
\hyperlink{gls:gaussian-binomial-at-minus-one}{Gaussian coefficients at \(q=-1\)} & \hyperlink{gls:gaussian-contraction}{Gaussian contraction} & \hyperlink{gls:gaussian-finite-difference}{Gaussian finite difference} \\
\hyperlink{gls:gaussian-integral}{Gaussian integral} & \hyperlink{gls:gaussian-binomial-pascal-recurrence}{Gaussian Pascal recurrence} & \hyperlink{gls:gaussian-polynomial}{Gaussian polynomial} \\
\hyperlink{gls:gaussian-quadrature}{Gaussian quadrature} & \hyperlink{gls:gaussian-rational}{Gaussian rational} & \hyperlink{gls:generalized-quantile-function}{Generalized quantile function} \\
\hyperlink{gls:generating-function}{Generating function} & \hyperlink{gls:genocchi-number}{Genocchi number} & \hyperlink{gls:genocchi-polynomial}{Genocchi polynomial} \\
\hyperlink{gls:geometric-bernoulli-polynomial}{Geometric Bernoulli polynomial} & \hyperlink{gls:geometric-interpolation-nodes}{Geometric interpolation nodes} & \hyperlink{gls:geometric-lagrange-interpolation}{Geometric Lagrange interpolation} \\
\hyperlink{gls:geometric-newton-interpolation}{Geometric Newton interpolation} & \hyperlink{gls:geometric-uniform-cdf}{Geometric uniform CDF} & \hyperlink{gls:geometric-uniform-density}{Geometric uniform density} \\
\hyperlink{gls:geometric-uniform-law}{Geometric uniform law} & \hyperlink{gls:geometric-uniform-random-series}{Geometric uniform random series} & \hyperlink{gls:germ}{Germ} \\
\hyperlink{gls:global-binary-series}{Global binary series} & \hyperlink{gls:global-differential-equation}{Global differential equation} & \hyperlink{gls:global-fabius-object}{Global Fabius object} \\
\hyperlink{gls:global-first-jet-reflection-law}{Global first-jet reflection law} & \hyperlink{gls:global-signed-extension}{Global signed extension} & \hyperlink{gls:global-uniform-spline-approximation}{Global uniform spline approximation} \\
\hyperlink{gls:gram-determinant}{Gram determinant} & \hyperlink{gls:gram-matrix}{Gram matrix} & \hyperlink{gls:gregory-coefficient}{Gregory coefficient} \\
\hyperlink{gls:grid-based-transseries}{Grid-based transseries} & \hyperlink{gls:grid-histogram}{Grid histogram} & \hyperlink{gls:grzegorczyk-computability}{Grzegorczyk computability} \\
\hyperlink{gls:hadamard-finite-part}{Hadamard finite part} & \hyperlink{gls:hahn-series}{Hahn series} & \hyperlink{gls:half-moment}{Half moment} \\
\hyperlink{gls:half-moment-numerator}{Half-moment numerator} & \hyperlink{gls:half-moment-numerator-g}{Half-moment numerator $G_n$} & \hyperlink{gls:half-moment-sequence-d}{Half-moment sequence $d_n$} \\
\hyperlink{gls:half-weighted-interval-indicator}{Half-weighted interval indicator} & \hyperlink{gls:hamming-weight}{Hamming weight} & \hyperlink{gls:hankel-moment-matrix}{Hankel moment matrix} \\
\hyperlink{gls:harmonic-number}{Harmonic number} & \hyperlink{gls:has-sum-and-tsum}{HasSum and tsum} & \hyperlink{gls:hausdorff-moment-sequence}{Hausdorff moment sequence} \\
\hyperlink{gls:head-tail-decomposition}{Head--tail decomposition} & \hyperlink{gls:heat-trace}{Heat trace} & \hyperlink{gls:heine-transformation}{Heine transformation} \\
\hyperlink{gls:herglotz-function}{Herglotz function} & \hyperlink{gls:highest-nonzero-derivative}{Highest nonzero dyadic derivative} & \hyperlink{gls:highest-set-bit-decomposition}{Highest-set-bit decomposition} \\
\hyperlink{gls:highest-set-bit-recursion}{Highest-set-bit recursion} & \hyperlink{gls:hilbert-space-projection}{Hilbert-space projection} & \hyperlink{gls:hilbert-transform}{Hilbert transform} \\
\hyperlink{gls:historical-notation-scope}{Historical notation scope} & \hyperlink{gls:hockey-stick-identity}{Hockey-stick identity} & \hyperlink{gls:holomorphic-function}{Holomorphic function} \\
\hyperlink{gls:holonomic-function}{Holonomic function} & \hyperlink{gls:homothety}{Homothety} & \hyperlink{gls:horner-scheme}{Horner scheme} \\
\hyperlink{gls:hypergeometric-term}{Hypergeometric term} & \hyperlink{gls:hypothesis-ledger}{Hypothesis ledger} & \hyperlink{gls:identity-theorem}{Identity theorem} \\
\hyperlink{gls:immutable-recovery-point}{Immutable recovery point} & \hyperlink{gls:improper-integral}{Improper integral} & \hyperlink{gls:inclusion-exclusion-principle}{Inclusion--exclusion principle} \\
\hyperlink{gls:independent-copy-equation}{Independent-copy equation} & \hyperlink{gls:independent-random-variables}{Independent random variables} & \hyperlink{gls:indicator-function}{Indicator function} \\
\hyperlink{gls:infinite-convolution}{Infinite convolution} & \hyperlink{gls:infinite-product}{Infinite product} & \hyperlink{gls:infinite-q-binomial-theorem}{Infinite \(q\)-binomial theorem} \\
\hyperlink{gls:infinite-q-pochhammer-product}{Infinite \(q\)-Pochhammer product} & \hyperlink{gls:infinite-random-series}{Infinite random series} & \hyperlink{gls:infinite-weighted-sum}{Infinite weighted sum} \\
\hyperlink{gls:inflection-point}{Inflection point} & \hyperlink{gls:information-density}{Information density} & \hyperlink{gls:injective-surjective-bijective}{Injective, surjective, and bijective} \\
\hyperlink{gls:integer-composition}{Integer composition} & \hyperlink{gls:integer-numerator-normalization}{Integer numerator normalization} & \hyperlink{gls:integer-partition}{Integer partition} \\
\hyperlink{gls:integer-valued-polynomial}{Integer-valued polynomial} & \hyperlink{gls:integrability}{Integrability} & \hyperlink{gls:integral-transform}{Integral transform} \\
\hyperlink{gls:integrality}{Integrality} & \hyperlink{gls:interior-support}{Interior of the support} & \hyperlink{gls:interpolating-polynomial}{Interpolating polynomial} \\
\hyperlink{gls:interval-partition}{Interval partition} & \hyperlink{gls:invariant-measure}{Invariant measure} & \hyperlink{gls:inverse-barnes-g-asymptotic}{Inverse Barnes-\(G\) asymptotic} \\
\hyperlink{gls:inverse-branch-closure}{Inverse-branch closure} & \hyperlink{gls:inverse-dyadic-value}{Inverse-dyadic value} & \hyperlink{gls:inverse-fabius-self-map-dynamics}{Inverse-Fabius self-map dynamics} \\
\hyperlink{gls:inverse-function}{Inverse function} & \hyperlink{gls:inverse-function-condition-number}{Inverse-function condition number} & \hyperlink{gls:inverse-gamma-asymptotic}{Inverse-gamma asymptotic} \\
\hyperlink{gls:inverse-lambert-expansion}{Inverse-Lambert expansion} & \hyperlink{gls:inverse-logarithmic-error}{Inverse-logarithmic error} & \hyperlink{gls:inverse-mellin-residue-expansion}{Inverse-Mellin residue expansion} \\
\hyperlink{gls:inverse-gap-modulus}{Inverse modulus from a gap bound} & \hyperlink{gls:inverse-power-table}{Inverse-power table} & \hyperlink{gls:inversion-formula}{Inversion formula} \\
\hyperlink{gls:fabius-isfabius-characterization}{IsFabius characterization} & \hyperlink{gls:iterated-derivative}{Iterated derivative} & \hyperlink{gls:iterated-prefix-kernel}{Iterated prefix kernel} \\
\hyperlink{gls:iterated-prefix-sum}{Iterated prefix sum} & \hyperlink{gls:jackson-fundamental-theorem}{Jackson fundamental theorem of calculus} & \hyperlink{gls:jackson-integration-by-parts}{Jackson integration by parts} \\
\hyperlink{gls:jackson-product-rule}{Jackson product rule} & \hyperlink{gls:jackson-q-derivative}{Jackson \(q\)-derivative} & \hyperlink{gls:jackson-q-integral}{Jackson \(q\)-integral} \\
\hyperlink{gls:j-fraction}{Jacobi continued fraction} & \hyperlink{gls:jacobi-cubic-identity}{Jacobi cubic identity} & \hyperlink{gls:jacobi-matrix}{Jacobi matrix} \\
\hyperlink{gls:jacobi-theta-function}{Jacobi theta function} & \hyperlink{gls:jacobi-triple-product}{Jacobi triple product} & \hyperlink{gls:jet}{Jet} \\
\hyperlink{gls:k-fold-convolution}{K-fold convolution} & \hyperlink{gls:k-fold-thue-morse-sum}{K-fold Thue-Morse sum} & \hyperlink{gls:k-function}{K-function (generalized hyperfactorial)} \\
\hyperlink{gls:kernel}{Kernel} & \hyperlink{gls:knot-multiplicity}{Knot multiplicity} & \hyperlink{gls:kolmogorov-distance}{Kolmogorov distance} \\
\hyperlink{gls:koopman-eigenfunction}{Koopman eigenfunction} & \hyperlink{gls:koopman-operator}{Koopman operator} & \hyperlink{gls:kummer-theorem}{Kummer's theorem} \\
\hyperlink{gls:l2-norm}{L2 norm} & \hyperlink{gls:lagrange-good-inversion}{Lagrange--Good multivariate inversion} & \hyperlink{gls:lagrange-interpolation}{Lagrange interpolation} \\
\hyperlink{gls:lagrange-inversion-formula}{Lagrange inversion formula} & \hyperlink{gls:lagrange-up-atom}{Lagrange up-atom} & \hyperlink{gls:lah-number}{Lah number} \\
\hyperlink{gls:lambert-polynomial}{Lambert polynomial} & \hyperlink{gls:lambert-w-branch-point}{Lambert \(W\) branch point} & \hyperlink{gls:lambert-w-function}{Lambert W function} \\
\hyperlink{gls:lambert-w-puiseux-coordinate}{Lambert \(W\) Puiseux coordinate} & \hyperlink{gls:laplace-transform}{Laplace transform} & \hyperlink{gls:large-deviation-principle}{Large-deviation principle} \\
\hyperlink{gls:laurent-expansion}{Laurent expansion} & \hyperlink{gls:laurent-log-series}{Laurent--logarithmic series} & \hyperlink{gls:law-pushforward-identity}{Law as a pushforward measure} \\
\hyperlink{gls:layer-cake-representation}{Layer-cake representation} & \hyperlink{gls:leading-coefficient-functional}{Leading-coefficient functional} & \hyperlink{gls:least-common-multiple}{Least common multiple} \\
\hyperlink{gls:least-lipschitz-constant}{Least Lipschitz constant} & \hyperlink{gls:least-squares-approximation}{Least-squares approximation} & \hyperlink{gls:least-squares-prediction-error}{Least-squares prediction error} \\
\hyperlink{gls:lebesgue-integral}{Lebesgue integral} & \hyperlink{gls:legendre-block}{Legendre block} & \hyperlink{gls:legendre-coefficient}{Legendre coefficient} \\
\hyperlink{gls:legendre-coefficient-parity-completion}{Legendre coefficient parity completion} & \hyperlink{gls:legendre-derivative-identity}{Legendre derivative identity} & \hyperlink{gls:legendre-energy-tail}{Legendre energy tail} \\
\hyperlink{gls:legendre-fenchel-transform}{Legendre--Fenchel transform} & \hyperlink{gls:legendre-gaunt-coefficient}{Legendre Gaunt coefficient} & \hyperlink{gls:legendre-linearization-coefficient}{Legendre linearization coefficient} \\
\hyperlink{gls:legendre-normalization}{Legendre normalization} & \hyperlink{gls:legendre-parity}{Legendre parity} & \hyperlink{gls:legendre-partial-sum}{Legendre partial sum} \\
\hyperlink{gls:legendre-polynomial}{Legendre polynomial} & \hyperlink{gls:legendres-factorial-valuation-formula}{Legendre's factorial-valuation formula} & \hyperlink{gls:legendre-superconvergence}{Legendre superconvergence} \\
\hyperlink{gls:leibniz-rule}{Leibniz rule} & \hyperlink{gls:levy-continuity-theorem}{Levy continuity theorem} & \hyperlink{gls:lipschitz-function}{Lipschitz function} \\
\hyperlink{gls:local-phase}{Local phase} & \hyperlink{gls:locally-finite-sum}{Locally finite sum} & \hyperlink{gls:locally-uniform-convergence}{Locally uniform convergence} \\
\hyperlink{gls:gamma-logarithm}{Log-gamma function} & \hyperlink{gls:log-periodic-asymptotic}{Log-periodic asymptotic} & \hyperlink{gls:log-periodic-function}{Log-periodic function} \\
\hyperlink{gls:log-squared-asymptotic}{Log-squared asymptotic} & \hyperlink{gls:logarithmic-delay-equation}{Logarithmic delay equation} & \hyperlink{gls:logarithmic-height}{Logarithmic height} \\
\hyperlink{gls:logarithmic-scale}{Logarithmic scale} & \hyperlink{gls:logarithmic-stieltjes-fixed-point}{Logarithmic Stieltjes fixed point} & \hyperlink{gls:lower-lambert-branch}{Lower Lambert branch} \\
\hyperlink{gls:lower-lambert-coordinate}{Lower Lambert coordinate} & \hyperlink{gls:lower-lambert-phase}{Lower-Lambert phase} & \hyperlink{gls:lower-order-term}{Lower-order term} \\
\hyperlink{gls:lucas-theorem}{Lucas's theorem} & \hyperlink{gls:maclaurin-series}{Maclaurin series} & \hyperlink{gls:mahler-equation}{Mahler equation} \\
\hyperlink{gls:mahler-function}{Mahler function} & \hyperlink{gls:major-arc}{Major arc} & \hyperlink{gls:mass-expansion}{Mass expansion} \\
\hyperlink{gls:mathematical-interface-theorem}{Mathematical interface theorem} & \hyperlink{gls:matrix-dilation}{Matrix dilation} & \hyperlink{gls:mean-value-theorem}{Mean value theorem} \\
\hyperlink{gls:with-density-measure}{Measure with density} & \hyperlink{gls:measure-zero-set}{Measure-zero set} & \hyperlink{gls:mellin-finite-part}{Mellin finite part} \\
\hyperlink{gls:mellin-inversion}{Mellin inversion} & \hyperlink{gls:mellin-pole-asymptotics}{Mellin-pole asymptotics} & \hyperlink{gls:mellin-strip}{Mellin strip} \\
\hyperlink{gls:mellin-transform}{Mellin transform} & \hyperlink{gls:meromorphic-continuation}{Meromorphic continuation} & \hyperlink{gls:mersenne-factor}{Mersenne factor} \\
\hyperlink{gls:minimum-mean-square-error}{Minimum mean-square error} & \hyperlink{gls:minkowski-sum-of-supports}{Minkowski sum of supports} & \hyperlink{gls:minor-arc}{Minor arc} \\
\hyperlink{gls:minor-arc-bound}{Minor-arc bound} & \hyperlink{gls:subspace-lattice-mobius-function}{M\"obius function of a subspace lattice} & \hyperlink{gls:partition-lattice-mobius-function}{M\"obius function of the partition lattice} \\
\hyperlink{gls:mod-gaussian-convergence}{Mod-Gaussian convergence} & \hyperlink{gls:moderate-deviation-regime}{Moderate-deviation regime} & \hyperlink{gls:module-docstring}{Module docstring} \\
\hyperlink{gls:modulus-of-continuity}{Modulus of continuity} & \hyperlink{gls:moment}{Moment} & \hyperlink{gls:moment-cumulant-transform}{Moment--cumulant transform} \\
\hyperlink{gls:moment-functional}{Moment functional} & \hyperlink{gls:moment-generating-function}{Moment-generating function} & \hyperlink{gls:moment-laurent-expansion}{Moment Laurent expansion} \\
\hyperlink{gls:moment-numerator}{Moment numerator} & \hyperlink{gls:moment-numerator-f}{Moment numerator $F_n$} & \hyperlink{gls:moment-recurrence}{Moment recurrence} \\
\hyperlink{gls:moment-sequence}{Moment sequence} & \hyperlink{gls:monic-orthogonal-polynomial}{Monic orthogonal polynomial} & \hyperlink{gls:monodromy}{Monodromy} \\
\hyperlink{gls:monotone-convergence-theorem}{Monotone convergence theorem} & \hyperlink{gls:monotone-function}{Monotone function} & \hyperlink{gls:monotone-order-isomorphism}{Monotone order isomorphism} \\
\hyperlink{gls:q-multifactorial-product}{Multiple \(q\)-Pochhammer notation} & \hyperlink{gls:multiplicative-periodic-fluctuation}{Multiplicative periodic fluctuation} & \hyperlink{gls:multiresolution-analysis}{Multiresolution analysis} \\
\hyperlink{gls:multiresolution-detail-operator}{Multiresolution detail operator} & \hyperlink{gls:mutual-information}{Mutual information} & \hyperlink{gls:namespace}{Namespace} \\
\hyperlink{gls:negative-index-bell-value}{Negative-index Bell value} & \hyperlink{gls:negative-laplace-decomposition}{Negative-Laplace decomposition} & \hyperlink{gls:negative-laplace-product}{Negative-Laplace product} \\
\hyperlink{gls:newton-identities}{Newton identities} & \hyperlink{gls:newton-iteration}{Newton iteration} & \hyperlink{gls:newton-series}{Newton series} \\
\hyperlink{gls:noise-floor}{Noise floor} & \hyperlink{gls:nome}{Nome} & \hyperlink{gls:fabius-non-elementarity}{Non-elementarity of the Fabius function} \\
\hyperlink{gls:inverse-fabius-non-elementarity}{Non-elementarity of the inverse Fabius function} & \hyperlink{gls:nonconstant-periodic-correction}{Nonconstant periodic correction} & \hyperlink{gls:noncrossing-partition}{Noncrossing partition} \\
\hyperlink{gls:nondegenerate-saddle}{Nondegenerate saddle} & \hyperlink{gls:normal-convergence}{Normal convergence} & \hyperlink{gls:normal-convergence-of-product}{Normal convergence of an infinite product} \\
\hyperlink{gls:q-pochhammer-normal-convergence}{Normal convergence of \(q\)-products} & \hyperlink{gls:normalized-reversion-coefficient}{Normalized reversion coefficient} & \hyperlink{gls:normalized-reversion-input}{Normalized reversion input} \\
\hyperlink{gls:normalized-saddle-mass}{Normalized saddle mass} & \hyperlink{gls:normalized-sinc-functions}{Normalized sinc functions} & \hyperlink{gls:normalized-volterra-operator}{Normalized Volterra operator} \\
\hyperlink{gls:normative-notation-contract}{Normative notation contract} & \hyperlink{gls:notation-collision}{Notation collision} & \hyperlink{gls:nowhere-analytic}{Nowhere analytic} \\
\hyperlink{gls:nowhere-analytic-fabius-iterate}{Nowhere-analytic Fabius iterate} & \hyperlink{gls:nowhere-analytic-inverse-iterate}{Nowhere-analytic inverse-Fabius iterate} & \hyperlink{gls:null-set}{Null set} \\
\hyperlink{gls:numerator-sequence}{Numerator sequence} & \hyperlink{gls:numerical-residual}{Numerical residual} & \hyperlink{gls:odd-cosine-series}{Odd cosine series} \\
\hyperlink{gls:odd-discrete-fourier-transform}{Odd discrete Fourier transform} & \hyperlink{gls:odd-part-of-an-integer}{Odd part of an integer} & \hyperlink{gls:one-periodic-function}{One-periodic function} \\
\hyperlink{gls:operation-count}{Operation count} & \hyperlink{gls:optimal-truncation}{Optimal truncation} & \hyperlink{gls:order-of-vanishing}{Order of vanishing} \\
\hyperlink{gls:ordered-composition}{Ordered composition} & \hyperlink{gls:ordered-composition-coefficient}{Ordered-composition coefficient} & \hyperlink{gls:ordinary-bell-polynomial}{Ordinary Bell polynomial} \\
\hyperlink{gls:ordinary-generating-function}{Ordinary generating function} & \hyperlink{gls:original-bounded-bridge}{Original--bounded Fabius bridge} & \hyperlink{gls:original-fabius-characterization}{Original compact-support characterization} \\
\hyperlink{gls:is-original-fabius}{Original Fabius predicate \(\mathrm{IsOriginalFabius}\)} & \hyperlink{gls:orthogonal-polynomial}{Orthogonal polynomial} & \hyperlink{gls:orthogonal-projection}{Orthogonal projection} \\
\hyperlink{gls:orthogonality}{Orthogonality} & \hyperlink{gls:orthonormal-polynomial}{Orthonormal polynomial} & \hyperlink{gls:oscillatory-correction}{Oscillatory correction} \\
\hyperlink{gls:p-recursive-sequence}{\(P\)-recursive sequence} & \hyperlink{gls:pade-approximant}{Padé approximant} & \hyperlink{gls:paley-wiener-theorem}{Paley--Wiener theorem} \\
\hyperlink{gls:parity}{Parity} & \hyperlink{gls:parity-power-series}{Parity-power series} & \hyperlink{gls:parseval-identity}{Parseval identity} \\
\hyperlink{gls:partial-exponential-bell-polynomial}{Partial exponential Bell polynomial} & \hyperlink{gls:partition-defect}{Partition defect} & \hyperlink{gls:partition-lattice}{Partition lattice} \\
\hyperlink{gls:partition-of-unity}{Partition of unity} & \hyperlink{gls:prime-power-pascal-row}{Pascal row modulo a prime power} & \hyperlink{gls:peano-kernel}{Peano kernel} \\
\hyperlink{gls:periodic-correction}{Periodic correction} & \hyperlink{gls:periodic-fourier-coefficient}{Periodic-correction Fourier coefficient} & \hyperlink{gls:periodic-correction-psi}{Periodic correction Psi} \\
\hyperlink{gls:periodic-mean}{Periodic mean} & \hyperlink{gls:periodicity}{Periodicity} & \hyperlink{gls:periodization}{Periodization} \\
\hyperlink{gls:perron-frobenius-operator}{Perron--Frobenius transfer operator} & \hyperlink{gls:perturbation-expansion}{Perturbation expansion} & \hyperlink{gls:phase-function}{Phase function} \\
\hyperlink{gls:physicists-hermite-polynomial}{Physicists' Hermite polynomial} & \hyperlink{gls:piecewise-polynomial}{Piecewise polynomial} & \hyperlink{gls:plancherel-theorem}{Plancherel theorem} \\
\hyperlink{gls:pochhammer-symbol}{Pochhammer symbol} & \hyperlink{gls:poincare-asymptotic-expansion}{Poincaré asymptotic expansion} & \hyperlink{gls:pointwise-convergence}{Pointwise convergence} \\
\hyperlink{gls:pointwise-equality}{Pointwise equality} & \hyperlink{gls:poisson-summation}{Poisson summation} & \hyperlink{gls:pole}{Pole} \\
\hyperlink{gls:polygamma-function}{Polygamma function} & \hyperlink{gls:polynomial-approximation}{Polynomial approximation} & \hyperlink{gls:polynomial-cell}{Polynomial cell} \\
\hyperlink{gls:polynomial-commutator}{Polynomial commutator with integration} & \hyperlink{gls:polynomial-deconvolution}{Polynomial deconvolution by the up-law} & \hyperlink{gls:polynomial-diagonal-stabilization}{Polynomial diagonal stabilization} \\
\hyperlink{gls:ring-polynomial-identity}{Polynomial identity over a ring} & \hyperlink{gls:polynomial-logarithmic-transseries}{Polynomial--logarithmic transseries} & \hyperlink{gls:polynomial-reproduction}{Polynomial reproduction} \\
\hyperlink{gls:polynomial-reproduction-order}{Polynomial reproduction order} & \hyperlink{gls:binomial-type-sequence}{Polynomial sequence of binomial type} & \hyperlink{gls:positive-compositional-iterate}{Positive compositional iterate} \\
\hyperlink{gls:positive-definite-kernel}{Positive-definite kernel} & \hyperlink{gls:positive-definite-moment-functional}{Positive-definite moment functional} & \hyperlink{gls:positive-half-line-fourier-energy}{Positive-half-line Fourier energy} \\
\hyperlink{gls:positive-part-kernel}{Positive-part kernel} & \hyperlink{gls:positive-rational-gap-sequence}{Positive rational gap lower-bound sequence} & \hyperlink{gls:post-render-source-union}{Post-render source union} \\
\hyperlink{gls:power-series}{Power series} & \hyperlink{gls:power-sum-symmetric-polynomial}{Power-sum symmetric polynomial} & \hyperlink{gls:powerset-expansion}{Powerset expansion} \\
\hyperlink{gls:prefix-polynomial}{Prefix polynomial} & \hyperlink{gls:prefix-sum-boundary-convention}{Prefix-sum boundary convention} & \hyperlink{gls:primitive-recursive-fold}{Primitive-recursive fold} \\
\hyperlink{gls:primitive-recursive-function}{Primitive-recursive function} & \hyperlink{gls:primitive-root-block}{Primitive-root block} & \hyperlink{gls:principal-branch}{Principal branch} \\
\hyperlink{gls:principal-lambert-phase}{Principal-Lambert phase} & \hyperlink{gls:probabilistic-representation}{Probabilistic representation} & \hyperlink{gls:probabilists-hermite-polynomial}{Probabilists' Hermite polynomial} \\
\hyperlink{gls:probability-density}{Probability density} & \hyperlink{gls:probability-measure}{Probability measure} & \hyperlink{gls:product-measure}{Product measure} \\
\hyperlink{gls:profile-generating-function}{Profile generating function} & \hyperlink{gls:proof-backed-exposition}{Proof-backed exposition} & \hyperlink{gls:proof-status-boundary}{Proof-status boundary} \\
\hyperlink{gls:prouhet-cancellation}{Prouhet cancellation} & \hyperlink{gls:prouhet-partition}{Prouhet partition} & \hyperlink{gls:prouhet-tarry-escott-problem}{Prouhet--Tarry--Escott problem} \\
\hyperlink{gls:provenance-ledger}{Provenance ledger} & \hyperlink{gls:public-declaration}{Public declaration} & \hyperlink{gls:public-facade}{Public facade} \\
\hyperlink{gls:puiseux-log-series}{Puiseux--logarithmic series} & \hyperlink{gls:puiseux-series}{Puiseux series} & \hyperlink{gls:pushforward-measure}{Pushforward measure} \\
\hyperlink{gls:pythagorean-identity}{Pythagorean identity} & \hyperlink{gls:q-bernstein-polynomial}{\(q\)-Bernstein polynomial} & \hyperlink{gls:q-beta-function}{\(q\)-beta function} \\
\hyperlink{gls:q-beta-integral}{\(q\)-beta integral} & \hyperlink{gls:q-binomial-coefficient}{q-binomial coefficient} & \hyperlink{gls:q-binomial-scalar-extension}{q-binomial scalar extension} \\
\hyperlink{gls:q-binomial-theorem}{q-binomial theorem} & \hyperlink{gls:q-calculus}{q-calculus} & \hyperlink{gls:q-catalan-number}{\(q\)-Catalan number} \\
\hyperlink{gls:q-digamma-function}{\(q\)-digamma function} & \hyperlink{gls:q-discrete-limit}{q-discrete limit} & \hyperlink{gls:q-factorial}{q-factorial} \\
\hyperlink{gls:q-finite-difference}{q-finite difference} & \hyperlink{gls:q-gamma-function}{\(q\)-gamma function} & \hyperlink{gls:q-gauss-summation}{\(q\)-Gauss summation} \\
\hyperlink{gls:q-integer}{q-integer} & \hyperlink{gls:q-kernel}{\(q\)-kernel of a sequence} & \hyperlink{gls:q-lucas-theorem}{\(q\)-Lucas theorem} \\
\hyperlink{gls:q-multinomial-coefficient}{\(q\)-multinomial coefficient} & \hyperlink{gls:q-pascal-summation}{\(q\)-Pascal summation} & \hyperlink{gls:q-pfaff-saalschutz-summation}{\(q\)-Pfaff--Saalschütz summation} \\
\hyperlink{gls:q-pochhammer-dissection}{\(q\)-Pochhammer dissection} & \hyperlink{gls:q-pochhammer-logarithmic-derivative}{\(q\)-Pochhammer logarithmic derivative} & \hyperlink{gls:q-pochhammer-order-derivative}{\(q\)-Pochhammer order derivative} \\
\hyperlink{gls:q-pochhammer-reversal}{\(q\)-Pochhammer reversal} & \hyperlink{gls:q-pochhammer-shift-law}{\(q\)-Pochhammer shift law} & \hyperlink{gls:q-pochhammer-symbol}{q-Pochhammer symbol} \\
\hyperlink{gls:q-product}{q-product} & \hyperlink{gls:q-richardson-extrapolation}{\(q\)-Richardson extrapolation} & \hyperlink{gls:q-series}{q-series} \\
\hyperlink{gls:q-shifted-factorial}{\(q\)-shifted factorial} & \hyperlink{gls:q-taylor-identity}{q-Taylor identity} & \hyperlink{gls:q-to-one-limit}{\(q\to1\) limit} \\
\hyperlink{gls:q-vandermonde-summation}{\(q\)-Vandermonde convolution} & \hyperlink{gls:quadratic-logarithmic-law}{Quadratic logarithmic law} & \hyperlink{gls:quadrature-weight}{Quadrature weight} \\
\hyperlink{gls:quantile-coupling}{Quantile coupling} & \hyperlink{gls:quantile-transport}{Quantile transport} & \hyperlink{gls:quantum-binomial-theorem}{Quantum binomial theorem} \\
\hyperlink{gls:quantum-multinomial-theorem}{Quantum multinomial theorem} & \hyperlink{gls:quantum-plane}{Quantum plane} & \hyperlink{gls:quotient-of-exponentials-approximation}{Quotient-of-exponentials approximation} \\
\hyperlink{gls:r-stirling-number}{\(r\)-Stirling number} & \hyperlink{gls:rademacher-sign}{Rademacher sign} & \hyperlink{gls:rademacher-sine-product}{Rademacher sine-product identity} \\
\hyperlink{gls:radius-asymptotic}{Radius asymptotic} & \hyperlink{gls:radon-nikodym-derivative}{Radon--Nikodym derivative} & \hyperlink{gls:radon-transform}{Radon transform} \\
\hyperlink{gls:ramification-index}{Ramification index} & \hyperlink{gls:random-series}{Random series} & \hyperlink{gls:rapid-decay}{Rapid decay} \\
\hyperlink{gls:rarefaction}{Rarefaction of a sequence} & \hyperlink{gls:rate-distortion-function}{Rate--distortion function} & \hyperlink{gls:rate-of-convergence}{Rate of convergence} \\
\hyperlink{gls:ratio-set}{Ratio set} & \hyperlink{gls:rational-legendre-coefficient}{Rational Legendre coefficient} & \hyperlink{gls:rational-recurrence}{Rational recurrence} \\
\hyperlink{gls:rationality}{Rationality} & \hyperlink{gls:real-analytic-function}{Real-analytic function} & \hyperlink{gls:q-reciprocal-base-duality}{Reciprocal-base \(q\)-duality} \\
\hyperlink{gls:reciprocal-gamma-function}{Reciprocal gamma function} & \hyperlink{gls:reciprocal-power-table}{Reciprocal-power table} & \hyperlink{gls:recurrence-relation}{Recurrence relation} \\
\hyperlink{gls:recursive-fast-name}{Recursive fast name} & \hyperlink{gls:recursive-modulus}{Recursive modulus} & \hyperlink{gls:reduced-dyadic-level}{Reduced dyadic level} \\
\hyperlink{gls:reference-tail}{Reference tail} & \hyperlink{gls:reference-weight}{Reference weight} & \hyperlink{gls:refinement-equation}{Refinement equation} \\
\hyperlink{gls:refinement-operator}{Refinement operator} & \hyperlink{gls:reflection-symmetry}{Reflection symmetry} & \hyperlink{gls:regular-sequence}{Regular sequence} \\
\hyperlink{gls:relative-entropy}{Relative entropy} & \hyperlink{gls:relative-error}{Relative error} & \hyperlink{gls:relative-saddle-mass}{Relative saddle mass} \\
\hyperlink{gls:relative-tail-logarithm}{Relative tail logarithm} & \hyperlink{gls:remainder-estimate}{Remainder estimate} & \hyperlink{gls:removable-singularity}{Removable singularity} \\
\hyperlink{gls:removable-value-convention}{Removable-value convention} & \hyperlink{gls:report-finite-fabius-approximant}{Report finite Fabius approximant} & \hyperlink{gls:representation-independence}{Representation independence} \\
\hyperlink{gls:representation-operator}{Representation operator} & \hyperlink{gls:reproducing-kernel-hilbert-space}{Reproducing-kernel Hilbert space} & \hyperlink{gls:reshetnikov-integers}{Reshetnikov integers} \\
\hyperlink{gls:residue}{Residue} & \hyperlink{gls:resolvent-order-hierarchy}{Resolvent-order hierarchy} & \hyperlink{gls:retired-alias}{Retired alias} \\
\hyperlink{gls:riemann-liouville-fractional-integral}{Riemann--Liouville fractional integral} & \hyperlink{gls:riemann-zeta-function}{Riemann zeta function} & \hyperlink{gls:riesz-basis}{Riesz basis} \\
\hyperlink{gls:rising-factorial}{Rising factorial} & \hyperlink{gls:rodrigues-formula}{Rodrigues formula} & \hyperlink{gls:rogers-szego-polynomial}{Rogers--Szegő polynomial} \\
\hyperlink{gls:root-of-unity-collision}{Root-of-unity collision} & \hyperlink{gls:rounding-to-nearest-dyadic}{Rounding to nearest dyadic} & \hyperlink{gls:rvachev-appell-polynomial}{Rvachev Appell polynomial} \\
\hyperlink{gls:rvachev-atom}{Rvachev atom} & \hyperlink{gls:rvachev-decoder}{Rvachev decoder} & \hyperlink{gls:rvachev-differential-difference-law}{Rvachev differential--difference law} \\
\hyperlink{gls:rvachev-polynomial-deconvolution}{Rvachev polynomial deconvolution} & \hyperlink{gls:rvachev-up-function}{Rvachev up-function} & \hyperlink{gls:saddle-central-arc}{Saddle central arc} \\
\hyperlink{gls:saddle-coefficient-recurrence}{Saddle coefficient recurrence} & \hyperlink{gls:saddle-equation}{Saddle equation} & \hyperlink{gls:saddle-hessian}{Saddle Hessian} \\
\hyperlink{gls:saddle-jet}{Saddle jet} & \hyperlink{gls:saddle-mass}{Saddle mass} & \hyperlink{gls:saddle-point}{Saddle point} \\
\hyperlink{gls:saddle-point-method}{Saddle-point method} & \hyperlink{gls:saddle-radius}{Saddle radius} & \hyperlink{gls:saddle-tail-bound}{Saddle tail bound} \\
\hyperlink{gls:saddlepoint-density-approximation}{Saddlepoint density approximation} & \hyperlink{gls:sampling-comb}{Sampling comb} & \hyperlink{gls:sampling-identity}{Sampling identity} \\
\hyperlink{gls:scalar-extension}{Scalar extension} & \hyperlink{gls:scaling-law}{Scaling law} & \hyperlink{gls:schroeder-equation}{Schröder equation} \\
\hyperlink{gls:schur-complement}{Schur complement} & \hyperlink{gls:schwartz-function}{Schwartz function} & \hyperlink{gls:second-order-eulerian-number}{Second-order Eulerian number} \\
\hyperlink{gls:second-saddle-correction}{Second saddle correction} & \hyperlink{gls:sectorial-asymptotic-expansion}{Sectorial asymptotic expansion} & \hyperlink{gls:self-convolution}{Self-convolution} \\
\hyperlink{gls:self-sampling-identity}{Self-sampling identity} & \hyperlink{gls:self-similarity}{Self-similarity} & \hyperlink{gls:semantic-api}{Semantic API} \\
\hyperlink{gls:sequential-computability}{Sequential computability} & \hyperlink{gls:series-reversion-coefficient}{Series-reversion coefficient} & \hyperlink{gls:sharp-asymptotic}{Sharp asymptotic} \\
\hyperlink{gls:sharp-constant}{Sharp constant} & \hyperlink{gls:sharp-derivative-bound}{Sharp derivative bound} & \hyperlink{gls:sharp-transfer}{Sharp transfer} \\
\hyperlink{gls:sheffer-sequence}{Sheffer sequence} & \hyperlink{gls:shifted-spline}{Shifted spline} & \hyperlink{gls:signed-global-series}{Signed global series} \\
\hyperlink{gls:signed-measure}{Signed measure} & \hyperlink{gls:signed-natural-numerator}{Signed-natural numerator} & \hyperlink{gls:signed-stirling-first-kind}{Signed Stirling numbers of the first kind} \\
\hyperlink{gls:signed-thue-morse-prefix}{Signed Thue--Morse prefix} & \hyperlink{gls:signed-weight-random-series}{Signed-weight random series} & \hyperlink{gls:sinc-function}{Sinc function} \\
\hyperlink{gls:sinc-lobe-sign}{Sinc-lobe sign} & \hyperlink{gls:sinc-product}{Sinc product} & \hyperlink{gls:small-argument-asymptotic}{Small-argument asymptotic} \\
\hyperlink{gls:smooth-compact-support}{Smooth compact support} & \hyperlink{gls:source-checkpoint}{Source checkpoint} & \hyperlink{gls:source-module}{Source module} \\
\hyperlink{gls:source-pdf-parity}{Source--PDF parity} & \hyperlink{gls:spectral-measure}{Spectral measure} & \hyperlink{gls:spectral-profile-inversion}{Spectral profile inversion} \\
\hyperlink{gls:spectral-zeta-function}{Spectral zeta function} & \hyperlink{gls:spline}{Spline} & \hyperlink{gls:square-summable-activation}{Square-summable activation profile} \\
\hyperlink{gls:stable-endpoint-evaluation}{Stable endpoint evaluation} & \hyperlink{gls:standardized-cumulant}{Standardized cumulant} & \hyperlink{gls:standardized-geometric-uniform-law}{Standardized geometric uniform law} \\
\hyperlink{gls:stars-and-bars}{Stars-and-bars principle} & \hyperlink{gls:stationary-point}{Stationary point} & \hyperlink{gls:steepest-descent}{Steepest descent} \\
\hyperlink{gls:steepest-descent-contour}{Steepest-descent contour} & \hyperlink{gls:stein-equation}{Stein equation} & \hyperlink{gls:stein-operator}{Stein operator} \\
\hyperlink{gls:stieltjes-constant}{Stieltjes constant} & \hyperlink{gls:stieltjes-inversion-formula}{Stieltjes inversion formula} & \hyperlink{gls:stieltjes-moment-sequence}{Stieltjes moment sequence} \\
\hyperlink{gls:stieltjes-pade-approximation}{Stieltjes--Padé approximation} & \hyperlink{gls:stieltjes-perron-inversion}{Stieltjes--Perron inversion} & \hyperlink{gls:stieltjes-transform}{Stieltjes transform} \\
\hyperlink{gls:stirling-formula}{Stirling formula} & \hyperlink{gls:stirling-numbers}{Stirling numbers} & \hyperlink{gls:stirling-second-kind}{Stirling numbers of the second kind} \\
\hyperlink{gls:stokes-line}{Stokes line} & \hyperlink{gls:stopped-primitive}{Stopped primitive} & \hyperlink{gls:strang-fix-condition}{Strang--Fix condition} \\
\hyperlink{gls:strict-interior-positivity}{Strict interior positivity} & \hyperlink{gls:strict-interior-range}{Strict interior range} & \hyperlink{gls:strict-unimodality}{Strict unimodality} \\
\hyperlink{gls:strip-of-convergence}{Strip of convergence} & \hyperlink{gls:sturm-liouville-equation}{Sturm-Liouville equation} & \hyperlink{gls:subgraph-fubini-identity}{Subgraph Fubini identity} \\
\hyperlink{gls:subspace-lattice}{Subspace lattice} & \hyperlink{gls:summability}{Summability} & \hyperlink{gls:superasymptotic-approximation}{Superasymptotic approximation} \\
\hyperlink{gls:supergraph-fubini-identity}{Supergraph Fubini identity} & \hyperlink{gls:support}{Support} & \hyperlink{gls:support-as-range}{Support-as-range principle} \\
\hyperlink{gls:family-support}{Support of an indexed family} & \hyperlink{gls:supremum-norm}{Supremum norm} & \hyperlink{gls:survival-function}{Survival function} \\
\hyperlink{gls:symmetric-mixed-difference}{Symmetric mixed difference} & \hyperlink{gls:symmetry}{Symmetry} & \hyperlink{gls:tail-estimate}{Tail estimate} \\
\hyperlink{gls:tail-flatness}{Tail flatness} & \hyperlink{gls:tanh-divided-slope}{Tanh divided slope} & \hyperlink{gls:taylor-block}{Taylor block} \\
\hyperlink{gls:taylor-polynomial}{Taylor polynomial} & \hyperlink{gls:taylor-series}{Taylor series} & \hyperlink{gls:terminating-basic-hypergeometric-series}{Terminating basic hypergeometric series} \\
\hyperlink{gls:terminating-taylor-series}{Terminating Taylor series} & \hyperlink{gls:q-series-termination-index}{Termination index} & \hyperlink{gls:termwise-differentiation}{Termwise differentiation} \\
\hyperlink{gls:termwise-integration}{Termwise integration} & \hyperlink{gls:test-function-space}{Test-function space} & \hyperlink{gls:theta-quasi-periodicity}{Theta quasi-periodicity} \\
\hyperlink{gls:three-term-recurrence}{Three-term recurrence for orthogonal polynomials} & \hyperlink{gls:thue-morse-bit}{Thue--Morse bit} & \hyperlink{gls:thue-morse-block-product}{Thue--Morse block product} \\
\hyperlink{gls:thue-morse-gamma-product}{Thue--Morse gamma product} & \hyperlink{gls:thue-morse-sequence}{Thue-Morse sequence} & \hyperlink{gls:thue-morse-sign}{Thue-Morse sign} \\
\hyperlink{gls:tightness}{Tightness} & \hyperlink{gls:toeplitz-lemma}{Toeplitz lemma} & \hyperlink{gls:toeplitz-row}{Toeplitz row} \\
\hyperlink{gls:tonelli-theorem}{Tonelli theorem} & \hyperlink{gls:topological-support}{Topological support} & \hyperlink{gls:total-variation}{Total variation} \\
\hyperlink{gls:total-variation-distance}{Total-variation distance} & \hyperlink{gls:totalized-definition}{Totalized definition} & \hyperlink{gls:touchard-polynomial}{Touchard polynomial} \\
\hyperlink{gls:transfer-theorem}{Transfer theorem} & \hyperlink{gls:translate-block}{Translate block} & \hyperlink{gls:translate-sum}{Translate sum} \\
\hyperlink{gls:translated-legendre-series}{Translated Legendre series} & \hyperlink{gls:transmonomial}{Transmonomial} & \hyperlink{gls:transseries}{Transseries} \\
\hyperlink{gls:transseries-composition}{Transseries composition} & \hyperlink{gls:transseries-reversion}{Transseries reversion} & \hyperlink{gls:triangle-admissibility}{Triangle admissibility} \\
\hyperlink{gls:triangular-cancellation}{Triangular cancellation} & \hyperlink{gls:triangular-recurrence}{Triangular recurrence} & \hyperlink{gls:triangular-reduction}{Triangular reduction} \\
\hyperlink{gls:truncated-power-basis}{Truncated-power basis} & \hyperlink{gls:truncation-error}{Truncation error} & \hyperlink{gls:two-adic-unit}{\(2\)-adic unit} \\
\hyperlink{gls:two-kernel}{\(2\)-kernel of a sequence} & \hyperlink{gls:two-pi-fourier-transform}{\(2\pi\)-frequency Fourier transform} & \hyperlink{gls:type-b-eulerian-number}{Type-B Eulerian number} \\
\hyperlink{gls:typed-map}{Typed map} & \hyperlink{gls:unconditional-convergence}{Unconditional convergence} & \hyperlink{gls:uniform-approximation}{Uniform approximation} \\
\hyperlink{gls:uniform-asymptotic-expansion}{Uniform asymptotic expansion} & \hyperlink{gls:uniform-convergence}{Uniform convergence} & \hyperlink{gls:uniform-distribution}{Uniform distribution} \\
\hyperlink{gls:uniform-legendre-convergence}{Uniform Legendre convergence} & \hyperlink{gls:uniform-norm}{Uniform norm} & \hyperlink{gls:uniform-random-variable}{Uniform random variable} \\
\hyperlink{gls:uniform-spline}{Uniform spline} & \hyperlink{gls:unimodality}{Unimodality} & \hyperlink{gls:unique-midpoint-inflection}{Unique midpoint inflection} \\
\hyperlink{gls:uniqueness-of-periodic-correction}{Uniqueness of the periodic correction} & \hyperlink{gls:fabius-translate-on-unit-interval}{Unit-interval Fabius--up translate} & \hyperlink{gls:unsigned-stirling-first-kind}{Unsigned Stirling numbers of the first kind} \\
\hyperlink{gls:up-synthesis-amplitude}{Up-synthesis amplitude} & \hyperlink{gls:valuation}{Valuation} & \hyperlink{gls:valuation-histogram}{Valuation histogram} \\
\hyperlink{gls:vandermonde-determinant}{Vandermonde determinant} & \hyperlink{gls:vanishing-moment}{Vanishing moment} & \hyperlink{gls:varentropy}{Varentropy} \\
\hyperlink{gls:variable-substitution}{Variable substitution} & \hyperlink{gls:variance}{Variance} & \hyperlink{gls:vertical-bromwich-line}{Vertical Bromwich line} \\
\hyperlink{gls:volterra-operator}{Volterra integral operator} & \hyperlink{gls:volterra-resolvent}{Volterra resolvent} & \hyperlink{gls:walsh-hadamard-transform}{Walsh--Hadamard transform} \\
\hyperlink{gls:wasserstein-distance}{Wasserstein distance} & \hyperlink{gls:weak-composition}{Weak composition} & \hyperlink{gls:weak-convergence}{Weak convergence} \\
\hyperlink{gls:weak-star-convergence}{Weak-star convergence} & \hyperlink{gls:weierstrass-m-test}{Weierstrass M-test} & \hyperlink{gls:weierstrass-product}{Weierstrass product} \\
\hyperlink{gls:weighted-path-expansion}{Weighted path expansion} & \hyperlink{gls:weighted-random-sum}{Weighted random sum} & \hyperlink{gls:weighted-uniform-law}{Weighted-uniform law} \\
\hyperlink{gls:well-based-support}{Well-based support} & \hyperlink{gls:well-founded-recursion}{Well-founded recursion} & \hyperlink{gls:well-poised-basic-hypergeometric-series}{Well-poised basic hypergeometric series} \\
\hyperlink{gls:wiener-chaos}{Wiener or orthogonal chaos} & \hyperlink{gls:wigner-three-j-symbol}{Wigner \(3j\)-symbol} & \hyperlink{gls:wright-omega-function}{Wright omega function} \\
\hyperlink{gls:zero-bias-transform}{Zero-bias transform} & \hyperlink{gls:zero-of-entire-function}{Zero of an entire function} & \hyperlink{gls:zero-run}{Zero run} \\
\hyperlink{gls:zeta-regularization}{Zeta regularization} & \hyperlink{gls:zeta-regularized-determinant}{Zeta-regularized determinant} & \hyperlink{gls:zonoid}{Zonoid} \\
\hyperlink{gls:zonotope}{Zonotope} &  &  \\
\end{longtable}\endgroup\newpage
\section{A}\label{sec:letter-a}

\glossaryentry{abel-equation}{Abel equation}{\statusfrontier}{Abel's equation
\[
 \psi(f(x))=\psi(x)+1
\]
conjugates iteration to translation.  It is suited to parabolic dynamics and
continuous iteration.  Solutions are not unique without normalization and
domain conditions; endpoint singularities of the Fabius inverse make a global
Abel coordinate a frontier problem.}

\glossaryentry{absolute-continuity}{Absolute continuity}{\statusclassical}{A function \(f:[a,b]\to\R\) is absolutely continuous if, for every \(\varepsilon>0\), some \(\delta>0\) makes \(\sum_j|f(b_j)-f(a_j)|<\varepsilon\) whenever the intervals \((a_j,b_j)\) are pairwise disjoint and have total length below \(\delta\). Equivalently, \(f\) has an integrable derivative almost everywhere and \(f(x)=f(a)+\int_a^x f'(t)\,\dd t\). Smooth functions are absolutely continuous on compact intervals. The concept underlies the passage between the Fabius cumulative distribution function and its density and justifies integral forms of the differential identities.}

\glossaryentry{absolute-convergence}{Absolute convergence}{\statusclassical}{A series \(\sum_{n\ge0}a_n\) is absolutely convergent when
\(\sum_{n\ge0}|a_n|<\infty\).  Absolute convergence permits rearrangement and
regrouping and is the usual route to interchanging sums with bounded linear
operations.  For series of functions one distinguishes pointwise absolute
convergence, uniform absolute convergence, and normal convergence on compact
sets.  In the current repository this language is used chiefly for transform
products, \(q\)-series, translate blocks, and analytic parameter families.  A
locally finite translate sum, such as the signed global Fabius construction at
a fixed real argument, needs no convergence theorem and should not be
misdescribed as an infinite absolutely convergent series.}

\glossaryentry{absolute-legendre-summability}{Absolute Legendre summability}{\statuslean}{A Legendre expansion is absolutely summable when
\(\sum_n|c_n|<\infty\).  Since \(|P_n(x)|\le1\) on \([-1,1]\), this implies
uniform convergence by the Weierstrass test.  Repeated integration by parts
and endpoint flatness give stronger-than-polynomial coefficient decay for the
up-function, hence absolute summability.}

\glossaryentry{activation-deficit}{Activation deficit}{\statuslean}{For a symmetric activation \(F\), an activation deficit measures the
distance from a plateau, such as \(F(x)\) near zero or \(1-F(x)\) near one.
Reflection identifies the two deficits, and the differential scaling law
relates a deficit at one scale to values at the next.  Deficit estimates feed
flatness bounds, tail inequalities, and comparison with standard activations.}

\glossaryentry{activation-profile}{Activation profile}{\statusmixed}{An activation profile is a monotone sigmoidal map used as a normalized
nonlinear transition between two plateaus.  The bounded Fabius function is a
\(C^\infty\) compact-transition activation with exactly constant exterior
tails and a symmetric derivative density.  Unlike logistic or error-function
activations, it is nowhere analytic on its transition interval and has exact
dyadic arithmetic.}

\glossaryentry{active-documentation-corpus}{Active documentation corpus}{\statusconvention}{The active documentation corpus is the set of current project-authored
Markdown and LaTeX sources outside the archival provenance tree.  It includes
the proof-backed exposition, the Lean walkthrough, the notation catalogue,
audits, and theorem-bearing frontier manuscripts.  Inclusion in the active
corpus means that shared notation and status conventions apply; it does not
mean that every mathematical assertion has already been formalized.}

\glossaryentry{gaussian-binomial-adjacent-ratio}{Adjacent Gaussian-binomial ratio}{\statuslean}{Where the denominator is nonzero,
\[
 \frac{\qbinom{n}{k+1}{q}}{\qbinom{n}{k}{q}}
 =\frac{1-q^{\,n-k}}{1-q^{\,k+1}}.
\]
The cross-multiplied polynomial identity remains meaningful at roots of unity
where the quotient itself may be undefined.  Ratio formulas are therefore
stored together with denominator hypotheses and separate cancellation lemmas.}

\glossaryentry{admissible-class}{Admissible class}{\statusclassical}{An admissible class is the collection of objects on which a defining operator is intended to act and within which existence or uniqueness is asserted. Conditions may include boundedness, continuity, prescribed endpoint values, symmetry, monotonicity, support restrictions, or differentiability. For the bounded Fabius function, a convenient admissible class consists of continuous maps with constant tails and the correct reflection law; the refinement or integration operator preserves this class and is contractive in the supremum norm. Stating the class explicitly prevents a local functional equation from being mistaken for a global characterization.}

\glossaryentry{fractional-affine-covariance}{Affine covariance of fractional integrals}{\statusmixed}{An affine change of variables transports a fractional integral with a power
of the scale factor:
\(I^\alpha(f\circ A)\) is related to
\((I^\alpha f)\circ A\) by \(|A'|^{-\alpha}\) in the orientation-preserving
case.  The terminal and side of integration must transform as well.  This
covariance connects Fabius and up normalizations at noninteger orders.}

\glossaryentry{affine-mixed-difference}{Affine mixed difference}{\statuslean}{An affine mixed difference evaluates a function at subset sums built from
affinely transformed increments.  Translation and scaling factors can be
pulled out according to degree, so the operator is suited to dyadic grids and
geometric interpolation.  The current finite-difference modules use explicit
powerset expansions to avoid ambiguous recursive conventions.}

\glossaryentry{affine-transformation}{Affine transformation}{\statusclassical}{An affine transformation of one real variable has the form \(x\mapsto ax+b\), with \(a\ne0\); it combines scaling, reflection when \(a<0\), and translation. Affine changes of variable transport support intervals, densities, derivatives, and Fourier transforms in controlled ways. The relations \(\Up(x)=F(1-|x|)\), \(F(x)=\Up(x-1)\) on \([0,1]\), and the derivative identity \(F'(x)=2\Up(2x-1)\) are built from such transformations. Dyadic self-similarity uses the special dilation \(x\mapsto2x\).}

\glossaryentry{algebraic-function-branch}{Algebraic function branch}{\statuslean}{An algebraic branch is a single-valued function \(y=f(x)\) satisfying a
nonzero polynomial relation \(P(x,y)=0\) on a domain where a branch has been
chosen.  Away from discriminant points, such a branch is analytic.  The
project formalizes this local analyticity principle and excludes the canonical
Fabius function from algebraic-branch representations on the open unit
interval.}

\glossaryentry{aliasing}{Aliasing}{\statusclassical}{Aliasing is the identification of distinct frequencies after sampling on a
lattice.  Frequencies differing by a dual-lattice multiple have the same
samples.  Exact polynomial reproduction exploits transform zeros to suppress
unwanted aliases; numerical Fourier reconstruction must instead bandlimit or
bound the aliased tail.}

\glossaryentry{all-orders-asymptotic-expansion}{All-orders asymptotic expansion}{\statusclassical}{An all-orders asymptotic expansion is a hierarchy of finite approximations valid to arbitrary prescribed order. Writing \(f(x)\sim\sum_{j\ge0}a_j\phi_j(x)\) in an asymptotic scale means that for every \(N\), subtracting the first \(N\) terms leaves \(o(\phi_{N-1}(x))\), or a stated sharper remainder. It does not generally assert convergence of the infinite formal series. For the Fabius endpoint, the repository proves expansions in inverse powers of the lower-Lambert coordinate \(\lambda(x)\), with smooth one-periodic coefficient functions that encode the surviving dyadic oscillation.}

\glossaryentry{almost-everywhere}{Almost everywhere}{\statusclassical}{A statement holds almost everywhere, abbreviated a.e., when the exceptional
set has measure zero.  Equality almost everywhere is the natural equality in
\(L^p\) spaces, but it is weaker than pointwise equality and does not preserve
chosen endpoint values.  The repository distinguishes a.e. identities arising
from integration from the pointwise functional and differential laws satisfied
by the canonical continuous Fabius and Rvachev representatives.}

\glossaryentry{almost-sure-convergence}{Almost-sure convergence}{\statusclassical}{Random variables \(X_n\) converge almost surely to \(X\) if \(\Prob(\{\omega:X_n(\omega)\to X(\omega)\})=1\). This is a pathwise mode of convergence, stronger than convergence in probability but not identical to convergence in mean. For an infinite weighted random sum \(\sum_j a_jU_j\), absolute summability of the deterministic weights together with boundedness of the independent uniforms gives almost-sure convergence directly. The resulting random variable furnishes the probabilistic representation of the Fabius distribution and its compactly supported density.}

\glossaryentry{amplitude}{Amplitude}{\statusclassical}{In an integral treated by the saddle-point or steepest-descent method, the amplitude is the factor multiplying the dominant exponential phase, as in \(I(\lambda)=\int A(z,\lambda)\exp(\lambda\Phi(z))\,\dd z\). The phase determines the saddle and Gaussian localization, while derivatives of the amplitude contribute to correction coefficients. In the Fabius Bromwich analysis, factors originating in the Laplace product, the Jacobian of the local change of variables, and reciprocal powers of the saddle coordinate form the amplitude. Separating amplitude from phase organizes the all-orders coefficient algebra.}

\glossaryentry{analysis-operator}{Analysis operator}{\statusclassical}{An analysis operator maps a function to coefficients, usually
\(Tf=(\langle f,\phi_k\rangle)_k\).  For an orthonormal basis it is an
isometry into \(\ell^2\); for a frame it is bounded with frame inequalities.
Exact moment functionals and sample-value maps are analysis operators of
different kinds and must not be conflated.}

\glossaryentry{analytic-continuation}{Analytic continuation}{\statusclassical}{Analytic continuation extends a holomorphic function from one domain to a larger connected domain by requiring agreement on their overlap. Uniqueness follows from the identity theorem: two holomorphic continuations that agree on a set with an accumulation point agree everywhere on the connected component. Mellin transforms and zeta expressions in the endpoint analysis are first defined in regions of absolute convergence and then continued meromorphically. The procedure must be distinguished from merely assigning a formal value to a divergent expression; poles and finite parts have to be tracked explicitly.}

\glossaryentry{analytic-germ}{Analytic germ}{\statusclassical}{The analytic germ of a function at a point records its behavior on some unspecified neighborhood of that point, with two representatives identified when they agree on a smaller neighborhood. A smooth function has an analytic germ at \(x_0\) precisely when it equals a convergent power series near \(x_0\). At dyadic interior points the Fabius formal Taylor series terminates, but the function does not agree locally with that polynomial; hence no analytic germ is present there. This distinction is central to proving nowhere analyticity despite exact finite Taylor data.}

\glossaryentry{analytic-locus}{Analytic locus}{\statusclassical}{The analytic locus of a smooth real function is the set of points at which it is real analytic on some neighborhood. It is open by definition. A constant tail is analytic at every point strictly outside the transition interval, whereas flatness at a boundary point need not imply analyticity there. The repository identifies the exact analytic loci of the bounded Fabius function, the compactly supported up-function, and the signed global extension, carefully separating exterior constant regions, support endpoints, and the dense set of dyadic points where Taylor series terminate without representing the function.}

\glossaryentry{analytic-moment}{Analytic moment}{\statusclassical}{An analytic moment is a moment regarded not only as an integral but as a Taylor coefficient of an analytic transform. If \(M(z)=\int\e^{zx}\,\mu(\dd x)\) is entire, then \(m_n=M^{(n)}(0)\). This viewpoint gives Cauchy estimates, generating-function identities, and scalar extension to complex arguments. The repository's analytic-moment modules connect the exact Fabius moment recurrences with the entire complex moment-generating function.}

\glossaryentry{analytic-zero-multiplicity}{Analytic zero multiplicity}{\statusclassical}{A holomorphic function has a zero of multiplicity \(r\) at \(z_0\) when
\(f(z)=(z-z_0)^rg(z)\) with \(g(z_0)\ne0\).  Equivalently its first
\(r-1\) derivatives vanish and the \(r\)-th does not.  Product
factorization makes multiplicities additive and supports exact jet-preserving
Fourier extrapolants.}

\glossaryentry{analyticity-at-a-point}{Analyticity at a point}{\statusclassical}{A real function is analytic at \(x_0\) when it agrees with a convergent power
series on some neighborhood of \(x_0\).  Possessing derivatives of all orders
at the point is insufficient.  For Fabius functions, terminating or flat jets
at dense dyadic configurations contradict neighborhood power-series
representations, giving nowhere-analytic rather than merely endpoint
nonanalytic behavior.}

\glossaryentry{angular-frequency-fourier-transform}{Angular-frequency Fourier transform}{\statusconvention}{The angular convention is
\(\mathcal F_\omega f(\omega)=\int f(x)e^{-\ii\omega x}\,\dd x\).
It is related to the \(2\pi\)-frequency transform by
\(\omega=2\pi\xi\).  Product factors, inversion constants, derivative
multipliers, and exponential type all change accordingly; formulas from the
two conventions cannot be mixed termwise.}

\glossaryentry{antidiagonal-polynomial}{Antidiagonal polynomial}{\statusfrontier}{An antidiagonal polynomial records values of a two-parameter repeated-sum
array along \(k-n=r\) or \(n+k=r\), depending on the indexing convention.
After the vanishing region is removed, many such diagonals are polynomial in
the remaining index.  Explicit diagonal formulas enable constant-time
evaluation far inside a repeated-summation table.}

\glossaryentry{appell-deconvolution}{Appell deconvolution}{\statuslean}{Given a random variable \(X\) with moment-generating function \(M_X(z)\), the
deconvolution Appell polynomials have generating function
\[
 \frac{e^{xz}}{M_X(z)}.
\]
They satisfy
\(\mathbb E[A_n(x+X)]=x^n\).  For the Rvachev law, reciprocal sinc or
moment-series coefficients produce exact polynomial reconstruction formulas.}

\glossaryentry{appell-sequence}{Appell sequence}{\statuslean}{A polynomial sequence \(A_n(x)\) is Appell when
\(A_n'(x)=nA_{n-1}(x)\), equivalently its exponential generating function is
\[
 A(z)e^{xz}
\]
for a series \(A\) with nonzero constant term.  Bernoulli and Euler
polynomials are examples.  Moment-based Appell sequences encode convolution
and deconvolution with the Rvachev law.}

\glossaryentry{approximate-identity}{Approximate identity}{\statusclassical}{An approximate identity is a family of kernels \(K_n\) whose mass is normalized, whose mass concentrates near the origin, and whose convolution with a suitable function tends to that function. In probability it is often realized by laws of random variables whose tails shrink to zero; in spline theory it appears through increasingly fine convolution kernels. The finite centered spline approximants to the Fabius density or distribution behave in this way: they are explicit piecewise-polynomial convolutions and converge uniformly after the support and centering are normalized correctly.}

\glossaryentry{approximation-theory}{Approximation theory}{\statusclassical}{Approximation theory studies how functions are represented or approximated by simpler classes such as polynomials, splines, rational functions, or translates of a fixed kernel. It asks for convergence, rates, optimality, stability, and constructive algorithms. The up-function belongs to Rvachev's atomic-function approach, while the repository develops both centered spline approximations with explicit uniform error and Fourier-Legendre polynomial approximations with least-squares optimality.}

\glossaryentry{arbitrary-phase-synthesis}{Arbitrary-phase synthesis}{\statuslean}{Arbitrary-phase synthesis reconstructs a polynomial for every real shift
\(\theta\) using samples at \((k+\theta)/M\) and correspondingly shifted
up-atoms.  Phase-uniform exactness is stronger than reconstruction on the
unshifted lattice and is the property classified by the composite-mesh
theorems.}

\glossaryentry{arbitrary-phase-up-synthesis}{Arbitrary-phase up synthesis}{\statuslean}{Arbitrary-phase synthesis permits the centers of the wide up-atoms to be
shifted by a common real phase while retaining exact polynomial reproduction
on the target interval.  The coefficients change through an explicitly
invertible interpolation system.  The phase is a geometric placement
parameter, not the periodic endpoint phase from the Lambert asymptotic.}

\glossaryentry{arithmetic-ray}{Arithmetic ray}{\statusfrontier}{An arithmetic ray is a sequence of frequencies or complex arguments
selected by a fixed arithmetic pattern, such as odd multiples times powers of
two.  Evaluating a dyadic product along such rays isolates valuation classes,
lobe signs, and extremal decay behavior.  Ray asymptotics need not describe
uniform behavior over the full real axis.}

\glossaryentry{associahedron}{Associahedron}{\statusclassical}{The associahedron is a convex polytope whose face structure encodes
parenthesizations or polygon dissections.  Indexing conventions differ:
\(K_n\) may denote dimension \(n-2\), \(n-1\), or triangulations of an
\(n\)-gon.  The notation catalogue replaces ambiguous historical \(K_n\)
symbols by a dimension-explicit form before formulas are merged.}

\glossaryentry{asymptotic-coefficient-jet}{Asymptotic coefficient jet}{\statusmixed}{An asymptotic coefficient jet is the finite list of phase and amplitude
derivatives needed to compute a chosen number of saddle corrections.  Bell
polynomials and Wick/Gaussian moments turn the jet into explicit coefficients.
The order of the required jet grows with the requested expansion depth.}

\glossaryentry{asymptotic-equivalence}{Asymptotic equivalence}{\statusclassical}{For nonzero comparison functions, \(f(x)\sim g(x)\) as \(x\to a\) means \(f(x)/g(x)\to1\). This multiplicative statement is stronger than equality of leading logarithmic orders. If \(\log f=\log g+o(1)\), then asymptotic equivalence follows after exponentiation; an \(O(1)\) logarithmic discrepancy generally does not suffice. In the Fabius endpoint problem, the periodic correction must be included in the exponential main term: omitting a nonconstant bounded function changes the ratio by a nonconvergent oscillatory factor and destroys the claimed equivalence.}

\glossaryentry{asymptotic-expansion}{Asymptotic expansion}{\statusclassical}{An asymptotic expansion expresses a function by successively smaller comparison terms near a limiting regime. If \(\phi_{j+1}=o(\phi_j)\), the notation \(f\sim\sum_{j\ge0}a_j\phi_j\) means that every finite truncation approximates \(f\) to the scale of its last retained term. Coefficients may be constants, polynomials, or periodic functions. The Fabius development uses logarithmic expansions, saddle expansions, and inverse-Lambert expansions; their remainders are stated quantitatively rather than inferred from a formal infinite series.}

\glossaryentry{asymptotic-scale}{Asymptotic scale}{\statusclassical}{An asymptotic scale is an ordered family \((\phi_j)\) with \(\phi_{j+1}(x)=o(\phi_j(x))\) in the limit under study. It supplies the notion of successive order needed for an asymptotic expansion. Powers of \(1/\lambda\), where \(\lambda\to\infty\), form the main scale in the corrected Fabius saddle analysis; powers of \(1/|\log x|\) are a coarser equivalent scale after expanding the lower Lambert coordinate. Logarithmic and log-periodic factors can modify the scale and must not be silently absorbed into constants.}

\glossaryentry{atom-of-a-measure}{Atom of a measure}{\statusclassical}{An atom is a measurable set of positive measure containing no smaller
measurable subset of intermediate positive measure; for a real law, a point
mass at \(a\) is measured by \(\Prob(X=a)\).  A distribution with a continuous
density is atomless.  Infinite weighted sums can lose atoms even when finite
approximants have piecewise-polynomial densities.}

\glossaryentry{atomic-function}{Atomic function}{\statusclassical}{In Rvachev's terminology, an atomic function is a compactly supported smooth solution of a functional-differential or refinement equation from which broader systems of localized basis functions can be generated. The up-function is the canonical example. Its support is finite, all derivatives vanish at the support boundary, and dilates and translates inherit exact algebraic relations. Atomic functions are used in approximation theory and numerical methods because they combine localization with very high regularity, although unlike polynomials or analytic kernels they are typically nowhere analytic in the interior.}

\glossaryentry{atomless-law}{Atomless law}{\statusclassical}{A probability law is atomless when every singleton has mass zero.
Continuity of its CDF is equivalent to atomlessness for real distributions.
The geometric uniform and Fabius laws are atomless and in fact have smooth
densities, allowing ordinary inverse functions rather than quantile intervals
on their strict support interiors.}

\glossaryentry{automatic-sequence}{Automatic sequence}{\statusclassical}{A sequence is \(k\)-automatic if a finite automaton reading the base-\(k\) digits of \(n\) outputs its \(n\)-th term. The Thue--Morse sequence is the fundamental \(2\)-automatic example: the automaton tracks the parity of the number of one bits. Automatic structure yields digit recurrences, finite block factorization, and efficient generation. These properties drive the Fabius signed translate and dyadic cancellation formulas.}

\glossaryentry{averaging-operator}{Averaging operator}{\statusclassical}{An averaging operator replaces a value by an integral or finite mean over translated or rescaled values. A typical probability operator is \((Tf)(x)=\int f(x-y)\,\mu(\dd y)\), namely convolution with a probability measure. Such operators preserve positivity and constants, contract the supremum norm, and smooth irregularities. The fixed-point construction of the Fabius law can be formulated with an operator induced by a scaled uniform random variable; iterating it corresponds to adding independent dyadic-scale summands and produces finite spline approximants.}

\section{B}\label{sec:letter-b}

\glossaryentry{b-spline}{B-spline}{\statusclassical}{A B-spline is a compactly supported piecewise polynomial obtained from repeated convolution of interval indicators, with degree and smoothness increasing under convolution. Uniform B-splines use equally spaced knots; centered forms are symmetric about the origin. Finite convolutions of the dyadically scaled uniform densities in the probabilistic Fabius construction are nonuniformly scaled box splines, and after a common-grid reformulation they yield explicit centered finite splines. These approximants have rational dyadic data, controlled support, and a certified uniform error.}

\glossaryentry{bailey-lemma}{Bailey lemma}{\statusfrontier}{The Bailey lemma maps one Bailey pair to another by multiplying and summing
with additional \(q\)-Pochhammer factors.  Iteration produces a Bailey chain
and many Rogers--Ramanujan type identities.  In the Fabius project it is a
frontier organizing principle for repeated triangular transforms, not a
substitute for checking convergence and denominator hypotheses.}

\glossaryentry{bailey-pair}{Bailey pair}{\statusfrontier}{A pair of sequences \((\alpha_n,\beta_n)\) is a Bailey pair relative to
\(a\) when \(\beta_n\) is a specified triangular \(q\)-hypergeometric transform
of \(\alpha_0,\ldots,\alpha_n\), commonly
\[
 \beta_n=\sum_{r=0}^n
 \frac{\alpha_r}{(q;q)_{n-r}(aq;q)_{n+r}}.
\]
The normalization varies.  Bailey pairs convert finite coefficient data into
families of product and series identities.}

\glossaryentry{balanced-basic-hypergeometric-series}{Balanced basic hypergeometric series}{\statusclassical}{A terminating basic hypergeometric series is balanced when the product of its
lower parameters is a prescribed power of \(q\) times the product of its upper
parameters, with the exact exponent determined by the \({}_r\phi_s\)
normalization.  Balance is what makes the \(q\)-Pfaff--Saalschütz product
evaluation possible.  It is stronger than mere termination.}

\glossaryentry{banach-fixed-point-theorem}{Banach fixed-point theorem}{\statusclassical}{The Banach fixed-point theorem states that a strict contraction \(T\) on a complete metric space has a unique fixed point. If \(d(Tx,Ty)\le qd(x,y)\) with \(0<q<1\), then the iterates \(T^n x_0\) converge geometrically to that fixed point and satisfy explicit a posteriori and a priori error bounds. In the Fabius construction, the ambient space is a closed class of bounded continuous functions with the supremum metric, and the defining averaging/refinement operator is contractive. The theorem simultaneously proves existence, uniqueness, and a certified approximation scheme.}

\glossaryentry{barnes-g-function}{Barnes \(G\)-function}{\statusfrontier}{Barnes' \(G\)-function satisfies
\(G(z+1)=\Gamma(z)G(z)\) and \(G(1)=1\).  It interpolates products of
factorials and has a large-\(z\) logarithmic expansion led by
\(\tfrac12z^2\log z\).  Inverting \(G\) therefore requires a
polynomial--logarithmic transseries on a square-root logarithmic scale.}

\glossaryentry{base-q-digit}{Base-\(q\) digit}{\statusmixed}{For an integer base \(q\ge2\), the \(j\)-th digit of \(n\) is
\(d_j(n)=\lfloor n/q^j\rfloor\bmod q\).  Then
\(n=\sum_{j\ge0}d_j(n)q^j\) with finite support.  Base-\(q\) digit formulas
generalize binary weight, support carry theorems, and enter \(q\)-Lucas and
prime-power valuation arguments.}

\glossaryentry{base-q-digit-sum}{Base-\(q\) digit sum}{\statusclassical}{The base-\(q\) digit sum is \(s_q(n)=\sum_jd_j(n)\).  For prime \(p\),
Legendre's formula connects it to \(v_p(n!)\), and Kummer's theorem connects
digit carries to binomial valuations.  At \(q=2\), its parity is the
Thue--Morse bit and its value is the Hamming weight.}

\glossaryentry{basic-hypergeometric-series}{Basic hypergeometric series}{\statusclassical}{A basic hypergeometric series is a \(q\)-series whose consecutive
coefficient ratio is rational in \(q^n\).  In standard notation,
\({}_r\phi_s\) is built from products \((a_i;q)_n/(b_j;q)_n\), the factor
\((q;q)_n\), and a convention-dependent power of
\((-1)^nq^{\binom n2}\).  Bilateral analogues sum over \(n\in\Z\).
The current repository fixes these normalizations explicitly and formalizes
selected transformation and summation theorems, including Heine,
\(q\)-Gauss, and terminating \(q\)-Pfaff--Saalschutz forms.  Parameter
denominators and convergence regions are part of every statement.}

\glossaryentry{basic-hypergeometric-r-s-phi}{Basic hypergeometric series \({}_r\phi_s\)}{\statuslean}{With one standard normalization,
\[
 {}_r\phi_s\!\left[\begin{matrix}a_1,\ldots,a_r\\
 b_1,\ldots,b_s\end{matrix};q,z\right]
 =\sum_{n\ge0}
 \frac{(a_1,\ldots,a_r;q)_n}{(q,b_1,\ldots,b_s;q)_n}
 \left((-1)^nq^{n(n-1)/2}\right)^{1+s-r}z^n .
\]
The exponent convention is part of the notation.  Denominator parameters,
termination, and convergence are stated explicitly in every theorem.}

\glossaryentry{bell-hessenberg-determinant}{Bell--Hessenberg determinant}{\statusfrontier}{A Bell--Hessenberg determinant is a lower-Hessenberg determinant whose
entries are a moment, cumulant, or coefficient sequence and whose superdiagonal
contains simple integer factors.  Expanding by the last row reproduces a Bell
recurrence.  The exact signs and factorial normalization depend on whether
ordinary or exponential series are used.}

\glossaryentry{bell-number}{Bell number}{\statusclassical}{The Bell number \(B_n\) counts set partitions of an \(n\)-element set:
\(B_n=\sum_k\genfrac{\{}{\}}{0pt}{}{n}{k}\).  Its exponential generating function
is \(\exp(e^z-1)\).  Bell numbers organize moments, composition derivatives,
determinants, and asymptotic reversion coefficients.}

\glossaryentry{bell-polynomial}{Bell polynomial}{\statusclassical}{Bell polynomials encode the combinatorics of differentiating a composite exponential or logarithm. The complete exponential Bell polynomial satisfies \(\exp(\sum_{m\ge1}x_m t^m/m!)=\sum_{n\ge0}B_n(x_1,\ldots,x_n)t^n/n!\); partial Bell polynomials refine the number of factors. They provide a compact form of the Faà di Bruno formula. In saddle expansions, they convert phase and amplitude jets into coefficients of the exponential series, while logarithmic coefficients can also be expressed through ordered compositions. The repository often spells out these finite sums to facilitate formal verification.}

\glossaryentry{bell-polynomial-asymptotic-coefficient}{Bell-polynomial asymptotic coefficient}{\statusmixed}{When an exponent or logarithm has a formal expansion, complete Bell
polynomials convert its coefficient sequence into the coefficients of the
exponential.  Partial Bell polynomials organize a fixed number of factors.
This gives finite, explicit formulas for saddle corrections, gamma derivative
jets, and moment--cumulant expansions.}

\glossaryentry{bernoulli-number}{Bernoulli number}{\statusclassical}{The Bernoulli numbers \(B_n\) are rational numbers defined by \(t/(\e^t-1)=\sum_{n\ge0}B_n t^n/n!\), with the convention \(B_1=-1/2\). They govern sums of powers, Euler--Maclaurin summation, and the Taylor expansions of trigonometric and Bose-type kernels. The Fabius arithmetic development derives Bernoulli recurrences related to moment sequences and reciprocal-sine products. Because two sign conventions for \(B_1\) occur in the literature, a formal development must fix the generating function rather than rely on the name alone.}

\glossaryentry{bernoulli-recurrence}{Bernoulli recurrence}{\statusclassical}{A Bernoulli recurrence is a triangular identity determining Bernoulli numbers from earlier values, classically \(\sum_{k=0}^{n}\binom{n+1}{k}B_k=0\) for \(n\ge1\). Equivalent recurrences arise by multiplying or differentiating the generating function \(t/(\e^t-1)\). In the repository, Bernoulli-type relations connect coefficients in reciprocal trigonometric expansions with Fabius moment recurrences. Triangularity means the newest coefficient appears with a nonzero scalar and can therefore be solved exactly over the rationals.}

\glossaryentry{berry-esseen-bound}{Berry--Esseen bound}{\statusfrontier}{A Berry--Esseen bound quantifies the Kolmogorov distance to a Gaussian by a
third absolute moment divided by a variance power, summed over independent
terms.  It gives an \(O(\epsilon)\)-type rate under the appropriate triangular
array normalization.  Symmetry may improve formal cumulant order but does not
automatically improve the classical absolute-moment bound.}

\glossaryentry{bessel-inequality}{Bessel inequality}{\statusclassical}{Bessel's inequality states that for an orthonormal system \((e_n)\) in a Hilbert space, \(\sum|\langle f,e_n\rangle|^2\le\|f\|^2\). Equality holds for a complete basis and becomes Parseval's identity. Applied to normalized Legendre polynomials, it bounds coefficient squares and the \(L^2\) error of partial sums. Stronger coefficient estimates are needed for the repository's absolute uniform convergence theorem.}

\glossaryentry{big-o-and-little-o}{Big-O and little-o notation}{\statusconvention}{The relation \(f(x)=\bigO(g(x))\) means that \(|f(x)|\le C|g(x)|\) near the specified limit for some constant \(C\); \(f(x)=\smallo(g(x))\) means \(f(x)/g(x)\to0\). Uniform qualifiers indicate that the constant or convergence is independent of auxiliary variables. These symbols describe error scales, not hidden equalities. In the corrected Fabius asymptotic, an \(\bigO(1/\lambda(x))\) logarithmic remainder is meaningful only after the nonconstant periodic main term has been retained.}

\glossaryentry{bijection}{Bijection}{\statusclassical}{A bijection is a map that is both injective and surjective, so every target element has exactly one preimage. A continuous strictly monotone function from an interval onto an interval is a bijection and admits a continuous inverse. The bounded Fabius function is constant outside \([0,1]\), but its restriction to \([0,1]\) is a strictly increasing bijection onto \([0,1]\). This fact licenses the inverse Fabius function and makes inverse-dyadic questions meaningful: for a dyadic ordinate one asks for the unique abscissa with that value.}

\glossaryentry{bilateral-basic-hypergeometric-series}{Bilateral basic hypergeometric series}{\statusmixed}{A bilateral series \({}_r\psi_s\) sums over all \(n\in\Z\) rather than
\(n\ge0\).  Negative-index \(q\)-Pochhammer symbols are interpreted through
reciprocal finite products.  Convergence typically requires an annulus rather
than a disk, and zeros of denominator factors must be excluded.  Bilateral
series appear in theta and dissection identities but are not interchangeable
with unilateral \({}_r\phi_s\) series.}

\glossaryentry{bilateral-laplace-transform}{Bilateral Laplace transform}{\statusclassical}{The bilateral Laplace transform of a suitable function \(f:\R\to\C\) is
\[
 \mathcal L_{\mathrm{bil}}f(s)=\int_{-\infty}^{\infty}e^{-sx}f(x)\,\dd x .
\]
It differs from the one-sided transform by integrating over the whole real
line, and its convergence domain is generally a vertical strip.  Compact
support makes the transform entire.  The distinction matters when passing
between the compactly supported up-function, the bounded CDF, and transforms
of global signed extensions.}

\glossaryentry{binary-block-factorization}{Binary block factorization}{\statusmixed}{When an index is split as \(2^m a+r\) with \(0\le r<2^m\), nonoverlap of
binary digits gives
\(\varepsilon_{2^ma+r}=\varepsilon_a\varepsilon_r\).
This block factorization separates coarse and fine scales and is the key
multiplicative identity behind complete-block cancellation and self-similar
prefix arrays.}

\glossaryentry{binary-digit-sum}{Binary digit sum}{\statusclassical}{The binary digit sum \(\wt(n)\), also called the Hamming weight or popcount, is the number of ones in the base-two expansion of a nonnegative integer \(n\). Its parity generates the Thue--Morse sign \(\varepsilon_n=(-1)^{\wt(n)}\). It obeys \(\wt(2n)=\wt(n)\) and \(\wt(2n+1)=\wt(n)+1\), which yield the fundamental two-block recursions. The Fabius development uses both the exact integer \(\wt(n)\) and its parity in dyadic cancellation, signed translate sums, and logarithmic identities.}

\glossaryentry{binary-expansion}{Binary expansion}{\statusclassical}{A binary expansion represents a nonnegative real number as \(x=\sum_{j}b_j2^j\) or, on the unit interval, as \(x=\sum_{j\ge1}b_j2^{-j}\) with digits \(b_j\in\{0,1\}\). Dyadic rationals have two infinite expansions, one terminating and one ending in infinitely many ones, so algorithms must impose a convention or prove representation independence. The global binary reduction in the repository peels off scale blocks according to the set bits and produces an absolutely convergent series whose value is independent of harmless encoding choices.}

\glossaryentry{binary-floor-digit-formula}{Binary floor-digit formula}{\statuslean}{A binary floor-digit formula expresses a bit, digit sum, or Thue--Morse
value entirely through floor functions.  Such formulas extend naturally to a
real variable between integer jumps and can be inserted into finite sums,
integrals, and interpolation identities.  Their discontinuity set and
endpoint convention must be stated explicitly.}

\glossaryentry{binary-reduction-series}{Binary reduction series}{\statusmixed}{A binary reduction series is a scale-by-scale expansion obtained by iterating an identity that reduces the argument according to its binary digits. In the Fabius setting, a Taylor block is extracted at each active binary scale, and the residual argument is shifted and rescaled. The finite reductions telescope; passing to the limit yields an absolutely convergent representation of the signed global extension. Unlike a mere digit algorithm, the construction is an analytic identity and includes explicit coefficients, convergence estimates, and a proof that the series agrees with the bounded Fabius function on \([0,1]\).}

\glossaryentry{binary-reflection-recursion}{Binary reflection recursion}{\statuslean}{A binary reflection recursion reduces an index in the upper half of a
power-of-two block to its reflected lower-half index, multiplying by a
Thue--Morse sign or adding a correction.  It is the finite-sequence analogue
of Fabius reflection and underlies fast prefix, dyadic-value, and automatic
kernel algorithms.}

\glossaryentry{binomial-coefficient}{Binomial coefficient}{\statusclassical}{The binomial coefficient \(\binom nk=n!/(k!(n-k)!)\) counts \(k\)-element subsets of an \(n\)-element set and is the coefficient of \(x^k\) in \((1+x)^n\). Pascal's recurrence and the binomial theorem make it ubiquitous in derivative scaling, moment recurrences, Taylor expansions, and finite differences. The Fabius formulas combine ordinary binomial coefficients with Gaussian \(q\)-binomial coefficients; keeping the two notations distinct is essential because the latter depend on a base parameter and encode geometric rather than equally spaced nodes.}

\glossaryentry{binomial-basis}{Binomial polynomial basis}{\statusclassical}{The polynomials \(\binom{x}{k}=x^{\underline{k}}/k!\) form a basis for
polynomials in \(x\).  Forward difference acts simply:
\(\Delta\binom{x}{k}=\binom{x}{k-1}\).  Integer-valued polynomials have
integer coefficients in this basis, making it natural for repeated prefix
sums and exact diagonal formulas.}

\glossaryentry{binomial-transform}{Binomial transform}{\statusclassical}{The binomial transform sends a sequence \((a_k)\) to \(b_n=\sum_{k=0}^{n}\binom nk a_k\), with signed variants obtained by inserting \((-1)^k\). It is inverted by the corresponding signed transform. Such transforms arise when moments are shifted, when an exponential generating function is multiplied by \(\e^z\), or when a variable is replaced by \(1-x\). In the Fabius arithmetic, the full moments and half moments are connected by finite even-index binomial sums produced by the symmetry and translation between the bounded and centered forms.}

\glossaryentry{biorthogonal-system}{Biorthogonal system}{\statusfrontier}{Two families \((\phi_n)\) and \((\psi_n)\) are biorthogonal when
\(\langle\phi_n,\psi_m\rangle=\delta_{nm}\).  Such pairs support coefficient
extraction without orthogonality of either family alone.  Appell
deconvolution and translated-kernel systems suggest natural biorthogonal
questions in the Fabius/Rvachev setting.}

\glossaryentry{bit-extraction}{Bit extraction}{\statuslean}{For \(n\in\N\), the \(j\)-th bit can be extracted by
\[
 b_j(n)=\left\lfloor\frac n{2^j}\right\rfloor
       -2\left\lfloor\frac n{2^{j+1}}\right\rfloor .
\]
This equals the remainder of \(\lfloor n/2^j\rfloor\) modulo two.  It links
Lean's bit-testing operations to floor formulas used in analytic sums and
makes binary digit identities available without choosing a string
representation.}

\glossaryentry{boolean-cube-difference}{Boolean-cube mixed difference}{\statusmixed}{For increments \(h_1,\ldots,h_m\), the mixed difference is an alternating
sum over the Boolean cube,
\[
 \Delta_{h_1}\cdots\Delta_{h_m}f(x)
 =\sum_{S\subseteq[m]}(-1)^{m-|S|}
 f\!\left(x+\sum_{j\in S}h_j\right).
\]
It annihilates polynomials of degree below \(m\) and is the continuous-node
analogue of Prouhet cancellation.}

\glossaryentry{boolean-cumulant}{Boolean cumulant}{\statusfrontier}{Boolean cumulants linearize Boolean convolution and are related to moments
through interval partitions.  Their generating functions satisfy a simple
reciprocal relation.  Free, Boolean, and classical cumulants share low-order
patterns but have different partition lattices and independence notions.}

\glossaryentry{borel-probability-measure}{Borel probability measure}{\statusclassical}{A Borel probability measure on a topological space assigns total mass one to the sigma-algebra generated by open sets. On \(\R\), it is determined by its cumulative distribution function and, when absolutely continuous, by a density. The law of the infinite dyadically weighted sum of independent uniform variables is a compactly supported Borel probability measure. Its distribution function is the bounded Fabius function, while a centered affine image has density equal to the Rvachev up-function after the chosen normalization.}

\glossaryentry{borel-summability}{Borel summability}{\statusfrontier}{A formal series is Borel summable in a direction when its Borel transform has
suitable analytic continuation and growth there and the directional Laplace
integral exists.  The resulting analytic function has the original series as
an asymptotic expansion.  Borel summability is stronger than formal
manipulability and is not claimed for a repository transseries unless an
analytic continuation theorem is supplied.}

\glossaryentry{borel-transform}{Borel transform}{\statusfrontier}{For a factorially divergent formal series
\(\sum_{n\ge0}a_nx^{-n-1}\), its formal Borel transform is
\(\sum_{n\ge0}a_n\xi^n/n!\).  If this series can be analytically continued and
Laplace transformed along a direction, the result is a Borel sum.  Singularities
of the Borel transform encode exponentially small sectors and Stokes jumps.}

\glossaryentry{bose-einstein-kernel}{Bose-Einstein kernel}{\statusclassical}{A Bose-Einstein kernel is a function built from \((\e^t-1)^{-1}\), often \(t/(\e^t-1)\) or \((\e^t-1)^{-1}-t^{-1}\), whose Mellin transform produces Gamma and zeta factors. Near zero it has a Bernoulli expansion and usually a pole that must be subtracted before integration. In the Fabius negative-Laplace analysis, a related dyadic harmonic sum is decomposed into polynomial-logarithmic terms plus a periodic fluctuation. The finite-part treatment of the kernel isolates the Mellin pole and makes the Gamma--zeta Fourier coefficients explicit.}

\glossaryentry{bose-finite-part-integral}{Bose finite-part integral}{\statusclassical}{A Bose finite-part integral is a regularized integral involving a Bose-Einstein kernel whose endpoint singular terms have been explicitly removed. If an integrand has an expansion \(a_{-m}t^{-m}+\cdots+a_{-1}t^{-1}+O(1)\), one subtracts the singular part, integrates the remainder, and adds the canonically determined finite contributions. In the repository this construction supplies constants and periodic corrections in the Mellin analysis of the negative-Laplace product. It is a Hadamard finite part, not an ordinary improper integral when the unsubtracted integral diverges.}

\glossaryentry{boundary-value-problem}{Boundary value problem}{\statusclassical}{A boundary value problem seeks a solution of a differential or functional-differential equation satisfying conditions at boundary points or on the boundary of a domain. Compact support can be viewed as an infinite collection of exterior boundary conditions, while endpoint flatness enforces smooth gluing to zero. Rvachev atomic functions were developed for nonclassical approximation methods in boundary value problems because their localized dilates and translates can encode domain geometry.}

\glossaryentry{bounded-fabius-candidate}{Bounded Fabius candidate}{\statuslean}{A bounded Fabius candidate is a function \(f:\R\to\R\) to which the
predicate \(\mathrm{IsFabius}\) may apply.  It has constant exterior tails,
reflection \(f(x)+f(1-x)=1\), regularity, and the left-half
differential-dilation law encoded by the predicate.  The word ``candidate''
keeps the variable function distinct from the canonical fixed point
\(F\).}

\glossaryentry{bounded-fabius-function}{Bounded Fabius function}{\statusmixed}{The bounded Fabius function is the normalization \(F:\R\to[0,1]\) that equals zero for \(x\le0\), one for \(x\ge1\), satisfies \(F(1-x)=1-F(x)\), and obeys \(F'(x)=2F(2x)\) on \([0,1/2]\). The adjective bounded distinguishes it from the signed global extension \(\mathcal F\), which satisfies the pure dilation equation everywhere but is not a CDF. It is the normalization used for dyadic values, inverse values, and computability.}

\glossaryentry{bounded-variation}{Bounded variation}{\statusclassical}{A real or complex function on an interval has bounded variation if the supremum of \(\sum_j|f(x_j)-f(x_{j-1})|\) over all partitions is finite. For a finite row of weights, total variation means \(\sum_j|\omega_{n,j}|\). Uniform bounded variation prevents cancellation from concealing unbounded coefficients and is the stability hypothesis in a Toeplitz convergence argument. The complex \(q\)-binomial discrete limit uses rows whose mass is one and whose total variation is uniformly bounded by an explicit finite product ratio.}

\glossaryentry{box-spline}{Box spline}{\statusfrontier}{A box spline is the density obtained by pushing forward uniform measure on a
cube by a linear map, or equivalently convolving uniform segment measures in
specified directions.  It is piecewise polynomial on a hyperplane arrangement
and supported on the associated zonotope.  The univariate finite sinc-prefix
splines are rank-one instances.}

\glossaryentry{bromwich-inversion-integral}{Bromwich inversion integral}{\statusclassical}{The Bromwich inversion integral recovers a function from its Laplace transform along a vertical line in the complex plane. In a common form, \(f(x)=(2\pi\ii)^{-1}\int_{c-\ii\infty}^{c+\ii\infty}\e^{sx}\mathcal Lf(s)\,\dd s\), where \(c\) lies to the right of all singularities. Analytic decay and contour justification are essential. For the Fabius endpoint, the contour is moved through a real negative saddle after rewriting the transform by an exact negative-Laplace product; the resulting integral supports quantitative saddle analysis.}

\glossaryentry{bromwich-saddle-method}{Bromwich saddle method}{\statusmixed}{The Bromwich saddle method applies steepest-descent ideas directly to a Laplace inversion contour. One first identifies a large real parameter and a saddle solving the derivative equation for the logarithmic phase, then deforms or localizes the vertical contour, expands the phase near the saddle, and bounds the tails and minor arcs. In the Fabius analysis the saddle is naturally expressed through the lower branch of Lambert \(W\). The method yields the main exponential scale, the log-periodic correction, the Gaussian mass, and systematic inverse-saddle corrections.}

\glossaryentry{bump-function}{Bump function}{\statusclassical}{A bump function is a \(C^\infty\) function with compact support. Nonzero bump functions cannot be real analytic on all of \(\R\), because an analytic function vanishing on an open set vanishes identically. Rvachev's up-function is a distinguished bump: it is even, nonnegative, supported on \([-1,1]\), strictly unimodal, infinitely flat at \(\pm1\), and satisfies an exact refinement-type differential equation. Its special algebraic structure makes it useful far beyond the generic existence of smooth cutoffs.}

\section{C}\label{sec:letter-c}

\glossaryentry{c-infinity-smoothness}{C-infinity smoothness}{\statusclassical}{A function is \(C^\infty\) if derivatives of every finite order exist and are continuous. Smoothness is weaker than real analyticity: the Taylor series may fail to converge to the function, or may even be identically zero at a flat point. The Fabius and up-functions are canonical examples of smooth but nowhere-analytic behavior on their nonconstant regions. Their differential scaling law gives exact formulas and sharp supremum bounds for all derivatives, while compact support forces infinite-order vanishing at the boundary.}

\glossaryentry{canonical-clamped-inverse}{Canonical clamped inverse}{\statusconvention}{A clamped inverse extends an inverse on its natural range by fixing the lower
endpoint below the range and the upper endpoint above it.  This gives a total
monotone map on \(\R\) while preserving the ordinary inverse on \([0,1]\).
Clamping is a convention, not analytic continuation; it creates constant
tails and nondifferentiable transition behavior unrelated to the interior
inverse.}

\glossaryentry{canonical-correlation}{Canonical correlation}{\statusfrontier}{Canonical correlation maximizes correlation between linear combinations of
two families of square-integrable variables.  It is determined by covariance
operators after quotienting null directions.  Pairwise correlation formulas
do not automatically solve the multi-family canonical-correlation problem.}

\glossaryentry{canonical-fabius-fixed-point}{Canonical Fabius fixed point}{\statuslean}{The canonical bounded Fabius function is obtained as the unique fixed point
of a refinement or integration operator on an admissible complete metric
space.  Its specification theorem packages the defining Fabius predicate,
while uniqueness lets later results quantify either over arbitrary witnesses
or specialize to the named function.  The fixed-point construction supplies a
mathematical object before exact arithmetic and transform representations are
developed.}

\glossaryentry{fabius-specification}{Canonical Fabius specification}{\statuslean}{The canonical specification theorem states that the named function
\(\operatorname{fabius}\) satisfies the abstract predicate
\(\mathrm{IsFabius}\).  It is the principal interface between the existence
construction and downstream theorems.  Rather than unfolding the fixed-point
definition, later proofs use the specification's tail, symmetry, derivative,
and smoothness consequences.}

\glossaryentry{canonical-fabius-inverse}{Canonical inverse Fabius function}{\statuslean}{The canonical inverse Fabius function is the inverse of
\(F:[0,1]\to[0,1]\), extended at endpoints according to the repository's
clamped convention.  On \((0,1)\) it is strictly increasing and continuous,
but endpoint flatness of \(F\) produces singular derivative behavior for the
inverse.  It is not the inverse of the signed global extension and should not
be assigned branches outside the declared range without a separate
construction.}

\glossaryentry{canonical-solution}{Canonical solution}{\statusclassical}{A canonical solution is a solution selected uniquely by natural normalization and regularity conditions, rather than by an arbitrary parameter choice. A differential or functional equation alone may admit many local or distributional solutions; support, symmetry, endpoint values, positivity, and smoothness can isolate one. The bounded Fabius function is canonical as the unique smooth distribution function with the stated tails, symmetry, and dyadic differential equation. Rvachev's up-function is the corresponding canonical compactly supported even density.}

\glossaryentry{caputo-fractional-derivative}{Caputo fractional derivative}{\statusclassical}{For \(m-1<\alpha\le m\), the Caputo derivative is
\({}^{C}D_{a+}^{\alpha}f=I_{a+}^{m-\alpha}f^{(m)}\).
It differs from the Riemann--Liouville derivative by explicit initial-value
terms and is often better adapted to classical boundary conditions.
Neither operator may be substituted for the other in Fabius fractional
extensions without checking endpoint jets.}

\glossaryentry{cardinal-interpolation}{Cardinal interpolation}{\statusclassical}{Cardinal interpolation uses a kernel \(L\) satisfying
\(L(k)=\delta_{0k}\) on an integer or scaled lattice, so that
\(f(x)\approx\sum_k f(kh)L(x/h-k)\).  Exactness depends on the function class
and kernel.  Rvachev atoms yield finite cardinal systems after polynomial
deconvolution and a sufficient dyadic mesh.}

\glossaryentry{carleman-condition}{Carleman condition}{\statusfrontier}{For a Hamburger moment sequence, Carleman's condition
\[
 \sum_{n\ge1}m_{2n}^{-1/(2n)}=\infty
\]
is sufficient for determinacy.  It is automatic for compactly supported
probability measures because \(m_{2n}^{1/(2n)}\) stays bounded.  The condition
also appears in ultradifferentiable and quasi-analytic discussions, where a
different derivative-weight sequence must be distinguished from moments.}

\glossaryentry{carleman-linearization}{Carleman linearization}{\statusfrontier}{Carleman linearization embeds nonlinear iteration into the linear action of
the Carleman matrix on monomials.  It can encode integer iterates and suggest
fractional iteration through matrix powers.  Infinite-dimensional domains,
spectral branches, and convergence of truncations are substantial additional
requirements.}

\glossaryentry{carleman-matrix}{Carleman matrix}{\statusfrontier}{For an analytic map \(f(z)=\sum_{n\ge0}a_nz^n\), its Carleman matrix is
\[
 C(f)_{mn}=[z^n]f(z)^m,\qquad m,n\ge0.
\]
Formally \(C(f\circ g)=C(f)C(g)\) when the coefficient sums are finite or
convergent.  A finite principal truncation is an approximation to the infinite
operator, not an exact linearization theorem.}

\glossaryentry{carry-count}{Carry count}{\statusclassical}{A carry count records how many digit positions produce carries when adding
integers in a fixed base.  Kummer's theorem identifies the number of base-\(p\)
carries in \(k+(n-k)\) with \(v_p\binom nk\).  Carry statistics also control
cyclotomic exponents in Gaussian coefficients and prime-power Pascal
patterns.}

\glossaryentry{catalan-number}{Catalan number}{\statusclassical}{The Catalan number
\(C_n=\frac1{n+1}\binom{2n}{n}\) counts binary trees, polygon triangulations,
Dyck paths, and many equivalent objects.  Its generating function is the
distinguished zero-constant solution of \(C(z)=1+zC(z)^2\), or after a shift a
quadratic inverse germ.  The repository also studies \(q\)-Catalan
polynomials and quarter-scaled Catalan germs.}

\glossaryentry{cauchy-repeated-integration-formula}{Cauchy formula for repeated integration}{\statusclassical}{For an integer \(n\ge1\),
\[
 I^nf(x)=\frac1{(n-1)!}\int_a^x(x-t)^{n-1}f(t)\,\dd t .
\]
The identity collapses an \(n\)-fold ordered integral to one integral over a
simplex volume.  It is the integer prototype of the Riemann--Liouville
fractional integral and explains the polynomial kernels in iterated Fabius
primitives.}

\glossaryentry{cauchy-integral-formula}{Cauchy integral formula}{\statusclassical}{If \(f\) is holomorphic inside and on a positively oriented simple closed contour \(\gamma\), then \(f^{(n)}(z_0)=\frac{n!}{2\pi\ii}\int_\gamma f(z)/(z-z_0)^{n+1}\,\dd z\). It implies analyticity, derivative bounds, and coefficient extraction. Transform functions in the Fabius development are entire, so Cauchy estimates can control Taylor coefficients and justify analytic scalar manipulations.}

\glossaryentry{cauchy-product}{Cauchy product}{\statusclassical}{The Cauchy product of \(\sum a_nz^n\) and \(\sum b_nz^n\) is \(\sum_{n\ge0}(\sum_{k=0}^{n}a_kb_{n-k})z^n\). Under absolute convergence it represents the product of the sums; for formal power series it is the definition of multiplication without analytic convergence. Moment and Bernoulli recurrences arise by equating Cauchy-product coefficients in generating-function identities. In formalized proofs, finite coefficient sums are preferable to informal series multiplication because every index range and normalization is explicit.}

\glossaryentry{cauchy-theorem}{Cauchy theorem}{\statusclassical}{Cauchy's theorem says that the integral of a holomorphic function around a closed contour is zero under standard domain hypotheses. Consequently, integrals along homotopic contours agree unless singularities are crossed. It underlies the deformation of Bromwich and Mellin contours used in the endpoint analysis. Quantitative tail bounds are still required because the contours are unbounded.}

\glossaryentry{cauchy-transform}{Cauchy transform of a measure}{\statuslean}{For a finite measure \(\mu\),
\[
 G_\mu(z)=\int_{\R}\frac{1}{z-x}\,\dd\mu(x),
 \qquad z\notin\supp\mu .
\]
The sign convention \(z-x\) is fixed.  Expanding at infinity gives the moment
series, while boundary imaginary parts recover the density under Stieltjes
inversion.}

\glossaryentry{cdf-layer-cake}{CDF layer-cake formula}{\statuslean}{A CDF layer-cake formula rewrites an integral involving \(F(x)\) or
\(1-F(x)\) as an expectation or a moment of the underlying random variable.
For example, integrating a survival function gives a first moment.  The
formalized versions state support and integrability hypotheses so that
boundary terms and totalized integrals are controlled.}

\glossaryentry{centered-appell-sequence}{Centered Appell sequence}{\statuslean}{A centered Appell sequence is normalized relative to a centered probability
law, usually so that applying its moment functional annihilates every positive
degree:
\(\mathbb E[A_n(x+X)]=x^n\).
Centering removes the first cumulant and imposes parity when the law is
symmetric.  The sequence acts as an exact polynomial deconvolution basis.}

\glossaryentry{centered-finite-spline}{Centered finite spline}{\statusmixed}{A centered finite spline in this development is the density or distribution obtained from a finite sum of centered, dyadically scaled uniform random variables, normalized so that its support is symmetric. It is a compactly supported piecewise polynomial computable from finitely many rational operations. As the number of summands grows, it converges uniformly to the Rvachev up-function or, after translation and integration, to the bounded Fabius function. The centering correction is essential: an uncentered approximation carries a systematic half-mesh displacement.}

\glossaryentry{centered-geometric-uniform-law}{Centered geometric uniform law}{\statuslean}{Centering subtracts the mean from a geometric uniform series, producing an
even law supported on a symmetric interval.  The characteristic function then
loses its phase and becomes a product of sinc factors.  Even cumulants remain,
odd moments vanish, and the \(q=1/2\) density is the up-function under the
project's normalization.}

\glossaryentry{centered-periodic-correction}{Centered periodic correction}{\statusmixed}{A centered periodic correction is a periodic function from which its mean value has been subtracted. If \(P(u+1)=P(u)\), then \(\Psi(u)=P(u)-\int_0^1P(v)\,\dd v\) has period one and zero average. Centering moves the constant Fourier mode into the nonoscillatory main constant and leaves only nonzero modes. In the corrected Fabius endpoint asymptotic, \(\Psi(\lambda(x))\) is a genuine nonconstant log-periodic term whose Fourier coefficients are explicit Gamma--zeta values.}

\glossaryentry{centered-sinc-partial-product}{Centered sinc partial product}{\statuslean}{A centered partial product is the phase-free Fourier product obtained after
centering each uniform summand.  Removing the phases makes it real and even on
the real axis and exposes sign changes only through sinc lobes.  It is the
transform counterpart of a centered finite spline approximant.}

\glossaryentry{centered-thue-morse-shell}{Centered Thue--Morse shell}{\statuslean}{The centered Thue--Morse shell is a standalone finite sinc-product model in
which dyadic lobe signs are encoded by centered binary data.  Its formal
theorems concern the product model itself; they do not automatically identify
every shell formula with a Fourier transform of a separately constructed
finite inverse object.}

\glossaryentry{central-binomial-valuation}{Central binomial valuation}{\statusmixed}{The valuation of \(\binom{2n}{n}\) is governed by carries in \(n+n\).
For \(p=2\),
\(v_2\binom{2n}{n}=s_2(n)\).  This exact bridge between central binomial
coefficients and binary weight appears in Legendre/Gaunt formulas,
denominator estimates, and \(q=-1\) specializations.}

\glossaryentry{central-gaussian-binomial-coefficient}{Central Gaussian binomial coefficient}{\statuslean}{A central Gaussian coefficient is \(\qbinom{2n}{n}{q}\), or more generally the
coefficient with \(k\) nearest \(n/2\).  It has degree \(n^2\), is palindromic,
and specializes at \(q=1\) to the central binomial coefficient.  The current
reduction theory factors central objects into cyclotomic and product pieces
suited to divisibility, asymptotic, and Catalan analyses.}

\glossaryentry{central-limit-theorem}{Central limit theorem}{\statusmixed}{A central limit theorem states that a centered, variance-normalized sequence
of random variables converges in distribution to a standard Gaussian.
For geometric-uniform families approaching the dense-scale boundary,
cumulants and characteristic functions yield an explicit CLT.  It is a weak
limit and does not alone imply density, entropy, or tail asymptotics.}

\glossaryentry{chaos-multi-index}{Chaos multi-index}{\statusfrontier}{A chaos multi-index \(\alpha=(\alpha_1,\alpha_2,\ldots)\) has finite support
and labels a tensor product of one-coordinate orthogonal polynomials.  Its
total degree is \(|\alpha|=\sum_j\alpha_j\).  The symbol and admissible zero
components are declared so it cannot be confused with a valuation or
asymptotic exponent.}

\glossaryentry{characteristic-function}{Characteristic function}{\statusclassical}{The characteristic function of a real random variable \(X\) is \(\varphi_X(t)=\E(\e^{\ii tX})\). It always exists, is positive definite and uniformly continuous, and determines the probability law uniquely. Independence turns sums into products: \(\varphi_{X+Y}=\varphi_X\varphi_Y\). For the centered dyadically weighted uniform series underlying the up-function, each summand contributes a sinc factor, so the characteristic function becomes an infinite sinc product. Under the repository's Fourier convention this is precisely the Fourier transform of the density.}

\glossaryentry{christoffel-darboux-kernel}{Christoffel--Darboux kernel}{\statusmixed}{For orthonormal polynomials \(p_0,\ldots,p_n\),
\[
 K_n(x,y)=\sum_{k=0}^{n}p_k(x)p_k(y).
\]
The Christoffel--Darboux identity rewrites this sum as a quotient involving
\(p_n,p_{n+1}\).  On the diagonal \(K_n(x,x)>0\); its reciprocal is the
Christoffel function.  The kernel is the reproducing kernel for polynomials of
degree at most \(n\) in the moment inner product.}

\glossaryentry{christoffel-function}{Christoffel function}{\statusmixed}{The Christoffel function is
\[
 \lambda_n(x)=K_n(x,x)^{-1}
 =\min\left\{\int |p|^2\,\mathrm d\mu:
 \deg p\le n,\ p(x)=1\right\}.
\]
It measures the cost of point evaluation and controls quadrature weights,
least-squares stability, and local polynomial concentration.  Its normalization
depends on whether the orthogonal basis is monic or orthonormal.}

\glossaryentry{closed-form}{Closed form}{\statusclassical}{A closed form is a finite expression assembled from an agreed collection of standard operations or special functions. The term has no universally fixed boundary: a finite sum involving moments and Thue--Morse signs may be considered explicit even when no elementary simplification exists. In the Fabius development, closed forms include exact dyadic finite sums, explicit recurrences, Gaussian \(q\)-binomial formulas, and Lambert-based asymptotic coordinates. The phrase should not imply that evaluation is numerically stable or computationally optimal.}

\glossaryentry{coefficient-extraction}{Coefficient extraction}{\statusclassical}{Coefficient extraction is denoted \([z^n]A(z)\) and returns the coefficient of \(z^n\) in a polynomial, formal power series, or convergent analytic series. It converts product identities into finite convolution sums and turns interpolation or generating-function statements into arithmetic formulas. In the \(q\)-binomial part of the Fabius theory, a weighted Thue--Morse block annihilates all lower-degree node polynomials, leaving a leading coefficient that equals the desired dyadic value after normalization.}

\glossaryentry{coercion}{Coercion and scalar cast}{\statusconvention}{A coercion transports an element along a canonical embedding, such as
\(\N\to\Z\to\Q\to\R\to\C\), or views a polynomial over a smaller coefficient
ring in a larger one.  A scalar-cast theorem states that an exact rational or
integer identity survives this transport.  The cast changes the ambient type,
not the mathematical normalization; explicit bridge lemmas keep arithmetic
proofs separate from analytic consequences.}

\glossaryentry{compact-support}{Compact support}{\statusclassical}{The support of a function is the closure of the set where it is nonzero; compact support means this closed set is compact. On \(\R\), such a function vanishes outside a bounded interval. The up-function has support exactly \([-1,1]\). Compact support implies that its Fourier transform extends to an entire function of exponential type, while \(C^\infty\) smoothness and endpoint flatness give rapid decay along the real axis. Support is topological: endpoints belong to it even though the function value there is zero.}

\glossaryentry{pascal-companion-row}{Companion Pascal row}{\statusfrontier}{A companion row is an auxiliary binomial or valuation row chosen so that a
recurrence, symmetry, or convolution becomes closed.  It may track neighboring
indices, reciprocals after denominator clearing, or carry corrections.
Because the term is report-local, its exact normalization is always printed
with the defining formula.}

\glossaryentry{compiled-tree}{Compiled tree}{\statusconvention}{The compiled tree is the set of Lean sources that successfully build under
the declared import graph and toolchain.  A theorem present only in a draft
file outside that graph is not part of the formal library.  Corpus census,
public-facade, and proof-backed-exposition claims refer to the compiled tree at
a recorded checkpoint.}

\glossaryentry{complementary-bell-number}{Complementary Bell number}{\statusfrontier}{A complementary Bell number is a signed or transformed Bell-family value
defined in the relevant report, often by replacing \(e^z-1\) with
\(1-e^z\) or by a binomial involution.  Because several conventions exist, the
defining generating function is authoritative.  The family appears in
derangement and inverse-coefficient asymptotics.}

\glossaryentry{complete-binary-block}{Complete binary block}{\statuslean}{A complete binary block is an interval of \(2^m\) consecutive integers
aligned at a multiple of \(2^m\).  Its lower \(m\) bits run through all binary
words exactly once.  Thue--Morse signs factor across aligned blocks, so
Prouhet cancellation and prefix recursions can be transferred from the initial
block to every aligned one.}

\glossaryentry{complete-exponential-bell-polynomial}{Complete exponential Bell polynomial}{\statusclassical}{The complete exponential Bell polynomial is
\(B_n(x_1,\ldots,x_n)=\sum_kB_{n,k}(x_1,\ldots)\).
It satisfies
\[
 \exp\!\left(\sum_{j\ge1}x_j\frac{t^j}{j!}\right)
 =\sum_{n\ge0}B_n(x_1,\ldots,x_n)\frac{t^n}{n!}.
\]
It converts cumulants to moments and reciprocal moment data to Appell
coefficients.}

\glossaryentry{complete-homogeneous-symmetric-polynomial}{Complete homogeneous symmetric polynomial}{\statuslean}{The complete homogeneous symmetric polynomial \(h_m(x_1,\ldots,x_r)\) is
the sum of all monomials of total degree \(m\).  Its generating function is
\(\prod_i(1-x_it)^{-1}\).  Evaluating it on a geometric progression produces
Gaussian-binomial expressions used in geometric Lagrange moment identities.}

\glossaryentry{complex-exponential-generating-function}{Complex exponential generating function}{\statusclassical}{An exponential generating function has the form \(A(z)=\sum_{n\ge0}a_nz^n/n!\). When \(z\in\C\) and the series converges, it is a complex analytic function; formally, it packages binomial convolutions especially efficiently. Moment generating functions are exponential generating functions of moments. For the Fabius and up distributions, complex EGFs connect the moment recurrences with products of elementary exponential factors and permit coefficient identities to be proved uniformly over \(\Q\), \(\R\), or \(\C\).}

\glossaryentry{complex-moment-generating-function}{Complex moment-generating function}{\statusclassical}{For a bounded real random variable \(X\), the complex moment-generating function is \(M_X(z)=\E(\e^{zX})\) for \(z\in\C\). Boundedness makes it entire, and \(M_X^{(n)}(0)=\E(X^n)\). Independence converts sums into products. Evaluating at imaginary arguments gives the characteristic function, while negative real arguments lead to Laplace asymptotics. The Fabius development uses this single entire object to relate moments, the sinc product, the negative-Laplace product, and Bromwich inversion.}

\glossaryentry{complex-shift}{Complex shift}{\statusclassical}{A complex shift replaces an argument or interpolation node by \(z+Q\) with a fixed \(Q\in\C\). Translation usually changes finite expressions, but a sufficiently high-order finite-difference functional annihilates all lower-degree contributions introduced by the shift. The repository proves that each normalized finite \(q\)-binomial Taylor block is independent of a common complex shift, even though individual summands depend on it. In the associated discrete-limit approximants, dependence may remain at finite level and disappears only in the limit.}

\glossaryentry{composite-mesh-sharpness}{Composite-mesh sharpness}{\statuslean}{Composite-mesh sharpness classifies all nonzero integer meshes achieving a
given universal reproduction degree.  The formal result equates exactness
through degree \(d\) with \(2^d\mid M\), hence identifies \(2^d\) as the
least mesh.  It is a whole polynomial-space theorem and does not assert
minimality for one particular target polynomial.}

\glossaryentry{computable-fabius-inverse}{Computable inverse Fabius function}{\statuslean}{Computability of the inverse means that from a name for
\(y\in[0,1]\) and a precision request one can produce a rational approximation
to \(F^{-1}(y)\) with certified error.  Monotonicity, computability of \(F\),
and an effective separation or modulus condition provide the algorithm.
This is a statement in computable analysis, not a practical complexity bound.}

\glossaryentry{computable-real}{Computable real}{\statusclassical}{A computable real is a real number supplied by an algorithm that, given a precision request \(p\), returns a rational or signed dyadic approximation within \(2^{-p}\). The algorithm must be uniform and terminating, and the error guarantee is part of the representation. A computable sequence requires one algorithm uniform in both the sequence index and the precision. The repository uses redundant pairs of natural-number numerators to keep the coding primitive recursive and to avoid relying on hidden integer arithmetic.}

\glossaryentry{computable-real-function}{Computable real function}{\statusclassical}{In the Grzegorczyk-style formulation used here, a real function is computable when it is sequentially computable and effectively uniformly continuous. Sequential computability means that it maps every computable real sequence to a computable sequence, uniformly in indices and requested precision. Effective uniform continuity supplies a recursive modulus controlling how accurate the input must be. The bounded Fabius function is proved computable by evaluating explicit finite splines and using a certified uniform approximation error together with the sharp global Lipschitz bound.}

\glossaryentry{computable-sequence}{Computable sequence}{\statusclassical}{A sequence \((x_i)_{i\ge0}\) of real or complex numbers is computable if a single algorithm, given \(i\) and a precision parameter \(p\), produces a rational approximation to \(x_i\) with a certified error such as \(2^{-p}\). Uniformity in \(i\) distinguishes a computable sequence from a list of individually computable numbers with unrelated algorithms. Moment sequences, dyadic values, and coefficients of spline approximants in the repository are represented by exact recurrences, yielding stronger information than approximation alone.}

\glossaryentry{conditional-asymptotic}{Conditional asymptotic theorem}{\statusfrontier}{A conditional asymptotic theorem derives a limiting formula under an explicit hypothesis not yet proved in the surrounding development, such as a parity assertion, a tail estimate, or a nonvanishing condition. Its logical status is weaker than an unconditional theorem even when the derivation is correct. Several repository modules retain conditional intermediate results as reusable transfer lemmas, while later modules discharge the hypotheses. The glossary distinguishes this proof-theoretic adjective from conditional convergence of series, which is a different analytic notion.}

\glossaryentry{conditional-law}{Conditional law}{\statusclassical}{A conditional law describes the distribution of a random variable after
conditioning on a sub-\(\sigma\)-algebra or event.  In geometric random series,
conditioning on the first uniform summand leaves an independent scaled tail,
which yields integral equations for the CDF and density.  Formal statements
avoid pointwise conditional densities when a regular conditional probability
has not been constructed.}

\glossaryentry{gaussian-binomial-continuity}{Continuity of Gaussian coefficients in the base}{\statuslean}{Because \(\qbinom{n}{k}{q}\) is a polynomial, it defines a continuous and entire
function of \(q\), including \(q=1\), \(q=0\), and roots of unity.  At \(q=1\)
it equals \(\binom nk\).  This polynomial continuation is the canonical
meaning at singular-looking bases; a quotient of \(q\)-factorials is only a
derived formula on its nonvanishing domain.}

\glossaryentry{continuous-gaussian-coefficient}{Continuous Gaussian coefficient}{\statusfrontier}{For \(0<q<1\) and suitable complex or positive real \(x,y\), a continuous
extension is
\[
 \qbinom{x}{y}{q}=
 \frac{\Gamma_q(x+1)}{\Gamma_q(y+1)\Gamma_q(x-y+1)}.
\]
It agrees with the Gaussian polynomial at integral \(x=n,y=k\) in range, but
away from integers it is meromorphic and branch-sensitive through \(q^z\).
It is used for inversion in the upper or lower parameter and for the
\(q\uparrow1\) bridge to the ordinary gamma quotient.}

\glossaryentry{continuous-linear-naturality}{Continuous-linear naturality}{\statuslean}{Continuous-linear naturality states that a convergent random or deterministic
series commutes with a continuous linear map:
\(T(\sum x_n)=\sum T(x_n)\).  It transports weighted-uniform constructions
between scalar and vector spaces, sends laws through affine projections, and
justifies termwise transform or coordinate calculations under the declared
summability mode.}

\glossaryentry{contour-deformation}{Contour deformation}{\statusclassical}{Contour deformation replaces an integration path in the complex plane by a homologous path on which the integral is easier to estimate, using Cauchy's theorem and accounting for any crossed singularities by residues. For an unbounded contour one must also control connecting arcs and convergence at infinity. In Bromwich saddle analysis, the original vertical line is localized or deformed toward a steepest-descent path through the saddle. The repository separates local central estimates from tail and minor-arc bounds so that the deformation is quantitatively justified.}

\glossaryentry{contour-integral}{Contour integral}{\statusclassical}{A contour integral \(\int_\gamma f(z)\,\dd z\) integrates a complex function along a parameterized curve \(\gamma\). Holomorphy allows path deformation, while poles contribute residues. Fourier inversion, Laplace inversion, Mellin inversion, and coefficient extraction all admit contour formulations. In the Fabius endpoint proof, the contour integral carries a large exponential phase determined by the negative-Laplace transform; its central portion yields the saddle expansion and its remote portions require uniform decay estimates.}

\glossaryentry{contraction-constant}{Contraction constant}{\statusclassical}{A contraction constant is a number \(q<1\) satisfying \(d(Tx,Ty)\le qd(x,y)\). It governs the geometric rate \(q^n\) of fixed-point iteration and the size of certified error bounds. In a Fabius fixed-point construction, the exact value depends on the chosen operator and function norm. Recording it turns uniqueness into an effective approximation theorem.}

\glossaryentry{contraction-mapping}{Contraction mapping}{\statusclassical}{A contraction mapping on a metric space satisfies \(d(Tx,Ty)\le qd(x,y)\) for a constant \(q<1\). Contractions are automatically continuous and, on complete spaces, have unique fixed points. Their iterates converge geometrically, giving explicit numerical error bounds. The Fabius defining operator is arranged to be a contraction on a space of bounded continuous functions with fixed tails and symmetry. This construction avoids assuming beforehand that the desired smooth function or probability law exists.}

\glossaryentry{convexity-and-concavity}{Convexity and concavity}{\statusclassical}{A function is convex on an interval when \(f(tx+(1-t)y)\le tf(x)+(1-t)f(y)\) for \(0\le t\le1\); it is concave when the inequality is reversed. For twice differentiable functions, \(f''\ge0\) implies convexity and \(f''\le0\) implies concavity. The bounded Fabius function is strictly convex on the left half of \([0,1]\), strictly concave on the right half, and has a unique inflection point at \(1/2\). This follows from the monotonicity and symmetry of its derivative, which is an affine copy of the up-function.}

\glossaryentry{convolution}{Convolution}{\statusclassical}{For integrable functions, \((f*g)(x)=\int_{\R}f(x-y)g(y)\,\dd y\); for measures, convolution is the law of the sum of independent random variables. Convolution is associative and commutative, preserves total mass, adds supports by Minkowski sum, and turns into multiplication under the Fourier transform. The Fabius probability model is an infinite convolution of scaled uniform laws. Finite convolutions are splines, while the limiting centered density is the up-function.}

\glossaryentry{convolution-measure}{Convolution of measures}{\statusclassical}{For finite Borel measures \(\mu,\nu\) on \(\R\), their convolution is \((\mu*\nu)(A)=\int\int\mathbf1_A(x+y)\,\mu(\dd x)\nu(\dd y)\). It is the pushforward of \(\mu\otimes\nu\) under addition. Characteristic functions multiply, and supports add. The finite and infinite convolution measures in the Fabius probability construction are defined at this measure level before smooth densities are identified.}

\glossaryentry{convolution-power}{Convolution power}{\statusclassical}{The \(k\)-fold convolution power \(f^{*k}\) is defined recursively by \(f^{*1}=f\) and \(f^{*(k+1)}=f^{*k}*f\). It represents the density of a sum of \(k\) independent identically distributed variables when \(f\) is a density. Refinement identities for the up-function lead to exact self-convolution and scaled convolution formulas. Fourier transformation converts a convolution power into an ordinary power of the transform, often making normalization and support calculations immediate.}

\glossaryentry{cornish-fisher-expansion}{Cornish--Fisher expansion}{\statusfrontier}{A Cornish--Fisher expansion reverses an Edgeworth CDF expansion to approximate
quantiles:
\[
 Q(u)=\Phi^{-1}(u)+\sum_{j\ge1}\epsilon^jq_j(\Phi^{-1}(u)).
\]
The coefficient polynomials are obtained by series reversion and Hermite
identities.  Uniformity is normally restricted to \(u\) away from \(0\) and
\(1\), where the Gaussian density stays bounded away from zero.}

\glossaryentry{correction-term}{Correction term}{\statusclassical}{A correction term is a subordinate contribution added to a leading approximation to reach a sharper error scale. Corrections can be algebraic powers, logarithms, periodic functions, or exponentially small terms. Their relative size must be judged in the variable actually governing the limit. In the Fabius endpoint expansion, the centered one-periodic fluctuation is part of the main exponent rather than a dispensable numerical decoration; subsequent saddle corrections occur in inverse powers of \(\lambda(x)\).}

\glossaryentry{cosine-series}{Cosine series}{\statusclassical}{A cosine series is a Fourier series involving only \(\cos(n\pi x/L)\), appropriate for even functions on a symmetric interval. If \(f\) is even, its sine coefficients vanish. The up-function admits a cosine expansion whose coefficients are samples of its Fourier transform; the sinc product forces all nonzero even-frequency coefficients in a particular normalization to vanish, leaving an odd cosine series. Smoothness gives rapid coefficient decay and hence strong convergence.}

\glossaryentry{coupling}{Coupling}{\statusclassical}{A coupling of probability measures \(\mu\) and \(\nu\) is a joint law with
those marginals.  Constructive couplings compare random variables on one
sample space and turn pointwise differences into metric bounds.  The natural
head--tail coupling of a geometric uniform series with its truncations controls
Wasserstein distance and supports inverse-distribution error estimates.}

\glossaryentry{covariance-kernel}{Covariance kernel}{\statusmixed}{For a centered family \(X_t\), the covariance kernel is
\(K(s,t)=\mathbb E[X_sX_t]\).  It is positive definite.  Normalizing the
geometric-uniform family yields a rational kernel in the parameters, from
which Gram determinants and projection errors can be computed exactly.}

\glossaryentry{critical-value}{Critical value}{\statusclassical}{A critical value is \(f(x_0)\) where \(f'(x_0)=0\) or, in a complex
setting, where the derivative of the map vanishes.  Local inverse branches can
meet or ramify there.  Critical values, not merely zeros and poles of the
forward function, determine the branch atlas of inverse \(q\)-special
functions.}

\glossaryentry{cumulant}{Cumulant}{\statusclassical}{Cumulants \(\kappa_n\) are the coefficients of the logarithm of a moment-generating or characteristic function: \(\log\E(\e^{zX})=\sum_{n\ge1}\kappa_nz^n/n!\) near zero. The first cumulant is the mean, the second is the variance, and higher cumulants measure non-Gaussian structure. Cumulants add for independent sums, making them natural for the infinite random-series representation. Endpoint saddle modules use tilted cumulants and their recurrences to control local phase derivatives and sharp dyadic asymptotics.}

\glossaryentry{cumulant-generating-function}{Cumulant-generating function}{\statusclassical}{The cumulant-generating function is \(K_X(z)=\log M_X(z)\), where \(M_X(z)=\E(\e^{zX})\), on a domain where a consistent logarithm is defined. Its derivatives at zero are cumulants. For products of moment-generating functions arising from independent summands, taking logarithms converts the product into a sum, revealing dyadic harmonic structure. The negative-real specialization for the Fabius law is denoted by a logarithmic Laplace function whose exact decomposition contains polynomial-logarithmic terms and a one-periodic remainder.}

\glossaryentry{cumulant-scaling-law}{Cumulant scaling law}{\statuslean}{If \(Y=aX+b\), then \(\kappa_1(Y)=a\kappa_1(X)+b\) and
\(\kappa_n(Y)=a^n\kappa_n(X)\) for \(n\ge2\).  Independent sums add cumulants.
For geometric uniform series this turns each cumulant into a geometric sum,
producing rational functions of \(q\) and Bernoulli-number factors.}

\glossaryentry{cumulative-distribution-function}{Cumulative distribution function}{\statusclassical}{The cumulative distribution function of a real random variable \(X\) is \(F(x)=\Prob(X\le x)\). It is nondecreasing, right-continuous, tends to zero at \(-\infty\), and tends to one at \(+\infty\). If the law has a continuous density \(f\), then \(F'(x)=f(x)\). The bounded Fabius function is a smooth CDF supported in the probabilistic sense on \([0,1]\); its derivative is a scaled translate of Rvachev's up-function.}

\glossaryentry{cyclotomic-inverse-branch}{Cyclotomic inverse branch}{\statusfrontier}{Near a root of unity \(\zeta\), a \(q\)-polynomial or meromorphic
\(q\)-function has a local expansion
\(Y-Y_0=c(q-\zeta)^m+\cdots\).  If \(m=1\), the inverse is analytic; if
\(m>1\), it has \(m\) Puiseux branches with leading power
\((Y-Y_0)^{1/m}\).  Cyclotomic factorization determines zeros, while derivative
cancellation determines additional critical ramification.}

\glossaryentry{cyclotomic-polynomial}{Cyclotomic polynomial}{\statusclassical}{The \(m\)-th cyclotomic polynomial \(\Phi_m(q)\in\Z[q]\) is the monic
polynomial whose roots are the primitive \(m\)-th roots of unity.  It satisfies
\[
 q^n-1=\prod_{d\mid n}\Phi_d(q).
\]
Cyclotomic factorization converts products of \(1-q^j\) into exact zero-order
and divisibility data.  This is the arithmetic backbone of root-of-unity
analysis for Gaussian coefficients and \(q\)-Catalan polynomials.}

\glossaryentry{cyclotomic-valuation}{Cyclotomic valuation}{\statuslean}{For a nonzero polynomial \(P(q)\), the cyclotomic valuation
\(v_{\Phi_m}(P)\) is the exponent of \(\Phi_m\) in its factorization.  For
products of \(1-q^j\), it is obtained by counting multiples of \(m\).
Subtracting numerator and denominator counts gives the exponent in a Gaussian
coefficient and decides whether a primitive \(m\)-th root is a zero, a
removable quotient singularity, or a nonzero value.}

\section{D}\label{sec:letter-d}

\glossaryentry{d-finite-function}{\(D\)-finite function}{\statusfrontier}{A formal power series is \(D\)-finite when it satisfies a nontrivial linear
differential equation with polynomial coefficients.  Over characteristic zero,
its coefficient sequence is \(P\)-recursive.  Smooth nonanalytic compactly
supported functions are not represented by one convergent power series on the
whole support, so local or transform formulations must be distinguished.}

\glossaryentry{decay-estimate}{Decay estimate}{\statusclassical}{A decay estimate bounds a function by a quantity tending to zero in a specified limit, often uniformly in auxiliary parameters. Polynomial, exponential, superpolynomial, and stretched-exponential decay have different analytic consequences. The Fourier transform of a compactly supported \(C^\infty\) function decays faster than every power on the real axis, while the Fabius endpoint value decays on a quadratic logarithmic scale. Quantitative decay estimates justify Fourier inversion, contour truncation, termwise operations, and numerical stopping rules.}

\glossaryentry{declaration-arity}{Declaration arity}{\statusconvention}{The arity of a definition or notation command is the number and order of its
explicit arguments.  Parameters suppressed by a local section are not part of
the printed command but remain part of the mathematical context.  Arity
auditing prevents, for example, a base \(q\), order \(\nu\), or evaluation
variable from being silently exchanged when formulas are migrated between
documents.}

\glossaryentry{declaration-inventory}{Declaration inventory}{\statusconvention}{A declaration inventory is an explicit list or counted summary of the public
definitions and theorems in a module or thematic tranche.  It lets a prose
document state exactly which formal objects support a section.  The current
README frequently records inventories as ``definitions + theorems'' and adds
non-claims to prevent a finite or formal result from being read as a stronger
analytic theorem.}

\glossaryentry{delay-differential-equation}{Delay differential equation}{\statusclassical}{A delay differential equation relates a derivative at one point to function values at transformed or earlier points, rather than only to values at the same point. The Fabius relation \(F'(x)=2F(2x)\) is more precisely an advanced dilation equation on the left half when read in \(x\), but after logarithmic change of variables it becomes a delay equation because multiplication of \(x\) by two becomes translation. Such equations encode self-similarity and can generate log-periodic phenomena.}

\glossaryentry{delta-series}{Delta series}{\statusclassical}{A delta series is a formal power series \(f(t)\) with zero constant term
and nonzero linear coefficient.  It has a compositional inverse and serves as
the substitution component of a Sheffer sequence or exponential Riordan
array.  Normalized reversion coefficients are formed from its higher
coefficients relative to the linear term.}

\glossaryentry{denominator-bound}{Denominator bound}{\statusclassical}{A denominator bound gives an explicit integer divisible by the reduced denominators of a family of rational quantities. It may be universal over all dyadic numerators at a fixed level or tailored to one sequence. In the Fabius arithmetic, exact finite formulas prove rationality at dyadic points, while valuation formulas identify powers of two and least-common-multiple products control odd prime factors. A denominator bound certifies exact storage and comparison even when it is not minimal.}

\glossaryentry{denominator-clearing}{Denominator clearing}{\statusmixed}{Denominator clearing multiplies a rational identity by a common nonzero
factor to obtain an integral or polynomial identity.  It enables valuation,
divisibility, parity, and semiring arguments.  The chosen factor should be
recorded: an oversized denominator may obscure sharp arithmetic even though it
preserves correctness.}

\glossaryentry{dense-analytic-locus}{Dense analytic locus}{\statuslean}{A function has dense analytic locus when every nonempty open subinterval
contains a point at which it is real analytic.  Elementary and algebraic
branches satisfy such a property under the repository's closure theorems.
Nowhere analyticity of the Fabius function and its inverse therefore yields
strong non-elementarity and non-algebraicity conclusions.}

\glossaryentry{derivative-scaling}{Derivative scaling}{\statusclassical}{Derivative scaling is the law obtained by repeatedly differentiating a dilation equation. From \(\mathcal F'(x)=2\mathcal F(2x)\), induction and the chain rule give \(\mathcal F^{(m)}(x)=2^{\binom{m+1}{2}}\mathcal F(2^mx)\). The triangular exponent comes from multiplying \(2^{1+2+\cdots+m}\). This identity yields exact dyadic derivatives, sharp uniform derivative bounds, terminating Taylor series at dyadic points, and high-order flatness estimates.}

\glossaryentry{derivative-value-distribution}{Derivative-value distribution}{\statuslean}{The derivative-value distribution treats \(F'(X)\), \(\Up'(X)\), or a
related derivative evaluated at a random argument as a random variable.  The
current formal development derives distributional symmetries and level-set
identities from reflection and monotonicity.  This viewpoint converts shape
information into pushforward measures and moments, rather than considering
derivative values only pointwise.}

\glossaryentry{determinant-bell-identity}{Determinant--Bell identity}{\statusfrontier}{A determinant--Bell identity expresses complete Bell polynomials,
reciprocal-series coefficients, or cumulants as structured determinants.
It packages a triangular recurrence into one closed algebraic object.  Such
identities are useful for exact arithmetic and symbolic verification but do
not automatically improve numerical conditioning.}

\glossaryentry{determinate-moment-problem}{Determinate moment problem}{\statusclassical}{A moment problem is determinate when at most one positive measure has the
given moments.  Compact support implies determinacy by polynomial density.
Carleman's condition provides a broader sufficient criterion on unbounded
support.  In the Fabius/Rvachev setting, compact support means that exact
moment identities determine the probability law, not merely its formal
generating series.}

\glossaryentry{differential-entropy}{Differential entropy}{\statusfrontier}{For a density \(f\),
\[
 h(f)=-\int f(x)\log f(x)\,\mathrm dx,
\]
when the positive and negative parts are defined, with \(0\log0=0\).
Unlike discrete entropy it can be negative and changes by \(\log|a|\) under
the scaling \(aX\).}

\glossaryentry{differential-equation}{Differential equation}{\statusclassical}{A differential equation imposes a relation among a function and its derivatives. Ordinary differential equations compare quantities at the same independent variable, while functional-differential equations may include dilated or translated arguments. The Fabius and up equations are of the latter type. Their global solution theory requires normalization and support or tail conditions; without these, the equation alone need not select the intended function. Repeated differentiation reveals the exact self-similar hierarchy of derivatives.}

\glossaryentry{differential-algebraic-function}{Differentially algebraic function}{\statusfrontier}{A function is differentially algebraic when it satisfies a polynomial
differential equation involving the function and finitely many derivatives.
Every \(D\)-finite function is differentially algebraic, but not conversely.
A non-elementarity theorem does not imply differential transcendence; these
are distinct levels of expressive complexity.}

\glossaryentry{differentially-transcendental-function}{Differentially transcendental function}{\statusfrontier}{A function is differentially transcendental over a base differential field if
it satisfies no nonzero algebraic differential equation over that field.
Proving this is much stronger than showing non-elementarity or
non-\(D\)-finiteness.  Fabius-related differential-transcendence statements are
frontier questions unless backed by a precise theorem.}

\glossaryentry{differentiation-under-the-integral-sign}{Differentiation under the integral sign}{\statusclassical}{Differentiation under the integral sign exchanges \(\frac{\dd}{\dd x}\int f(x,t)\,\dd t\) with \(\int \partial_x f(x,t)\,\dd t\). It is justified by hypotheses such as continuity of the derivative and an integrable dominating function uniform in \(x\). The method derives smoothness and differential identities for convolution and transform representations. In the Fabius setting, compact support and rapid decay often provide the needed domination, but endpoint or complex-contour integrals require separate estimates.}

\glossaryentry{digit-character}{Digit character}{\statusfrontier}{A digit character assigns a root of unity to a digit statistic, for example
\(\omega^{s_q(n)}\).  The binary choice \(\omega=-1\) gives the
Thue--Morse sign.  Character orthogonality decomposes digit-restricted sums and
suggests multibase analogues of Prouhet cancellation, rarefaction, and
geometric interpolation identities.}

\glossaryentry{digit-multiplicity-profile}{Digit multiplicity profile}{\statusfrontier}{A digit multiplicity profile records how often each digit occurs in a base
expansion or across a finite block.  Complete blocks have uniform positional
profiles, while constrained or carry-conditioned sets do not.  Such profiles
feed weighted digit characters, rarefaction matrices, and valuation
statistics.}

\glossaryentry{digit-power-sum}{Digit power sum}{\statusfrontier}{A digit power sum is a finite sum of powers of digit values or digit-weighted
positions, such as \(\sum_jd_j(n)^r\) or
\(\sum_jj^r d_j(n)\).  These statistics refine the ordinary digit sum and can
parameterize weighted automatic sequences, carry distributions, and higher
derivatives of digit generating functions.}

\glossaryentry{digitwise-factorization}{Digitwise factorization}{\statusclassical}{Digitwise factorization is an identity in which a sum indexed by binary words factors into independent contributions from each bit position. The basic Thue--Morse polynomial satisfies \(\sum_{r=0}^{2^k-1}\varepsilon_r z^r=\prod_{j=0}^{k-1}(1-z^{2^j})\). Choosing a bit either contributes \(1\) or contributes \(-z^{2^j}\), so expanding the product enumerates all binary integers and records the parity of their digit sum. This factorization is the engine behind Prouhet cancellation and several dyadic formulas.}

\glossaryentry{dirac-comb}{Dirac comb}{\statusclassical}{The Dirac comb of spacing \(h>0\) is the tempered distribution
\(\Sha_h=\sum_{k\in\Z}\delta_{kh}\).  Under a fixed Fourier convention it
transforms to a reciprocal-spacing comb with a normalization factor.
Multiplying by a function samples it, while convolving periodizes it.  Poisson
summation is the equality of these two Fourier descriptions.}

\glossaryentry{dirac-measure}{Dirac measure}{\statusclassical}{The Dirac measure \(\delta_a\) places unit mass at \(a\):
\(\int f\,\dd\delta_a=f(a)\).  Convolution with it translates,
\(f*\delta_a=f(\,\cdot-a)\), and its \(2\pi\)-Fourier transform is
\(e^{-2\pi\ii a\xi}\).  Finite atomic measures encode spline difference
operators, sampling combs, and distributional derivatives of piecewise
polynomials throughout the repository.}

\glossaryentry{dirichlet-eta-function}{Dirichlet eta function}{\statusclassical}{The Dirichlet eta function is \(\eta(s)=\sum_{n\ge1}(-1)^{n-1}n^{-s}\) for \(\Re s>0\), with analytic continuation to an entire function. It satisfies \(\eta(s)=(1-2^{1-s})\zeta(s)\). Alternating sums involving powers or reciprocal powers can therefore be expressed in zeta terms with a dyadic correction factor. Eta-type combinations appear naturally when Thue--Morse or parity decompositions separate even and odd indices.}

\glossaryentry{discrete-antiderivative}{Discrete antiderivative}{\statusclassical}{For a sequence \(a_n\), a discrete antiderivative is a sequence \(A_n\)
with \(\Delta A_n=A_{n+1}-A_n=a_n\).  Choosing \(A_0=0\) gives the prefix sum.
Iterating produces binomial-kernel convolutions and raises polynomial degree,
just as continuous integration does.  Boundary constants are part of the
definition.}

\glossaryentry{discrete-b-spline}{Discrete B-spline}{\statusclassical}{A discrete B-spline is a finite sequence or grid function obtained by repeated discrete convolution of interval indicators, serving as the lattice analogue of a continuous B-spline. Iterated prefix sums are an equivalent construction. After normalization and rescaling, such sequences approximate compactly supported spline densities. The repository identifies corrected Thue--Morse grid arrays with histogram or discrete-spline data, allowing a combinatorial recurrence to be connected to the analytic up-function through a discrete limit.}

\glossaryentry{discrete-convolution}{Discrete convolution}{\statusclassical}{For sequences \(a,b\), their discrete convolution is \((a*b)_n=\sum_k a_kb_{n-k}\), with the sum restricted to indices where both terms are defined. It models the distribution of sums of integer-valued random variables and is the coefficient rule for products of generating functions. Iterated convolution of finite box sequences produces discrete B-splines. In Thue--Morse limit arguments, discrete convolution and prefix summation reveal cancellation and smooth the signed data before scaling.}

\glossaryentry{discrete-convolution-power-kernel}{Discrete convolution-power kernel}{\statusclassical}{The \(r\)-fold convolution power of the constant-one sequence has
coefficient \(\binom{n+r-1}{r-1}\).  Consequently repeated prefix summation is
convolution with a binomial kernel.  This observation turns nested sums into a
single finite sum and links discrete antiderivatives to weak compositions.}

\glossaryentry{discrete-fourier-transform}{Discrete Fourier transform}{\statusclassical}{For \(N\) samples \(a_0,\ldots,a_{N-1}\), the discrete Fourier transform is
\(\widehat a_k=\sum_{n=0}^{N-1}a_ne^{-2\pi\ii nk/N}\).
It diagonalizes cyclic convolution and converts periodic digit or valuation
data into spectral modes.  Normalization may be placed in the forward,
inverse, or both transforms and is always declared.}

\glossaryentry{discrete-limit}{Discrete limit}{\statusclassical}{A discrete limit theorem shows that normalized finite arrays, sums, or grid functions converge to a continuous object as the mesh tends to zero or the degree tends to infinity. Convergence may be pointwise, uniform, in \(L^p\), or weak. The Fabius directory contains several discrete-limit constructions: spline grids converging to the up-function, complex-shift \(q\)-binomial rows converging to the signed global extension, and Thue--Morse histograms converging after exact normalization. Finite identities and stability bounds are kept separate from the limiting argument.}

\glossaryentry{distribution-law}{Distribution law}{\statusclassical}{The distribution law of a random variable is the probability measure it induces on its value space. Two random variables can be defined on different probability spaces yet have the same law. Laws are compared through cumulative distribution functions, moments when determinate, or characteristic functions. The Fabius function is most naturally defined as the CDF of an infinite weighted sum of independent uniform variables; this law-based viewpoint makes positivity, symmetry, and convolution identities immediate.}

\glossaryentry{division-free-identity}{Division-free identity}{\statusmixed}{A division-free identity is written using multiplication and polynomial
operations after all denominators have been cleared.  It remains meaningful at
specializations where a quotient formula would display \(0/0\), and it is
well-suited to integral, modular, and semiring proofs.  Current Gaussian
adjacent-ratio, moment-numerator, and cyclotomic arguments commonly prove a
division-free statement first and derive quotient forms under nonvanishing
hypotheses.}

\glossaryentry{division-free-recurrence}{Division-free recurrence}{\statusmixed}{A division-free recurrence computes each new term using additions and multiplications of previously known integers, with any apparent rational factors already shown integral. Such a recurrence is valuable for exact arithmetic because it avoids repeated fraction reduction and exposes positivity or divisibility. The repository gives division-free recurrences for normalized moment numerators and denominator sequences by clearing triangular factors such as factorials and Mersenne products. The proof must verify that every displayed quotient is an integer before the recurrence can be called division-free.}

\glossaryentry{dominated-convergence-theorem}{Dominated convergence theorem}{\statusclassical}{The dominated convergence theorem states that if measurable functions \(f_n\to f\) almost everywhere and \(|f_n|\le g\) for an integrable \(g\), then \(\int f_n\to\int f\). It is a principal tool for interchanging limits with integrals. In the Fabius theory it supports limits of finite convolutions, differentiation or expansion of transform integrals, and passage from finite random sums to the infinite law. Uniform or local domination must be checked explicitly when a parameter varies.}

\glossaryentry{double-factorial}{Double factorial}{\statusclassical}{The double factorial is \(n!!=n(n-2)(n-4)\cdots\), terminating at \(1\) or \(2\), with \((-1)!!=0!!=1\). Thus \((2m-1)!!=(2m)!/(2^m m!)\). Odd double factorials are Gaussian moments and appear in saddle-point integrals and in the normalization of Reshetnikov's integer sequence. Their product structure isolates powers of two from odd factors, which is useful in valuation and integrality arguments.}

\glossaryentry{dyadic-argument}{Dyadic argument}{\statusmixed}{A dyadic argument is a real input of the form \(a/2^n\), usually with integers \(a,n\ge0\). Reducing \(a\) to odd numerator identifies the exact dyadic level. The Fabius differential scaling maps dyadic points to integers after finitely many doublings, causing all sufficiently high derivatives to vanish there. Consequently the formal Taylor series terminates and exact rational evaluation becomes possible by finite recurrences or Thue--Morse sums.}

\glossaryentry{dyadic-common-denominator}{Dyadic common denominator}{\statusmixed}{A dyadic common denominator is an explicit integer \(D_n\) that clears the denominators of all Fabius values on a grid of level \(n\), or of an associated family of normalized moments. Its construction combines factorials, powers of two, and least common multiples of odd factors. It is stronger than proving each value rational separately because it supports uniform exact tables and modular arguments. The repository develops both recurrence and closed-product controls for such denominators.}

\glossaryentry{dyadic-derivative-cell}{Dyadic derivative cell}{\statuslean}{A dyadic derivative cell is an open interval between adjacent dyadic grid
points on which a sufficiently high derivative of a finite spline
approximant, or the exact derivative relation inherited from the Fabius
scaling law, has a simple polynomial or signed form.  Cell language separates
interior formulas from knot behavior and records the level at which a rational
point is reduced.}

\glossaryentry{dyadic-derivative-sign-law}{Dyadic derivative sign law}{\statuslean}{At a reduced dyadic point, a distinguished highest nonzero derivative of
the Fabius function has sign determined by a Thue--Morse or binary-weight
parity.  Lower and higher derivative behavior follows from the scaling law and
cell structure.  This exact sign law links analytic jets to automatic
sequences and supplies noncancellation data in nowhere-analyticity arguments.}

\glossaryentry{direction-profile}{Dyadic direction profile}{\statusfrontier}{A dyadic direction profile records how many projected segment weights occur at
each dyadic scale.  It determines the integer zero multiplicities, cumulants,
and comb precision of the corresponding one-dimensional product.  Recovering
the profile from spectral data is an inverse problem in discrete convexity.}

\glossaryentry{dyadic-factorized-product-extrapolant}{Dyadic factorized-product extrapolant}{\statusfrontier}{A dyadic factorized-product extrapolant has the form
\(\widehat\Phi_{N,r}=\Phi_N\widehat T_{N,r}\), where
\(\widehat T_{N,r}\) approximates the omitted tail.  Because the finite factor
is retained exactly, included integer zeros, their multiplicities, and the
real lobe sign are preserved.  The hat denotes extrapolation of a transform,
not a second Fourier transform.}

\glossaryentry{dyadic-grid}{Dyadic grid}{\statusmixed}{The dyadic grid of level \(n\) is \(2^{-n}\Z\), or its intersection with a bounded interval. It is nested under refinement: every level-\(n\) point also belongs to level \(n+1\). Functional equations with dilation by two interact especially well with this grid. Fabius values and derivatives on it admit finite exact formulas, while spline approximants and numerical verification suites use the same grid for certified comparison.}

\glossaryentry{dyadic-harmonic-sum}{Dyadic harmonic sum}{\statusmixed}{A dyadic harmonic sum has the form \(H(s)=\sum_{j\in\Z\text{ or }j\ge0}h(s/2^j)\). Mellin transformation turns scaling by \(2^{-j}\) into a geometric series factor \((1-2^{-z})^{-1}\). Its imaginary lattice of poles produces one-periodic functions of \(\log_2s\) or a related phase. The logarithm of the negative-Laplace Fabius product is such a sum.}

\glossaryentry{dyadic-name}{Dyadic name}{\statusmixed}{A dyadic name for a real number is a sequence of dyadic rationals \(a_p/2^p\) converging with a prescribed fast error, commonly \(|x-a_p/2^p|\le2^{-p}\). A signed-dyadic name may encode the numerator as a pair of naturals \((a_p^+,a_p^-)\). Names are representations used in computable analysis: an algorithm operates on approximations with guarantees rather than on an inaccessible exact real. The Fabius computability proof constructs such names uniformly from finite spline values.}

\glossaryentry{dyadic-rational}{Dyadic rational}{\statusmixed}{A dyadic rational is a rational number whose reduced denominator is a power of two, equivalently an integer multiple of \(2^{-n}\). Dyadic rationals form a dense subring of \(\R\) and are exactly the finite binary fractions. They are privileged for the Fabius function because the dilation equation repeatedly doubles the argument; after finitely many steps a dyadic point reaches an integer, producing terminating derivative and evaluation formulas. Every bounded Fabius value at a dyadic argument is rational.}

\glossaryentry{dyadic-refinement-invariance}{Dyadic refinement invariance}{\statusmixed}{Dyadic refinement invariance is the property that representing the same dyadic rational on a finer grid does not change the value produced by a formula or algorithm. For example, \(a/2^n=(2a)/2^{n+1}\) must yield identical outputs. This is nontrivial when formulas contain level-dependent moments, factorials, or sign sums. The repository proves compatibility directly, ensuring that exact evaluators depend on the rational number itself rather than on a noncanonical numerator-denominator pair.}

\glossaryentry{dyadic-spine}{Dyadic spine}{\statuslean}{The dyadic spine is the dense set of dyadic points and their iterated
preimages used to propagate exact derivative information through scaling and
composition.  At these points the Fabius derivatives have finite-order
termination or explicitly signed nonzero values.  Nowhere-analyticity proofs
use the spine to force incompatible jets arbitrarily close to any proposed
analytic neighborhood.}

\glossaryentry{dyadic-taylor-evaluator}{Dyadic Taylor evaluator}{\statusmixed}{The dyadic Taylor evaluator is a terminating exact algorithm for \(F(a/2^n)\). It uses a Taylor identity whose derivatives at a dyadic center are known through the global scaling law and vanish above a finite order. A binary recursion removes a highest set bit or folds the numerator into a smaller block, so the number of recursive calls is tied to the binary weight rather than to the denominator size. All arithmetic is rational and the termination proof is explicit.}

\glossaryentry{dyadic-valuation}{Dyadic valuation}{\statusmixed}{The dyadic, or \(2\)-adic, valuation \(\vtwo(r)\) of a nonzero rational \(r\) is the exponent of two in its prime factorization: write \(r=2^ku/v\) with odd integers \(u,v\), then \(\vtwo(r)=k\). It extends the integer valuation and satisfies \(\vtwo(rs)=\vtwo(r)+\vtwo(s)\) and \(\vtwo(r+s)\ge\min(\vtwo(r),\vtwo(s))\). The repository proves an exact formula for \(\vtwo(F(2^{-n}))\), thereby determining the power of two in the reduced denominator.}

\glossaryentry{dyadic-denominator-lcm}{Dyadic-value denominator LCM}{\statuslean}{At level \(n\), the dyadic-value denominator LCM is the least common
multiple of reduced denominators of the selected exact Fabius values on the
odd grid.  It is a positive integer with explicit boundary conventions.  The
quantity summarizes global denominator complexity at a level without claiming
that every individual denominator attains it.}

\glossaryentry{dyadic-zero-classification}{Dyadic zero classification}{\statuslean}{For the up transform, a nonzero integer frequency \(m\) is a zero in every
factor whose scale divides out the odd part of \(m\).  The number of vanishing
factors gives an explicit multiplicity controlled by \(v_2(m)\).  Formal
modules distinguish analytic zero order of an entire function from a
\(2\)-adic valuation of an integer.}

\section{E}\label{sec:letter-e}

\glossaryentry{edgeworth-expansion}{Edgeworth expansion}{\statusfrontier}{An Edgeworth expansion corrects a Gaussian density or CDF by Hermite
polynomials whose coefficients are standardized cumulants.  A rigorous theorem
states the normalization, small parameter, number of terms, uniformity region,
and remainder norm.  The formal correction polynomial by itself need not be a
nonnegative density.}

\glossaryentry{effective-flatness}{Effective flatness}{\statusclassical}{Effective flatness is a quantitative version of infinite-order vanishing at an endpoint. Instead of merely asserting \(F^{(m)}(0)=0\) for every \(m\), it gives explicit bounds such as \(F(x)\le 2^{\binom{n+1}{2}}x^n/n!\) when \(2^nx\le1\). Such inequalities provide computable rates, justify endpoint truncations, and show decay faster than every fixed power. The repository derives both mean-value and iterated-integral forms from sharp derivative scaling.}

\glossaryentry{effective-modulus}{Effective modulus}{\statusclassical}{An effective modulus converts a qualitative continuity or convergence statement into a recursive rate. For uniform continuity, a function \(d(n)\) may guarantee that \(|x-y|<1/d(n)\) implies \(|f(x)-f(y)|<1/n\). For convergence, a modulus gives an index after which an error tolerance is met. The bounded Fabius function has an explicit modulus derived from its least Lipschitz constant, while the spline approximation has a direct precision-to-level rule. These data are essential in computable analysis.}

\glossaryentry{effective-uniform-continuity}{Effective uniform continuity}{\statusclassical}{A function is effectively uniformly continuous when it has a recursively computable modulus valid over its whole domain. In the convention used by the repository, a recursive positive integer \(d(n)\) satisfies \(|x-y|<1/d(n)\Rightarrow|f(x)-f(y)|<1/n\) for every positive \(n\). This strengthens ordinary uniform continuity by exposing a usable rate. For the constant-tail bounded Fabius function, the global Lipschitz bound yields the primitive-recursive choice \(d(n)=2n\).}

\glossaryentry{elementary-function-class}{Elementary-function class}{\statuslean}{The elementary-function class is generated from rational functions,
exponentials, logarithms, powers, trigonometric functions, and permitted
compositions under a declared real or complex branch discipline.  The formal
definition used in the project is designed to prove a dense analytic-locus
theorem.  Since the Fabius function is nowhere analytic, it cannot belong to
that class on any nontrivial open interval.}

\glossaryentry{elementary-logarithmic-asymptotic}{Elementary logarithmic asymptotic}{\statusclassical}{An elementary logarithmic asymptotic rewrites a special-function coordinate, such as the lower Lambert variable, in nested elementary logarithms. For small \(x\), \(W_{-1}(-c x)\) can be expanded using \(\log x\), \(\log|\log x|\), and inverse powers of \(|\log x|\). Substitution makes the scale more familiar but also lengthens the formula and can obscure the periodic phase. The repository derives this form only after establishing the exact Lambert-coordinate expansion, so error propagation remains controlled.}

\glossaryentry{elementary-symmetric-polynomial}{Elementary symmetric polynomial}{\statusclassical}{The elementary symmetric polynomial \(e_m(x_1,\ldots,x_r)\) sums squarefree
monomials over \(m\)-element subsets.  Its generating function is
\(\prod_i(1+x_it)\).  Finite \(q\)-binomial expansions are evaluations on
geometric progressions.}

\glossaryentry{empty-sum-product-conventions}{Empty sum, empty product, and zero-power conventions}{\statusconvention}{A sum over an empty finite index set is \(0\), a product over one is \(1\),
and a zeroth convolution or operator power is the identity.  In the
natural-number polynomial identities used by the project, \(x^0=1\), including
the formal order-zero case \(0^0=1\).  These neutral-element conventions are
not cosmetic: they make boundary formulas at \(n=0\), empty compositions,
finite \(q\)-products, and triangular recurrences genuinely total.}

\glossaryentry{endpoint-asymptotic}{Endpoint asymptotic}{\statusclassical}{An endpoint asymptotic describes a function as its argument approaches the boundary of its support or domain, here \(x\downarrow0\) for the bounded Fabius function. Infinite flatness makes ordinary Taylor expansion useless: every derivative at the endpoint vanishes. The correct analysis instead uses Laplace transforms, a saddle whose location diverges, a lower-Lambert coordinate, and a log-periodic correction. Symmetry transports the result to the opposite endpoint.}

\glossaryentry{endpoint-convention}{Endpoint convention}{\statusconvention}{An endpoint convention specifies values at the boundary of an interval or
support.  It matters for pointwise refinement equations, half-weighted
indicators, totalized inverses, and clamped CDFs even when it does not change a
Lebesgue integral.  The repository states whether an interval is open or
closed, whether a density vanishes at support endpoints, and whether an inverse
uses left, right, or canonical endpoint values.}

\glossaryentry{endpoint-derivative-jet}{Endpoint derivative jet}{\statuslean}{The endpoint derivative jet is the sequence
\((f^{(n)}(a))_{n\ge0}\) at a support or interval endpoint.  For \(F\) at
zero and for \(\Up\) at \(\pm1\), the nonconstant derivatives vanish to all
orders.  Formal theorems package these values because they control extension
by zero, repeated integration, inverse asymptotics, and nonanalyticity.}

\glossaryentry{endpoint-flatness}{Endpoint flatness}{\statusclassical}{A smooth function is flat at an endpoint if all derivatives of positive order vanish there. For a compactly supported \(C^\infty\) function extended by zero, flatness at the support boundary is necessary for global smoothness. Flatness implies \(f(x)=o(x^N)\) for every fixed \(N\) under suitable one-sided conditions, but it does not imply local analyticity. The Fabius and up-functions are infinitely flat at their transition endpoints, and the repository supplies explicit bounds quantifying that fact.}

\glossaryentry{entire-function}{Entire function}{\statusclassical}{An entire function is holomorphic on all of \(\C\). Its Taylor series at any point has infinite radius of convergence, though its growth can vary greatly. The Fourier transform of the compactly supported up-function is entire and is represented by a locally uniformly convergent infinite product of sinc factors. Its zeros, growth, and real-axis decay encode partitions of unity, sampling identities, and inversion formulas. Entirety follows from the product or from compact support together with standard Fourier-Laplace theory.}

\glossaryentry{entire-function-exponential-type}{Entire function of exponential type}{\statusclassical}{An entire function \(F\) is of exponential type \(\tau\) if roughly \(|F(z)|\le C_\varepsilon\exp((\tau+\varepsilon)|z|)\) for every \(\varepsilon>0\). Fourier transforms of functions supported in a compact interval have exponential type controlled by the support radius. For the up-function supported on \([-1,1]\), this growth is compatible with rapid decay on the real axis. The distinction between complex-plane growth and real-axis decay is important: an entire function can grow exponentially off the real axis while decreasing superpolynomially along it.}

\glossaryentry{entire-product-zero-ledger}{Entire-product zero ledger}{\statusconvention}{A zero ledger lists every zero of an entire product, its multiplicity,
possible cancellations, and the factor responsible.  It prevents a
factorization valid only away from zeros from being used as a global quotient.
The current q-product and Fourier-product modules pair product identities with
separate zero-order and meromorphic-pole statements.}

\glossaryentry{entropy-power}{Entropy power}{\statusfrontier}{For an absolutely continuous real random variable with differential entropy
\(h(X)\), its entropy power is
\(N(X)=(2\pi e)^{-1}e^{2h(X)}\).
It equals the variance for a Gaussian and satisfies the entropy-power
inequality under independence.  Compact support guarantees neither finite
Fisher information nor smooth parameter dependence at every finite spline
stage.}

\glossaryentry{equality-in-distribution}{Equality in distribution}{\statusclassical}{Random variables \(X\) and \(Y\) are equal in distribution, written
\(X\stackrel d=Y\), when their pushforward probability measures coincide.
They need not be equal on a common sample space.  Independent-copy fixed-point
equations for geometric uniform laws and the Fabius distribution are naturally
distributional; after uniqueness and density theorems they induce pointwise
identities for the canonical CDF and density.}

\glossaryentry{equality-of-measures}{Equality of measures}{\statusclassical}{Measures \(\mu\) and \(\nu\) are equal when they assign the same value to
every measurable set.  It is often enough to prove equality on a generating
\(\pi\)-system or by equality of integrals against an appropriate test class.
A density identity holding almost everywhere yields equality of the associated
absolutely continuous measures, while pointwise endpoint discrepancies in the
density remain immaterial to that measure equality.}

\glossaryentry{equicontinuity}{Equicontinuity}{\statusclassical}{A family \(\mathcal F\) of functions is equicontinuous if one neighborhood size works simultaneously for every member of the family at each point; uniform equicontinuity uses one size over the whole domain. Equicontinuity, boundedness, and compactness interact through the Arzela--Ascoli theorem. In approximation arguments for finite Fabius splines or convolutions, a common Lipschitz bound gives equicontinuity and helps upgrade pointwise or weak convergence to uniform convergence on compact sets.}

\glossaryentry{euler-maclaurin-formula}{Euler-Maclaurin formula}{\statusclassical}{The Euler--Maclaurin formula connects sums and integrals by endpoint derivatives and Bernoulli numbers. In one form, \(\sum_{n=a}^{b}f(n)=\int_a^bf(t)\,\dd t+(f(a)+f(b))/2+\sum_{k\ge1}B_{2k}(f^{(2k-1)}(b)-f^{(2k-1)}(a))/(2k)!+R\). It is a principal tool for extracting constants, zeta values, and asymptotic corrections from harmonic sums. The Fabius Mellin and saddle analysis uses Euler--Maclaurin-type evaluations with explicit remainder control rather than as a purely formal expansion.}

\glossaryentry{euler-mascheroni-constant}{Euler-Mascheroni constant}{\statusclassical}{The Euler--Mascheroni constant is \(\gamma=\lim_{n\to\infty}(\sum_{k=1}^{n}1/k-\log n)\). It is the constant term in the Laurent expansion \(\zeta(s)=1/(s-1)+\gamma+\cdots\) and satisfies \(\Gamma'(1)=-\gamma\). Finite-part Mellin calculations and expansions of Gamma-zeta products naturally produce \(\gamma\). In an endpoint asymptotic, it belongs to the nonoscillatory constant after the periodic component has been centered.}

\glossaryentry{euler-pentagonal-theorem}{Euler pentagonal theorem}{\statusmixed}{Euler's pentagonal theorem is
\[
 (q;q)_\infty=\sum_{n\in\Z}(-1)^n q^{n(3n-1)/2},
 \qquad |q|<1.
\]
It is a specialization of the Jacobi triple product and gives the sparse
coefficient pattern of the Euler product.  The two generalized pentagonal
exponents arise from positive and negative \(n\).}

\glossaryentry{euler-first-q-exponential}{Euler's first \(q\)-exponential}{\statuslean}{For \(|q|<1\),
\[
 e_q(z)=\sum_{n\ge0}\frac{z^n}{(q;q)_n}
       =\frac1{(z;q)_\infty},\qquad |z|<1
\]
in the unscaled convention.  A classical-limit convention often replaces
\(z\) by \((1-q)z\).  The radius and scaling must be stated because several
functions called \(e_q\) coexist in the literature.}

\glossaryentry{euler-second-q-exponential}{Euler's second \(q\)-exponential}{\statuslean}{Euler's companion is
\[
 E_q(z)=\sum_{n\ge0}\frac{q^{n(n-1)/2}z^n}{(q;q)_n}
       =(-z;q)_\infty .
\]
It is entire in \(z\) for \(|q|<1\), and \(e_q(z)E_q(-z)=1\) on the common
domain.  The two functions correspond to the two signs in the finite
\(q\)-binomial theorem.}

\glossaryentry{eulerian-number}{Eulerian number}{\statusclassical}{The Eulerian number
\(\genfrac{\langle}{\rangle}{0pt}{}{n}{k}\) counts permutations of \(n\) with
\(k\) descents under the chosen indexing convention.  It forms the
coefficients of the Eulerian polynomial and transforms powers into rational
generating functions of the form
\(\sum_{m\ge0}m^nz^m\).}

\glossaryentry{eulerian-polynomial}{Eulerian polynomial}{\statusclassical}{The Eulerian polynomial packages Eulerian numbers, commonly as
\(A_n(t)=\sum_k\genfrac{\langle}{\rangle}{0pt}{}{n}{k} t^k\).
A shift in the descent index changes the displayed power of \(t\), so the
project prints the convention.  These polynomials occur in finite-difference,
polylogarithmic, and rational generating-function formulas.}

\glossaryentry{exact-arithmetic}{Exact arithmetic}{\statusclassical}{Exact arithmetic manipulates integers, rationals, algebraic data, or symbolic products without rounding error. It contrasts with floating-point evaluation, where cancellation and underflow can erase tiny endpoint values. The dyadic Fabius algorithms use integer signs, binomial coefficients, moment numerators, factorials, and known denominators to return exact rational results. These values serve both as mathematical theorems and as high-precision reference data for validating asymptotic or numerical procedures.}

\glossaryentry{exact-dyadic-derivative}{Exact dyadic derivative}{\statusmixed}{At a reduced interior dyadic point \(x_0=a/2^n\) with odd \(a\), the Fabius derivative of order \(n\) has the exact magnitude \(2^{\binom{n+1}{2}}\) and a sign determined by the Thue--Morse sequence; all higher derivatives vanish. This follows from global derivative scaling because \(2^nx_0=a\) is an integer where the signed extension takes a sign value. The result proves that the terminating Taylor polynomial has degree exactly \(n\).}

\glossaryentry{exact-rational-moment-table}{Exact rational moment table}{\statuslean}{An exact rational moment table lists moments or energy partial sums as
reduced fractions obtained symbolically from proved recurrences.  It is not a
floating-point reconstruction.  Such tables provide regression tests,
denominator data, and certified input for polynomial, Appell, Legendre, and
asymptotic calculations.}

\glossaryentry{exact-topological-support}{Exact topological support}{\statuslean}{An exact support theorem proves equality rather than only containment.
For the up-function the topological support is precisely \([-1,1]\); for its
density law the support is the same closed interval.  Endpoint membership
comes from every neighborhood containing positive interior values even though
the function itself vanishes at the endpoints.}

\glossaryentry{expansive-matrix}{Expansive matrix}{\statusclassical}{A square matrix is expansive when every eigenvalue has modulus greater than
one.  Its inverse contracts, so powers \(A^{-n}\) generate a shrinking
multivariate scale.  Expansiveness is the matrix analogue of a scalar
refinement dilation factor exceeding one.}

\glossaryentry{expectation}{Expectation}{\statusclassical}{The expectation \(\E[X]\) of an integrable random variable is its probability-weighted mean, defined as a Lebesgue integral. More generally, \(\E[g(X)]\) integrates \(g\) against the law of \(X\). Moments, characteristic functions, and moment-generating functions are expectations of \(X^n\), \(\e^{\ii tX}\), and \(\e^{zX}\). In the Fabius probability representation, linearity of expectation and independence translate the infinite weighted sum into exact moment recurrences and transform products.}

\glossaryentry{nonclaim}{Explicit non-claim}{\statusconvention}{An explicit non-claim states a plausible strengthening that a theorem or
module does not establish.  Examples include lack of a convergence theorem for
a formal identity, lack of a target-specific minimality result, or absence of
complex branch continuation.  Non-claims are part of rigorous documentation:
they bound the semantic reach of names that might otherwise invite
overinterpretation.}

\glossaryentry{explicit-uniform-error}{Explicit uniform error}{\statusclassical}{An explicit uniform error bound has the form \(\sup_{x\in D}|f_n(x)-f(x)|\le E_n\) with a fully specified \(E_n\to0\). It is stronger than an existence statement and directly yields a stopping rule. The centered finite-spline approximation to the bounded Fabius function satisfies a dyadic error bound independent of \(x\). This estimate is used both for certified numerical evaluation and for constructing a computable real function with a recursive precision modulus.}

\glossaryentry{exponential-formula}{Exponential formula}{\statusclassical}{The combinatorial exponential formula states that labeled assemblies of
connected components have an exponential generating function equal to the
exponential of the connected-component series.  Algebraically it is the Bell
polynomial identity.  Moments versus cumulants and arbitrary partitions versus
connected blocks are its probabilistic manifestation.}

\glossaryentry{exponential-generating-function}{Exponential generating function}{\statusclassical}{For a sequence \((a_n)\), the exponential generating function is \(A(z)=\sum_{n\ge0}a_nz^n/n!\). The factorial denominator makes multiplication encode binomial convolution: the coefficient of \(z^n/n!\) in \(A(z)B(z)\) is \(\sum_k\binom nk a_kb_{n-k}\). Moment-generating functions are canonical examples. Fabius moment and half-moment recurrences become compact product or functional equations when expressed in this form.}

\glossaryentry{exponential-riordan-array}{Exponential Riordan array}{\statuslean}{For \(g(0)\ne0\) and \(f(z)=f_1z+\cdots\) with \(f_1\ne0\), the exponential
Riordan array has entries
\[
 R_{n,k}=\frac{n!}{k!}[z^n]g(z)f(z)^k.
\]
It sends an exponential generating function \(A\) to
\(g(z)A(f(z))\).  Composition and inversion of such arrays encode Sheffer
basis changes and series reversion.}

\glossaryentry{tilted-measure}{Exponentially tilted measure}{\statusmixed}{For a random variable with finite MGF,
\[
 \frac{\mathrm d\Prob_t}{\mathrm d\Prob}
 =\exp(tX-\Lambda(t))
\]
defines the exponentially tilted law.  Under it, the mean is
\(\Lambda'(t)\) and variance \(\Lambda''(t)\).  Saddlepoint and
large-deviation arguments choose the tilt that makes the rare value typical.}

\glossaryentry{extended-nonnegative-reals}{Extended nonnegative reals}{\statusclassical}{The extended nonnegative reals are \([0,\infty]=\R_{\ge0}\cup\{\infty\}\).
They support monotone limits and nonnegative integrals without first proving
finiteness.  Addition and order extend in the expected way, while subtraction
and multiplication involving \(0\cdot\infty\) require explicit conventions or
avoidance.  Lean measure-theoretic statements often live first in this type
before finite moment or norm hypotheses convert them to ordinary real
equalities.}

\glossaryentry{extension-by-zero}{Extension by zero}{\statusconvention}{Extension by zero takes a function defined on a subset and declares it to vanish outside that subset. For a function on a closed interval to extend as a \(C^m\) function, its derivatives through order \(m\) must match zero at the boundary; \(C^\infty\) extension therefore requires flatness. Rvachev's up-function is naturally a compactly supported global function, while finite spline pieces are often first defined on their support and then extended by zero. Endpoint normalization determines whether the extension has the claimed regularity.}

\glossaryentry{exterior-constant-tails}{Exterior constant tails}{\statuslean}{The bounded Fabius function is identically \(0\) on the nonpositive
half-line and identically \(1\) on \([1,\infty)\).  These constant tails make
it a total CDF-style function on \(\R\), simplify convolution and inverse
statements, and distinguish it from the signed global extension.  Derivatives
of every positive order vanish on both exterior regions.}

\section{F}\label{sec:letter-f}

\glossaryentry{faa-di-bruno-formula}{Faa di Bruno formula}{\statusclassical}{The Faa di Bruno formula is the higher-order chain rule for a composite \(f\circ g\). One form is \(\frac{\dd^n}{\dd x^n}f(g(x))=\sum_{k=1}^{n}f^{(k)}(g(x))B_{n,k}(g'(x),\ldots,g^{(n-k+1)}(x))\), where \(B_{n,k}\) are partial Bell polynomials. It organizes the derivatives of exponentials, logarithms, and inverse functions. In the all-orders saddle and inverse-Lambert algebra, the formula explains why coefficients are finite polynomials in lower-order phase jets, even when the Lean proof unfolds them through recurrences.}

\glossaryentry{faa-di-bruno-partition-propagation}{Fa\`a di Bruno partition propagation}{\statusfrontier}{Fa\`a di Bruno partition propagation tracks the first potentially nonzero
derivative of a composition by summing over set partitions or Bell-polynomial
indices.  When lower derivatives vanish, most partitions disappear and a
small surviving family determines the new jet.  This mechanism drives the
dyadic-spine proofs for compositional iterates and also appears in formal
series reversion.}

\glossaryentry{cdf-density-bridge}{Fabius CDF--density bridge}{\statuslean}{The bounded Fabius function is a cumulative distribution function, and its
derivative is an affine rescaling of the up-density:
\[
 F'(x)=2\,\Up(2x-1).
\]
The formula is pointwise for the canonical smooth representatives and both
sides vanish outside \([0,1]\).  It transports support, positivity,
smoothness, moments, and transform identities between the CDF and density
pictures.}

\glossaryentry{fabius-distribution}{Fabius distribution}{\statusmixed}{The Fabius distribution is the compactly supported probability law whose cumulative distribution function is the bounded Fabius function. It can be realized by an infinite sum of independent uniform random variables with geometrically decreasing weights, after the normalization used in the repository. Its law is symmetric about \(1/2\), has a smooth density on \((0,1)\), and is infinitely flat at the support endpoints. Centering and rescaling the density yields Rvachev's up-function.}

\glossaryentry{fabius-fractional-shift}{Fabius fractional shift}{\statusfrontier}{A Fabius fractional shift is a proposed or derived continuation of the
integer derivative/primitive ladder to a real or complex order using
fractional integration, Mellin coordinates, or operator semigroups.  It must
specify terminal, branch, normalization, and function space.  Different
choices need not agree, so the term denotes a framework rather than one
canonical globally formalized object.}

\glossaryentry{fabius-function}{Fabius function}{\statusmixed}{The bounded Fabius function is the unique \(C^\infty\) map \(F:\R\to[0,1]\) with \(F=0\) on \(( -\infty,0]\), \(F=1\) on \([1,\infty)\), reflection law \(F(1-x)=1-F(x)\), and differential equation \(F'(x)=2F(2x)\) on \([0,1/2]\). It is a strictly increasing distribution function on \([0,1]\), strictly convex then strictly concave, and nowhere analytic on its nonconstant interval. Its dyadic values are rational, and its endpoint decay contains a nonconstant log-periodic correction.}

\glossaryentry{fabius-recurrence-sequence}{Fabius recurrence sequence}{\statusmixed}{A Fabius recurrence sequence is one of the exact rational or integer sequences introduced to encode moments, inverse-power values, dyadic numerators, or normalized endpoint values. Such sequences are usually triangular: the \(n\)-th term is determined from earlier terms by binomial coefficients, factorials, and factors \(2^j-1\). The phrase is repository-specific rather than universally standardized, so the defining recurrence and normalization must accompany it. Different source papers use related symbols for moments, half moments, numerators, and Reshetnikov integers.}

\glossaryentry{fabius-fold-map}{Fabius--Rvachev fold map}{\statuslean}{The fold map sends a bounded CDF-style function to an even compactly
supported profile by
\[
 u(x)=F(1-|x|).
\]
For the canonical objects this is Rvachev's up-function.  The absolute value
folds the two sides of the unit interval onto the origin, transfers endpoint
flatness to both support endpoints, and turns complement reflection of \(F\)
into even symmetry and unimodality of \(u\).}

\glossaryentry{fabius-self-map-dynamics}{Fabius self-map dynamics}{\statusfrontier}{On \([0,1]\), the bounded Fabius function is an increasing self-map with
fixed endpoints and a symmetry-determined midpoint value.  Iteration raises
questions about fixed points, convergence rates, conjugacies, Carleman
linearization, and the propagation of nonanalyticity.  The frontier literature
treats these dynamics while keeping them distinct from pointwise powers and
from the signed global extension.}

\glossaryentry{factorial-normalization}{Factorial normalization}{\statusconvention}{Factorial normalization records whether a coefficient family is encoded by
\(\sum a_nz^n\) or \(\sum a_nz^n/n!\).  The two conventions change
convolution, Bell-polynomial, determinant, and reversion formulas.  A notation
such as \(f_n\) is not transported between ordinary and exponential series
without the corresponding \(n!\) factors.}

\glossaryentry{falling-factorial}{Falling factorial}{\statusclassical}{The falling factorial is
\(x^{\underline n}=x(x-1)\cdots(x-n+1)\), with
\(x^{\underline0}=1\).  It is the natural basis for forward differences and
counts injections when \(x\) is a nonnegative integer.  Signed Stirling
numbers of the first kind convert it to monomials.}

\glossaryentry{fast-signed-dyadic-name}{Fast signed-dyadic name}{\statusmixed}{A fast signed-dyadic name represents a real number at precision \(p\) by a pair \((c^+,c^-)\in\N^2\), interpreted as \((c^+-c^-)/2^p\), with error at most \(2^{-p}\). The adjective fast refers to the guaranteed exponential error rate in the precision index, not necessarily to computational complexity. This redundant coding avoids negative integers inside primitive-recursive constructions and is extensionally equivalent to ordinary signed dyadic approximation. It is the representation used in the repository's computable-analysis layer.}

\glossaryentry{filter-convergence}{Filter convergence}{\statusclassical}{Filter convergence expresses limits without choosing a particular sequence.
The notation \(f\to L\) along a filter includes ordinary limits, one-sided
limits, limits at infinity, and uniform convergence on sets.  In Lean this is
the common language behind \texttt{Tendsto}, \texttt{TendstoUniformly}, and
\texttt{TendstoUniformlyOn}.  The project uses it to keep endpoint, parameter,
series, and approximation limits typed and to prevent an unstated limiting
regime from entering a theorem.}

\glossaryentry{finite-difference}{Finite difference}{\statusclassical}{The forward finite-difference operator is \((\Delta_hf)(x)=f(x+h)-f(x)\); its \(n\)-fold iterate is an alternating binomial sum. It annihilates polynomials of degree below \(n\) and extracts the leading coefficient of a degree-\(n\) polynomial up to \(n!h^n\). Gaussian \(q\)-differences replace equally spaced nodes by geometric nodes and ordinary binomial weights by \(q\)-binomial weights. The Fabius dyadic formulas combine such leading-coefficient extraction with Thue--Morse cancellation.}

\glossaryentry{finite-difference-table}{Finite-difference table}{\statusclassical}{A finite-difference table arranges a sequence and its successive forward
differences in triangular form.  Constant \(d\)-th differences characterize
polynomials of degree at most \(d\).  In exact Fabius and repeated
Thue--Morse computations, tables reveal diagonal polynomiality, integrality,
and the first level where cancellation fails.}

\glossaryentry{finite-field}{Finite field}{\statusclassical}{A finite field \(\mathbb F_q\) exists exactly when \(q\) is a prime
power and is unique up to isomorphism.  Vector subspaces over \(\mathbb F_q\)
are counted by Gaussian binomial coefficients.  The base \(q\) here is a field
cardinality, not an arbitrary complex nome, although the resulting counting
polynomial can later be evaluated elsewhere.}

\glossaryentry{finite-part-integral}{Finite-part integral}{\statusclassical}{A finite-part integral assigns a canonical finite value to an integral with a controlled endpoint divergence by subtracting the singular asymptotic terms before taking the limit. For example, if \(f(t)=a_{-1}/t+O(1)\), the finite part of \(\int_0^1f(t)\,\dd t\) removes \(a_{-1}\log(1/\varepsilon)\) from the cutoff integral. It differs from a Cauchy principal value, which relies on symmetric cancellation. Mellin and Bose-kernel calculations in the Fabius asymptotic use Hadamard finite parts to isolate constants and periodic terms.}

\glossaryentry{finite-q-binomial-theorem}{Finite \(q\)-binomial theorem}{\statuslean}{The finite theorem is
\[
 (z;q)_n=\sum_{k=0}^{n}(-1)^kq^{k(k-1)/2}\qbinom{n}{k}{q}z^k.
\]
Equivalently, \(\prod_{j=0}^{n-1}(1+zq^j)\) has coefficients
\(q^{k(k-1)/2}\qbinom{n}{k}{q}\).  The sign convention depends on whether the
product uses \(1-zq^j\) or \(1+zq^j\); the glossary records both through this
single normalized identity.}

\glossaryentry{finite-q-pochhammer-product}{Finite \(q\)-Pochhammer product}{\statuslean}{For \(n\in\N\),
\[
 (a;q)_n=\prod_{j=0}^{n-1}(1-aq^j),\qquad (a;q)_0=1.
\]
The empty-product convention is part of the definition.  As a polynomial in
\(a\), it has degree \(n\); as a polynomial in \(q\), its degree depends on the
specialization of \(a\).  The formal development proves shift, splitting,
reversal, and scalar-extension laws and uses this product as the basic atom of
Gaussian coefficients and basic hypergeometric terms.}

\glossaryentry{finite-shifted-up-synthesis}{Finite shifted-up synthesis}{\statuslean}{A finite shifted-up synthesis represents a polynomial on a prescribed compact
interval by a finite linear combination of sufficiently wide translates and
dilates of the up-function.  On that interval every selected atom lies in one
polynomial cell of its global differential structure, reducing the problem to
a finite Vandermonde or finite-difference system.  The equality is local to the
declared interval, not global on \(\R\).}

\glossaryentry{finite-sinc-product}{Finite sinc product}{\statuslean}{A finite sinc product
\[
 \Phi_N(\xi)=\prod_{j=0}^{N}\operatorname{sinc}_{\pi}(\xi/2^j)
\]
is the Fourier transform of a finite convolution of centered interval
densities.  Its inverse transform is a compactly supported spline with
explicit knots and finite smoothness.  As \(N\) grows, the products converge
normally on compact sets to the entire up transform.}

\glossaryentry{finite-spline-approximation}{Finite spline approximation}{\statusclassical}{A finite spline approximation replaces the limiting Fabius law by the distribution of a finite partial random sum, whose density is a compactly supported piecewise polynomial. Every approximant can be evaluated using finite rational arithmetic on a dyadic grid. Centering the support removes a half-step bias, and coupling the omitted tail gives an explicit uniform error. These approximants are simultaneously probabilistic truncations, convolution splines, and constructive witnesses for computability.}

\glossaryentry{first-defect-monomial}{First-defect monomial}{\statuslean}{For a mesh \(M\), the first-defect monomial has degree \(v_2(M)+1\).
A formal theorem supplies a phase at which its shifted Rvachev reconstruction
fails.  This proves sharpness of the universal degree bound rather than merely
failure of a proof technique.}

\glossaryentry{first-nonzero-shell}{First nonzero shell}{\statusfrontier}{In a sum organized by weight or defect, the first nonzero shell is the
lowest layer that survives all parity, moment, or partition cancellations.
Identifying it gives the leading asymptotic or the highest nonzero derivative.
Subsequent shells supply systematic corrections.}

\glossaryentry{first-saddle-correction}{First saddle correction}{\statusmixed}{The first saddle correction is the coefficient of \(1/\lambda\) after the leading Gaussian saddle approximation. It combines the quartic phase term, the square of the cubic phase term, derivatives of the amplitude, and any first-order periodic dependence of the transform. In logarithmic form it becomes \(A_1(\lambda(x))/\lambda(x)\). The repository derives it explicitly and proves an error one order smaller.}

\glossaryentry{fisher-information}{Fisher information}{\statusfrontier}{For a positive differentiable density \(f\),
\[
 J(f)=\int \frac{(f'(x))^2}{f(x)}\,\dd x
\]
when finite.  It measures translation sensitivity and enters de Bruijn-type
entropy expansions.  Endpoint flatness can make the quotient delicate even
when both numerator and density vanish; finiteness requires proof.}

\glossaryentry{fixed-point}{Fixed point}{\statusclassical}{A fixed point of a map \(T\) is an element \(x\) satisfying \(T(x)=x\). Functional equations often become fixed-point equations on spaces of functions. A unique fixed point can be obtained from contraction, order, compactness, or monotone-iteration arguments. The Fabius distribution is constructed as the fixed point of an averaging/refinement operator; the operator identity is equivalent to the probabilistic self-similarity of the underlying random series and ultimately to the differential equation.}

\glossaryentry{fixed-scale-up-synthesis}{Fixed-scale up synthesis}{\statuslean}{In fixed-scale synthesis all atoms share one dilation and differ only by
translation.  A finite consecutive set of centers spans polynomials up to a
degree determined by the scale and active cells.  Keeping one scale simplifies
implementation and exposes Toeplitz or finite-difference structure in the
coefficient matrix.}

\glossaryentry{flat-function}{Flat function}{\statusclassical}{A smooth function is flat at a point \(x_0\) if \(f^{(n)}(x_0)=0\) for every \(n\ge0\), or, in some conventions, for every positive \(n\) after the value is separated. The model \(\e^{-1/x^2}\) extended by zero is flat but not analytic. The Fabius function is flat at its transition endpoints relative to its constant tails, and the up-function is flat at \(\pm1\). Effective flatness estimates quantify the otherwise qualitative vanishing of every derivative.}

\glossaryentry{flat-jet}{Flat jet}{\statusclassical}{A function has a flat jet at \(x_0\) when every derivative at \(x_0\)
vanishes.  Its Taylor series is then identically zero, although the function
need not vanish in any neighborhood.  The Fabius function has flat endpoint
jets, and the up-function is flat at both support endpoints; these are
canonical examples separating \(C^\infty\) smoothness from analyticity.}

\glossaryentry{focused-import}{Focused import}{\statusconvention}{A focused import exposes a small thematic slice of the library rather than
the complete root facade.  It is useful for faster builds, dependency audits,
and frontier-facing algebraic experiments.  A focused module name describes
what has been formalized there, but does not license conclusions absent from
its public declarations--for example analytic convergence of a merely formal
germ or approximation bounds for a finite identity.}

\glossaryentry{formal-asymptotic-series}{Formal asymptotic series}{\statusclassical}{A formal asymptotic series is an ordered coefficient expression manipulated
algebraically without asserting that the infinite sum converges.  Its analytic
meaning is supplied by a remainder theorem: every finite truncation approximates
the target to a declared order in a stated limit.  The repository keeps formal
coefficient identities separate from realized Fabius, Lambert, gamma, and
\(q\)-function asymptotics.}

\glossaryentry{formal-germ}{Formal germ}{\statusmixed}{A formal germ is a power series considered algebraically near a chosen base
point, without an assertion that it converges to an analytic function.  It can
satisfy composition, inversion, or algebraic equations coefficient by
coefficient.  The quarter-Catalan and dyadic-rescaling constructions in the
current tree are formal germs; their uniqueness in a power-series ring is
distinct from existence of a global analytic branch.}

\glossaryentry{formal-logarithm}{Formal logarithm}{\statusclassical}{For a formal power series \(A(z)=1+U(z)\) with zero constant term in \(U\), its formal logarithm is \(\log A=\sum_{r\ge1}(-1)^{r-1}U^r/r\). Every coefficient depends on only finitely many coefficients of \(U\), so no analytic convergence is required. In the saddle expansion, the normalized mass series is converted into logarithmic correction coefficients by this operation. Ordered compositions give an explicit coefficient formula for the formal logarithm.}

\glossaryentry{formal-normal-form}{Formal normal form}{\statusmixed}{A formal normal form is a canonical arrangement of terms in a polynomial,
power series, transseries, or recurrence after algebraic equivalence has been
resolved.  Typical choices order monomials, collect equal exponents, and
separate leading from residual blocks.  Normal forms make coefficient
extraction and series reversal deterministic; they do not assert analytic
convergence unless a separate theorem supplies it.}

\glossaryentry{formal-power-series}{Formal power series}{\statusclassical}{A formal power series \(\sum_{n\ge0}a_nz^n\) is an algebraic object whose coefficients are manipulated without assigning a numerical value to \(z\) or requiring convergence. Addition is coefficientwise and multiplication is the Cauchy product. Formal differentiation, composition under suitable constant-term conditions, exponentiation, and logarithms are defined coefficientwise. The repository uses formal series to certify moment recurrences, saddle jets, and all-orders coefficient identities before separate analytic estimates justify any asymptotic interpretation.}

\glossaryentry{formal-versus-analytic-identity}{Formal versus analytic identity}{\statusconvention}{A formal identity is proved by coefficient algebra in a series ring; an
analytic identity additionally requires convergence and equality of the
resulting functions on a domain.  The same printed series may support both
interpretations under different hypotheses.  The repository repeatedly marks
this boundary in Catalan--Gaussian filters, complex-order products, inverse
series, and transseries so that convergence is never inferred from formal
manipulation alone.}

\glossaryentry{formal-versus-realized-transseries}{Formal versus realized transseries}{\statusconvention}{A formal transseries identity is an algebraic equality of ordered supports and
coefficients.  A realized asymptotic identity additionally names a function,
domain, branch, sector, and remainder bound.  The repository marks this
boundary explicitly: formal inversion may guide a conjecture, but it is not an
analytic theorem until realization and error transfer are proved.}

\glossaryentry{formalization-crosswalk}{Formalization crosswalk}{\statusconvention}{A formalization crosswalk maps a mathematical notion or displayed formula to
its Lean name, module, hypotheses, and exact conclusion.  It prevents prose
terminology from being mistaken for a theorem of greater generality.  This
glossary supplies a selected crosswalk for major objects; the root README and
module docstrings remain authoritative when exact declaration names or source
checkpoints matter.}

\glossaryentry{fourier-coefficient}{Fourier coefficient}{\statusclassical}{For a one-periodic integrable function \(P\), its \(k\)-th complex Fourier coefficient is \(\widehat P(k)=\int_0^1P(u)\e^{-2\pi\ii ku}\,\dd u\). The coefficient \(k=0\) is the mean; nonzero coefficients measure oscillation. In the corrected Fabius endpoint theory, the centered periodic fluctuation has coefficients proportional to \(-\Gamma(-\chi_k)\zeta(1-\chi_k)/\log2\), where \(\chi_k=2\pi\ii k/\log2\). Their nonvanishing proves that the fluctuation is genuinely nonconstant.}

\glossaryentry{fourier-energy}{Fourier energy}{\statuslean}{Fourier energy is an \(L^2\) norm such as
\(\int_{\R}|\widehat f(\xi)|^2\,\dd\xi\).  Plancherel identifies it with the
spatial \(L^2\) energy of \(f\).  For \(\Up\), the positive-half-line sinc
product integral, the square integral of \(\Up\), and a rational Legendre
coefficient series are formally connected.}

\glossaryentry{fourier-inversion}{Fourier inversion}{\statusclassical}{Fourier inversion reconstructs a function from its Fourier transform under suitable integrability or regularity assumptions. With the convention \(\widehat f(\xi)=\int f(x)\e^{-2\pi\ii x\xi}\,\dd x\), one has \(f(x)=\int\widehat f(\xi)\e^{2\pi\ii x\xi}\,\dd\xi\) when both sides are appropriately integrable, with standard variants at discontinuities. The up-function's rapid transform decay makes inversion straightforward and supports cosine expansions, derivative formulas, and sampling identities.}

\glossaryentry{fourier-legendre-coefficient}{Fourier-Legendre coefficient}{\statusclassical}{The Fourier-Legendre coefficient of \(f\in L^2[-1,1]\) is the scalar projection onto a Legendre polynomial: \(a_n=(2n+1)\int_{-1}^{1}f(x)P_n(x)\,\dd x/2\) in the standard normalization. These coefficients minimize mean-square error among polynomial approximations when truncated. The repository derives exact coefficients for the up-function, exploiting moments and symmetry, and proves more than \(L^2\) convergence: the resulting series converges absolutely and uniformly.}

\glossaryentry{fourier-legendre-series}{Fourier-Legendre series}{\statusclassical}{A Fourier-Legendre series expands a function on \([-1,1]\) in the orthogonal basis \((P_n)\) of Legendre polynomials. It is analogous to a trigonometric Fourier series but adapted to polynomial approximation with unit weight. For even functions only even degrees occur. The up-function's series has exact arithmetic coefficients, least-squares optimality of partial sums, and absolute uniform convergence established through derivative and coefficient estimates. It should not be confused with the ordinary Fourier transform of the same function.}

\glossaryentry{fourier-mode}{Fourier mode}{\statusclassical}{A Fourier mode is a basic exponential \(u\mapsto\e^{2\pi\ii ku}\) or its sine-cosine components at integer frequency \(k\). Periodic functions decompose into such modes. Under a dyadic dilation equation, complex exponents with imaginary parts \(2\pi k/\log2\) recur because multiplication by two advances the logarithmic phase by one period. The nonzero modes of the Fabius periodic correction are expressed by Gamma--zeta coefficients and cannot be replaced by a constant.}

\glossaryentry{fourier-series}{Fourier series}{\statusclassical}{A Fourier series represents a periodic function as \(\sum_{k\in\Z}c_k\e^{2\pi\ii kx}\), or equivalently by sine and cosine terms. Convergence may be in \(L^2\), pointwise, uniform, or absolute, depending on regularity. Smooth periodic functions have rapidly decaying coefficients. The repository uses Fourier series both for the endpoint's one-periodic correction and for periodicizations of the up-function; the two applications have different variables and should not be conflated.}

\glossaryentry{fourier-shell}{Fourier shell}{\statusfrontier}{A Fourier shell groups frequencies or product factors with the same dyadic
valuation, radial size, or newly introduced zero set.  Shell decomposition
separates coarse zeros from fine tail corrections and is useful for
factorized approximation, spectral counting, and Thue--Morse sign analysis.}

\glossaryentry{fourier-transform}{Fourier transform}{\statusclassical}{The Fourier transform converts translation into multiplication by a phase, convolution into multiplication, and differentiation into multiplication by frequency. Under the repository's convention, \(\widehat f(z)=\int_{\R}f(x)\e^{-2\pi\ii xz}\,\dd x\). For the up-function it equals the entire product \(\prod_{n\ge0}\sinc(\pi z/2^n)\). Zeros of the first factors yield partition and sampling identities, while compact support and smoothness determine complex growth and real-axis decay.}

\glossaryentry{fourier-convention}{Fourier-transform convention}{\statusconvention}{A Fourier-transform convention specifies the signs and factors of \(2\pi\) in the transform and inverse transform. The reference article uses \(\widehat f(z)=\int f(x)\e^{-2\pi\ii xz}\,\dd x\), for which the centered interval density has a \(\sinc(\pi z)\) transform. Sampling frequencies, derivative multipliers, and Poisson-summation constants all depend on this choice. Stating the convention prevents normalization errors.}

\glossaryentry{transform-of-a-distribution}{Fourier transform of a distribution}{\statusclassical}{For a tempered distribution \(T\), the Fourier transform is defined by
duality against Schwartz test functions.  Dirac masses, combs, polynomial
factors, and derivatives then transform by exact algebraic rules even when no
ordinary integral exists.  Distributional formulas must be distinguished from
pointwise identities between functions.}

\glossaryentry{fractional-cdf-layer-cake}{Fractional CDF layer-cake formula}{\statuslean}{Fractional layer-cake identities weight level sets by powers such as
\(t^{\alpha-1}\) or \((x-a)^{\alpha-1}\).  They connect fractional moments,
Riemann--Liouville integrals, and CDF tails through Beta and Gamma constants.
In the Fabius corpus they provide a systematic extension of integer moment
relations to real orders under positivity and integrability assumptions.}

\glossaryentry{fractional-semigroup-law}{Fractional integration semigroup law}{\statusmixed}{Under appropriate integrability hypotheses,
\(I^\alpha I^\beta=I^{\alpha+\beta}\) for positive orders.  The proof is a
Beta-integral computation justified by Tonelli or Fubini.  This law makes the
fractional order a continuous analogue of an iteration count and underlies
operator interpolation between integer primitive levels.}

\glossaryentry{fractional-iterate}{Fractional iterate}{\statusfrontier}{A fractional iterate \(f^{\circ t}\) extends integer composition so that
\(f^{\circ(s+t)}=f^{\circ s}\circ f^{\circ t}\) and
\(f^{\circ1}=f\).  Constructing one requires a chosen conjugacy or semigroup
and is generally nonunique.  The notation is never interpreted as a pointwise
power \(f(x)^t\).}

\glossaryentry{fractional-order-lowering}{Fractional order lowering}{\statusmixed}{Order lowering applies a suitable derivative or resolvent identity to reduce
a generalized integral order.  It is valid only after endpoint terms and
regularity are controlled.  The repository's real Stieltjes-resolvent
hierarchies and fractional Fabius formulas state the admissible range
explicitly rather than treating order as a purely formal index.}

\glossaryentry{fractional-order-raising}{Fractional order raising}{\statusmixed}{Order raising composes a fractional integral of order \(\alpha\) with an
ordinary or fractional integral to obtain order \(\alpha+\beta\).  At the
kernel level it is the Beta convolution of positive-part powers.  Fabius
primitive ladders use this rule to interpolate between discrete integration
levels and to transfer support and scaling relations.}

\glossaryentry{fractional-part}{Fractional part}{\statusconvention}{The fractional part of a real number is \(\{x\}=x-\lfloor x\rfloor\in[0,1)\). Any one-periodic function \(P(x)\) can be viewed as a function of \(\{x\}\). In dyadic asymptotics, the relevant phase is often the fractional part of a logarithmic or Lambert coordinate; changing \(x\) by a dyadic scale shifts that coordinate by approximately an integer. The periodic correction records the residual dependence on this phase.}

\glossaryentry{frame-bound}{Frame bound}{\statusfrontier}{A family \((\phi_k)\) in a Hilbert space has frame bounds \(A,B>0\) when
\[
 A\|f\|^2\le\sum_k|\langle f,\phi_k\rangle|^2\le B\|f\|^2.
\]
Frames permit redundancy.  Shifted compactly supported kernels may reproduce
polynomials while failing a global frame inequality, so approximation order
and frame stability are distinct properties.}

\glossaryentry{free-cumulant}{Free cumulant}{\statusfrontier}{Free cumulants \(r_n\) are related to moments by summing over noncrossing
partitions:
\[
 m_n=\sum_{\pi\in NC(n)}\prod_{B\in\pi}r_{|B|}.
\]
They linearize free additive convolution.  They are not classical cumulants,
whose partition formula uses all set partitions.}

\glossaryentry{free-infinite-divisibility}{Free infinite divisibility}{\statusfrontier}{A probability law is freely infinitely divisible when, for every \(n\), it is
the free additive convolution of \(n\) identical laws.  Conditional
positivity of the free-cumulant sequence gives algebraic tests.  Failure of a
finite Hankel necessary condition proves nondivisibility; passing all tested
blocks is not a proof of divisibility.}

\glossaryentry{fubini-theorem}{Fubini theorem}{\statusclassical}{Fubini's theorem permits the order of integration in a product space to be exchanged when the integrand is absolutely integrable. Tonelli's theorem covers nonnegative measurable integrands without prior integrability. These results justify factorization of expectations for independent variables, convolution formulas, and rearrangement of transform integrals. Signed Thue--Morse expressions require absolute bounds before Fubini can be invoked.}

\glossaryentry{full-logarithmic-expansion}{Full logarithmic expansion}{\statusclassical}{The full logarithmic expansion is the all-orders expansion of \(\log F(x)\), rather than of \(F(x)\) itself. In the repository it has the form \(\log F(x)=\mathcal M(x)+\sum_{j=1}^{N-1}A_j(\lambda(x))/\lambda(x)^j+\bigO(\lambda(x)^{-N})\), where the \(A_j\) are smooth one-periodic functions. Taking the logarithm separates multiplicative saddle corrections into additive coefficients and makes the uniqueness and size of the periodic term transparent.}

\glossaryentry{full-moment-sequence-c}{Full moment sequence $c_n$}{\statusclassical}{The sequence \(c_n\) in the repository denotes a normalized family of full centered moments of the up-function, with exact normalization fixed by the displayed integral and generating function in the source. Symmetry means only even powers contribute to the principal sequence. The triangular recurrence for \(c_n\) feeds the dyadic closed formula, Fourier-Legendre coefficients, and moment numerator arithmetic. Because symbols vary between papers, \(c_n\) should always be read with the repository's normalization.}

\glossaryentry{functional-differential-equation}{Functional-differential equation}{\statusclassical}{A functional-differential equation combines derivatives with values at transformed arguments, for example \(f'(x)=2f(2x)\) or an equation involving translates \(f(ax+b)\). Such equations are common in refinement theory, delay equations, and self-similar probability laws. They generally require boundary, support, or normalization data for uniqueness. The Fabius, up, and signed global forms satisfy related but not identical global equations, so the domain and chosen normalization are part of the statement.}

\glossaryentry{functional-equation}{Functional equation}{\statusclassical}{A functional equation constrains a function through values at different arguments, such as reflection, translation, dilation, or convolution identities. The equations \(F(1-x)=1-F(x)\), \(\mathcal F'(x)=2\mathcal F(2x)\), and the up-function refinement law are central examples. Functional equations often generate recurrences on special grids and transform products after Fourier or Laplace transformation. Local validity does not automatically imply a global identity.}

\glossaryentry{fundamental-theorem-of-calculus}{Fundamental theorem of calculus}{\statusclassical}{The fundamental theorem of calculus links differentiation and integration: a continuous function \(g\) has an antiderivative \(G(x)=\int_a^xg(t)\,\dd t\), and a continuously differentiable \(F\) satisfies \(F(b)-F(a)=\int_a^bF'(t)\,\dd t\). Iterating this identity with vanishing endpoint derivatives gives factorially improved flatness bounds. It also converts the up-density into the Fabius cumulative distribution and supports moment integration by parts.}

\section{G}\label{sec:letter-g}

\glossaryentry{gamma-derivative-tower}{Gamma-derivative tower}{\statuslean}{A gamma-derivative tower is a recursively organized family of derivatives of
a product or quotient built from gamma and reciprocal-gamma factors.  Bell
polynomials in polygamma values express the \(n\)-th derivative.  The current
Thue--Morse gamma modules track zeros, poles, jets, and exact ratios rather than
differentiate a meromorphic expression informally.}

\glossaryentry{gamma-function}{Gamma function}{\statusclassical}{The Gamma function extends the factorial by \(\Gamma(s)=\int_0^\infty t^{s-1}\e^{-t}\,\dd t\) for \(\Re s>0\) and meromorphic continuation elsewhere. It satisfies \(\Gamma(s+1)=s\Gamma(s)\) and has simple poles at the nonpositive integers. Gamma factors arise in Mellin transforms and Gaussian saddle integrals. In the Fabius periodic correction, \(\Gamma(-\chi_k)\) combines with a zeta value at a purely imaginary dyadic frequency.}

\glossaryentry{gamma-ratio}{Gamma ratio}{\statusmixed}{A gamma ratio has the form
\(\Gamma(z+a)/\Gamma(z+b)\).  For large \(z\) in a sector it admits a
power expansion beginning with \(z^{a-b}\), with coefficients expressed by
Bernoulli polynomials.  Ratios often cancel the dominant Stirling exponential
and are therefore more stable than separate gamma evaluations.}

\glossaryentry{gamma-tower-ratio}{Gamma-tower ratio}{\statuslean}{A gamma-tower ratio compares consecutive levels of a recursively defined
Thue--Morse gamma product.  Binary block factorization reduces it to a smaller
set of shifted gamma factors and makes induction possible.  The ratio is
stated as a meromorphic identity together with explicit nonvanishing
conditions or a cross-multiplied entire form.}

\glossaryentry{gamma-zeta-coefficient}{Gamma-zeta coefficient}{\statusclassical}{A Gamma-zeta coefficient is a Fourier coefficient expressed as a product of values of the Gamma and Riemann zeta functions. For the centered Fabius periodic correction, the nonzero coefficient at frequency \(k\) has the form \(-\Gamma(-\chi_k)\zeta(1-\chi_k)/\log2\), with \(\chi_k=2\pi\ii k/\log2\). Mellin inversion explains this form: poles repeated along a vertical arithmetic progression become Fourier modes in a logarithmic variable. Exponential Gamma decay makes the periodic fluctuation extremely smooth and numerically small.}

\glossaryentry{gartner-ellis-theorem}{Gärtner--Ellis theorem}{\statusfrontier}{The Gärtner--Ellis theorem derives an LDP from convergence of scaled
log-moment-generating functions under regularity such as essential smoothness
or, in compact settings, appropriate exposed-point control.  Merely computing
a pointwise log-MGF limit does not automatically supply the full lower bound.}

\glossaryentry{gauss-lobatto-quadrature}{Gauss--Lobatto quadrature}{\statusfrontier}{A Gauss--Lobatto rule includes both endpoints of an interval and chooses the
remaining nodes for maximal polynomial precision.  With \(n\) total nodes its
classical precision is \(2n-3\).  For Fabius/Rvachev weighted rules, endpoint
flatness changes numerical behavior but not the algebraic moment conditions.}

\glossaryentry{gauss-radau-quadrature}{Gauss--Radau quadrature}{\statusfrontier}{A Gauss--Radau rule prescribes one endpoint and chooses the remaining nodes
to achieve degree \(2n-2\) exactness with \(n\) nodes.  It can be obtained from
a modified orthogonality measure or a rank-one Jacobi-matrix perturbation.
This differs from a dyadic-comb cubature, whose nodes are fixed arithmetically
rather than optimized.}

\glossaryentry{gaussian-binomial-bounds}{Gaussian-binomial bounds}{\statuslean}{For real \(0\le q<1\), positivity gives
\[
 1\le \qbinom{n}{k}{q}\le \frac1{(q;q)_\infty},
\]
with sharper finite bounds obtained from the product representation.  For
\(q>1\), palindromy transfers estimates to \(q^{-1}\) and contributes the
factor \(q^{k(n-k)}\).  Such bounds provide domination for fixed-column limits,
Toeplitz rows, and \(q\)-interpolation kernels.}

\glossaryentry{gaussian-binomial-coefficient}{Gaussian binomial coefficient}{\statusclassical}{The Gaussian, or \(q\)-binomial, coefficient is defined by \[ \qbinom nkq=\frac{(q;q)_n}{(q;q)_k(q;q)_{n-k}}, \qquad 0\le k\le n. \] It is a polynomial in \(q\) with nonnegative integer coefficients and specializes to \(\binom nk\) as \(q\to1\). At \(q=1/2\), normalized Gaussian rows act as finite-difference weights on geometric nodes. The Fabius formulas use them to extract leading coefficients, prove translation invariance, and construct complex discrete limits.}

\glossaryentry{gaussian-binomial-palindromy}{Gaussian binomial palindromy}{\statuslean}{The Gaussian polynomial satisfies
\[
 \qbinom{n}{k}{q}=q^{k(n-k)}\qbinom{n}{k}{q^{-1}}.
\]
Thus its coefficient sequence is palindromic.  This is polynomial reciprocity,
not merely equality of two rational expressions for \(q\ne0\); after
multiplication by the displayed power it extends through \(q=0\).  The identity
also controls large-\(q\) asymptotics and reciprocal-base inverse branches.}

\glossaryentry{gaussian-binomial-polynomial}{Gaussian binomial polynomial}{\statuslean}{For \(0\le k\le n\),
\[
 \qbinom{n}{k}{q}=\frac{(q;q)_n}{(q;q)_k(q;q)_{n-k}}
\]
is a polynomial in \(q\) with nonnegative integer coefficients, despite the
displayed quotient.  Its degree is \(k(n-k)\), its constant and leading
coefficients are one, and it is symmetric under \(k\leftrightarrow n-k\).
The formal development proves polynomial structure before using rational
quotient presentations at exceptional bases.}

\glossaryentry{gaussian-binomial-at-minus-one}{Gaussian coefficients at \(q=-1\)}{\statuslean}{Evaluation at \(q=-1\) is governed by parity:
zeros caused by even cyclotomic factors may cancel between numerator and
denominator.  The surviving value reduces to an ordinary binomial coefficient
with halved indices in the compatible parity cases and vanishes in the
remaining odd case.  Current modules also track the first derivative at
\(-1\), which requires the polynomial realization rather than substitution
into the quotient formula.}

\glossaryentry{gaussian-contraction}{Gaussian contraction}{\statusclassical}{A Gaussian contraction is the pairing operation used to integrate products
against a centered Gaussian law.  Wick's rule expresses an even Gaussian
moment as a sum over pairings and makes every odd moment vanish.  In
saddle-point expansions the same combinatorics contracts powers of the local
integration variable after the quadratic phase has been normalized; it
organizes higher correction coefficients and connects them with Bell
polynomials and cumulants.}

\glossaryentry{gaussian-finite-difference}{Gaussian finite difference}{\statusmixed}{A Gaussian finite difference is a \(q\)-analogue of an ordinary finite difference, with geometrically spaced nodes and Gaussian binomial weights. Its finite \(q\)-binomial theorem implies annihilation of node polynomials below a prescribed degree and an exact response at the leading degree. In the Fabius dyadic evaluation, the base \(q=1/2\) matches the powers of two in the dilation equation, while a second Thue--Morse cancellation acts inside binary blocks.}

\glossaryentry{gaussian-integral}{Gaussian integral}{\statusclassical}{The Gaussian integral is \(\int_{-\infty}^{\infty}\e^{-t^2/2}\,\dd t=\sqrt{2\pi}\). Its even moments are \(\int t^{2m}\e^{-t^2/2}\,\dd t=\sqrt{2\pi}(2m-1)!!\), while odd moments vanish. These identities produce the universal coefficients in saddle mass expansions. The reference Gaussian weight is obtained by scaling according to the second derivative of the Fabius phase at the saddle.}

\glossaryentry{gaussian-binomial-pascal-recurrence}{Gaussian Pascal recurrence}{\statuslean}{Two equivalent \(q\)-Pascal recurrences are
\[
 \qbinom{n+1}{k}{q}=\qbinom{n}{k}{q}+q^{n+1-k}\qbinom{n}{k-1}{q}
 =q^k\qbinom{n}{k}{q}+\qbinom{n}{k-1}{q}.
\]
Boundary values are \(1\) at \(k=0,n\) and \(0\) outside the natural range.
The two forms correspond to adjoining a terminal or initial geometric node and
are used in induction, summation, and interpolation proofs.}

\glossaryentry{gaussian-polynomial}{Gaussian polynomial}{\statusclassical}{Gaussian polynomial is another name for a Gaussian \(q\)-binomial coefficient, emphasizing that \(\qbinom nkq\) is genuinely a polynomial in \(q\). More generally, the term may refer to finite products and combinations built from these coefficients. Repository modules on Gaussian polynomial contraction establish identities and bounds for rows at \(q=1/2\), including normalization and total-variation estimates required by the Toeplitz limit.}

\glossaryentry{gaussian-quadrature}{Gaussian quadrature}{\statusmixed}{An \(n\)-node Gaussian quadrature rule for a positive measure integrates all
polynomials of degree at most \(2n-1\) exactly.  Its nodes are the zeros of the
degree-\(n\) orthogonal polynomial and its weights are positive.  Exact
Rvachev moments determine such rules algebraically; numerical node
certification is a separate analytic or interval-arithmetic task.}

\glossaryentry{gaussian-rational}{Gaussian rational}{\statusclassical}{A Gaussian rational is a complex number \(a+b\ii\) with \(a,b\in\Q\), that is, an element of \(\Q(\ii)\). It is the fraction field of the Gaussian integers after allowing rational denominators. Exact complex \(q\)-binomial approximants can be evaluated in this field when the shift parameter is Gaussian rational. Scalar-extension lemmas then interpret rational polynomial identities over \(\C\) without redoing the arithmetic proof.}

\glossaryentry{generalized-quantile-function}{Generalized quantile function}{\statusclassical}{For a CDF \(F\), the left-continuous generalized inverse is
\(Q(u)=\inf\{x:F(x)\ge u\}\).  It exists even when \(F\) has jumps or flat
parts.  For the continuous strictly increasing Fabius CDF on \([0,1]\), this
quantile agrees with the ordinary inverse.  The generalized formulation is
useful when transferring approximation bounds from CDFs to inverse
distributions.}

\glossaryentry{generating-function}{Generating function}{\statusclassical}{A generating function encodes a sequence as coefficients of a formal or analytic series. Ordinary generating functions use \(\sum a_nz^n\); exponential generating functions use \(\sum a_nz^n/n!\). Products encode convolutions, derivatives encode weighted index shifts, and functional equations encode recurrences. Fabius moment sequences, Thue--Morse signs, Bernoulli coefficients, and saddle corrections all have useful generating functions, but their convergence domains and normalizations differ.}

\glossaryentry{genocchi-number}{Genocchi number}{\statusclassical}{Genocchi numbers are defined by a generating function such as
\[
 \frac{2t}{e^t+1}=\sum_{n\ge0}G_n\frac{t^n}{n!}.
\]
They are closely related to Bernoulli numbers and vanish in one parity beyond
the initial terms.  Sign and index conventions vary and are therefore stated
explicitly in the corpus.}

\glossaryentry{genocchi-polynomial}{Genocchi polynomial}{\statusclassical}{Genocchi polynomials insert \(e^{xt}\) into the Genocchi generating
function.  They form an Appell sequence and satisfy translation and reflection
identities analogous to Bernoulli and Euler polynomials.  They occur in the
broader combinatorial-coefficient and inversion materials.}

\glossaryentry{geometric-bernoulli-polynomial}{Geometric Bernoulli polynomial}{\statusfrontier}{For a normalized geometric-uniform law with moment-generating function
\(M_q(t)\), the geometric Bernoulli polynomials are defined by
\[
 \frac{e^{xt}}{M_q(t)}
 =\sum_{n\ge0}B_{n,q}^{\mathrm{geom}}(x)\frac{t^n}{n!}.
\]
They are monic Appell polynomials that invert averaging by the law.  At limiting
parameters they connect to classical Bernoulli-type families after the stated
normalization.}

\glossaryentry{geometric-interpolation-nodes}{Geometric interpolation nodes}{\statusclassical}{Geometric interpolation nodes form a progression such as \(-1,-2,\ldots,-2^n\), with constant ratio rather than constant spacing. Lagrange denominators on these nodes factor into \(q\)-Pochhammer products, so the leading-coefficient functional is expressed by Gaussian binomial weights with \(q=1/2\). This is the interpolation mechanism behind the finite Fabius \(q\)-binomial formulas. A common translation of all nodes changes lower-degree terms but leaves the extracted leading coefficient unchanged.}

\glossaryentry{geometric-lagrange-interpolation}{Geometric Lagrange interpolation}{\statusmixed}{For distinct geometric nodes \(x_j=q^j\), the Lagrange basis is
\[
 L_j(x)=\prod_{m\ne j}\frac{x-q^m}{q^j-q^m}.
\]
Factoring the denominator yields powers of \(q\), signs, and finite
\(q\)-Pochhammer products.  At \(q=1/2\), these weights explain the Gaussian
coefficients and Mersenne factors in exact Fabius evaluation.}

\glossaryentry{geometric-newton-interpolation}{Geometric Newton interpolation}{\statuslean}{Newton interpolation on nodes \(1,q,q^2,\ldots\) uses divided differences
adapted to a geometric comb.  The node polynomial is a rescaled
\(q\)-Pochhammer product, and the interpolation coefficients are expressible
through Gaussian factors.  It is the \(q\)-analogue of equally spaced Newton
differences and supplies exact formulas for dyadic and reciprocal-power
tables.}

\glossaryentry{geometric-uniform-cdf}{Geometric uniform CDF}{\statuslean}{The geometric uniform CDF is
\(F_q(x)=\Prob(X_q\le x)\) for the normalized geometric uniform law.
The current formalization proves monotonicity, measurability, continuity,
reflection, endpoint values, integral representations, and smoothness in the
positive-contraction regime.  Exterior formulas use the support normalization
and should not be extrapolated to signed or complex \(q\).}

\glossaryentry{geometric-uniform-density}{Geometric uniform density}{\statuslean}{The geometric uniform density is the continuous compactly supported
derivative of the geometric uniform CDF when \(0<q<1\).  It is nonnegative,
reflective about the support midpoint, integrates to one, and realizes the law
as a measure with density.  At \(q=1/2\), an affine rescaling yields the
Rvachev up-density.}

\glossaryentry{geometric-uniform-law}{Geometric uniform law}{\statuslean}{The geometric uniform law is the distribution of a normalized geometric
uniform random series.  Its support, reflection symmetry, transform product,
moments, cumulants, CDF, and smooth density are developed uniformly in the
ratio \(q\) under explicit hypotheses such as \(0<q<1\).  The Fabius law is a
distinguished \(q=1/2\) member after affine normalization.}

\glossaryentry{geometric-uniform-random-series}{Geometric uniform random series}{\statuslean}{For \(|q|<1\), a geometric uniform random series is a sum
\[
 X_q=\sum_{j\ge0}a q^j U_j
\]
of independent uniform variables with a normalization \(a\) chosen to fix the
support or total weight.  Almost-sure convergence follows from bounded
geometric tails.  At the dyadic specialization the law is affinely related to
the Fabius distribution and the centered density to Rvachev's up-function.}

\glossaryentry{germ}{Germ}{\statusclassical}{A germ at a point is an equivalence class of functions that agree on some neighborhood of that point. Germs let local properties be discussed without fixing a particular neighborhood. A smooth germ records all nearby values, whereas a formal Taylor series records only derivatives at the point; the latter does not determine a smooth germ. The terminating Taylor series of the Fabius function at a dyadic point is therefore compatible with a nonpolynomial, nowhere-analytic germ.}

\glossaryentry{global-binary-series}{Global binary series}{\statusclassical}{The global binary reduction formula evaluates the signed global
Fabius extension by peeling an argument through its binary scales and adding
finite Taylor blocks determined by exact dyadic data.  For a fixed dyadic input
the process terminates; for more general encodings the accompanying theorem
specifies the limiting or summability mode.  Its principal value is structural:
it converts the global functional equation into digitwise arithmetic while
preserving signs, endpoint conventions, and the distinction between the
bounded and global functions.}

\glossaryentry{global-differential-equation}{Global differential equation}{\statusclassical}{The global differential equation is \(\mathcal F'(x)=2\mathcal F(2x)\) for every real \(x\), satisfied by the signed global extension \(\mathcal F\). The bounded CDF does not satisfy this equation on all of \(\R\) because of its constant tails, and the compactly supported up-function satisfies a folded variant. Keeping this distinction explicit prevents invalid differentiation or scaling outside the natural interval. Repeated differentiation of the global equation yields exact derivative scaling.}

\glossaryentry{global-fabius-object}{Global Fabius object}{\statuslean}{The repository's \(\operatorname{globalFabius}\) or equivalent exported
object denotes the signed extension of the bounded function to all real
arguments.  It agrees with \(F\) on the unit interval but follows the
unclamped global differential law elsewhere.  Names containing ``global'' are
therefore semantically significant: substituting the bounded function outside
\([0,1]\) can change both values and signs.}

\glossaryentry{global-first-jet-reflection-law}{Global first-jet reflection law}{\statuslean}{A first-jet reflection law couples both function values and first
derivatives at reflected points.  For the bounded Fabius normalization,
complement reflection gives \(F(x)+F(1-x)=1\), and differentiation gives the
corresponding equality of derivatives after accounting for the chain-rule
sign.  The current development packages global first-jet consequences so that
shape and inverse arguments need not rederive them locally.}

\glossaryentry{global-signed-extension}{Global signed extension}{\statusmixed}{The global signed extension, written \(\mathcal F\) here and
implemented by the repository's global Fabius object, is the real-valued
extension that agrees with the bounded Fabius function on \([0,1]\) and obeys
\(\mathcal F'(x)=2\mathcal F(2x)\) on all of \(\R\).  Outside the unit interval
it oscillates rather than remaining clamped; for example the documented global
evaluator gives \(\mathcal F(3)=-1\).  It is constructed through exact
translation and reduction identities, so it must not be confused with the
bounded CDF-style function or with the even compactly supported up-function.}

\glossaryentry{global-uniform-spline-approximation}{Global uniform spline approximation}{\statuslean}{The global uniform spline approximation extends the finite spline model to
all real arguments in the normalization appropriate to the bounded or signed
global function.  The formal error theorem gives a uniform dyadic bound and a
filter-convergence theorem to the global Fabius object.  This analytic
approximation is stronger than the coarse estimate required merely to certify
computability.}

\glossaryentry{gram-determinant}{Gram determinant}{\statusclassical}{The Gram determinant is \(\det G\).  It equals the squared volume of the
parallelotope spanned by the vectors and vanishes exactly under linear
dependence.  Ratios of consecutive Hankel--Gram determinants give squared
norms and recurrence coefficients of monic orthogonal polynomials.}

\glossaryentry{gram-matrix}{Gram matrix}{\statuslean}{For vectors or functions \(v_0,\ldots,v_{N-1}\) in an inner-product space,
the Gram matrix is \(G_{ij}=\langle v_i,v_j\rangle\).  It is Hermitian positive
semidefinite and is positive definite exactly when the family is linearly
independent.  For monomials under a moment inner product, the Gram matrix is
the Hankel moment matrix.}

\glossaryentry{gregory-coefficient}{Gregory coefficient}{\statusfrontier}{Gregory coefficients are the coefficients in
\(z/\log(1+z)\) or a sign-shifted equivalent convention.  They appear in
numerical integration, reciprocal logarithms, and series reversion.  Because
the literature alternates signs and starting indices, the defining generating
function is part of the notation.}

\glossaryentry{grid-based-transseries}{Grid-based transseries}{\statusfrontier}{A grid-based transseries has support contained in a finitely generated
multiplicative grid of monomials.  It can often be computed like a multivariate
formal power series while retaining an asymptotic order.  Not every well-based
transseries is grid based, but explicit Lambert and polynomial--logarithmic
examples usually are.}

\glossaryentry{grid-histogram}{Grid histogram}{\statusmixed}{A grid histogram is a piecewise-constant or discrete array representation of mass on a regular mesh. Its values may be cumulative counts, normalized bin heights, or samples of a spline. In the Thue--Morse part of the development, a corrected finite grid is identified exactly with a normalized histogram produced by iterated discrete smoothing. This identity bridges combinatorial sign sums and weak or uniform convergence to the continuous Fabius-related density.}

\glossaryentry{grzegorczyk-computability}{Grzegorczyk computability}{\statusclassical}{Grzegorczyk computability for real functions is an early constructive framework based on computable real sequences and effective uniform continuity. A function must transform fast rational names effectively and possess a recursive continuity modulus. This avoids treating arbitrary real inputs as finite data. The repository proves the bounded Fabius function computable in this sense by combining primitive-recursive evaluation of centered finite splines with a uniform approximation theorem and the modulus \(d(n)=2n\).}

\section{H}\label{sec:letter-h}

\glossaryentry{hadamard-finite-part}{Hadamard finite part}{\statusclassical}{The Hadamard finite part regularizes an integral or divergent expression by subtracting all explicitly known singular terms in its cutoff expansion and retaining the constant term. It handles power and logarithmic divergences and is local to the singular endpoint. Unlike analytic continuation, it is defined directly from asymptotic subtraction, although the two methods often agree when both apply. The Fabius Mellin/Bose analysis uses finite parts at the Mellin origin to isolate the constant and periodic components of a dyadic harmonic sum.}

\glossaryentry{hahn-series}{Hahn series}{\statusmixed}{For an ordered abelian exponent group \(\Gamma\), a Hahn series is
\[
 \sum_{\gamma\in S}a_\gamma t^\gamma
\]
with well-ordered support \(S\subseteq\Gamma\) in the chosen direction.
Laurent and Puiseux series are special cases.  Hahn power--log series let the
coefficients \(a_\gamma\) be logarithmic polynomials and support real exponent
scales used in generalized reversion.}

\glossaryentry{half-moment}{Half moment}{\statusmixed}{A half moment is a moment integrated over one half of a symmetric support, with normalization chosen in the source. For the up-function or bounded Fabius law, the sequence \(d_n\) packages integrals over a half interval and is linked by symmetry to the full even moments \(c_k\). A key bridge identifies normalized half moments with endpoint values \(F(2^{-n})\). This converts analytic integrals into exact rational arithmetic and makes parity information about numerator sequences yield dyadic valuation theorems.}

\glossaryentry{half-moment-numerator}{Half-moment numerator}{\statusmixed}{A half-moment numerator is an integer sequence obtained by clearing the canonical factorial and power-of-two denominators from the rational half moments. The repository uses a dedicated symbol, often \(G_n\) or a related normalization depending on the source layer. Its recurrence is triangular and division-free after the correct scaling. Parity and positivity properties of these numerators are the arithmetic input for exact endpoint valuations and Reshetnikov integer formulas.}

\glossaryentry{half-moment-numerator-g}{Half-moment numerator $G_n$}{\statusmixed}{The symbol \(G_n\) is used in the arithmetic layer for an integer normalization of the half moments, with the exact scaling fixed in the repository. Its recurrence is derived from the rational half-moment recurrence after clearing factorial, dyadic, and Mersenne denominators. Parity and divisibility results for \(G_n\) are structural inputs to denominator theorems. It should not be confused with the transform sometimes also denoted \(G(z)\) in analytic sections.}

\glossaryentry{half-moment-sequence-d}{Half-moment sequence $d_n$}{\statusmixed}{The sequence \(d_n\) is the repository's normalized half-moment sequence. It is related to the full moments by a finite even-index binomial transform and to endpoint values by an exact inverse-power bridge of the form \(F(2^{-n})\) equals a specified factorial and power-of-two normalization times \(d_n\). Arithmetic properties of \(d_n\), especially the oddness of a normalized numerator, yield the exact \(2\)-adic valuation and Reshetnikov integers.}

\glossaryentry{half-weighted-interval-indicator}{Half-weighted interval indicator}{\statusconvention}{A half-weighted interval indicator equals \(1\) in the interior of an
interval, \(1/2\) at its two endpoints, and \(0\) outside.  It is useful when
adjacent closed cells are summed, because shared endpoints then receive total
weight one.  It is a separate object from the ordinary indicator and appears
in exact periodization, spline, and sampling formulas where endpoint values
are pointwise data rather than negligible measure-zero choices.}

\glossaryentry{hamming-weight}{Hamming weight}{\statusclassical}{The Hamming weight of a binary word or nonnegative integer is the number of one bits. For integers it is the same as the binary digit sum \(\wt(n)\). Its parity determines the Thue--Morse sign and its exact value counts how many active scales occur in a binary recursion. Algorithms that remove one set bit per step therefore have complexity proportional to Hamming weight rather than to the numerical size of the numerator.}

\glossaryentry{hankel-moment-matrix}{Hankel moment matrix}{\statusmixed}{The \(N\times N\) Hankel moment matrix is
\[
 H_N=(m_{i+j})_{0\le i,j<N}.
\]
Its entries are constant on anti-diagonals.  Positivity of
\(\int p(x)^2\,\mathrm d\mu(x)\) implies positive semidefiniteness; strict
positivity through degree \(N-1\) makes \(H_N\) positive definite.  Its
determinants control existence and normalization of monic orthogonal
polynomials.}

\glossaryentry{harmonic-number}{Harmonic number}{\statusclassical}{The \(n\)-th harmonic number is \(H_n=\sum_{k=1}^{n}1/k\), with \(H_0=0\). Harmonic numbers arise in derivatives of Gamma functions, coefficients of logarithms of factorial products, and finite-part expansions. Generalized harmonic numbers use \(\sum k^{-r}\). In saddle and Mellin coefficient formulas, they may appear when expanding Gamma ratios or zeta corrections near integer arguments.}

\glossaryentry{has-sum-and-tsum}{HasSum and tsum}{\statusconvention}{In Lean, \texttt{HasSum f a} records a proof that the series with terms
\(f(n)\) converges to \(a\), whereas \texttt{tsum f} denotes the library's
totalized infinite sum.  A theorem named \texttt{*\_hasSum} retains the
summability witness and a theorem named \texttt{*\_eq\_tsum} states the
resulting equality.  The distinction is important in the repository because a
formal expression may exist as a totalized value even when the intended
analytic theorem still requires a summability hypothesis.}

\glossaryentry{hausdorff-moment-sequence}{Hausdorff moment sequence}{\statusmixed}{A sequence is a Hausdorff moment sequence when
\(m_n=\int_0^1x^n\,\mathrm d\mu(x)\) for a finite positive measure.
Hausdorff's criterion is complete monotonicity of the finite differences:
\((-1)^k\Delta^k m_n\ge0\).  The representing measure on \([0,1]\) is unique.
Bounded Fabius moments and normalized profile sequences fit this compact
moment setting.}

\glossaryentry{head-tail-decomposition}{Head--tail decomposition}{\statuslean}{A head--tail decomposition splits an infinite random series at level \(N\)
into a finite partial sum and an independent remainder.  The head has an
explicit spline density, while the tail is confined to a short interval or is
a scaled copy of the whole law.  This yields uniform CDF sandwiches,
self-similar recurrences, and computable approximation errors.}

\glossaryentry{heat-trace}{Heat trace}{\statusclassical}{For a nonnegative operator with discrete eigenvalues \(\lambda_j\), the
heat trace is
\(\Theta(t)=\sum_j e^{-t\lambda_j}\) when trace class.  Its small-\(t\)
asymptotics encode spectral growth, while Mellin transformation connects it to
the spectral zeta function.  Frontier work applies analogous constructions to
dyadically structured zero spectra.}

\glossaryentry{heine-transformation}{Heine transformation}{\statuslean}{Heine's transformation rewrites a \({}_2\phi_1\) series as another
\({}_2\phi_1\) multiplied by an infinite-product quotient.  One standard form is
\[
 {}_2\phi_1(a,b;c;q,z)
 =\frac{(b,az;q)_\infty}{(c,z;q)_\infty}
 {}_2\phi_1(c/b,z;az;q,b).
\]
Its validity requires the stated convergence and nonvanishing conditions;
other equivalent forms permute parameters and should not be conflated.}

\glossaryentry{herglotz-function}{Herglotz function}{\statuslean}{A Herglotz, or Pick--Nevanlinna, function maps the upper half-plane into
itself.  With the \(z-x\) Cauchy-transform convention, a positive measure may
instead produce a function with negative imaginary part, so a minus sign is
sometimes inserted.  The repository fixes this orientation before applying
representation and boundary-value theorems.}

\glossaryentry{highest-nonzero-derivative}{Highest nonzero dyadic derivative}{\statuslean}{At a fixed dyadic point, scaling and flatness can force all derivatives
above a finite order to vanish while leaving one distinguished top derivative
nonzero.  The order and sign are computed from the reduced dyadic level and
binary data.  This terminating jet is a central local invariant in exact
Taylor evaluation and dyadic-spine nonanalyticity proofs.}

\glossaryentry{highest-set-bit-decomposition}{Highest-set-bit decomposition}{\statuslean}{Every positive integer has a unique decomposition
\(n=2^m+r\) with \(0\le r<2^m\).  Recursions based on the highest set bit
process one nonzero binary digit at a time and expose complete aligned blocks.
Their depth is bounded by bit length and often more closely by Hamming weight.}

\glossaryentry{highest-set-bit-recursion}{Highest-set-bit recursion}{\statusclassical}{A highest-set-bit recursion decomposes a positive integer \(a\) as \(2^m+r\) with \(0\le r<2^m\), processes the leading binary block, and recurses on a smaller integer. For Fabius dyadic evaluation, reflection and Taylor reduction turn this decomposition into an exact arithmetic algorithm. Each recursive call removes at least one set bit, providing a transparent termination measure and an operation count controlled by the binary representation.}

\glossaryentry{hilbert-space-projection}{Hilbert-space projection}{\statuslean}{For a closed subspace \(V\) of a Hilbert space, every \(x\) has a unique
orthogonal projection \(P_Vx\) minimizing \(\|x-v\|\).  The residual is
orthogonal to \(V\), and Pythagoras splits the squared norm.  Finite spans of
standardized geometric-uniform random variables give exact least-squares
predictors in the current corpus.}

\glossaryentry{hilbert-transform}{Hilbert transform}{\statusclassical}{The Hilbert transform is the principal-value singular integral
\[
 Hf(x)=\frac1\pi\operatorname{p.v.}\int_{\R}\frac{f(t)}{x-t}\,\dd t .
\]
It is the real boundary part conjugate to a Cauchy transform and multiplies
Fourier modes by \(-\ii\,\operatorname{sgn}\xi\) under a standard convention.
Compact support does not remove the principal-value issue inside the support.}

\glossaryentry{historical-notation-scope}{Historical notation scope}{\statushistorical}{A historical notation scope is a clearly marked region in which a quoted or
transcribed source retains its original symbols.  It preserves bibliographic
fidelity without changing the active project's notation.  The scope ends with
the quotation or transcription; surrounding commentary returns to canonical
commands.  A historical marker is therefore provenance metadata, not a blanket
waiver from notation consistency.}

\glossaryentry{hockey-stick-identity}{Hockey-stick identity}{\statusclassical}{The hockey-stick identity is \(\sum_{j=r}^{n}\binom jr=\binom{n+1}{r+1}\). Its name comes from the shape of the selected entries in Pascal's triangle. Repeated prefix summation and discrete B-spline formulas often reduce to this identity, as do sums over ordered compositions with a final free part. The repository uses it to collapse nested finite sums arising from Thue--Morse cancellation and spline integration.}

\glossaryentry{holomorphic-function}{Holomorphic function}{\statusclassical}{A holomorphic function is complex differentiable at every point of an open domain. Complex differentiability implies analytic power-series expansion and strong rigidity properties such as the identity theorem and Cauchy's integral formula. Entire functions are holomorphic on all of \(\C\). The complex moment-generating function and Fourier transform of the up-function are holomorphic objects; contour and Mellin arguments rely on their domains, singularities, and growth estimates.}

\glossaryentry{holonomic-function}{Holonomic function}{\statusfrontier}{In one variable, ``holonomic'' is commonly synonymous with \(D\)-finite.
In several variables it refers to a system with minimal-dimensional
characteristic variety.  Closure under sums, products, derivatives, and
definite integration makes holonomic methods powerful, but composition and
functional inversion have restricted closure properties.}

\glossaryentry{homothety}{Homothety}{\statusclassical}{A homothety is a dilation about a fixed center, \(x\mapsto c+a(x-c)\). Homothetic copies of a compactly supported function are rescaled and shifted versions that preserve shape while changing support. Rvachev's atomic-function theory builds approximation systems from homothetic copies of the up-function. The Fabius dilation equation is centered at zero and uses ratio two, while folding and translating connect it to support-centered copies.}

\glossaryentry{horner-scheme}{Horner scheme}{\statusclassical}{Horner's scheme evaluates a polynomial \(a_0+a_1x+\cdots+a_nx^n\) as \(a_0+x(a_1+x(\cdots+x a_n))\). It uses \(n\) multiplications and \(n\) additions and reduces intermediate expression growth. Exact dyadic formulas involving moment polynomials can be implemented in Horner form to avoid repeated exponentiation. In rational arithmetic it also helps control numerator and denominator sizes, although bit complexity still depends on coefficient growth.}

\glossaryentry{hypergeometric-term}{Hypergeometric term}{\statusclassical}{A sequence term \(a_n\) is hypergeometric when the ratio \(a_{n+1}/a_n\) is a rational function of \(n\). Factorials, binomial coefficients, and double factorials are basic examples. Their sums can often be manipulated by telescoping or symbolic recurrence methods. The Fabius arithmetic formulas combine hypergeometric prefactors with recursively defined moment numerators; the \(q\)-analogues replace rational dependence on \(n\) by rational dependence on \(q^n\).}

\glossaryentry{hypothesis-ledger}{Hypothesis ledger}{\statusconvention}{A hypothesis ledger lists the assumptions attached to a family of formulas:
domains, positivity, nonvanishing denominators, branch choices, convergence
regions, and endpoint exclusions.  The current \(q\)-series and inverse
materials use ledgers because superficially similar identities can require
different parameter regimes.  A glossary definition summarizes the object;
the ledger determines when a particular theorem may be applied.}

\section{I}\label{sec:letter-i}

\glossaryentry{identity-theorem}{Identity theorem}{\statusclassical}{The identity theorem says that two holomorphic functions on a connected domain that agree on a set with an accumulation point agree everywhere. A real-analytic analogue holds on intervals. It is used conceptually in nowhere-analytic arguments: if a smooth Fabius germ were analytic and agreed with a terminating Taylor polynomial on a neighborhood, overlap and continuation would force polynomial behavior incompatible with the global scaling and shape.}

\glossaryentry{immutable-recovery-point}{Immutable recovery point}{\statusconvention}{An immutable recovery point is a commit or archived artifact retained so
that an earlier source state can be reconstructed exactly.  It differs from an
active document, which should receive current notation and corrections.
Recovery points preserve scholarly provenance and allow regressions to be
diagnosed without forcing obsolete material back into the live corpus.}

\glossaryentry{improper-integral}{Improper integral}{\statusclassical}{An improper integral is defined as a limit when the interval is unbounded or the integrand is singular at an endpoint. Convergence requires the cutoff integrals to approach a finite limit. Absolute convergence is stronger and permits more transformations. Laplace, Mellin, and Fourier integrals in the Fabius theory are often improper; when a singular Mellin term prevents ordinary convergence, a finite-part integral rather than an improper integral is used.}

\glossaryentry{inclusion-exclusion-principle}{Inclusion--exclusion principle}{\statusclassical}{Inclusion--exclusion computes the size or measure of a union by alternating
over intersections.  Algebraically it is a Boolean-cube M\"obius inversion.
Finite-difference and Thue--Morse subset sums share the same alternating-sign
architecture, although their summands represent function values rather than
event intersections.}

\glossaryentry{independent-copy-equation}{Independent-copy equation}{\statuslean}{An independent-copy equation is a distributional fixed point such as
\[
 X\stackrel d=aU+qX',
\]
where \(X'\) is independent of \(U\) and has the same law as \(X\).
It encodes the first term and rescaled tail of a geometric random series.
Taking transforms, moments, or CDFs produces the corresponding analytic
functional equations, while uniqueness requires an appropriate class of
probability laws.}

\glossaryentry{independent-random-variables}{Independent random variables}{\statusclassical}{Random variables are independent when probabilities of joint events factor, equivalently when expectations of suitable products factor. Independence implies that characteristic and moment-generating functions multiply under addition and that probability laws convolve. The Fabius distribution arises from an infinite sum of independent uniform variables at dyadically decreasing scales. This representation explains the infinite product transform and produces finite convolution splines by truncation.}

\glossaryentry{indicator-function}{Indicator function}{\statusconvention}{The indicator of a set \(A\) is \(\mathbf1_A(x)=1\) for \(x\in A\) and \(0\) otherwise. Normalized interval indicators are uniform probability densities. Their repeated convolutions are splines, and their Fourier transforms are phase factors times sinc functions. The finite approximants to the up-function begin with scaled interval indicators, so elementary support and mass calculations propagate through convolution.}

\glossaryentry{infinite-convolution}{Infinite convolution}{\statusclassical}{An infinite convolution is a limit of finite convolutions of probability measures or functions. For probability laws it represents the distribution of an infinite sum of independent random variables, provided the partial sums converge. Weak convergence is often the natural initial topology, while additional smoothness and tail control may yield densities and uniform convergence. The Fabius law is the infinite convolution of dyadically scaled uniform laws; its Fourier transform is the corresponding infinite product of sinc factors.}

\glossaryentry{infinite-product}{Infinite product}{\statusclassical}{An infinite product \(\prod_{n\ge0}(1+u_n)\) converges to a nonzero limit when the partial products converge and, under standard conditions, when \(\sum|u_n|<\infty\). Locally uniform convergence of holomorphic factors produces a holomorphic limit unless zeros accumulate pathologically. The up-function transform is an entire sinc product. Its zero structure, logarithm, and negative-real specialization drive the Fourier and endpoint theories.}

\glossaryentry{infinite-q-binomial-theorem}{Infinite \(q\)-binomial theorem}{\statuslean}{For \(|q|<1\) and \(|z|<1\),
\[
 \frac{(az;q)_\infty}{(z;q)_\infty}
 =\sum_{n=0}^{\infty}\frac{(a;q)_n}{(q;q)_n}z^n.
\]
The domain stated here guarantees absolute convergence; analytic continuation
may enlarge it only after poles and zeros are controlled.  Euler's two
\(q\)-exponential series are special cases.}

\glossaryentry{infinite-q-pochhammer-product}{Infinite \(q\)-Pochhammer product}{\statuslean}{When \(|q|<1\),
\[
 (a;q)_\infty=\prod_{j=0}^{\infty}(1-aq^j).
\]
The product is an analytic function of \(a\), locally uniformly on \(\C\);
its zeros are exactly \(a=q^{-j}\) unless repeated or cancelled in a quotient.
The repository separates the product-valued definition, normal-convergence
theorems, zero-order statements, and logarithmic-derivative formulas so that
division by the product is never justified merely by a formal factorization.}

\glossaryentry{infinite-random-series}{Infinite random series}{\statusclassical}{An infinite random series is a sum \(\sum_{n\ge0}X_n\) of random variables defined as a limit of partial sums in a chosen mode, such as almost surely or in \(L^2\). If \(X_n\) are uniformly bounded by an absolutely summable deterministic sequence, convergence is immediate for every outcome. The Fabius variable is a weighted series of independent uniforms with geometric weights, so it has compact support and an entire moment-generating function.}

\glossaryentry{infinite-weighted-sum}{Infinite weighted sum}{\statusclassical}{An infinite weighted sum has the form \(\sum a_nx_n\), with convergence governed by the weights and the size of the terms. In probability, independent bounded variables with \(\sum|a_n|<\infty\) give almost-sure absolute convergence. Dyadic weights make the tail comparable to the next scale and create exact self-similarity. This is the probabilistic source of the Fabius dilation equation and of the finite spline truncations.}

\glossaryentry{inflection-point}{Inflection point}{\statusclassical}{An inflection point is a point where the local convexity changes from convex to concave or conversely. A zero of the second derivative alone is not sufficient unless a sign change or equivalent convexity statement is proved. The bounded Fabius function has a unique inflection point at \(x=1/2\): its derivative is strictly increasing on the left half and strictly decreasing on the right half. Symmetry forces the location, while strict unimodality of the up-density gives uniqueness.}

\glossaryentry{information-density}{Information density}{\statusfrontier}{When the joint law is absolutely continuous with respect to the product of
marginals, the information density is the log Radon--Nikodym derivative
\[
 \imath_{X;Y}(x,y)=
 \log\frac{\mathrm dP_{XY}}{\mathrm d(P_XP_Y)}(x,y).
\]
Its expectation is mutual information and its variance is an information
dispersion quantity.}

\glossaryentry{injective-surjective-bijective}{Injective, surjective, and bijective}{\statusclassical}{A map is injective if equal outputs imply equal inputs, surjective if every target value is attained, and bijective if both hold. Strict monotonicity gives injectivity on an interval; continuity and endpoint values often give surjectivity onto the interval between them. The Fabius restriction \(F:[0,1]\to[0,1]\) is bijective, so it has a well-defined inverse. Globally, the constant tails destroy injectivity, which is why inverse statements always specify the restricted domain.}

\glossaryentry{integer-composition}{Integer composition}{\statusclassical}{A composition of a positive integer \(n\) is an ordered tuple of positive integers \((m_1,\ldots,m_r)\) with sum \(n\). Unlike a partition, order matters. Weak compositions allow zero parts. Compositions index coefficients of powers and logarithms of formal series: the coefficient of \(z^n\) in \(U(z)^r\) is a sum over \(r\)-part compositions of \(n\). The all-orders Fabius saddle coefficients use such sums explicitly.}

\glossaryentry{integer-numerator-normalization}{Integer numerator normalization}{\statuslean}{An integer numerator normalization chooses a prescribed denominator
\(D_n\) and defines \(N_n=D_nm_n\) for a rational moment or value \(m_n\).
The normalization is useful when \(N_n\) obeys a division-free recurrence,
has controlled parity, or factors naturally.  Different valid choices produce
different numerator sequences and are not interchangeable.}

\glossaryentry{integer-partition}{Integer partition}{\statusclassical}{An integer partition of \(n\) is a multiset of positive integers whose sum is \(n\), with order ignored. Partitions classify multiplicity patterns in Bell-polynomial and Faà di Bruno expansions. They are distinct from partitions of unity, which are families of functions summing to one, and from ordered compositions. In all-orders coefficient formulas one may pass between partition-indexed multiplicities and composition-indexed ordered products depending on which is easier to formalize.}

\glossaryentry{integer-valued-polynomial}{Integer-valued polynomial}{\statusclassical}{An integer-valued polynomial is a rational-coefficient polynomial taking
integer values at every integer.  Such polynomials are exactly the finite
integer combinations of \(\binom{x}{k}\).  Repeated-sum diagonals often belong
to this ring even when their monomial-basis coefficients have denominators.}

\glossaryentry{integrability}{Integrability}{\statusclassical}{Integrability means that a function has a well-defined integral in the relevant sense, usually Lebesgue or Riemann. Absolute integrability, \(\int|f|<\infty\), ensures stable Fourier transforms and interchange of sums and integrals. Compactly supported continuous functions are absolutely integrable, as are the up-function and its derivatives. Mellin kernels may fail to be integrable at zero until singular terms are subtracted.}

\glossaryentry{integral-transform}{Integral transform}{\statusclassical}{An integral transform maps \(f\) to a new function \(Tf(s)=\int K(s,x)f(x)\,\dd x\) using a kernel \(K\). Fourier, Laplace, Mellin, and moment transforms emphasize different structures: translation, exponential decay, scaling, and polynomial moments. The Fabius theory moves among all four. Each transform has its own convergence region and inversion theorem, so identities must be transferred with explicit hypotheses.}

\glossaryentry{integrality}{Integrality}{\statusclassical}{An integrality theorem proves that a rationally defined expression is in fact an integer. Such results often require divisibility, valuation, or recurrence arguments that are invisible in the original quotient. The repository proves positivity and oddness of normalized Reshetnikov integers and constructs division-free numerator recurrences. Integrality is stronger than rationality and enables modular or parity conclusions that feed back into exact denominator formulas.}

\glossaryentry{interior-support}{Interior of the support}{\statusclassical}{The interior of the support is the largest open set contained in the support. For the up-function, the support is \([-1,1]\) and its interior is \((-1,1)\), where the function is strictly positive. Boundary points belong to the topological support even though the value and every derivative vanish there. Distinguishing support from its interior is important in statements about positivity, analyticity, and zero sets.}

\glossaryentry{interpolating-polynomial}{Interpolating polynomial}{\statusclassical}{Given distinct nodes \(x_0,\ldots,x_n\) and values \(y_j\), the interpolating polynomial is the unique polynomial of degree at most \(n\) satisfying \(p(x_j)=y_j\). Lagrange's formula writes it as a sum of basis polynomials. On geometric nodes, its leading coefficient is a Gaussian-weighted functional of the data. The Fabius \(q\)-binomial identities identify this leading coefficient after Thue--Morse cancellation with an exact dyadic value.}

\glossaryentry{interval-partition}{Interval partition}{\statusfrontier}{An interval partition has every block equal to a consecutive interval in the
natural order.  Such partitions index Boolean moment--cumulant relations.
They are fewer than noncrossing partitions and depend on the linear order.}

\glossaryentry{invariant-measure}{Invariant measure}{\statusmixed}{A measure \(\mu\) is invariant under \(T\) when \(T_\#\mu=\mu\), equivalently
\(\int f\circ T\,\mathrm d\mu=\int f\,\mathrm d\mu\) for bounded measurable
\(f\).  An invariant density is a fixed point of the corresponding transfer
operator.  A functional equation for a Fabius-related function does not by
itself define an invariant dynamical measure.}

\glossaryentry{inverse-barnes-g-asymptotic}{Inverse Barnes-\(G\) asymptotic}{\statusfrontier}{The inverse of Barnes \(G\) at large positive values is obtained by reverting
its quadratic-logarithmic expansion.  The leading coordinate is comparable to
\(\sqrt{\log y/\log\log y}\), after which nested logarithmic corrections are
determined recursively.  Complex inverse branches require additional
monodromy and critical-value analysis.}

\glossaryentry{inverse-branch-closure}{Inverse-branch closure}{\statuslean}{Inverse-branch closure is the class generated from elementary functions by
allowing inverses on intervals or domains where a function is injective and an
analytic branch exists.  The repository uses an explicitly typed version to
state that the inverse Fabius function is not obtained even after adjoining
such inverse branches to the elementary class.  Branch domains and
nonvanishing derivatives are indispensable hypotheses.}

\glossaryentry{inverse-dyadic-value}{Inverse-dyadic value}{\statusmixed}{An inverse-dyadic value is a number \(x\in[0,1]\) satisfying \(F(x)=a/2^n\) for a dyadic ordinate. Existence and uniqueness follow from the bijection \(F:[0,1]\to[0,1]\). Closed forms are generally harder than direct dyadic evaluation because the argument need not itself be dyadic. The repository studies compositions and exact identities at selected inverse-dyadic values, using symmetry and functional equations rather than assuming a simple rational form.}

\glossaryentry{inverse-fabius-self-map-dynamics}{Inverse-Fabius self-map dynamics}{\statusfrontier}{The inverse Fabius function is likewise an increasing self-map of the unit
interval.  Its iterates reverse the time direction of the bounded Fabius
dynamics wherever ordinary inversion applies.  Endpoint flatness of \(F\)
becomes extreme expansion for the inverse, so local derivative formulas,
fractional iteration proposals, and numerical conditioning require separate
endpoint analysis.}

\glossaryentry{inverse-function}{Inverse function}{\statusclassical}{If \(f:I\to J\) is bijective, its inverse \(f^{-1}:J\to I\) satisfies \(f^{-1}(f(x))=x\) and \(f(f^{-1}(y))=y\). For a continuous strictly monotone map between intervals, the inverse is continuous; if \(f'\ne0\), it is differentiable locally with derivative \(1/f'\). The Fabius inverse is defined on \([0,1]\) using the strictly increasing bounded function. Near the endpoints its behavior is governed by inverting the corrected flat asymptotic.}

\glossaryentry{inverse-function-condition-number}{Inverse-function condition number}{\statusmixed}{For a differentiable scalar map \(y=f(x)\) with \(f'(x)\ne0\), the local
absolute sensitivity of the inverse is \(|1/f'(x)|\).  A relative condition
number also scales by \(|y/x|\) when meaningful.  Near a critical point or a
flat Fabius endpoint, inverse evaluation becomes ill conditioned even though
the inverse remains continuous.}

\glossaryentry{inverse-gamma-asymptotic}{Inverse-gamma asymptotic}{\statusfrontier}{An inverse-gamma asymptotic solves \(\Gamma(x)=y\) or
\(\log\Gamma(x)=\log y\) for large \(y\).  Stirling's formula gives a first
Lambert-type approximation; polynomial--logarithmic reversion supplies
successive corrections.  Since gamma is not globally injective on \(\C\),
the branch and real monotonicity region are part of the result.}

\glossaryentry{inverse-lambert-expansion}{Inverse-Lambert expansion}{\statusmixed}{An inverse-Lambert expansion rewrites the solution of an equation such as \(\lambda2^{-\lambda}=x\) in elementary logarithmic terms. Exactly, \(\lambda=-W_{-1}(-x\log2)/\log2\). Iterating the relation \(\lambda\log2=|\log x|+\log\lambda\) gives nested logarithms and a full asymptotic series. The repository formalizes the coefficient algebra and remainder bounds needed to transfer the sharp saddle expansion from \(\lambda\) to powers of \(1/|\log x|\).}

\glossaryentry{inverse-logarithmic-error}{Inverse-logarithmic error}{\statusclassical}{An inverse-logarithmic error is a remainder bounded by a negative power of a growing logarithm, for example \(\bigO(1/|\log x|)\) as \(x\downarrow0\). Such errors decay very slowly and are therefore sensitive to bounded oscillatory terms: a missing nonconstant periodic contribution is not swallowed by them. The corrected Fabius asymptotic attains an inverse-logarithmic error only after evaluating the periodic fluctuation at the exact lower-Lambert phase.}

\glossaryentry{inverse-mellin-residue-expansion}{Inverse-Mellin residue expansion}{\statusmixed}{Shifting a Mellin inversion contour and summing crossed residues yields an
asymptotic expansion.  The remaining contour supplies the error term and
requires growth bounds.  Repeated poles create logarithmic factors; arithmetic
pole lattices create periodic fluctuations.  Formal residue lists alone do not
justify the contour shift.}

\glossaryentry{inverse-gap-modulus}{Inverse modulus from a gap bound}{\statuslean}{An inverse modulus converts a positive lower bound on
\(F(y)-F(x)\) for separated rational or dyadic inputs into a modulus of
continuity for \(F^{-1}\).  The generic formal interface assumes a computable
sequence of positive rational gap bounds; constructing a particular sharp
sequence is a separate task.  This separation makes the logical content of
inverse computability explicit.}

\glossaryentry{inverse-power-table}{Inverse-power table}{\statusclassical}{An inverse-power table records exact or formally generated
coefficients in expansions involving \(x^{-1},x^{-2},\ldots\), usually after a
large coordinate has been selected.  In saddle and Lambert analyses it stores
the algebra needed to divide, compose, exponentiate, or reverse truncated
series without repeatedly expanding from first principles.  The table is an
indexed coefficient family, not the reciprocal of a single scalar.  The former
headword ``inverse-power value'' is retained only as a historical alias.}

\glossaryentry{inversion-formula}{Inversion formula}{\statusclassical}{An inversion formula reconstructs an original object from a transform or cumulative operation. Fourier, Laplace, and Mellin transforms each have inversion integrals; finite differences invert prefix summation; binomial inversion undoes a binomial transform. In the Fabius directory, inversion formulas recover the up-density from its sinc product and the endpoint function from a Bromwich contour. Each inversion requires its own convergence and regularity hypotheses.}

\glossaryentry{fabius-isfabius-characterization}{IsFabius characterization}{\statusmixed}{\texttt{IsFabius} is the Lean predicate collecting the axioms that characterize the bounded Fabius function: constant tails, reflection symmetry, differentiability or smoothness, and the left-half differential equation in the exact repository formulation. Mathematically, it defines a property of candidate functions rather than a new function. Existence supplies one witness and uniqueness proves that any two witnesses agree, allowing later theorems to be stated abstractly for every function satisfying the characterization.}

\glossaryentry{iterated-derivative}{Iterated derivative}{\statusclassical}{The iterated derivative \(f^{(m)}\) is obtained by differentiating \(m\) times. For a self-similar differential equation, each iteration introduces both a coefficient and a chain-rule scaling. The signed Fabius extension satisfies the closed formula \(\mathcal F^{(m)}(x)=2^{\binom{m+1}{2}}\mathcal F(2^mx)\). This turns an apparently analytic hierarchy into exact arithmetic at dyadic points and yields the sharp derivative suprema.}

\glossaryentry{iterated-prefix-kernel}{Iterated prefix kernel}{\statuslean}{The \(r\)-fold prefix-sum operator has kernel
\[
 S^r a(n)=\sum_{j=0}^{n-1}\binom{n-1-j+r-1}{r-1}a(j)
\]
for \(r\ge1\), with the chosen exclusive-prefix convention.  Alternative
inclusive conventions shift the binomial factor.  This kernel makes repeated
Thue--Morse summation accessible to coefficient extraction and polynomial
diagonal analysis.}

\glossaryentry{iterated-prefix-sum}{Iterated prefix sum}{\statusclassical}{For a sequence \(a_n\), its prefix sum is \(A_n=\sum_{k\le n}a_k\); iterating this operation produces higher discrete antiderivatives. Repeated prefix sums of an indicator sequence yield binomial coefficients and discrete B-splines. When applied to Thue--Morse signs, low-order sums vanish over complete dyadic blocks by Prouhet cancellation, and the first nonzero iterated sum has a structured polynomial form. This connects binary combinatorics to spline limits.}

\section{J}\label{sec:letter-j}

\glossaryentry{jackson-fundamental-theorem}{Jackson fundamental theorem of calculus}{\statuslean}{Under convergence and endpoint hypotheses,
\[
 D_q\!\left(\int_0^x f(t)\,\mathrm d_qt\right)=f(x),
 \qquad
 \int_0^a D_qf(t)\,\mathrm d_qt=f(a)-f(0).
\]
The proof is a telescoping geometric series.  This is a \(q\)-calculus theorem,
not a direct consequence of the ordinary fundamental theorem, and the behavior
at the accumulation point must be controlled.}

\glossaryentry{jackson-integration-by-parts}{Jackson integration by parts}{\statuslean}{The integration-by-parts identity is
\[
 \int_0^a f(t)D_qg(t)\,\mathrm d_qt
 =f(a)g(a)-f(0)g(0)
 -\int_0^a g(qt)D_qf(t)\,\mathrm d_qt .
\]
Absolute convergence sufficient for rearrangement is part of the hypothesis.
The shift \(g(qt)\) encodes the nonlocal product rule.}

\glossaryentry{jackson-product-rule}{Jackson product rule}{\statuslean}{The Jackson derivative obeys
\[
 D_q(fg)(x)=g(x)D_qf(x)+f(qx)D_qg(x)
\]
and equivalently
\(D_q(fg)=fD_qg+g(q\,\cdot)D_qf\).
The shifted factor is not optional.  Iteration produces Gaussian-binomial
Leibniz formulas whose powers of \(q\) record how shifts pass through
derivatives.}

\glossaryentry{jackson-q-derivative}{Jackson \(q\)-derivative}{\statuslean}{For \(q\ne1\) and \(x\ne0\),
\[
 D_q f(x)=\frac{f(x)-f(qx)}{(1-q)x}.
\]
At \(x=0\) a value is supplied only under an explicit limiting hypothesis.
The operator satisfies \(D_qx^n=[n]_qx^{n-1}\) and a shifted product rule.
It approaches the ordinary derivative as \(q\to1\) under suitable regularity,
but it is a finite-difference operator at fixed \(q\).}

\glossaryentry{jackson-q-integral}{Jackson \(q\)-integral}{\statuslean}{For \(0<q<1\) and \(a>0\),
\[
 \int_0^a f(t)\,\mathrm d_qt
 =(1-q)a\sum_{n\ge0}q^n f(aq^n),
\]
when the series converges absolutely.  It samples on a geometric comb and is
therefore naturally paired with \(D_q\).  The endpoint \(0\) is approached
through the accumulating nodes rather than through a continuum of points.}

\glossaryentry{j-fraction}{Jacobi continued fraction}{\statusmixed}{The Stieltjes or Cauchy transform of a moment functional has a Jacobi
continued fraction,
\[
 \frac{1}{z-\alpha_0-\displaystyle\frac{\beta_1}
 {z-\alpha_1-\displaystyle\frac{\beta_2}{\ddots}}}.
\]
Its coefficients are the three-term recurrence data.  Finite truncations are
rational functions whose poles are Gaussian nodes and whose residues are
quadrature weights.}

\glossaryentry{jacobi-cubic-identity}{Jacobi cubic identity}{\statuslean}{Jacobi's cubic identity is a theta/product relation obtained by a
three-dissection of the triple product.  It organizes coefficients by residues
modulo three and yields cubic theta series.  The current formalization treats
it as an exact \(q\)-series theorem with absolute-convergence hypotheses, not
as a heuristic rearrangement of bilateral sums.}

\glossaryentry{jacobi-matrix}{Jacobi matrix}{\statusmixed}{The Jacobi matrix is the tridiagonal matrix representing multiplication by
\(x\) in an orthonormal polynomial basis.  Its diagonal entries are the
recurrence coefficients \(\alpha_n\), and its positive off-diagonal entries
are square roots of \(\beta_n\).  Finite principal blocks have eigenvalues
equal to Gaussian quadrature nodes.}

\glossaryentry{jacobi-theta-function}{Jacobi theta function}{\statusmixed}{A Jacobi theta function is a doubly infinite \(q\)-series with
quasi-periodicity and an equivalent infinite-product representation.  One
multiplicative normalization is
\[
 \theta(z;q)=(z;q)_\infty(q/z;q)_\infty(q;q)_\infty.
\]
Different classical \(\vartheta_j\) notations include powers of \(q\), signs,
and additive variables; every use must state its normalization.}

\glossaryentry{jacobi-triple-product}{Jacobi triple product}{\statuslean}{For \(|q|<1\) and \(z\ne0\), a standard form is
\[
 \sum_{n\in\Z}(-1)^n q^{n(n-1)/2}z^n
 =(z;q)_\infty(q/z;q)_\infty(q;q)_\infty.
\]
It identifies a bilateral series with a theta product.  Sign and exponent
shifts vary among conventions; the displayed normalization is the one from
which quasi-periodicity follows directly.}

\glossaryentry{jet}{Jet}{\statusclassical}{The \(m\)-jet of a smooth function at a point is the finite list of derivatives through order \(m\), often encoded as its Taylor polynomial modulo terms of degree \(m+1\). Jets capture the local data needed for finite-order approximation without asserting convergence of the full Taylor series. Saddle expansions use phase and amplitude jets at a moving critical point; formal recurrences combine these jets into all-orders coefficients. At a dyadic Fabius point the infinite jet is eventually zero, yet the function is not locally polynomial.}

\section{K}\label{sec:letter-k}

\glossaryentry{k-fold-convolution}{K-fold convolution}{\statusclassical}{A \(K\)-fold convolution combines \(K\) functions or measures by repeated convolution. For probability laws it is the distribution of a sum of \(K\) independent variables; for compact supports, the support is the Minkowski sum of the individual supports. K-fold identities involving up-functions and Thue--Morse coefficients generalize the basic self-convolution relation. Transform methods reduce them to products or powers and make normalization transparent.}

\glossaryentry{k-fold-thue-morse-sum}{K-fold Thue-Morse sum}{\statusmixed}{A \(K\)-fold Thue--Morse sum is a multiple discrete convolution or signed sum built from the sequence \(\varepsilon_n=(-1)^{\wt(n)}\). Its generating function is a power of the finite Thue--Morse product, so high-order zeros at one translate into vanishing moments and Prouhet-type cancellation. In the repository such sums are related to convolution powers, discrete splines, and scaled limits of the Fabius/up functions.}

\glossaryentry{k-function}{K-function (generalized hyperfactorial)}{\statusfrontier}{The K-function is a continuous extension of the hyperfactorial satisfying a
functional equation related to powers \(z^z\) and Barnes \(G\).  Its
normalization varies in the literature.  Large-argument inversion combines a
quadratic logarithmic leading term with Stirling-type corrections and must
state the selected real or complex branch.}

\glossaryentry{kernel}{Kernel}{\statusclassical}{A kernel is a function that defines an integral, convolution, summation, or transformation operator. Examples include the Fourier exponential, the Laplace exponential, interval indicators, and Bose-Einstein kernels. Kernel properties such as positivity, normalization, support, symmetry, and decay determine the operator's behavior. The Fabius development uses probability kernels for the fixed point, sinc factors as transform kernels, and regularized Mellin kernels for endpoint asymptotics.}

\glossaryentry{knot-multiplicity}{Knot multiplicity}{\statusclassical}{The multiplicity of a spline knot is the number of repeated breakpoint
factors at that location.  Higher multiplicity reduces smoothness across the
knot.  For unequal-scale box splines and finite sinc prefixes, coincident signed
subset sums create multiplicities that must be counted before differentiability
or polynomial degree is inferred.}

\glossaryentry{kolmogorov-distance}{Kolmogorov distance}{\statusclassical}{For real distribution functions \(F\) and \(G\), the Kolmogorov distance is
\[
 d_K(F,G)=\sup_{x\in\R}|F(x)-G(x)|.
\]
It is the natural norm for uniform CDF approximation.  A deterministic bound
on the omitted tail of a monotone random sum yields a Kolmogorov sandwich after
shifting the finite-sum CDF, while a direct density bound may give a sharper
estimate.}

\glossaryentry{koopman-eigenfunction}{Koopman eigenfunction}{\statusfrontier}{A nonzero observable \(f\) is a Koopman eigenfunction when
\(f\circ T=\lambda f\).  Products of eigenfunctions multiply eigenvalues when
the function space is an algebra.  Boundary behavior and the choice of space
can exclude formally plausible eigenfunctions, so spectral claims always name
their ambient domain.}

\glossaryentry{koopman-operator}{Koopman operator}{\statusfrontier}{For a measurable self-map \(T:X\to X\), the Koopman operator is
\[
 (\mathcal K_Tf)(x)=f(Tx)
\]
on a declared function space.  It is linear even when \(T\) is nonlinear.
Its spectrum depends on the Banach or Hilbert space and invariant measure.
Fabius Koopman models remain frontier objects unless those data and a spectral
theorem are explicitly supplied.}

\glossaryentry{kummer-theorem}{Kummer's theorem}{\statusclassical}{For a prime \(p\), Kummer's theorem states that
\(v_p\binom nk\) equals the number of carries when adding \(k\) and \(n-k\)
in base \(p\).  Equivalently, it is
\((s_p(k)+s_p(n-k)-s_p(n))/(p-1)\).  The theorem converts valuation questions
into digit combinatorics and complements Lucas's congruence theorem.}

\section{L}\label{sec:letter-l}

\glossaryentry{l2-norm}{L2 norm}{\statusclassical}{For a square-integrable function on a measure space, the \(L^2\) norm is \(\|f\|_2=(\int|f|^2)^{1/2}\). It comes from the inner product \(\langle f,g\rangle=\int f\overline g\), so orthogonality and Pythagorean identities are available. Fourier-Legendre partial sums are orthogonal projections in \(L^2[-1,1]\) and therefore give the unique least-squares polynomial approximants. Uniform convergence proved later is stronger than convergence in this norm.}

\glossaryentry{lagrange-good-inversion}{Lagrange--Good multivariate inversion}{\statusfrontier}{Good's multivariate inversion formula extracts coefficients of a solution to
\(w_i=z_i\phi_i(\mathbf w)\) through a determinant correcting the Jacobian of
the substitution.  It generalizes Lagrange inversion and is effective for
coupled exponentially small variables in oscillatory inverse problems.  The
full determinant factor is essential.}

\glossaryentry{lagrange-interpolation}{Lagrange interpolation}{\statusclassical}{Lagrange interpolation writes the unique degree-at-most-\(n\) polynomial through distinct data \((x_j,y_j)\) as \(p(x)=\sum_{j=0}^{n}y_j\ell_j(x)\), where \(\ell_j(x)=\prod_{m\ne j}(x-x_m)/(x_j-x_m)\). Its leading coefficient is \(\sum_j y_j/\prod_{m\ne j}(x_j-x_m)\). For geometric nodes \(-2^j\), these denominators simplify to \(q\)-Pochhammer products and Gaussian binomial weights, revealing the interpolation structure of exact Fabius dyadic formulas.}

\glossaryentry{lagrange-inversion-formula}{Lagrange inversion formula}{\statusmixed}{If \(w=z\phi(w)\) with \(\phi(0)\ne0\), then
\[
 [z^n]H(w(z))=\frac1n[t^{n-1}]H'(t)\phi(t)^n
\]
for \(n\ge1\).  In particular it gives coefficients of the inverse of a
power series with nonzero linear term.  Formal and analytic versions share the
coefficient formula but have different convergence hypotheses.}

\glossaryentry{lagrange-up-atom}{Lagrange up-atom}{\statuslean}{A Lagrange up-atom is a normalized translate/dilate of the up-function chosen
so that, on a target interval, its active polynomial piece realizes one
Lagrange cardinal basis polynomial.  Finite sums then interpolate arbitrary
data at the selected nodes.  Outside the interval the atom remains a smooth
compactly supported function rather than the global polynomial.}

\glossaryentry{lah-number}{Lah number}{\statusclassical}{The unsigned Lah number
\(L(n,k)=\binom{n-1}{k-1}n!/k!\) counts partitions of an \(n\)-element set
into \(k\) nonempty linearly ordered lists.  Lah numbers convert rising and
falling factorial bases.  Signed variants insert \((-1)^{n-k}\).}

\glossaryentry{lambert-polynomial}{Lambert polynomial}{\statusmixed}{The Lambert polynomials \(L_n(u)\) package the logarithmic coefficient blocks
in the large-\(x\) expansion of Wright omega or Lambert \(W\).  In the current
zero-based normalization \(L_0(u)=-u\), and for \(n\ge1\) the degree is \(n\)
with leading coefficient \(1/n\).  A differential-integral recurrence
generates all orders.}

\glossaryentry{lambert-w-branch-point}{Lambert \(W\) branch point}{\statusmixed}{The map \(w\mapsto we^w\) has derivative zero at \(w=-1\), giving the
critical value \(-e^{-1}\).  The branches \(W_0\) and \(W_{-1}\) meet there.
With
\(p=\sqrt{2(1+ez)}\), both admit convergent local Puiseux series whose sign
selects the branch.}

\glossaryentry{lambert-w-function}{Lambert W function}{\statusmixed}{The Lambert \(W\) function is the multivalued inverse of \(w\mapsto w\e^w\): it satisfies \(W(z)\e^{W(z)}=z\). On \((-\e^{-1},0)\) there are two real branches, \(W_0\in(-1,0)\) and \(W_{-1}\in(-\infty,-1)\). The Fabius endpoint saddle diverges as \(x\downarrow0\), so the lower branch \(W_{-1}\), not the principal branch, gives the relevant coordinate. Its nested-log expansion translates the exact saddle result into elementary logarithms.}

\glossaryentry{lambert-w-puiseux-coordinate}{Lambert \(W\) Puiseux coordinate}{\statusmixed}{Near \(z=-e^{-1}\), the signed square-root coordinate
\[
 p=\pm\sqrt{2(1+ez)}
\]
linearizes the leading branch separation:
\(W(z)=-1+p-\frac13p^2+\cdots\).
The sign and square-root cut must be stated.  Ordinary Taylor expansion in
\(z+e^{-1}\) cannot represent either inverse branch.}

\glossaryentry{laplace-transform}{Laplace transform}{\statusclassical}{The Laplace transform of a function or measure on a half-line is \(\mathcal Lf(s)=\int_0^\infty\e^{-sx}f(x)\,\dd x\), on the region where the integral converges. For a compactly supported probability law it extends to an entire moment-generating function after a sign change. Convolution becomes multiplication. The Fabius transform factors into elementary terms from the independent uniform summands, and its large negative-real behavior is the starting point for the corrected endpoint asymptotic.}

\glossaryentry{large-deviation-principle}{Large-deviation principle}{\statusfrontier}{A sequence satisfies an LDP with speed \(a_n\) and rate \(I\) when probabilities
of sets decay like \(e^{-a_n\inf I}\), with separate open-set lower and
closed-set upper bounds.  The rate is often the Legendre--Fenchel transform of
a limiting log-MGF.  An LDP describes exponential tails, not central Gaussian
fluctuations.}

\glossaryentry{laurent-expansion}{Laurent expansion}{\statusclassical}{A Laurent expansion represents a holomorphic function on an annulus as \(\sum_{n=-\infty}^{\infty}a_n(z-z_0)^n\). Finitely many negative powers characterize a pole; the coefficient of \((z-z_0)^{-1}\) is the residue. Mellin transforms in the Fabius analysis have Laurent expansions at singular points such as the origin or \(s=1\). The finite part is the constant coefficient after the principal polar terms are removed.}

\glossaryentry{laurent-log-series}{Laurent--logarithmic series}{\statusmixed}{A Laurent--logarithmic series has terms \(t^n p_n(\log t)\), where \(n\in\Z\)
and each \(p_n\) is a polynomial.  It models inverse problems in which
logarithms enter at every algebraic order.  Formal differentiation changes both
the power and the logarithmic coefficient; analytic realization requires a
logarithm branch.}

\glossaryentry{law-pushforward-identity}{Law as a pushforward measure}{\statusclassical}{The law of a random variable \(X:\Omega\to S\) is the pushforward
\(X_\#\Prob\), defined by
\(\operatorname{Law}(X)(B)=\Prob(X\in B)\).  An almost-sure equality of random
variables implies equality of laws, and a measurable map sends laws by another
pushforward.  The repository uses this identity to pass from random-series
recursions to CDF and density statements.}

\glossaryentry{layer-cake-representation}{Layer-cake representation}{\statusclassical}{For a nonnegative measurable function \(f\),
\[
 f(x)=\int_0^\infty \mathbf1_{\{f(x)>t\}}\,\dd t
\]
pointwise, and Tonelli gives
\(\int f=\int_0^\infty\mu\{f>t\}\,\dd t\).
Weighted variants yield moment and fractional-power formulas.  The project
uses layer-cake arguments for CDFs, densities, inverse functions, and level-set
geometry.}

\glossaryentry{leading-coefficient-functional}{Leading-coefficient functional}{\statusclassical}{On polynomials of degree at most \(n\), the leading-coefficient functional sends \(p(x)=a_nx^n+\cdots\) to \(a_n\). Evaluating \(p\) at \(n+1\) distinct nodes gives another coordinate system, so the functional becomes a weighted sum of node values. For geometric nodes the weights are normalized Gaussian binomial coefficients. The Fabius \(q\)-calculus shows that this functional, followed by Thue--Morse block cancellation and exact normalization, returns the desired dyadic value.}

\glossaryentry{least-common-multiple}{Least common multiple}{\statusclassical}{The least common multiple \(\lcm(a_1,\ldots,a_m)\) is the smallest positive integer divisible by each \(a_j\). Prime valuations satisfy \(v_p(\lcm(a_j))=\max_jv_p(a_j)\). LCM products give economical common denominators for families of rational moment or dyadic values. In the Fabius arithmetic layer, they control odd prime factors separately from the exactly known power of two.}

\glossaryentry{least-lipschitz-constant}{Least Lipschitz constant}{\statusmixed}{The least Lipschitz constant of \(f\) is \(\sup_{x\ne y}|f(x)-f(y)|/|x-y|\), possibly infinite. For a continuously differentiable function on an interval it equals \(\sup|f'|\). The bounded Fabius function is globally Lipschitz with optimal constant \(2\): the derivative bound gives the upper estimate, and equality is attained at the midpoint because \(F'(1/2)=2\). Optimality matters for the sharp effective modulus of continuity.}

\glossaryentry{least-squares-approximation}{Least-squares approximation}{\statusclassical}{A least-squares approximation minimizes the squared \(L^2\) error \(\int|f-p|^2\) over a finite-dimensional approximation space. In an orthogonal basis, the minimizer is obtained by truncating the corresponding Fourier coefficients. The Legendre partial sum of degree \(n\) is therefore the unique best polynomial approximation to the up-function in \(L^2[-1,1]\). This optimality is metric-specific and does not by itself imply best uniform approximation.}

\glossaryentry{least-squares-prediction-error}{Least-squares prediction error}{\statuslean}{The least-squares prediction error of \(Y\) from a subspace \(V\subset L^2\)
is
\[
 \mathbb E[(Y-P_VY)^2].
\]
For standardized geometric-uniform variables constructed from common digits,
the current exact formula factors this error into squared hyperbolic-distance
terms \(\prod_j((q-q_j)/(1-qq_j))^2\).  The common probability coupling is part
of the statement.}

\glossaryentry{lebesgue-integral}{Lebesgue integral}{\statusclassical}{The Lebesgue integral integrates measurable functions by measuring level sets and is stable under limit theorems such as monotone and dominated convergence. It naturally handles probability measures and almost-everywhere equivalence. Moments, expectations, convolution measures, and transform integrals in the Fabius development are most naturally Lebesgue integrals, even when continuity and compact support make them agree with classical Riemann integrals.}

\glossaryentry{legendre-block}{Legendre block}{\statuslean}{A Legendre block is a finite polynomial assembled from a consecutive or
parity-restricted range of Legendre modes.  Blocks are used to isolate scale,
parity, or a translate contribution before summing an infinite expansion.
Exact rational coefficients may be computed from moments, while analytic
estimates control the block norm.}

\glossaryentry{legendre-coefficient}{Legendre coefficient}{\statusclassical}{A Legendre coefficient is the component of a function along the Legendre polynomial \(P_n\) under the inner product on \([-1,1]\). With \(\int_{-1}^{1}P_n^2=2/(2n+1)\), the coefficient is \((2n+1)\int fP_n/2\). Symmetry annihilates odd coefficients for an even function. The Fabius directory expresses the up-function coefficients exactly through moments and proves bounds strong enough for absolute uniform convergence.}

\glossaryentry{legendre-coefficient-parity-completion}{Legendre coefficient parity completion}{\statuslean}{If \(a_r\) records only the even coefficient \(c_{2r}\), its parity-completed
sequence is \(c_{2r}=a_r\), \(c_{2r+1}=0\).  The completion is a new sequence,
not a change of index convention.  It allows full-series theorems to consume
an efficiently stored even sequence without losing the parity proof.}

\glossaryentry{legendre-derivative-identity}{Legendre derivative identity}{\statusmixed}{A useful identity is
\[
 (2n+1)P_n(x)=P_{n+1}'(x)-P_{n-1}'(x),\qquad n\ge1.
\]
It transfers derivatives from a smooth function to polynomial combinations
under integration by parts.  Repeated use, together with endpoint flatness,
gives rapid decay of Legendre coefficients of the up-function.}

\glossaryentry{legendre-energy-tail}{Legendre energy tail}{\statuslean}{With classical Legendre normalization, the squared \(L^2\) error after degree
\(N\) is
\[
 \sum_{n>N}\frac{2}{2n+1}|c_n|^2.
\]
This is the Legendre energy tail.  Exact rational coefficients give exact
finite energy sums; decay estimates certify the infinite remainder.  For even
Rvachev data only even indices contribute.}

\glossaryentry{legendre-fenchel-transform}{Legendre--Fenchel transform}{\statusclassical}{For a function \(\Lambda\), its convex conjugate is
\[
 \Lambda^\ast(x)=\sup_t\{tx-\Lambda(t)\}.
\]
When \(\Lambda\) is a differentiable steep log-MGF, the maximizing \(t\)
solves \(\Lambda'(t)=x\).  This transform supplies large-deviation rate
functions for normalized geometric-uniform laws.}

\glossaryentry{legendre-gaunt-coefficient}{Legendre Gaunt coefficient}{\statuslean}{The Legendre Gaunt coefficient is the triple product
\[
 G_{\ell mn}=\int_{-1}^{1}P_\ell(x)P_m(x)P_n(x)\,\dd x.
\]
It vanishes unless the triangle inequalities hold and \(\ell+m+n\) is even.
When admissible, it has an exact factorial formula equivalent to a squared
Wigner \(3j\)-symbol.  The repository proves rational and real-cast forms.}

\glossaryentry{legendre-linearization-coefficient}{Legendre linearization coefficient}{\statuslean}{The product of two Legendre polynomials expands as
\[
 P_\ell(x)P_m(x)=\sum_n L_{\ell m}^{\,n}P_n(x),
\qquad
 L_{\ell m}^{\,n}=\frac{2n+1}{2}G_{\ell mn}.
\]
Only triangle-admissible indices of the correct parity occur.  These
coefficients turn products and nonlinear energies into finite convolutions of
Legendre coefficient sequences.}

\glossaryentry{legendre-normalization}{Legendre normalization}{\statusconvention}{The classical Legendre polynomial \(P_n\) is normalized by
\(P_n(1)=1\) and
\[
 \int_{-1}^{1}P_n(x)P_m(x)\,\dd x
 =\frac{2}{2n+1}\delta_{nm}.
\]
This is neither monic nor orthonormal.  Every coefficient formula in the
project includes the factor \((2n+1)/2\) appropriate to this normalization.}

\glossaryentry{legendre-parity}{Legendre parity}{\statuslean}{Legendre polynomials satisfy \(P_n(-x)=(-1)^nP_n(x)\).  Therefore an even
function has zero odd Legendre coefficients and an odd function has zero even
coefficients.  The Rvachev up-function is even, so its full coefficient
sequence is the parity completion of an even-indexed rational sequence.}

\glossaryentry{legendre-partial-sum}{Legendre partial sum}{\statuslean}{For coefficients \(c_n\), the degree-\(N\) partial sum is
\(S_N(x)=\sum_{n=0}^{N}c_nP_n(x)\).  It is the \(L^2\)-orthogonal projection
onto polynomials of degree at most \(N\) when \(c_n\) are Fourier--Legendre
coefficients.  Uniform convergence requires more than Hilbert-space
projection and is proved from absolute coefficient bounds.}

\glossaryentry{legendre-polynomial}{Legendre polynomial}{\statusclassical}{The Legendre polynomial \(P_n\) is the degree-\(n\) polynomial normalized by \(P_n(1)=1\) and orthogonal to lower-degree polynomials on \([-1,1]\) with unit weight. Rodrigues' formula is \(P_n(x)=\frac1{2^nn!}\frac{\dd^n}{\dd x^n}(x^2-1)^n\), and it solves \(((1-x^2)P_n')'+n(n+1)P_n=0\). These properties support integration by parts, coefficient recurrences, and least-squares approximation of the up-function.}

\glossaryentry{legendres-factorial-valuation-formula}{Legendre's factorial-valuation formula}{\statusclassical}{For a prime \(p\) and \(n\in\N\), Legendre's formula is
\[
 v_p(n!)=\sum_{j\ge1}\left\lfloor\frac{n}{p^j}\right\rfloor
       =\frac{n-s_p(n)}{p-1},
\]
where \(s_p(n)\) is the base-\(p\) digit sum.  The sum is finite.  It converts
factorial quotients into digit and carry data and underlies valuation proofs
for binomial, multinomial, Gaussian-specialization, and denominator
expressions in the arithmetic parts of the project.}

\glossaryentry{legendre-superconvergence}{Legendre superconvergence}{\statusfrontier}{Superconvergence is accuracy at selected points or for selected functionals
that is higher than the global truncation order predicts.  Symmetry, endpoint
flatness, and cancellation of leading Legendre tails can produce it for
Fabius/Rvachev expansions.  A superconvergence claim must state the point set,
norm or functional, and asymptotic scale.}

\glossaryentry{leibniz-rule}{Leibniz rule}{\statusclassical}{The Leibniz rule for the \(n\)-th derivative of a product is \((fg)^{(n)}=\sum_{k=0}^{n}\binom nk f^{(k)}g^{(n-k)}\). It is the differential analogue of binomial convolution. Formalized moment, transform, and saddle calculations repeatedly use it to expand products with exact finite index ranges. For exponential products it combines naturally with Bell polynomials and cumulants.}

\glossaryentry{levy-continuity-theorem}{Levy continuity theorem}{\statusclassical}{Levy's continuity theorem states that pointwise convergence of characteristic functions to a function continuous at zero is equivalent to weak convergence of the corresponding probability measures. It is a standard route from finite convolution products to an infinite probability law. In the Fabius construction, the explicit convergence of the sinc products can identify the weak limit and its characteristic function.}

\glossaryentry{lipschitz-function}{Lipschitz function}{\statusclassical}{A function is Lipschitz with constant \(L\) when \(|f(x)-f(y)|\le L|x-y|\) for all relevant \(x,y\). Lipschitz continuity is stronger than uniform continuity and implies absolute continuity on intervals. A bounded derivative gives a Lipschitz bound by the mean value theorem. The Fabius function has least global constant \(2\), while related spline approximants inherit uniform Lipschitz control.}

\glossaryentry{local-phase}{Local phase}{\statusclassical}{The local phase is the Taylor expansion of the exponent in a contour integral around its saddle after the linear term has vanished. In a scaled variable \(t\), it begins with a negative quadratic Gaussian term, followed by cubic and higher corrections whose coefficients are phase jets. The Fabius saddle modules isolate the local phase from the global periodic factor, prove a central approximation, and bound the remainder on a radius where the expansion is uniformly valid.}

\glossaryentry{locally-finite-sum}{Locally finite sum}{\statusclassical}{A sum of functions \(\sum_jf_j(x)\) is locally finite if every point has a neighborhood meeting the support of only finitely many summands. Such a sum is unambiguously defined and inherits continuity or smoothness termwise without convergence estimates. Translate sums of compactly supported up-functions are often locally finite. The signed global Fabius extension also has presentations where only finitely many translated support pieces contribute at a fixed point, alongside other genuinely infinite absolutely convergent series.}

\glossaryentry{locally-uniform-convergence}{Locally uniform convergence}{\statusclassical}{A sequence converges locally uniformly on a domain if it converges uniformly on every compact subset. For holomorphic functions, locally uniform convergence preserves holomorphy and permits termwise differentiation. Infinite products are also commonly controlled locally uniformly. The sinc product for the up-function's Fourier transform and the negative-Laplace logarithmic series are established on compact subsets of their natural domains before global growth or tail estimates are considered.}

\glossaryentry{gamma-logarithm}{Log-gamma function}{\statusmixed}{The log-gamma function is a branch of \(\log\Gamma(z)\) on a domain avoiding
gamma poles and chosen cuts.  On the positive real axis it is the real
logarithm of \(\Gamma(x)\).  Stirling expansions are naturally stated for
\(\log\Gamma\), while inverse gamma problems first invert this large
polynomial--logarithmic scale and then exponentiate if needed.}

\glossaryentry{log-periodic-asymptotic}{Log-periodic asymptotic}{\statuslean}{A log-periodic asymptotic contains a periodic function of a logarithmic or
Lambert-type scale, such as \(\Psi(\log_b(1/x))\).  Multiplying \(x\) by a
fixed base shifts the phase by an integer and leaves the oscillation invariant.
In the Fabius endpoint law, the fluctuation is nonconstant and must remain in
the exponential main term to obtain ratio asymptotics.}

\glossaryentry{log-periodic-function}{Log-periodic function}{\statusclassical}{A function of \(x>0\) is log-periodic with ratio \(b>1\) if it is periodic in \(\log_b x\), equivalently \(Q(bx)=Q(x)\). Such oscillations arise in discrete-scale-invariant problems because scaling by \(b\) advances the logarithmic phase by one. The Fabius endpoint correction is one-periodic in a lower-Lambert coordinate rather than exactly in \(\log_2x\), but it has the same dyadic origin. A log-periodic term can be tiny yet prevent convergence to a constant asymptotic ratio.}

\glossaryentry{log-squared-asymptotic}{Log-squared asymptotic}{\statusclassical}{A log-squared asymptotic has a dominant term proportional to \((\log(1/x))^2\), often in the logarithm of a very flat function. For the Fabius endpoint, \(\log F(x)\) has a leading negative quadratic logarithmic scale, modified by lower nested-log, constant, and periodic contributions. This explains decay faster than every power of \(x\) but slower than typical \(\exp(-c/x^\alpha)\) scales. The exact coefficient is determined by the dyadic dilation factor \(\log2\).}

\glossaryentry{logarithmic-delay-equation}{Logarithmic delay equation}{\statusclassical}{A logarithmic change of variables converts multiplicative scaling into additive translation. If \(x=\e^{-t}\), then \(f(2x)\) becomes a shifted function of \(t-\log2\). The resulting relation is a delay or advance equation with fixed step \(\log2\). This viewpoint explains why periodic functions of a logarithmic coordinate appear as homogeneous or residual solutions in the Fabius endpoint analysis.}

\glossaryentry{logarithmic-height}{Logarithmic height}{\statusconvention}{Logarithmic height records the deepest iterated logarithm needed by a
transmonomial, while exponential height records nested exponentials.  The
height is structural, not the power of a logarithm.  Declaring it helps choose
a closed transseries field before composition or inversion.}

\glossaryentry{logarithmic-scale}{Logarithmic scale}{\statusclassical}{A logarithmic scale uses \(L=\log(1/x)\), \(\log L\), and inverse powers of these quantities to organize a limit as \(x\downarrow0\). It is appropriate when the decisive saddle grows only logarithmically. The Fabius expansion is most exact in \(\lambda=-W_{-1}(-x\log2)/\log2\), but \(\lambda\) itself has a nested-log expansion. Distinguishing the exact coordinate from its elementary logarithmic approximation prevents phase errors.}

\glossaryentry{logarithmic-stieltjes-fixed-point}{Logarithmic Stieltjes fixed point}{\statuslean}{A logarithmic Stieltjes representation integrates
\(\log(z-x)\) or its derivative against a compactly supported law and satisfies
a scale-recursive fixed-point equation.  The repository proves a real
logarithmic identity for \(z>1\).  No unmentioned complex logarithm branch or
cut continuation is part of that theorem.}

\glossaryentry{lower-lambert-branch}{Lower Lambert branch}{\statusconvention}{The lower real branch \(W_{-1}\) of Lambert \(W\) is defined on \([-\e^{-1},0)\), takes values in \(( -\infty,-1]\), and satisfies \(W_{-1}(z)\e^{W_{-1}(z)}=z\). As \(z\uparrow0\), it tends to \(-\infty\). This branch solves the large positive saddle equation in the Fabius endpoint problem. Using the principal branch would produce a bounded coordinate and the wrong asymptotic regime.}

\glossaryentry{lower-lambert-coordinate}{Lower Lambert coordinate}{\statusmixed}{For small \(x>0\), the repository defines the lower Lambert coordinate by \[ \lambda(x)=-\frac{W_{-1}(-x\log2)}{\log2}, \qquad \lambda(x)2^{-\lambda(x)}=x. \] It is the exact real saddle parameter, tends to infinity as \(x\downarrow0\), and advances by approximately one under dyadic rescaling. The periodic correction and all inverse-power coefficients are most naturally evaluated at \(\lambda(x)\), not at a truncated logarithmic surrogate.}

\glossaryentry{lower-lambert-phase}{Lower-Lambert phase}{\statuslean}{The lower-Lambert phase is the large positive coordinate obtained from the
\(W_{-1}\) solution of the endpoint saddle equation.  It converts the implicit
balance between a small argument and a logarithm into a monotone parameter
tending to infinity.  Fabius endpoint expansions are organized in inverse
powers of this phase, with a separate period-one fluctuation.}

\glossaryentry{lower-order-term}{Lower-order term}{\statusclassical}{A lower-order term is asymptotically smaller than the leading term in a specified scale. The phrase is meaningful only after the comparison variable and mode of comparison are fixed. A bounded periodic term is lower order than a quadratic logarithm but is not lower order than an \(o(1)\) remainder in the exponent. The Fabius correction illustrates why informal claims that a term is negligible must be matched to the claimed error scale.}

\glossaryentry{lucas-theorem}{Lucas's theorem}{\statusclassical}{Let \(p\) be prime and write \(n=\sum n_jp^j\) and
\(k=\sum k_jp^j\).  Lucas's theorem states
\[
 \binom nk\equiv\prod_j\binom{n_j}{k_j}\pmod p .
\]
It detects binomial coefficients nonzero modulo \(p\) digit by digit.  The
repository uses the binary specialization together with carry theorems and
\(q\)-Lucas analogues to analyze Pascal rows, parity masks, cyclotomic
specializations, and Thue--Morse structure.}

\section{M}\label{sec:letter-m}

\glossaryentry{maclaurin-series}{Maclaurin series}{\statusclassical}{A Maclaurin series is a Taylor series centered at zero. When the function is analytic there, \[ f(z)=\sum_{n\ge0}\frac{f^{(n)}(0)}{n!}z^n. \] Moment-generating and Fourier transforms of compactly supported laws have entire Maclaurin series whose coefficients are moments. By contrast, the bounded Fabius function itself has the zero Taylor series at the flat endpoint \(0\) but is positive immediately to the right, a standard demonstration that smoothness does not imply analyticity.}

\glossaryentry{mahler-equation}{Mahler equation}{\statusmixed}{A Mahler equation relates a function at \(z\) to its values at
\(z^q,z^{q^2},\ldots\), for example
\(\sum_{j=0}^rp_j(z)f(z^{q^j})=0\).  Generating functions of automatic and
regular sequences often satisfy such equations.  The reciprocal-Gamma and
geometric-product parts of the current tree also contain dilation equations
of Mahler type.}

\glossaryentry{mahler-function}{Mahler function}{\statusmixed}{A Mahler function is an analytic or formal function satisfying a Mahler
equation for an integer dilation \(q\ge2\).  Its arithmetic values,
transcendence, and singularities are constrained by the functional equation.
The Fabius corpus invokes this viewpoint for digit sequences, \(q\)-products,
and scale-recursive transforms, while distinguishing formal and analytic
solutions.}

\glossaryentry{major-arc}{Major arc}{\statusclassical}{In a contour or oscillatory integral, a major arc is the neighborhood of the dominant saddle or stationary point where the main contribution is extracted. Its width is chosen to make a local Taylor expansion accurate while retaining nearly all Gaussian mass. The complement is the minor arc or tail. Fabius saddle modules prove separate central estimates on the major arc and exponential or comparison bounds outside it, then balance the chosen radius to obtain a uniform remainder.}

\glossaryentry{mass-expansion}{Mass expansion}{\statusclassical}{A mass expansion is the asymptotic series for the normalized integral remaining after the leading exponential value and Gaussian scale have been factored out. It typically begins \(1+a_1/\lambda+a_2/\lambda^2+\cdots\). The coefficients combine phase and amplitude jets and Gaussian moments. Taking the formal logarithm of this series gives the additive coefficients in \(\log F(x)\). The repository provides both recurrences and explicit ordered-composition formulas for these saddle mass coefficients.}

\glossaryentry{mathematical-interface-theorem}{Mathematical interface theorem}{\statusmixed}{An interface theorem translates between representations without exposing
the internal construction of either side.  Typical interfaces identify a
rational recurrence with a real integral, a formal coefficient with an
analytic transform derivative, or a distributional fixed point with a CDF
equation.  Such theorems stabilize downstream work: later modules can use the
mathematical contract even if implementation details change.}

\glossaryentry{matrix-dilation}{Matrix dilation}{\statusfrontier}{For an expansive real matrix \(A\), the dilation \(D_Af(x)=f(Ax)\) scales
Fourier variables by \(A^{-\mathsf T}\) and volume by \(|\det A|^{-1}\).
A matrix refinement equation combines this dilation with lattice translates.
Scalar dyadic theorems do not automatically extend because anisotropy and
lattice geometry enter.}

\glossaryentry{mean-value-theorem}{Mean value theorem}{\statusclassical}{The mean value theorem states that a continuous function on \([a,b]\), differentiable on \((a,b)\), satisfies \(f(b)-f(a)=f'(c)(b-a)\) for some \(c\). Consequently \(|f(b)-f(a)|\le\sup|f'|\,|b-a|\). Iterating such estimates with the exact derivative scaling yields one form of effective Fabius flatness. It also proves Lipschitz continuity and turns derivative maxima into sharp global constants.}

\glossaryentry{with-density-measure}{Measure with density}{\statuslean}{Given a measure \(\mu\) and a nonnegative measurable function \(f\), the
measure \(f\,\mu\), written \(\mu.\mathrm{withDensity}(f)\) in Lean, satisfies
\(\int g\,\dd(f\mu)=\int gf\,\dd\mu\) under the standard extended-integral
conventions.  The geometric-uniform modules identify their probability law
with Lebesgue measure weighted by the explicit density.}

\glossaryentry{measure-zero-set}{Measure-zero set}{\statusclassical}{A set has Lebesgue measure zero if it can be covered by intervals of arbitrarily small total length. Statements holding outside such a set hold almost everywhere. Countable sets, including the dyadic rationals, have measure zero even when dense. In analytic-locus discussions this distinction matters: exceptional behavior on every dyadic point may be topologically dense yet measure-theoretically negligible. The Fabius nowhere-analytic result is stronger than an almost-everywhere statement.}

\glossaryentry{mellin-finite-part}{Mellin finite part}{\statusclassical}{A Mellin finite part is the constant term obtained when a Mellin transform has a pole at the point of interest. One subtracts the polar Laurent terms or equivalently regularizes the original endpoint integral. In the Fabius negative-Laplace decomposition, the Mellin origin carries the polynomial-logarithmic main terms, while shifted imaginary poles generate the periodic correction. Correct extraction of the finite part fixes the endpoint constant and the centering of the periodic function.}

\glossaryentry{mellin-inversion}{Mellin inversion}{\statusclassical}{Mellin inversion reconstructs a function from
\[
 M(s)=\int_0^\infty f(x)x^{s-1}\,\dd x
\]
by
\[
 f(x)=\frac{1}{2\pi\ii}\int_{c-\ii\infty}^{c+\ii\infty}
 M(s)x^{-s}\,\dd s,
\]
under suitable hypotheses.  Shifting the contour across poles produces
asymptotic expansions.  Periodically spaced imaginary poles yield Fourier
series in \(\log x\).  This is the conceptual source of the Gamma--zeta
periodic correction in the Fabius Laplace logarithm.}

\glossaryentry{mellin-pole-asymptotics}{Mellin-pole asymptotics}{\statusmixed}{In Mellin inversion, a pole of the transformed integrand contributes a power
or power--logarithm term obtained from its residue.  A vertical lattice of
poles with equal real part contributes a periodic function of \(\log x\).
This is the mechanism by which dyadic harmonic sums produce log-periodic
Fabius corrections.}

\glossaryentry{mellin-strip}{Mellin strip}{\statusclassical}{The fundamental strip of a Mellin transform
\(\int_0^\infty f(x)x^{s-1}\,\dd x\) is the vertical strip in which the
integral converges absolutely.  Meromorphic continuation beyond the strip is a
separate theorem.  Inversion contours begin inside a valid strip before they
are shifted across poles.}

\glossaryentry{mellin-transform}{Mellin transform}{\statusclassical}{The Mellin transform is \(M_f(s)=\int_0^\infty f(x)x^{s-1}\,\dd x\). It converts multiplicative scaling into multiplication by powers and is the Fourier transform in logarithmic coordinates. Harmonic sums over dyadic scales become products with \((1-2^{-s})^{-1}\), whose imaginary lattice of poles creates one-periodic functions. The Fabius endpoint analysis applies Mellin methods to a regularized Bose-type kernel derived from the negative-Laplace product.}

\glossaryentry{meromorphic-continuation}{Meromorphic continuation}{\statusclassical}{Meromorphic continuation extends a function beyond its initial domain while allowing isolated poles. Gamma and zeta functions are standard examples. Mellin transforms often converge only in a vertical strip but continue meromorphically through functional identities or subtraction of endpoint asymptotics. In the Fabius analysis, locating the poles and computing their residues is more important than merely knowing continuation exists, because the pole at zero and the imaginary dyadic lattice determine the main and periodic terms.}

\glossaryentry{mersenne-factor}{Mersenne factor}{\statusclassical}{A Mersenne factor is an integer of the form \(2^m-1\). Products of odd Mersenne factors arise when the half-base \(q\)-Pochhammer symbol \((1/2;1/2)_n=\prod_{j=1}^{n}(1-2^{-j})\) is cleared of powers of two. These products appear in moment denominators, Gaussian binomial normalizations, and Reshetnikov integers. Their oddness makes them especially convenient for separating \(2\)-adic valuations from odd divisibility.}

\glossaryentry{minimum-mean-square-error}{Minimum mean-square error}{\statusfrontier}{The MMSE of estimating \(X\) from \(Y\) is
\[
 \operatorname{mmse}(X|Y)=
 \mathbb E[(X-\mathbb E[X|Y])^2].
\]
The conditional expectation is the unique \(L^2\) optimizer.  Restricting the
estimator to a finite linear span gives a larger least-squares projection error,
which must be distinguished from the unrestricted MMSE.}

\glossaryentry{minkowski-sum-of-supports}{Minkowski sum of supports}{\statusclassical}{For subsets \(A,B\subset\R\), their Minkowski sum is \(A+B=\{a+b:a\in A,b\in B\}\). The support of a convolution satisfies \(\supp(f*g)\subseteq\supp f+\supp g\), with equality under positivity and mild nondegeneracy. Finite sums of scaled uniform variables therefore have support equal to the sum of their interval supports. This calculation fixes the centering and endpoint locations of the spline approximants.}

\glossaryentry{minor-arc}{Minor arc}{\statusclassical}{The minor arc is the portion of an oscillatory or contour integral outside the central saddle neighborhood. A successful saddle-point proof shows that its contribution is smaller than the target remainder, often by a strict loss in the real part of the phase or by transform decay. The Fabius development treats central, tail, and reference-weight estimates in separate modules, preventing a local Gaussian expansion from being mistaken for a complete contour proof.}

\glossaryentry{minor-arc-bound}{Minor-arc bound}{\statuslean}{A minor-arc bound shows that oscillatory or remote parts of a contour are
smaller than the main saddle contribution, often exponentially so.  It uses
nonresonance, product decay, or a quantitative phase gap.  The terminology is
analogous to the circle method but here applies to the Bromwich/Fourier
localization.}

\glossaryentry{subspace-lattice-mobius-function}{M\"obius function of a subspace lattice}{\statusclassical}{If \(U\subseteq Z\) and
\(d=\dim_{\mathbb F_q}(Z/U)\), then the subspace-lattice M\"obius function is
\[
 \mu(U,Z)=(-1)^d q^{\binom d2}.
\]
This is the \(q\)-analogue of the Boolean-lattice sign and is central to
finite-field inversion formulas.}

\glossaryentry{partition-lattice-mobius-function}{M\"obius function of the partition lattice}{\statusclassical}{For partitions \(\pi\le\sigma\), the partition-lattice M\"obius function
is the incidence-algebra inverse of the constant-one zeta function.  In
particular,
\(\mu(\hat0,\hat1)=(-1)^{n-1}(n-1)!\).  M\"obius inversion converts moments
to cumulants and connected combinatorial structures to arbitrary assemblies.}

\glossaryentry{mod-gaussian-convergence}{Mod-Gaussian convergence}{\statusfrontier}{Mod-Gaussian convergence factors a characteristic or moment-generating
function into a diverging Gaussian term and a residual analytic limit.  It
refines the central limit theorem and can yield moderate deviations and local
limits.  The normalization and domain of locally uniform convergence are part
of the definition.}

\glossaryentry{moderate-deviation-regime}{Moderate-deviation regime}{\statusfrontier}{Moderate deviations lie between central-limit fluctuations and full
large-deviation scale.  The deviation size grows relative to the standard
deviation but remains smaller than the macroscopic scale.  Cumulant bounds or
mod-Gaussian convergence can provide asymptotics there; neither a fixed-order
Edgeworth expansion nor an LDP alone covers the entire intermediate range.}

\glossaryentry{module-docstring}{Module docstring}{\statusconvention}{A module docstring is the opening documentation block that states a Lean
file's purpose, conventions, imported prerequisites, and theorem boundaries.
In a mathematically dense repository it functions as a local contract: it
should identify normalization, endpoint behavior, parameter hypotheses, and
known limitations before individual declarations are read.  Documentation
audits test the presence of these headers separately from theorem comments.}

\glossaryentry{modulus-of-continuity}{Modulus of continuity}{\statusclassical}{A modulus of continuity is a nondecreasing function \(\omega(\delta)\to0\) such that \(|f(x)-f(y)|\le\omega(|x-y|)\). A Lipschitz function has \(\omega(\delta)=L\delta\). An effective modulus must be computable in a specified representation. For the Fabius function, the exact least Lipschitz constant gives both a sharp analytic modulus and a simple recursive modulus used in computability.}

\glossaryentry{moment}{Moment}{\statusclassical}{The \(n\)-th moment of a random variable or finite measure is \(\E[X^n]\) or \(\int x^n\,\mu(\dd x)\), when finite. Moments encode location, scale, and shape; for compactly supported laws all moments exist and determine the law. The Fabius development uses full moments of the centered up-density and half moments on one side of the support. Their triangular recurrences and generating functions control exact dyadic values and Legendre coefficients.}

\glossaryentry{moment-cumulant-transform}{Moment--cumulant transform}{\statuslean}{Moments and cumulants are related by set partitions:
\[
 m_n=\sum_{\pi\in\Pi_n}\prod_{B\in\pi}\kappa_{|B|}.
\]
The inverse uses the M\"obius function of the partition lattice.  Complete
exponential Bell polynomials encode both directions.  Current moment,
Appell, and standardized-law modules use these transforms with exact rational
coefficients.}

\glossaryentry{moment-functional}{Moment functional}{\statuslean}{A moment functional is a linear map \(\mathcal M\) on polynomials determined by
\(\mathcal M(x^n)=m_n\).  When \(m_n=\int x^n\,\mathrm d\mu(x)\), it is
integration against a measure, but the algebraic definition does not require a
measure.  The Rvachev moment functional permits exact rational construction of
orthogonal polynomials, Legendre coefficients, and Appell deconvolution before
casting to real integrals.}

\glossaryentry{moment-generating-function}{Moment-generating function}{\statusclassical}{The moment-generating function is \(M_X(t)=\E(\e^{tX})\), defined near zero when the exponential moment exists. Its Taylor coefficients are moments. For bounded \(X\), it is entire. Independence turns sums into products, while the logarithm gives cumulants. The Fabius random series therefore has a product MGF whose imaginary restriction is the sinc characteristic function and whose negative-real restriction drives endpoint analysis.}

\glossaryentry{moment-laurent-expansion}{Moment Laurent expansion}{\statuslean}{For compactly supported \(\mu\),
\[
 G_\mu(z)=\sum_{n\ge0}\frac{m_n}{z^{n+1}}
\]
for \(|z|\) larger than the support radius, where
\(m_n=\int x^n\,\dd\mu(x)\).  The expansion is a convergent Laurent series,
not merely formal.  It transfers exact moment recurrences to the resolvent.}

\glossaryentry{moment-numerator}{Moment numerator}{\statusmixed}{A moment numerator is an integer sequence obtained by multiplying a rational moment by a canonical denominator built from factorials, powers of two, and Mersenne factors. In the repository the full-moment numerators and half-moment numerators are distinct sequences, often denoted \(F_n\) and \(G_n\) in the arithmetic paper's conventions. Their integer recurrences expose parity, positivity, and divisibility information hidden in the raw rational moments.}

\glossaryentry{moment-numerator-f}{Moment numerator $F_n$}{\statusmixed}{The arithmetic sequence \(F_n\) is a canonical integer numerator obtained from the full moments after clearing their prescribed denominator. It is distinct from the function \(F(x)\), despite sharing the letter in source conventions. Its triangular recurrence and divisibility properties organize exact moment arithmetic. The glossary uses the phrase ``moment numerator \(F_n\)'' whenever ambiguity with the bounded Fabius function is possible.}

\glossaryentry{moment-recurrence}{Moment recurrence}{\statusclassical}{A moment recurrence expresses a moment of order \(n\) in terms of lower moments, often by expanding a distributional fixed-point equation. If \(X\) has the same law as a scaled combination of \(X\) and an independent elementary variable, applying \(\E[(\cdot)^n]\) and the binomial theorem yields a triangular relation. The Fabius full and half moment recurrences are rational and exact, and after normalization become division-free integer recurrences.}

\glossaryentry{moment-sequence}{Moment sequence}{\statusclassical}{A moment sequence is a sequence \(m_n=\int x^n\,\mu(\dd x)\) associated with a measure. Not every sequence is a moment sequence; positivity constraints and growth conditions are involved. Compact support makes the moment problem determinate. The repository's sequence \(c_n\) records even centered moments, while \(d_n\) records normalized half moments. Their arithmetic normalizations are used throughout the dyadic and Legendre layers.}

\glossaryentry{monic-orthogonal-polynomial}{Monic orthogonal polynomial}{\statuslean}{The monic orthogonal polynomial \(\pi_n\) has degree \(n\), leading
coefficient one, and is orthogonal to every polynomial of degree \(<n\) under
the chosen moment functional.  Positive definiteness makes it unique.
Determinant formulas and Gram--Schmidt construct it entirely from
\(m_0,\ldots,m_{2n}\).}

\glossaryentry{monodromy}{Monodromy}{\statusfrontier}{Monodromy is the permutation of analytic branches produced by continuation
around a singular value.  For an algebraic critical point it cycles Puiseux
branches; for logarithmic inverses it may add multiples of \(2\pi i\).
A branch label is therefore incomplete without a base point and continuation
domain.}

\glossaryentry{monotone-convergence-theorem}{Monotone convergence theorem}{\statusclassical}{The monotone convergence theorem states that if \(0\le f_n\uparrow f\) pointwise, then \(\int f_n\uparrow\int f\), allowing the value \(+\infty\). It is useful for increasing approximations by positive densities or distribution functions. Fabius finite probability constructions often allow stronger dominated or uniform convergence, but monotone convergence remains a foundational tool for nonnegative integral limits.}

\glossaryentry{monotone-function}{Monotone function}{\statusclassical}{A function is nondecreasing when \(x<y\) implies \(f(x)\le f(y)\), and strictly increasing when the inequality is strict. Every cumulative distribution function is nondecreasing. The bounded Fabius function is strictly increasing on \([0,1]\), because its density is positive in the interior, but it is constant on the exterior tails. The up-function is not globally monotone; it increases to the origin and then decreases.}

\glossaryentry{monotone-order-isomorphism}{Monotone order isomorphism}{\statuslean}{A monotone order isomorphism is a bijection preserving and reflecting order,
with a monotone inverse.  The restriction of the canonical Fabius function to
\([0,1]\) has this property.  Order-theoretic inverse lemmas yield continuity,
endpoint values, and interval mapping without differentiability, and they
remain valid in settings where analytic inverse-function theorems are
inapplicable.}

\glossaryentry{q-multifactorial-product}{Multiple \(q\)-Pochhammer notation}{\statusconvention}{A compact product such as
\[
 (a_1,\ldots,a_r;q)_n=\prod_{\ell=1}^r(a_\ell;q)_n
\]
is multiple \(q\)-Pochhammer notation.  It abbreviates a product of independent
factors and introduces no multivariate coupling.  The base \(q\) and length
\(n\) are shared; each upper parameter retains its own zero set.}

\glossaryentry{multiplicative-periodic-fluctuation}{Multiplicative periodic fluctuation}{\statusclassical}{A multiplicative periodic fluctuation is a bounded factor whose value depends periodically on a logarithmic phase, typically \(\exp(\Psi(\lambda(x)))\). In logarithmic asymptotics it appears additively as \(\Psi\). Because it need not approach a limit, it can prevent a ratio from converging even when all coarse logarithmic orders agree. The corrected Fabius endpoint formula contains exactly such a factor, with very small but explicitly nonzero Fourier modes.}

\glossaryentry{multiresolution-analysis}{Multiresolution analysis}{\statusfrontier}{A multiresolution analysis is a nested sequence of closed spaces
\(V_j\) compatible with dilation and translation, whose union is dense and
intersection is trivial.  Orthogonal or biorthogonal detail spaces describe
the increments.  Finite nested interpolation spaces are a discrete analogue
but do not constitute an infinite MRA until density and stability are proved.}

\glossaryentry{multiresolution-detail-operator}{Multiresolution detail operator}{\statuslean}{For nested finite interpolation projectors \(P_{q-1},P_q\), the detail
operator is \(\Delta_q=P_q-P_{q-1}\), with \(\Delta_0=P_0\).
Under compatible nesting these details are idempotent on the terminal
coefficient space and mutually annihilating.  Their ranks equal the increments
in polynomial dimension.}

\glossaryentry{mutual-information}{Mutual information}{\statusfrontier}{Mutual information is
\[
 I(X;Y)=D_{\rm KL}(\mathcal L(X,Y)\|
 \mathcal L(X)\otimes\mathcal L(Y)).
\]
It measures dependence and is invariant under measurable bijections of either
variable.  Exact geometric-digit decompositions can turn it into a finite sum
of scale contributions.}

\section{N}\label{sec:letter-n}

\glossaryentry{namespace}{Namespace}{\statusconvention}{A namespace groups related Lean declarations and prevents name collisions.
The mathematical namespace \texttt{Fabius} contains functions, arithmetic
sequences, transform objects, and approximation operators that would otherwise
have generic names.  Opening a namespace shortens notation locally but does
not change a declaration's full identity; crosswalks use sufficiently complete
names to remain unambiguous.}

\glossaryentry{negative-index-bell-value}{Negative-index Bell value}{\statusfrontier}{Negative-index Bell values in the corpus are defined by an explicit
absolutely convergent Roman-type continuation, not by counting partitions of
a negative set.  A representative formula expresses \(B_{-k}\) through
negative-index first-kind Stirling values or
\(e^{-1}\sum_{j\ge1}(j^k j!)^{-1}\).  The definition must precede any use of
the familiar symbol \(B_n\) at negative \(n\).}

\glossaryentry{negative-laplace-decomposition}{Negative-Laplace decomposition}{\statusmixed}{The negative-Laplace decomposition analyzes \(\log M(-s)\) for large \(s>0\), where \(M\) is the Fabius moment-generating function. An exact dyadic product is converted into a harmonic sum of the kernel \(\kappa(u)=\log((1-\e^{-u})/u)\). Mellin analysis then separates explicit polynomial-logarithmic terms, a constant, a smooth one-periodic correction, and a decaying remainder. This decomposition is the global input for the Bromwich saddle method.}

\glossaryentry{negative-laplace-product}{Negative-Laplace product}{\statusmixed}{The negative-Laplace product is the exact factorization of the Fabius transform on the negative real axis, for example \(G(-s)=\prod_{j\ge1}(1-\e^{-s/2^j})/(s/2^j)\). Each factor is the Laplace transform of a uniform summand after normalization. The product is positive for \(s>0\), so its real logarithm is unambiguous. Its dyadic scale structure produces the Mellin denominator \(1-2^{-z}\) and hence the periodic correction.}

\glossaryentry{newton-identities}{Newton identities}{\statusclassical}{Newton identities relate power sums to elementary or complete homogeneous
symmetric polynomials.  For example,
\(me_m=\sum_{j=1}^m(-1)^{j-1}e_{m-j}p_j\).
They convert product data into logarithmic coefficients and underlie moment,
cumulant, and \(q\)-product recurrences.}

\glossaryentry{newton-iteration}{Newton iteration}{\statusclassical}{Newton iteration solves \(h(x)=0\) by \(x_{n+1}=x_n-h(x_n)/h'(x_n)\). Near a simple root it converges quadratically under standard hypotheses. Stable endpoint evaluation of inverse Fabius values or saddle coordinates can use Newton's method with monotone bracketing and exact asymptotic initial guesses. The lower-Lambert equation \(\lambda2^{-\lambda}=x\) may also be solved numerically this way when a special-function implementation is unavailable.}

\glossaryentry{newton-series}{Newton series}{\statusclassical}{A Newton series expands a function on integers as
\[
 f(x)=\sum_{k\ge0}\binom{x}{k}\Delta^kf(0),
\]
formally or under analytic growth conditions.  Finite sums give ordinary
polynomial interpolation on consecutive nodes.  The expansion is the discrete
counterpart of a Taylor series and interacts directly with prefix summation.}

\glossaryentry{noise-floor}{Noise floor}{\statusfrontier}{The numerical noise floor is the error level below which rounding, cancellation, underflow, or implementation effects dominate the signal being measured. The Fabius periodic correction has very small amplitude, so ordinary floating-point evaluation can obscure it. Exact rational endpoint values, arbitrary precision, and subtraction of the full nonoscillatory main term are required before interpreting residual plots. A flat residual at the noise floor does not establish absence of a mathematical term.}

\glossaryentry{nome}{Nome}{\statusconvention}{A nome is the base \(q\) used in theta, elliptic, and modular products,
typically with \(0<|q|<1\).  Some literature uses \(q=e^{\pi i\tau}\), while
other sources use \(e^{2\pi i\tau}\).  This factor-of-two choice changes every
theta exponent and modular transformation, so the project declares it before
using theta notation.}

\glossaryentry{fabius-non-elementarity}{Non-elementarity of the Fabius function}{\statuslean}{The non-elementarity theorem states that the canonical bounded Fabius
function cannot be represented by the project's elementary-function class on
the open unit interval, and consequently not globally in the corresponding
sense.  The proof is qualitative: elementary expressions have a dense
analytic locus, whereas the Fabius function is nowhere analytic.  It does not
depend on failure to find a particular closed form.}

\glossaryentry{inverse-fabius-non-elementarity}{Non-elementarity of the inverse Fabius function}{\statuslean}{The inverse non-elementarity theorem applies the same analytic-locus
obstruction to the canonical inverse on \((0,1)\).  The result is strengthened
to the class obtained by adjoining permitted inverse branches of elementary
functions.  Endpoint singularity alone would not prove this statement; the
key input is nonanalyticity at every interior point.}

\glossaryentry{nonconstant-periodic-correction}{Nonconstant periodic correction}{\statusfrontier}{A nonconstant periodic correction is a periodic term with at least one nonzero Fourier mode. In the Fabius endpoint asymptotic, every relevant Gamma--zeta coefficient is explicit, and nonvanishing proves that the correction cannot be absorbed into a constant. Its numerical amplitude is extremely small because Gamma values decay rapidly on vertical lines, but asymptotic correctness depends on logical size relative to the stated remainder, not on visual prominence.}

\glossaryentry{noncrossing-partition}{Noncrossing partition}{\statusfrontier}{A partition of \(\{1,\ldots,n\}\) is noncrossing if no four indices
\(a<b<c<d\) have \(a,c\) in one block and \(b,d\) in another.  Noncrossing
partitions form a Catalan-counted lattice and index free moment--cumulant
relations.}

\glossaryentry{nondegenerate-saddle}{Nondegenerate saddle}{\statusclassical}{A saddle \(z_0\) is nondegenerate when \(\Phi''(z_0)\ne0\).  The local
exponent is then quadratic to first order, producing a Gaussian factor and a
half-power scale in the large parameter.  Degenerate saddles require different
normal forms such as Airy or higher catastrophe integrals.}

\glossaryentry{normal-convergence}{Normal convergence}{\statusclassical}{A series of functions converges normally on a domain if it converges absolutely and uniformly on every compact subset, usually by the Weierstrass M-test. Normal convergence of holomorphic functions permits termwise differentiation and yields a holomorphic sum. Infinite transform products and Fourier series of the smooth periodic correction are controlled in this manner. The term is stronger than pointwise convergence and is especially useful for complex-analytic families.}

\glossaryentry{normal-convergence-of-product}{Normal convergence of an infinite product}{\statusmixed}{An infinite product \(\prod(1+u_n(z))\) converges normally on a domain when
\(\sum|u_n(z)|\) converges uniformly on every compact subset after finitely
many exceptional factors.  The limit is holomorphic, and zeros arise from
factor zeros unless the limit collapses identically.  This controls sinc,
\(q\)-Pochhammer, theta, and reciprocal-Gamma products in the repository.}

\glossaryentry{q-pochhammer-normal-convergence}{Normal convergence of \(q\)-products}{\statuslean}{A family of analytic products converges normally on a domain when it converges
uniformly on every compact subset.  For \(|q|<1\), the estimate
\(\sum_j |a q^j|<\infty\) locally uniformly in \(a\) gives normal convergence
of \((a;q)_\infty\).  Normal convergence implies holomorphy of the limit and
permits termwise differentiation after the usual nonvanishing or logarithmic
localization hypotheses.}

\glossaryentry{normalized-reversion-coefficient}{Normalized reversion coefficient}{\statusfrontier}{For a delta series
\(f(z)=f_1z+f_2z^2/2!+\cdots\), a normalized reversion input is typically
\(f_{j+1}/((j+1)f_1)\).  This isolates the linear scale before Bell-polynomial
or partition formulas are applied.  The exact normalization is printed because
ordinary-series conventions use different factorial factors.}

\glossaryentry{normalized-reversion-input}{Normalized reversion input}{\statusconvention}{A normalized reversion input is
\[
 a_j^{\rm rev}=\frac{f_{j+1}}{(j+1)f_1},
 \qquad j\ge1,
\]
for an exponential-series forward map \(f\).  These quantities enter compact
Bell-polynomial inverse formulas.  They are neither ordinary Taylor
coefficients nor the inverse coefficients themselves.}

\glossaryentry{normalized-saddle-mass}{Normalized saddle mass}{\statusmixed}{The normalized saddle mass is what remains of a saddle integral after factoring out the exponential at the saddle, the leading amplitude, and the Gaussian normalization. It tends to one and has a formal expansion \(1+\sum_{j\ge1}a_j\lambda^{-j}\). This normalization isolates universal Gaussian moments from problem-specific phase jets. Its logarithm supplies the additive all-orders corrections \(A_j\) in the Fabius endpoint formula.}

\glossaryentry{normalized-sinc-functions}{Normalized sinc functions}{\statusconvention}{The radian sinc is \(\operatorname{sinc}_{\mathrm{rad}}z=\sin z/z\),
and the normalized sinc is
\(\operatorname{sinc}_{\pi}x=\sin(\pi x)/(\pi x)\), both with removable
value \(1\) at zero.  Their zeros occur at \(\pi\Z\setminus\{0\}\) and
\(\Z\setminus\{0\}\), respectively.  The notation catalogue forbids a bare
ambiguous sinc in normative formulas.}

\glossaryentry{normalized-volterra-operator}{Normalized Volterra operator}{\statuslean}{A normalized Volterra operator includes the Gamma or factorial factor that
makes its order parameter additive, such as
\[
 (I^\alpha f)(x)=\frac1{\Gamma(\alpha)}
 \int_a^x(x-t)^{\alpha-1}f(t)\,\dd t .
\]
This normalization gives a clean semigroup law and matches ordinary repeated
integration at positive integers.}

\glossaryentry{normative-notation-contract}{Normative notation contract}{\statusconvention}{The normative notation contract is the shared LaTeX command file together
with the mathematical catalogue that defines each command's semantics.  It
fixes function normalizations, Fourier conventions, branch choices, and
collision rules across active documents.  Changing a command's arity or
meaning requires a contract-version update and synchronized catalogue entries;
historical sources are handled by explicit exceptions rather than silent
reuse.}

\glossaryentry{notation-collision}{Notation collision}{\statusconvention}{A notation collision occurs when the same glyph can denote different objects
in nearby topics--for example \(F\) as the bounded Fabius function, a
numerator sequence, or a generic map; \(P_n\) as a Legendre polynomial, a
probability, or a Pad\'e denominator; and \(\mu\) as a measure, mean, or
M\"obius function.  The project resolves collisions with semantic commands,
superscripts, or mandatory local declarations rather than relying on context
alone.}

\glossaryentry{nowhere-analytic}{Nowhere analytic}{\statusclassical}{A smooth function is nowhere analytic on a set if it is not real analytic at any point of that set. This is stronger than having isolated nonanalytic points. The Fabius function is nowhere analytic on its nonconstant interval: dyadic points are dense, their formal Taylor series terminate, and if the function were analytic near any point, analytic continuation and overlap with such a dyadic neighborhood would force local polynomial behavior incompatible with strict shape and scaling. The exterior constant tails remain analytic away from the endpoints.}

\glossaryentry{nowhere-analytic-fabius-iterate}{Nowhere-analytic Fabius iterate}{\statusfrontier}{A positive Fabius iterate \(F^{\circ n}\) is asserted in the dedicated
frontier development to be nonanalytic at every point of \([0,1]\).
The proof architecture follows dyadic points through the chain rule and
Fa\`a di Bruno partitions, locating a finite first nonzero derivative that
cannot be reconciled with an analytic germ.  This status is distinct from the
separately formalized nonanalyticity of \(F\) itself.}

\glossaryentry{nowhere-analytic-inverse-iterate}{Nowhere-analytic inverse-Fabius iterate}{\statusfrontier}{The corresponding frontier theorem treats
\((F^{-1})^{\circ n}\) for every positive \(n\).  Monotone conjugacy and
dyadic-image structure transport singular jets through successive inverse
branches.  Care is required near endpoints, where ordinary derivative
formulas may diverge; the result concerns real analyticity, not mere failure of
a bounded first derivative.}

\glossaryentry{null-set}{Null set}{\statusclassical}{A null set is a set of measure zero. Altering an integrable function on a null set does not change its Lebesgue integral or its \(L^p\) equivalence class, but it can change pointwise properties such as continuity or exact functional equations. Formalizations that reason with continuous representatives cannot silently replace functions almost everywhere. The Fabius theory usually works with canonical continuous or smooth versions rather than equivalence classes.}

\glossaryentry{numerator-sequence}{Numerator sequence}{\statusclassical}{A numerator sequence records the integer numerators obtained after multiplying rational quantities by a prescribed denominator. The choice of denominator is part of the definition and can be designed to expose recurrence, parity, or divisibility. The Fabius directory contains several such sequences for moments, half moments, and dyadic endpoint values. Similar names in different papers need a notation crosswalk because their normalizations are not interchangeable.}

\glossaryentry{numerical-residual}{Numerical residual}{\statusfrontier}{A numerical residual is the discrepancy obtained by substituting a computed approximation into an equation or by subtracting a predicted approximation from trusted data. A small residual is evidence, not by itself a proof, because conditioning and cancellation matter. Exact rational dyadic values provide unusually strong reference data for testing the corrected Fabius asymptotic. Comparing residuals across scales reveals the periodic oscillation and the order of the first omitted saddle term.}

\section{O}\label{sec:letter-o}

\glossaryentry{odd-cosine-series}{Odd cosine series}{\statusclassical}{An odd cosine series is a cosine expansion containing only odd harmonics \(\cos((2m+1)\pi x)\). It arises when even-frequency Fourier coefficients vanish. For the up-function, zeros of the sinc product force this pattern in the chosen periodization. Despite the adjective odd, the resulting function is even in \(x\); odd refers to the harmonic indices, not parity of the function.}

\glossaryentry{odd-discrete-fourier-transform}{Odd discrete Fourier transform}{\statusfrontier}{An odd DFT restricts or twists the discrete Fourier transform to odd
residue classes, often modulo a power of two.  It isolates primitive dyadic
frequencies and avoids modes annihilated by even symmetry.  The exact node and
normalization convention is report-specific.}

\glossaryentry{odd-part-of-an-integer}{Odd part of an integer}{\statusclassical}{For nonzero \(n\in\Z\), the odd part is \(n/2^{v_2(n)}\), an odd integer
up to sign.  It retains the unit information after dyadic scale is removed.
Odd-part recurrences are often more stable than raw numerator recurrences and
make parity statements immediate.}

\glossaryentry{one-periodic-function}{One-periodic function}{\statusclassical}{A function \(P\) is one-periodic when \(P(u+1)=P(u)\) for all \(u\). It is determined by its values on any interval of length one and has Fourier frequencies in \(\Z\). The endpoint correction and every all-orders coefficient \(A_j\) in the Fabius logarithmic expansion are one-periodic functions of the lower-Lambert coordinate. Centering sets the mean of the leading correction to zero.}

\glossaryentry{operation-count}{Operation count}{\statusclassical}{An operation count estimates the number of arithmetic or logical steps performed by an algorithm as a function of input size. For integer inputs, bit length and Hamming weight are more informative than numerical magnitude. The fast dyadic Fabius evaluator makes one main recursive reduction per set bit and evaluates finite Taylor blocks with bounded-degree arithmetic. Exact bit complexity also depends on growth of numerators and denominators, so a scalar operation count is only one layer of analysis.}

\glossaryentry{optimal-truncation}{Optimal truncation}{\statusmixed}{For a divergent asymptotic series, optimal truncation stops near the least
term.  The resulting error is often exponentially smaller than any fixed
algebraic order.  The optimal index depends on the large parameter and on
coefficient growth; it is not obtained by summing the formal series to
infinity.}

\glossaryentry{order-of-vanishing}{Order of vanishing}{\statusclassical}{A function has a zero of order \(m\) at \(x_0\) if its first \(m-1\) derivatives vanish there and the \(m\)-th derivative does not. Infinite-order vanishing is flatness. Finite Thue--Morse generating polynomials have high-order zeros at \(z=1\), encoding Prouhet cancellation of low powers. At an interior dyadic point, the Fabius Taylor polynomial has exact highest degree determined by the first nonzero scaled integer value, not by a zero of the function itself.}

\glossaryentry{ordered-composition}{Ordered composition}{\statusclassical}{An ordered composition of \(n\) is a sequence of positive integers summing to \(n\). It indexes terms produced by expanding powers of a series or taking a formal logarithm. The explicit Fabius saddle mass coefficients sum over compositions of \(2j\), with additional subset choices selecting phase or amplitude contributions; the logarithmic coefficients sum over compositions of \(j\). Ordered indexing avoids symmetry factors but produces more terms than partition notation.}

\glossaryentry{ordered-composition-coefficient}{Ordered-composition coefficient}{\statusmixed}{An ordered-composition coefficient sums products over tuples of positive
integers with prescribed total.  It is the coefficient form of powers of a
series and provides a direct alternative to Bell-polynomial notation.
Ordered compositions are useful in formalization because every index range is
finite and elementary.}

\glossaryentry{ordinary-bell-polynomial}{Ordinary Bell polynomial}{\statusclassical}{An ordinary Bell polynomial uses ordinary rather than exponential
generating-function normalization.  Factorials therefore move between its
arguments and coefficients relative to the exponential Bell family.  The
project names the normalization explicitly whenever determinant or reversion
formulas could otherwise differ by products of factorials.}

\glossaryentry{ordinary-generating-function}{Ordinary generating function}{\statusclassical}{The ordinary generating function of \((a_n)\) is \(A(z)=\sum_{n\ge0}a_nz^n\). Its product corresponds to ordinary discrete convolution. The finite Thue--Morse generating polynomial factors digitwise, and iterated prefix sums correspond to division by powers of \(1-z\) in a formal setting. Ordinary and exponential generating functions occur side by side in the repository and encode different convolution laws.}

\glossaryentry{original-bounded-bridge}{Original--bounded Fabius bridge}{\statuslean}{The original--bounded bridge proves that the compact-support predicate for
\(\Up\), together with the appropriate right-tail condition, is equivalent
to the bounded \(\mathrm{IsFabius}\) predicate after folding.  It reconciles
the historical differential-difference definition with the modern CDF-style
normalization and makes uniqueness transferable in both directions.}

\glossaryentry{original-fabius-characterization}{Original compact-support characterization}{\statuslean}{The original characterization identifies the normalized compactly supported
solution of
\[
 u'(x)=2\bigl(u(2x+1)-u(2x-1)\bigr)
\]
together with the appropriate smoothness, support, and value at the origin.
The formal development proves existence and uniqueness and then connects this
solution to the folded bounded Fabius function.  The characterization is
global; a local differential identity alone is not sufficient for uniqueness.}

\glossaryentry{is-original-fabius}{Original Fabius predicate \(\mathrm{IsOriginalFabius}\)}{\statuslean}{The original compact-support predicate characterizes the Rvachev-style
function directly: smoothness, compact support, normalization, and the
differential-difference law are packaged without first mentioning the bounded
CDF.  Current bridge theorems convert between this predicate for
\(\Up\) and the bounded predicate \(\mathrm{IsFabius}\), with the required
tail condition stated separately.  It is a proposition about a candidate
function, not an additional canonical function.}

\glossaryentry{orthogonal-polynomial}{Orthogonal polynomial}{\statusclassical}{A sequence \((p_n)\) is orthogonal with respect to a measure \(\mu\) if \(\int p_np_m\,\dd\mu=0\) for \(m\ne n\), with \(\deg p_n=n\). Orthogonal polynomials satisfy three-term recurrences and provide spectral bases for approximation. Legendre polynomials are orthogonal for unit weight on \([-1,1]\) and are used to expand the up-function. Orthogonality makes coefficient extraction and least-squares optimality immediate.}

\glossaryentry{orthogonal-projection}{Orthogonal projection}{\statusclassical}{The orthogonal projection of a Hilbert-space vector \(f\) onto a closed subspace \(V\) is the unique \(p\in V\) with \(f-p\) orthogonal to \(V\). It minimizes \(\|f-p\|\). The degree-\(n\) Fourier-Legendre partial sum is the orthogonal projection onto polynomials of degree at most \(n\). The corresponding Pythagorean identity quantifies exactly how the squared error decreases when another Legendre coefficient is included.}

\glossaryentry{orthogonality}{Orthogonality}{\statusclassical}{Vectors or functions are orthogonal when their inner product is zero. In \(L^2[-1,1]\), \(\langle f,g\rangle=\int f\overline g\). Orthogonality eliminates cross terms in norm calculations and yields the Pythagorean identity. The Legendre basis separates the polynomial components of the up-function, while Fourier exponentials separate periodic modes of the endpoint correction.}

\glossaryentry{orthonormal-polynomial}{Orthonormal polynomial}{\statusclassical}{An orthonormal polynomial \(p_n\) satisfies
\(\langle p_n,p_m\rangle=\delta_{nm}\) and has a declared positive leading
coefficient.  It differs from the monic polynomial by division by its norm.
Legendre polynomials in the repository are usually classically normalized by
\(P_n(1)=1\), so a factor \(\sqrt{(2n+1)/2}\) is required for orthonormality on
\([-1,1]\).}

\glossaryentry{oscillatory-correction}{Oscillatory correction}{\statusclassical}{An oscillatory correction varies rather than tending monotonically to a constant, often through sine, cosine, or a periodic function of a slowly varying phase. Its mean may be zero while its pointwise effect persists indefinitely. The Fabius correction is smooth and one-periodic in \(\lambda(x)\), so sequences approaching the endpoint through different phases have different normalized limits. The term is asymptotically indispensable despite its tiny amplitude.}

\section{P}\label{sec:letter-p}

\glossaryentry{p-recursive-sequence}{\(P\)-recursive sequence}{\statusfrontier}{A sequence \((a_n)\) is \(P\)-recursive when
\[
 \sum_{j=0}^{r}p_j(n)a_{n+j}=0
\]
for polynomial coefficients \(p_j\), eventually in \(n\).  Its ordinary
generating function is \(D\)-finite in characteristic zero.  Failure of a
guessed recurrence at finite order is evidence, not a proof of
non-\(P\)-recursiveness.}

\glossaryentry{pade-approximant}{Padé approximant}{\statusmixed}{A Padé approximant \([m/n]\) to a series is a rational function \(P_m/Q_n\)
whose Taylor expansion agrees through the maximal prescribed order, with
\(Q_n(0)=1\).  For Stieltjes transforms, diagonal and near-diagonal Padé
approximants inherit positivity and interlacing properties.  The project uses
moment determinants and continued fractions to organize certified rational
approximations.}

\glossaryentry{paley-wiener-theorem}{Paley--Wiener theorem}{\statusclassical}{Paley--Wiener theorems relate compact support to holomorphic growth of a
Fourier transform.  For a smooth function supported in \([-R,R]\), the
transform extends to an entire function bounded by rapidly decaying powers on
horizontal directions times \(e^{2\pi R|\operatorname{Im}z|}\) under the
\(2\pi\)-frequency convention.  Conversely, suitable entire functions of this
exponential type arise from compactly supported data.  This framework explains
the simultaneous entireness and rapid real-axis decay of the up transform.}

\glossaryentry{parity}{Parity}{\statusclassical}{Parity distinguishes even integers from odd integers. It is additive modulo two and governs the Thue--Morse sign through the parity of binary digit sum. Parity assertions about normalized Fabius moment numerators are stronger than nonvanishing and determine exact \(2\)-adic valuations. In finite \(q\)-binomial sums, parity also controls sign patterns and even-odd term separation.}

\glossaryentry{parity-power-series}{Parity-power series}{\statusclassical}{The parity-power series is a signed global Fabius representation whose summands are organized by binary parity data and powers arising from finite Taylor blocks. A corrected version in the repository includes the precise zeroth term and domain convention needed for global validity. It converges absolutely. The name is local to the development and should be understood through its displayed summands rather than as a standard class of power series.}

\glossaryentry{parseval-identity}{Parseval identity}{\statusclassical}{Parseval's identity equates the \(L^2\) norm of a function with the sum of squared coefficients in a complete orthonormal basis. For Fourier series it is \(\|f\|_2^2=\sum|c_k|^2\); for Legendre series the norm factors reflect polynomial normalization. It justifies energy estimates and coefficient square summability. In the Fabius context it complements the stronger absolute-uniform convergence results for the up-function expansion.}

\glossaryentry{partial-exponential-bell-polynomial}{Partial exponential Bell polynomial}{\statusclassical}{The partial exponential Bell polynomial \(B_{n,k}(x_1,x_2,\ldots)\)
sums over set partitions of total size \(n\) with \(k\) blocks, weighting a
block of size \(j\) by \(x_j\).  It is the coefficient engine in
Fa\`a di Bruno's formula and in fixed-block moment--cumulant transforms.}

\glossaryentry{partition-defect}{Partition defect}{\statusfrontier}{A partition defect measures how far an index tuple falls short of a target
weight, length, or admissibility condition.  Organizing a recurrence by defect
produces shells whose first nonzero layer often determines the leading
coefficient.  The term is used in frontier coefficient and derivative
expansions rather than as a single universal invariant.}

\glossaryentry{partition-lattice}{Partition lattice}{\statusclassical}{The partition lattice \(\Pi_n\) orders set partitions of
\(\{1,\ldots,n\}\) by refinement.  Its minimum is the discrete singleton
partition and its maximum is the one-block partition.  Intervals factor over
blocks, and its M\"obius function inverts the moment--cumulant relation.}

\glossaryentry{partition-of-unity}{Partition of unity}{\statusclassical}{A partition of unity is a family of nonnegative functions whose sum equals one, often locally finite and subordinate to a covering. Translates of the up-function satisfy exact partition identities because the Fourier transform vanishes at nonzero integers under an appropriate normalization. Scaled versions yield further partitions. These identities are analytically analogous to cardinal spline partitions and are useful in approximation and sampling.}

\glossaryentry{prime-power-pascal-row}{Pascal row modulo a prime power}{\statusfrontier}{A Pascal row modulo \(p^r\) is the sequence
\((\binom nk\bmod p^r)_{0\le k\le n}\).  Its zero pattern is controlled by
carry counts and higher congruence data, not Lucas's theorem alone.
Histograms and recursive block decompositions connect these rows to automatic
sequences and Fabius denominator questions.}

\glossaryentry{peano-kernel}{Peano kernel}{\statusmixed}{If a linear functional annihilates polynomials below degree \(m\), Peano's
theorem may represent its error as
\[
 L(f)=\int K(t)f^{(m)}(t)\,\dd t.
\]
The kernel \(K\) determines sign, sharp constants, and extremizers.  Repeated
integration and spline approximation in the project use positive Peano kernels
to obtain certified convergence rather than only formal Taylor cancellation.}

\glossaryentry{periodic-correction}{Periodic correction}{\statusmixed}{A periodic correction is a periodic function added to an otherwise smooth asymptotic main term to account for discrete-scale effects. In the Fabius endpoint theory it arises from the nonzero imaginary poles of \((1-2^{-s})^{-1}\) in Mellin space. Its centered form has zero mean and explicit Gamma--zeta Fourier coefficients. Uniqueness means that no different continuous one-periodic function can satisfy the same sharp remainder estimate.}

\glossaryentry{periodic-fourier-coefficient}{Periodic-correction Fourier coefficient}{\statuslean}{For a one-periodic correction \(\Psi\), its Fourier coefficients are
\[
 \widehat\Psi(k)=\int_0^1\Psi(u)e^{-2\pi iku}\,\dd u.
\]
The Fabius Mellin analysis expresses nonzero modes through explicit
Gamma--zeta values.  Centering sets the \(k=0\) coefficient to zero and moves
its mean into the nonoscillatory constant.}

\glossaryentry{periodic-correction-psi}{Periodic correction Psi}{\statusmixed}{The symbol \(\Psi\) denotes the centered smooth one-periodic function in the corrected Fabius endpoint expansion. Its mean over one period is zero, and its nonzero Fourier coefficients are explicit Gamma--zeta values. It is evaluated at the exact lower-Lambert coordinate \(\lambda(x)\). The centering convention moves the zero Fourier mode into the endpoint constant, making \(\Psi\) unique under the stated sharp remainder.}

\glossaryentry{periodic-mean}{Periodic mean}{\statusclassical}{The mean of a one-periodic integrable function \(P\) is \(\int_0^1P(u)\,\dd u\). It is the zero-frequency Fourier coefficient. Subtracting it centers the function; adding it to the constant term keeps the full expression unchanged. Repository modules compute the mean of the raw Mellin periodic term so that the endpoint constant and centered oscillation have a canonical normalization.}

\glossaryentry{periodicity}{Periodicity}{\statusconvention}{A function has period \(T>0\) if \(f(x+T)=f(x)\). The smallest positive period, when it exists, is the fundamental period. Periodicity can arise from translation symmetry or from multiplicative scaling after a logarithmic change of variables. The Fabius asymptotic uses period one in the lower-Lambert coordinate, corresponding to the base-two structure of the original equation.}

\glossaryentry{periodization}{Periodization}{\statusclassical}{The periodization of \(f\) with period \(h\) is
\(\sum_{k\in\Z}f(x+kh)\) when the sum is meaningful.  Compact support makes
it locally finite.  Poisson summation expresses its Fourier coefficients as
samples of \(\widehat f\).  Integer zeros of the up transform make several
periodizations exactly constant.}

\glossaryentry{perron-frobenius-operator}{Perron--Frobenius transfer operator}{\statusfrontier}{A transfer operator \(\mathcal P_T\) is dual to composition:
\[
 \int(\mathcal P_Tf)g\,\mathrm d\mu
 =\int f(g\circ T)\,\mathrm d\mu .
\]
When \(T\) is nonsingular it transports densities through inverse branches and
Jacobian factors.  It is not simply composition by \(T^{-1}\), especially when
\(T\) is noninjective.}

\glossaryentry{perturbation-expansion}{Perturbation expansion}{\statusclassical}{A perturbation expansion expresses a solution or integral as a formal or asymptotic series in a small parameter. One expands the phase, amplitude, saddle location, or recurrence around a limiting model and solves coefficient equations successively. The inverse-saddle parameter \(1/\lambda\) is the perturbation parameter in the Fabius all-orders analysis. Formal coefficient algebra and analytic remainder estimates are separate components of a valid perturbation proof.}

\glossaryentry{phase-function}{Phase function}{\statusclassical}{The phase function is the exponent controlling an oscillatory or Laplace-type integral, commonly \(\Phi(z)\) in \(\int A(z)\e^{\lambda\Phi(z)}\,\dd z\). Saddles satisfy \(\Phi'(z)=0\); the real part of \(\Phi\) determines concentration and decay. In the Fabius Bromwich integral the phase incorporates the linear inversion term and the logarithm of the negative-Laplace product. Its exact saddle equation leads to the lower Lambert coordinate.}

\glossaryentry{physicists-hermite-polynomial}{Physicists' Hermite polynomial}{\statusconvention}{The physicists' Hermite polynomial is
\(H_n^{\rm phys}(x)=(-1)^ne^{x^2}\frac{\mathrm d^n}{\mathrm dx^n}e^{-x^2}\).
It differs from the probabilists' normalization by
\(H_n^{\rm phys}(x)=2^{n/2}\operatorname{He}_n(\sqrt2\,x)\).
The two families must not share an undecorated \(H_n\) in probability
asymptotics.}

\glossaryentry{piecewise-polynomial}{Piecewise polynomial}{\statusclassical}{A piecewise-polynomial function agrees with a polynomial on each interval of a finite or locally finite partition. Splines impose smoothness conditions at the knots. Finite convolutions of uniform densities are piecewise polynomial, and their degrees increase with the number of factors. The centered Fabius approximants are therefore exactly evaluable on each dyadic cell, which is crucial for primitive-recursive computation.}

\glossaryentry{plancherel-theorem}{Plancherel theorem}{\statusclassical}{Plancherel's theorem extends the Fourier transform to a unitary map on
\(L^2\), up to the selected normalization.  It states equality of spatial and
frequency energies and turns convolution into multiplication where defined.
The theorem underlies exact up-function energy identities and \(L^2\)
approximation estimates.}

\glossaryentry{pochhammer-symbol}{Pochhammer symbol}{\statusclassical}{The ordinary Pochhammer symbol \((a)_n=a(a+1)\cdots(a+n-1)\) is the rising factorial, with \((a)_0=1\). It appears in hypergeometric series and Gamma ratios: \((a)_n=\Gamma(a+n)/\Gamma(a)\). The \(q\)-Pochhammer symbol is a different multiplicative object \((a;q)_n=\prod_{j=0}^{n-1}(1-aq^j)\). The Fabius \(q\)-calculus mainly uses the latter, so the base parameter should always be displayed.}

\glossaryentry{poincare-asymptotic-expansion}{Poincaré asymptotic expansion}{\statusclassical}{For an asymptotic scale \(\phi_0\gg\phi_1\gg\cdots\),
\[
 f(x)\sim\sum_{n\ge0}a_n\phi_n(x)
\]
means
\(f-\sum_{n<N}a_n\phi_n=o(\phi_{N-1})\) for every fixed \(N\).
The statement is termwise and finite at each order.  It neither promises
convergence nor specifies behavior when the truncation index grows with
\(x\).}

\glossaryentry{pointwise-convergence}{Pointwise convergence}{\statusclassical}{A sequence \(f_n\) converges pointwise to \(f\) if \(f_n(x)\to f(x)\) for each fixed \(x\). The index after which an approximation is accurate may depend on \(x\), unlike uniform convergence. Pointwise convergence alone does not preserve continuity or justify interchange with integration. The repository's stronger spline and Legendre results supply uniform convergence, while some discrete-limit formulas are first established pointwise and then strengthened through bounded rows or compactness.}

\glossaryentry{pointwise-equality}{Pointwise equality}{\statusclassical}{Two functions are pointwise equal when \(f(x)=g(x)\) for every input in the
declared domain.  This is stronger than equality almost everywhere and differs
from equality of distributions of random variables.  Functional equations,
support statements for continuous representatives, and exact dyadic
evaluations in the Fabius development are pointwise unless a theorem
explicitly uses a measure-theoretic equality.}

\glossaryentry{poisson-summation}{Poisson summation}{\statusclassical}{Poisson summation relates samples of a function to samples of its Fourier transform: under suitable decay, \(\sum_{n\in\Z}f(x+n)=\sum_{k\in\Z}\widehat f(k)\e^{2\pi\ii kx}\). If all nonzero integer transform samples vanish, the periodized sum is constant. Applying this to the up-function's sinc product yields exact partitions of unity and sampling identities. Smooth compact support supplies the regularity needed for classical justification.}

\glossaryentry{pole}{Pole}{\statusclassical}{A pole of order \(m\) at \(z_0\) is an isolated singularity where \((z-z_0)^mf(z)\) extends holomorphically and nonvanishingly. The Laurent series has finitely many negative powers. Mellin contour shifts collect residues at poles; a real pole generates polynomial-logarithmic main terms, while imaginary translates generate periodic modes. The dyadic factor \((1-2^{-s})^{-1}\) creates the relevant arithmetic lattice of poles.}

\glossaryentry{polygamma-function}{Polygamma function}{\statusmixed}{The order-\(m\) polygamma function is
\(\psi^{(m)}(z)=\frac{\mathrm d^{m+1}}{\mathrm dz^{m+1}}\log\Gamma(z)\).
For \(m\ge1\) and away from poles,
\[
 \psi^{(m)}(z)=(-1)^{m+1}m!\sum_{k\ge0}(z+k)^{-m-1}.
\]
These functions provide the derivative jets of gamma ratios and their
asymptotic coefficients.}

\glossaryentry{polynomial-approximation}{Polynomial approximation}{\statusclassical}{Polynomial approximation seeks polynomials close to a function in a chosen norm. Weierstrass' theorem gives uniform approximation for continuous functions on compact intervals, while orthogonal projections give least-squares approximants. The up-function's Fourier-Legendre series provides explicit polynomial approximants with exact coefficients and proven absolute uniform convergence. This is distinct from the finite spline approximation, whose pieces are polynomial but whose global approximant is not a single polynomial.}

\glossaryentry{polynomial-cell}{Polynomial cell}{\statuslean}{A polynomial cell is a connected interval between consecutive knots on which
a spline or finite sinc-prefix inverse transform is represented by one fixed
polynomial.  Cell boundaries depend on the dilation and translate parameters.
Local synthesis arguments choose atoms so the target interval stays within
specified cells, avoiding hidden changes of formula.}

\glossaryentry{polynomial-commutator}{Polynomial commutator with integration}{\statuslean}{For multiplication by \(x\) and an integration operator \(I\), the
commutator \(xI-Ix\) is another integral operator with an extra kernel factor.
Iterating this relation moves polynomials across Volterra or fractional
integrals and yields finite binomial identities.  Such commutators organize
moment recurrences and polynomial reproduction formulas.}

\glossaryentry{polynomial-deconvolution}{Polynomial deconvolution by the up-law}{\statuslean}{Polynomial deconvolution assigns to a polynomial \(P\) the unique
polynomial \(D(P)\) satisfying
\[
 \int_{\R}D(P)(x-t)\Up(t)\,\dd t=P(x).
\]
Moment Appell polynomials make this map explicit.  The formal operator is a
linear, injective, degree-preserving map with computable rational coefficients
for rational input.}

\glossaryentry{polynomial-diagonal-stabilization}{Polynomial diagonal stabilization}{\statusfrontier}{Polynomial diagonal stabilization is the phenomenon that, for fixed offset
from a natural boundary, entries of an iterated Thue--Morse summation array
eventually agree with a polynomial in the iteration order.  Degree,
factorization, and denominator patterns reflect the binomial prefix kernel.
A rigorous statement specifies the exact index region and whether the formula
is an identity or an experimentally inferred continuation.}

\glossaryentry{ring-polynomial-identity}{Polynomial identity over a ring}{\statusclassical}{A polynomial identity over a ring means equality in \(R[x]\), hence equality
of all coefficients.  It is stronger than agreement at finitely many points
unless an interpolation theorem applies, and it remains valid under every ring
homomorphism.  Gaussian, cyclotomic, \(q\)-Catalan, and quantum-plane results
often live first in a polynomial or noncommutative ring before evaluation at a
numeric base.}

\glossaryentry{polynomial-logarithmic-transseries}{Polynomial--logarithmic transseries}{\statusmixed}{A polynomial--logarithmic transseries at infinity is an expansion
\[
 X+\sum_{n\ge0}X^{-n}p_n(\log X),
\]
or a related Laurent form, with polynomial coefficient blocks.  Lambert
\(W\), Wright's omega function, inverse gamma-type functions, and logarithmic
series reversions live in this class.  Multiplication, division, composition,
and reversion are coefficientwise finite under a triangular order.}

\glossaryentry{polynomial-reproduction}{Polynomial reproduction}{\statusclassical}{An approximation or summation scheme reproduces polynomials up to degree \(m\) if it returns every such polynomial exactly. Finite-difference annihilation and partition-of-unity identities are dual manifestations of reproduction. Gaussian rows on geometric nodes reproduce the leading-coefficient functional after annihilating lower degrees. In spline and Toeplitz limits, exact mass and low-moment conditions control the approximation order.}

\glossaryentry{polynomial-reproduction-order}{Polynomial reproduction order}{\statuslean}{The reproduction order of a mesh-generator pair is the largest \(d\) for
which every polynomial of degree at most \(d\) is reconstructed exactly from
the prescribed translate sum.  In the current Rvachev comb classification, a
nonzero integer mesh \(M\) reproduces exactly through degree \(v_2(M)\) in the
universal shifted sense.}

\glossaryentry{binomial-type-sequence}{Polynomial sequence of binomial type}{\statusclassical}{A polynomial sequence \(p_n(x)\) is of binomial type when
\(p_n(x+y)=\sum_k\binom nkp_k(x)p_{n-k}(y)\).
Its exponential generating function is \(e^{x f(t)}\) for a delta series
\(f\).  Falling factorials and ordinary powers are basic examples.}

\glossaryentry{positive-compositional-iterate}{Positive compositional iterate}{\statusmixed}{For a self-map \(f\), its positive compositional iterates are
\(f^{\circ n}=f\circ\cdots\circ f\) for \(n\ge1\); \(f^{\circ0}\) is the
identity.  In the Fabius program the bounded function and its inverse are
iterated on the invariant interval \([0,1]\).  Composition exponents must not
be confused with pointwise powers.}

\glossaryentry{positive-definite-kernel}{Positive-definite kernel}{\statusmixed}{A Hermitian kernel \(K\) is positive definite when
\[
 \sum_{i,j}\overline c_i c_jK(x_i,x_j)\ge0
\]
for every finite choice of points and scalars.  Such a kernel determines a
reproducing-kernel Hilbert space.  Translation-invariant positive definiteness
is characterized by Bochner's theorem through a positive spectral measure.}

\glossaryentry{positive-definite-moment-functional}{Positive-definite moment functional}{\statusmixed}{A moment functional \(\mathcal M\) is positive definite when
\(\mathcal M(p\overline p)>0\) for every nonzero polynomial \(p\).
This ensures nonsingular Hankel matrices and a unique monic orthogonal
polynomial of every degree.  For the positive Rvachev density, the analytic
integral supplies positivity; the rational moment calculus supplies the same
Gram entries exactly.}

\glossaryentry{positive-half-line-fourier-energy}{Positive-half-line Fourier energy}{\statuslean}{Because the up transform is even in modulus, its full real-axis squared
mass is twice the integral over \((0,\infty)\).  The formal
Fabius-square-energy constant equals this positive-half-line mass under the
fixed Fourier convention.  A rescaled sinc-product integral gives an
equivalent representation with its Jacobian displayed.}

\glossaryentry{positive-part-kernel}{Positive-part kernel}{\statusclassical}{The positive-part power
\((x-t)_+^\alpha=\max(x-t,0)^\alpha\) packages a one-sided integration kernel
on the whole line.  It converts interval limits into support of the kernel and
is convenient for translation, convolution, splines, and fractional
integration.  Endpoint exponents and the convention at zero must be stated
when \(\alpha\le0\).}

\glossaryentry{positive-rational-gap-sequence}{Positive rational gap lower-bound sequence}{\statuslean}{A positive rational gap sequence supplies computable numbers
\(g(n)>0\) such that points separated by a prescribed mesh have Fabius values
separated by at least \(g(n)\).  Positivity guarantees injective numerical
decoding, while rationality allows exact verification.  The current inverse
computability module treats such a sequence as an input hypothesis rather than
claiming a universal closed form for it.}

\glossaryentry{post-render-source-union}{Post-render source union}{\statusconvention}{A post-render source union is material merged into a LaTeX source after the
last retained PDF was built.  The phrase warns that page counts describe the
PDF, whereas the source contains additional sections or corrections.  It is
especially relevant to rapidly growing frontier, integration, representation,
and \(q\)-series syntheses.  Such material is active documentation even before
a matching PDF is committed.}

\glossaryentry{power-series}{Power series}{\statusclassical}{A power series \(\sum a_n(z-z_0)^n\) converges inside a disk determined by its radius of convergence and defines a holomorphic function there. Within the disk it may be differentiated and integrated termwise. A formal power series has the same algebra but no convergence claim. The Fabius function's Taylor series at dyadic points is finite yet fails to represent the function locally, whereas its transform power series is entire and analytically valid.}

\glossaryentry{power-sum-symmetric-polynomial}{Power-sum symmetric polynomial}{\statusclassical}{The power sum is \(p_m(x_1,\ldots,x_r)=\sum_i x_i^m\).
Power sums linearize disjoint unions of alphabets and are related to elementary
and complete homogeneous functions by Newton identities.  Cumulants of
weighted independent sums are power sums of the weights times component
cumulants.}

\glossaryentry{powerset-expansion}{Powerset expansion}{\statusclassical}{A powerset expansion sums over all subsets of a finite index set.  Products
of terms \(1+u_j\), iterated finite differences, and independent Bernoulli
choices all expand this way.  In formal proofs a finite-set formulation makes
duplicate indices, empty cases, and sign exponents explicit.}

\glossaryentry{prefix-polynomial}{Prefix polynomial}{\statusmixed}{A prefix polynomial packages a finite signed prefix as
\(P_N(z)=\sum_{n<N}\varepsilon_n z^n\), or packages repeated prefix values as a
polynomial in the endpoint \(N\).  Complete blocks factor by
\(\prod_{j<m}(1-z^{2^j})\).  Derivatives at \(z=1\) encode Prouhet moments and
orders of cancellation.}

\glossaryentry{prefix-sum-boundary-convention}{Prefix-sum boundary convention}{\statusconvention}{An exclusive prefix uses \(\sum_{j<n}a_j\) and vanishes at \(n=0\);
an inclusive prefix uses \(\sum_{j\le n}a_j\).  Iteration shifts binomial
kernels and diagonal formulas, so the convention must be fixed before
comparing repeated-summation results.  The current project states boundary
values explicitly rather than relying on informal array indexing.}

\glossaryentry{primitive-recursive-fold}{Primitive-recursive fold}{\statuslean}{A primitive-recursive fold evaluates a finite list or bounded range by
repeatedly applying an update operation whose iteration count is itself
primitive recursive.  The Fabius computability certificates use such folds to
evaluate exact spline or dyadic data.  A large unrestricted bound may make the
certificate inefficient while still proving membership in the intended
recursion class.}

\glossaryentry{primitive-recursive-function}{Primitive-recursive function}{\statusclassical}{Primitive-recursive functions are built from zero, successor, and projections using composition and primitive recursion. They form a robust class of total computable functions with explicit structural termination, though they do not include every computable function. The repository expresses the finite-spline evaluator, signed-natural arithmetic, precision selection, and continuity modulus in primitive-recursive terms. This supplies a constructive proof of Fabius computability stronger than an informal numerical algorithm.}

\glossaryentry{primitive-root-block}{Primitive-root block}{\statuslean}{A primitive-root block groups factors whose exponents form a complete residue
system modulo the order \(d\) of a root of unity.  At a primitive \(d\)-th root,
products over a complete nonzero block collapse to explicit constants, while
the zero residue records vanishing order.  Blocking isolates the singular
factor before taking limits and avoids invalid termwise substitution into
\(0/0\) quotients.}

\glossaryentry{principal-branch}{Principal branch}{\statusconvention}{A principal branch is a conventional single-valued choice from a multivalued complex function, such as the principal logarithm or \(W_0\) branch of Lambert \(W\). Branch choices determine cuts, arguments, and continuity domains. The Fabius saddle requires the nonprincipal real branch \(W_{-1}\); using the principal branch would solve a different root of the same equation. Complex shifts and logarithms in transform formulas likewise require consistent branch declarations.}

\glossaryentry{principal-lambert-phase}{Principal-Lambert phase}{\statusmixed}{A principal-Lambert phase uses the branch \(W_0\) in an inverse or saddle
coordinate.  It is appropriate when the implicit solution remains near zero
or grows on the principal sheet.  It must not replace the \(W_{-1}\) phase at
a small positive Fabius endpoint, where the relevant solution tends to
\(-\infty\).}

\glossaryentry{probabilistic-representation}{Probabilistic representation}{\statusclassical}{A probabilistic representation identifies a deterministic function or sequence with a probability, expectation, density, or law. It can make positivity, monotonicity, normalization, and convolution structure transparent. The Fabius function is the CDF of an infinite dyadic weighted sum of independent uniforms; the up-function is the density of a centered affine version. Transform products and finite spline approximants follow directly from this representation.}

\glossaryentry{probabilists-hermite-polynomial}{Probabilists' Hermite polynomial}{\statusmixed}{The probabilists' Hermite polynomial is
\[
 \operatorname{He}_n(x)=(-1)^ne^{x^2/2}
 \frac{\mathrm d^n}{\mathrm dx^n}e^{-x^2/2}.
\]
It satisfies
\((\phi(x)\operatorname{He}_n(x))'=-\phi(x)\operatorname{He}_{n+1}(x)\)
for the standard normal density \(\phi\).  This is the normalization used in
Edgeworth and Cornish--Fisher formulas.}

\glossaryentry{probability-density}{Probability density}{\statusclassical}{A probability density \(f\) is a nonnegative integrable function with total integral one; probabilities are \(\Prob(X\in A)=\int_Af(x)\,\dd x\). Its cumulative distribution is \(F(x)=\int_{-\infty}^{x}f(t)\,\dd t\). Under the chosen affine normalization, the up-function is the smooth compactly supported density associated with the Fabius distribution. Strict positivity in the interior gives strict monotonicity of the CDF.}

\glossaryentry{probability-measure}{Probability measure}{\statusclassical}{A probability measure is a countably additive measure of total mass one. It provides the abstract framework for random variables and distributions. Pushforwards, products, and convolutions construct new laws. The infinite product probability space carrying independent uniform variables and the pushforward under their weighted sum gives a rigorous construction of the Fabius measure.}

\glossaryentry{product-measure}{Product measure}{\statusclassical}{Given probability spaces \((\Omega_n,\mu_n)\), their product measure is the canonical measure on the Cartesian product under which coordinate random variables are independent and have laws \(\mu_n\). Infinite product measures are constructed from consistent finite-dimensional distributions. The Fabius random series can be defined on a product of unit intervals, with each coordinate carrying a uniform law. Absolute summability of weights makes the sum measurable and convergent.}

\glossaryentry{profile-generating-function}{Profile generating function}{\statusfrontier}{For scale multiplicities \(c_h\), the profile generating function is
\(C(z)=\sum_{h\ge0}c_hz^h\).  Evaluation at \(z=4^{-n}\) encodes normalized
even cumulants, while substitution at \(2^{-s}\) appears in the spectral
Dirichlet transform.  Its coefficient nonnegativity reflects geometric
multiplicity.}

\glossaryentry{proof-backed-exposition}{Proof-backed exposition}{\statuslean}{A proof-backed exposition is project prose whose mathematical claims are
required to have counterparts in the compiled Lean development.  Such a
document may reorganize proofs and suppress implementation detail, but it must
not introduce unproved analytic strengthening.  In this repository the primary
Fabius--Rvachev article has this role, unlike the frontier volumes, which
explicitly mix proved, derived, conjectural, and exploratory material.}

\glossaryentry{proof-status-boundary}{Proof-status boundary}{\statusconvention}{A proof-status boundary separates Lean-proved results, classical background,
frontier derivations, conjectures, refuted claims, and historical quotations.
Crossing that boundary requires a proof or an explicit change of label.
Numerical evidence, a formal power-series identity, or a manuscript derivation
does not by itself become a theorem in the compiled library.  The present
glossary uses status badges to preserve this distinction.}

\glossaryentry{prouhet-cancellation}{Prouhet cancellation}{\statusclassical}{Prouhet cancellation is the vanishing of signed power sums generated by the Thue--Morse sequence. For a complete block of length \(2^m\), \(\sum_{n=0}^{2^m-1}\varepsilon_nn^r=0\) for \(0\le r<m\). It follows from the order-\(m\) zero at \(z=1\) of \(\prod_{j<m}(1-z^{2^j})\). This annihilation removes lower-degree polynomial contributions in dyadic Fabius sums and underlies the exact leading-coefficient extraction.}

\glossaryentry{prouhet-partition}{Prouhet partition}{\statusclassical}{A Prouhet partition divides a finite set of integers into classes whose low
power sums agree.  The binary Thue--Morse partition of
\(\{0,\ldots,2^m-1\}\) into even and odd digit-sum classes has equal moments
through degree \(m-1\).  Root-of-unity digit characters give higher-base
generalizations.}

\glossaryentry{prouhet-tarry-escott-problem}{Prouhet--Tarry--Escott problem}{\statusclassical}{The Prouhet--Tarry--Escott problem asks for disjoint multisets with equal
power sums through a prescribed degree.  Splitting a complete binary block by
the Thue--Morse parity gives Prouhet's classical solution: signed power sums
vanish below the block level.  These cancellations drive finite-difference,
prefix-sum, and exact dyadic identities in the project.}

\glossaryentry{provenance-ledger}{Provenance ledger}{\statusconvention}{A provenance ledger records where a formula entered the corpus, how its
notation was normalized, whether it was corrected, and which formal
declarations support it.  It is especially valuable when multiple frontier
drafts are merged or when a historical paper uses incompatible symbols.
Provenance is not proof, but it makes proof obligations and source corrections
auditable.}

\glossaryentry{public-declaration}{Public declaration}{\statusconvention}{A public declaration is a definition, theorem, structure, instance, or other
named Lean object exported without private visibility through the relevant
facade.  The repository audit counts declarations after taking a union, so a
theorem imported through several modules is counted once.  A declaration count
is a coverage diagnostic, not a claim that every declaration deserves a
separate headword.}

\glossaryentry{public-facade}{Public facade}{\statusconvention}{A public facade is an import module that re-exports the declarations intended
for downstream use.  The root import \texttt{FabiusFunction} is the complete
facade for this development, while smaller thematic imports reduce dependency
footprint.  Corpus census numbers refer to the union visible through the
facade, not to every private helper definition or local theorem in every source
file.}

\glossaryentry{puiseux-log-series}{Puiseux--logarithmic series}{\statusmixed}{A Puiseux--logarithmic series permits
\(t^{r/s}p_{r}(\log t)\).  It combines algebraic ramification with logarithmic
corrections and is the natural local class when an inverse has a multiple
critical point and the forward function also contains logarithms.  A fixed
denominator \(s\) is part of the formal ambient ring.}

\glossaryentry{puiseux-series}{Puiseux series}{\statusmixed}{A Puiseux series is a Laurent series in a fractional power
\(t^{1/s}\), so its exponents lie in \(s^{-1}\Z\).  It represents local
algebraic inverse branches near a ramification point after a branch of the
fractional power is chosen.  The Lambert \(W\) branch point and cyclotomic
critical points naturally produce square-root or higher Puiseux expansions.}

\glossaryentry{pushforward-measure}{Pushforward measure}{\statusclassical}{For a measurable map \(T:X\to Y\) and a measure \(\mu\) on \(X\), the pushforward \(T_\#\mu\) is defined by \(T_\#\mu(A)=\mu(T^{-1}(A))\). It is the law of \(T(X)\) when \(X\) has law \(\mu\). The Fabius measure is the pushforward of an infinite product of uniform measures under the weighted-sum map. Affine pushforwards relate the bounded distribution and centered up-density.}

\glossaryentry{pythagorean-identity}{Pythagorean identity}{\statusclassical}{In an inner-product space, if \(u\perp v\) then \(\|u+v\|^2=\|u\|^2+\|v\|^2\). For an orthogonal projection \(p\) of \(f\), one has \(\|f-q\|^2=\|f-p\|^2+\|p-q\|^2\) for every \(q\) in the subspace. This proves least-squares optimality and quantifies the error reduction of Fourier-Legendre truncation for the up-function.}

\section{Q}\label{sec:letter-q}

\glossaryentry{q-bernstein-polynomial}{\(q\)-Bernstein polynomial}{\statusfrontier}{A \(q\)-Bernstein basis deforms the classical Bernstein basis by replacing
binomial coefficients and powers with Gaussian and \(q\)-shifted factors.
Several inequivalent definitions occur.  In geometric-comb approximation it
preserves a chosen positive partition of unity and converges to the classical
basis as \(q\to1\), subject to the declared normalization.}

\glossaryentry{q-beta-function}{\(q\)-beta function}{\statusmixed}{The \(q\)-beta function is
\[
 B_q(x,y)=\frac{\Gamma_q(x)\Gamma_q(y)}{\Gamma_q(x+y)}.
\]
For positive real arguments and \(0<q<1\), it has a Jackson-integral
representation after choosing the matching normalization.  It is symmetric in
\(x,y\), and its logarithmic derivatives are differences of \(q\)-digamma
functions.  Pole and branch conventions are inherited from \(\Gamma_q\).}

\glossaryentry{q-beta-integral}{\(q\)-beta integral}{\statuslean}{A \(q\)-beta integral evaluates a Jackson integral of a product of powers and
\(q\)-shifted factors in terms of \(q\)-gamma functions.  It is the geometric
comb analogue of Euler's beta integral.  The formal theorem states positivity,
convergence, and endpoint conditions explicitly; the symbol
\(\mathrm d_qt\) cannot be replaced by ordinary Lebesgue measure.}

\glossaryentry{q-binomial-coefficient}{q-binomial coefficient}{\statusclassical}{The \(q\)-binomial coefficient is the Gaussian coefficient \(\qbinom nkq=(q;q)_n/((q;q)_k(q;q)_{n-k})\). It satisfies a \(q\)-Pascal recurrence and counts subspaces of a finite vector space when \(q\) is a prime power, while algebraically it is a polynomial in \(q\). The Fabius directory uses the analytic value \(q=1/2\), where a normalized row gives interpolation and Toeplitz weights. The parameter is a scalar base, not a probability.}

\glossaryentry{q-binomial-scalar-extension}{q-binomial scalar extension}{\statusconvention}{Scalar extension transports a polynomial or finite-sum identity proved over \(\Q\) or \(\R\) into a larger field such as \(\C\). Since Gaussian binomial coefficients at \(q=1/2\) are rational, their identities can be mapped through the canonical embedding. The repository separates this algebraic step from analytic convergence, allowing the complex-shift \(q\)-binomial Taylor identity to reuse an exact real polynomial theorem.}

\glossaryentry{q-binomial-theorem}{q-binomial theorem}{\statusclassical}{The \(q\)-binomial theorem comprises a finite polynomial identity
and an infinite analytic identity.  One finite form is
\[
 \prod_{j=0}^{n-1}(1+zq^j)
 =\sum_{k=0}^{n}q^{\binom{k}{2}}\qbinom{n}{k}{q}z^k .
\]
For \(|q|<1\) and \(|z|<1\), the infinite form expresses a quotient of
\(q\)-Pochhammer products as a basic hypergeometric series.  The repository
formalizes terminating identities, scalar and polynomial forms, reciprocal
Euler specializations, and the infinite theorem with explicit summability
hypotheses.  Finite algebraic statements must not be silently upgraded to
analytic convergence statements.}

\glossaryentry{q-calculus}{q-calculus}{\statusclassical}{\(q\)-calculus is a deformation of ordinary calculus in which additive increments are replaced by multiplicative scaling by \(q\). Its basic objects include \(q\)-integers, \(q\)-factorials, \(q\)-differences, Gaussian binomial coefficients, and \(q\)-Pochhammer symbols. The Fabius problem has a natural half-base \(q=1/2\) because every differential scaling step doubles the argument. Here \(q\)-calculus is finite and algebraic rather than a replacement for the ordinary derivative defining \(F\).}

\glossaryentry{q-catalan-number}{\(q\)-Catalan number}{\statuslean}{One common normalization is
\[
 C_n(q)=\frac{1}{[n+1]_q}\qbinom{2n}{n}{q}.
\]
The displayed quotient is in fact a polynomial, but this requires a
divisibility theorem.  Different \(q\)-Catalan families exist, so the
normalization must be stated.  The project uses the Gaussian quotient form and
keeps its cyclotomic divisibility proof separate from combinatorial
interpretations.}

\glossaryentry{q-digamma-function}{\(q\)-digamma function}{\statusmixed}{The \(q\)-digamma function is
\(\psi_q(z)=\frac{\mathrm d}{\mathrm dz}\log\Gamma_q(z)\).
For \(0<q<1\),
\[
 \psi_q(z)=-\log(1-q)+(\log q)
 \sum_{n\ge0}\frac{q^{n+z}}{1-q^{n+z}},
\]
away from poles.  It controls monotonicity, convexity, and parameter inversion
for \(\Gamma_q\), and tends to the ordinary digamma function as \(q\uparrow1\)
under standard fixed-\(z\) hypotheses.}

\glossaryentry{q-discrete-limit}{q-discrete limit}{\statusmixed}{The complex \(q\)-discrete limit is a sequence of finite Gaussian-binomial combinations of complex-shift spline approximants that converges to the signed global Fabius extension. After reindexing, each approximant is a Toeplitz row with total mass one and uniformly bounded variation, sampling only sufficiently advanced splines. A finite row can depend on the complex shift \(Q\); the limit is independent of fixed \(Q\). This limit is not termwise identical to the global \(q\)-series even though both converge to the same function.}

\glossaryentry{q-factorial}{q-factorial}{\statusclassical}{The \(q\)-integer is \([n]_q=(1-q^n)/(1-q)\), and the \(q\)-factorial is \([n]_q!=\prod_{j=1}^{n}[j]_q\). It is related to the \(q\)-Pochhammer product by \([n]_q!=(q;q)_n/(1-q)^n\). Gaussian binomial coefficients are ratios of \(q\)-factorials. At \(q=1/2\), clearing powers of two converts these products into odd Mersenne products.}

\glossaryentry{q-finite-difference}{q-finite difference}{\statusclassical}{A \(q\)-finite difference compares values at multiplicatively related points, often through \(D_qf(x)=(f(x)-f(qx))/((1-q)x)\), or through a finite Gaussian-weighted node sum. The repository mainly uses the latter leading-coefficient functional on geometric nodes. Like an ordinary \(n\)-th difference, it annihilates polynomials of degree below \(n\), but the weights and nodes encode powers of two.}

\glossaryentry{q-gamma-function}{\(q\)-gamma function}{\statusmixed}{For \(0<q<1\),
\[
 \Gamma_q(z)=(1-q)^{1-z}\frac{(q;q)_\infty}{(q^z;q)_\infty},
 \qquad q^z=e^{z\log q}.
\]
It satisfies \(\Gamma_q(1)=1\) and
\(\Gamma_q(z+1)=[z]_q\Gamma_q(z)\).
Its meromorphic pole lattice comes from the denominator.  A separate
reciprocal-base normalization is required for \(q>1\); the two conventions
must not be mixed inside an inverse or asymptotic formula.}

\glossaryentry{q-gauss-summation}{\(q\)-Gauss summation}{\statuslean}{Under \(|c/(ab)|<1\) and denominator nonvanishing,
\[
 {}_2\phi_1\!\left[\begin{matrix}a,b\\c\end{matrix};q,\frac{c}{ab}\right]
 =\frac{(c/a,c/b;q)_\infty}{(c,c/(ab);q)_\infty}.
\]
It is the basic-hypergeometric analogue of Gauss's ordinary hypergeometric
summation.  The infinite products encode the complete value; truncating them
without an error estimate changes the theorem into an approximation.}

\glossaryentry{q-integer}{q-integer}{\statusclassical}{The \(q\)-integer \([n]_q=1+q+\cdots+q^{n-1}=(1-q^n)/(1-q)\) deforms the ordinary integer \(n\), recovered as \(q\to1\). Products of \(q\)-integers form \(q\)-factorials. Although the Fabius formulas are often written directly with \((q;q)_n\) and Gaussian coefficients, \(q\)-integers explain the parallel with ordinary binomial algebra and the appearance of factors \(2^j-1\) at half base.}

\glossaryentry{q-kernel}{\(q\)-kernel of a sequence}{\statusclassical}{The \(q\)-kernel replaces powers of two by powers of an integer base
\(q\ge2\).  Finiteness characterizes \(q\)-automatic sequences, and finite
module generation characterizes \(q\)-regular sequences.  Kernel recursion is
the automata-theoretic counterpart of digitwise decomposition and Mahler
functional equations.}

\glossaryentry{q-lucas-theorem}{\(q\)-Lucas theorem}{\statuslean}{At a primitive \(d\)-th root of unity \(\zeta\), writing
\(n=n_1d+n_0\) and \(k=k_1d+k_0\) with \(0\le n_0,k_0<d\), a \(q\)-Lucas
formula reduces
\[
 \qbinom{n}{k}{\zeta}
\]
to an ordinary binomial factor involving \(n_1,k_1\) times a small Gaussian
coefficient \(\qbinom{n_0}{k_0}{\zeta}\).  The theorem is a root-of-unity digit
law and is distinct from the ordinary Lucas theorem modulo a prime.}

\glossaryentry{q-multinomial-coefficient}{\(q\)-multinomial coefficient}{\statuslean}{For nonnegative \(k_1,\ldots,k_r\) with sum \(n\),
\[
 \genfrac{[}{]}{0pt}{}{n}{k_1,\ldots,k_r}_{q}
 =\frac{(q;q)_n}{(q;q)_{k_1}\cdots(q;q)_{k_r}}.
\]
It is a polynomial with nonnegative integer coefficients and counts flags or
weighted words in standard interpretations.  Iterating Gaussian binomial
factorizations proves its polynomiality and yields quantum multinomial
expansions for suitably \(q\)-commuting variables.}

\glossaryentry{q-pascal-summation}{\(q\)-Pascal summation}{\statuslean}{A \(q\)-Pascal summation is a finite identity obtained by summing a
Gaussian Pascal recurrence, often a \(q\)-weighted analogue of the
hockey-stick identity.  Typical forms collapse a triangular Gaussian sum to a
single Gaussian coefficient.  Since several exponent placements are
equivalent after reindexing, the exact power of \(q\) is included in the
formula rather than hidden in the name.}

\glossaryentry{q-pfaff-saalschutz-summation}{\(q\)-Pfaff--Saalschütz summation}{\statuslean}{A terminating balanced \({}_3\phi_2\) sum has the form
\[
 {}_3\phi_2\!\left[\begin{matrix}a,b,q^{-N}\\
 c,abq^{1-N}/c\end{matrix};q,q\right]
 =\frac{(c/a,c/b;q)_N}{(c,c/(ab);q)_N}.
\]
The balance relation among the lower parameters is essential.  The project
proves the finite identity by a telescoping or inductive certificate suitable
for formal verification.}

\glossaryentry{q-pochhammer-dissection}{\(q\)-Pochhammer dissection}{\statuslean}{For an integer \(m\ge1\), splitting indices by residue class gives
\[
 (a;q)_\infty=\prod_{r=0}^{m-1}(aq^r;q^m)_\infty
\]
when \(|q|<1\), with the analogous finite formula after the terminal lengths
are distributed among residue classes.  Dissection is the product-side
counterpart of digit decomposition.  It supports root-of-unity analysis,
theta identities, and multiscale factorizations of Fabius-related products.}

\glossaryentry{q-pochhammer-logarithmic-derivative}{\(q\)-Pochhammer logarithmic derivative}{\statuslean}{Away from the zeros of \((a;q)_\infty\),
\[
 \frac{\partial}{\partial a}\log(a;q)_\infty
 =-\sum_{j\ge0}\frac{q^j}{1-aq^j}.
\]
The right side is a Lambert series and converges locally normally on compact
sets avoiding the zero lattice.  Higher derivatives are obtained by
differentiating the summands.  The formula is meromorphic, not an everywhere
holomorphic identity, because each product zero becomes a pole.}

\glossaryentry{q-pochhammer-order-derivative}{\(q\)-Pochhammer order derivative}{\statusmixed}{A continuous-order extension replaces the integer length \(n\) by a complex
parameter, commonly through
\[
 (a;q)_z=\frac{(a;q)_\infty}{(aq^z;q)_\infty},
 \qquad |q|<1,
\]
after fixing a branch of \(q^z\).  Differentiation with respect to \(z\) then
produces a logarithmic Lambert series.  This object agrees with the finite
product at nonnegative integers but is branch-dependent for complex \(q\);
the repository records the branch and pole lattice explicitly.}

\glossaryentry{q-pochhammer-reversal}{\(q\)-Pochhammer reversal}{\statuslean}{For nonzero \(a,q\),
\[
 (a;q)_n=(-a)^nq^{n(n-1)/2}
 \bigl(a^{-1}q^{1-n};q\bigr)_n.
\]
This reverses the order of the factors and exposes reciprocal-base symmetry.
Specializing \(a=q^{-N}\) makes termination explicit in basic hypergeometric
series.  The powers of \(-a\) and \(q\) are essential; omitting either changes
both the degree and the zero divisor.}

\glossaryentry{q-pochhammer-shift-law}{\(q\)-Pochhammer shift law}{\statuslean}{The elementary shift identity is
\[
 (a;q)_{m+n}=(a;q)_m(aq^m;q)_n.
\]
For the infinite product this becomes
\((a;q)_\infty=(a;q)_m(aq^m;q)_\infty\) when \(|q|<1\).
It is the multiplicative analogue of cutting a series into a finite head and a
tail.  Many quotient identities, contiguous relations, and numerical
evaluators are obtained by applying the law before cancellation.}

\glossaryentry{q-pochhammer-symbol}{q-Pochhammer symbol}{\statusclassical}{For a base \(q\), parameter \(a\), and natural length \(n\), the
finite \(q\)-Pochhammer symbol is
\[
 (a;q)_n=\prod_{j=0}^{n-1}(1-aq^j),\qquad (a;q)_0=1.
\]
When \(|q|<1\), the infinite symbol \((a;q)_\infty\) is the normally
convergent product over \(j\ge0\).  The current tree treats finite and infinite
forms, dissection, base reversal, analytic dependence on \(a\), logarithmic
derivatives, and branch-dependent complex order.  Every generalized-order
formula states its logarithm branch and denominator nonvanishing conditions;
the notation is never interpreted by informal substitution alone.}

\glossaryentry{q-product}{q-product}{\statusclassical}{A \(q\)-product is a finite or infinite product with geometrically progressing factors, typically a \(q\)-Pochhammer symbol. Its zeros and ratios encode \(q\)-recurrences and basic hypergeometric coefficients. The Fabius half-base products link geometric interpolation denominators to Mersenne factors. They should be distinguished from the independent sinc product of the Fourier transform, although both reflect the same dyadic scale hierarchy.}

\glossaryentry{q-richardson-extrapolation}{\(q\)-Richardson extrapolation}{\statusfrontier}{A \(q\)-Richardson scheme cancels leading error powers in data sampled at
geometrically shrinking scales \(h,qh,q^2h,\ldots\).  Its weights are
Lagrange weights at the target scale \(0\) and therefore have explicit
\(q\)-product forms.  Stability depends on the base, order, and error model;
exact cancellation does not by itself imply numerical robustness.}

\glossaryentry{q-series}{q-series}{\statusclassical}{A \(q\)-series is a series whose coefficients or exponents are organized by powers of a base \(q\), often using \(q\)-Pochhammer symbols. It may be formal, terminating, or analytically convergent. The global Fabius \(q\)-binomial series is an absolutely convergent sum of finite Taylor blocks expressed through Gaussian rows and a fixed complex shift. The convergence comes from binary reduction estimates, not merely from formal \(q\)-series algebra.}

\glossaryentry{q-shifted-factorial}{\(q\)-shifted factorial}{\statusclassical}{\(q\)-shifted factorial is a standard synonym for the
\(q\)-Pochhammer symbol.  The word ``factorial'' emphasizes the recurrence
\((a;q)_{n+1}=(1-aq^n)(a;q)_n\); it does not mean the special product
\([n]_q!\).  The glossary keeps a separate entry because both phrases occur in
the literature and their collision is a common source of normalization errors.}

\glossaryentry{q-taylor-identity}{q-Taylor identity}{\statusmixed}{The repository's \(q\)-Taylor identity states that a finite Gaussian-binomial expression with arbitrary common complex shift \(Q\) equals a real Taylor block embedded in \(\C\). Translation invariance follows because the order-\(n\) geometric finite difference annihilates every lower-degree term generated by expanding the shift. This identity converts local derivative data into the finite summands of the global complex \(q\)-series.}

\glossaryentry{q-to-one-limit}{\(q\to1\) limit}{\statusmixed}{A \(q\to1\) limit recovers a classical object only after its scaling is fixed:
\([n]_q\to n\), \(\qbinom{n}{k}{q}\to\binom nk\), and a suitably scaled
\(q\)-exponential tends to \(e^z\).  Infinite products often require
renormalization or logarithmic asymptotics.  The direction of approach and
uniformity in other parameters are part of the theorem.}

\glossaryentry{q-vandermonde-summation}{\(q\)-Vandermonde convolution}{\statusmixed}{The \(q\)-Vandermonde identity convolves two Gaussian rows with a quadratic
or bilinear power of \(q\) to obtain a larger Gaussian coefficient.  One form
is
\[
 \qbinom{m+n}{r}{q}=\sum_k q^{(m-k)(r-k)}
 \qbinom{m}{k}{q}\qbinom{n}{r-k}{q}.
\]
Different but equivalent exponents correspond to reversing one row.  It is the
geometric-node analogue of the ordinary Vandermonde convolution.}

\glossaryentry{quadratic-logarithmic-law}{Quadratic logarithmic law}{\statusclassical}{A quadratic logarithmic law states that the leading behavior of \(\log F(x)\) is a negative constant times \((\log(1/x))^2\). It is a coarse consequence of the sharp endpoint expansion and ignores nested logs and periodic fluctuations. Such a law proves superpolynomial flatness and gives the correct broad decay class. It is insufficient for asymptotic equivalence or accurate numerical evaluation because lower terms in the exponent remain large or bounded but nonconvergent.}

\glossaryentry{quadrature-weight}{Quadrature weight}{\statusclassical}{In a rule
\(\int f\,\mathrm d\mu\approx\sum_jw_jf(x_j)\), the numbers \(w_j\) are
quadrature weights.  Exactness through a polynomial degree gives moment
equations \(\sum_jw_jx_j^k=m_k\).  Positive weights yield stability and a
probabilistic interpretation; signed extrapolation weights can be exact but
ill-conditioned.}

\glossaryentry{quantile-coupling}{Quantile coupling}{\statusclassical}{For real laws with generalized inverses \(Q_\mu,Q_\nu\), taking one uniform
\(U\) and setting \(X=Q_\mu(U)\), \(Y=Q_\nu(U)\) gives the monotone quantile
coupling.  It is optimal for one-dimensional Wasserstein costs.  In the Fabius
setting the ordinary inverse on \([0,1]\) supplies the quantile exactly.}

\glossaryentry{quantile-transport}{Quantile transport}{\statuslean}{Quantile transport maps a reference probability \(u\in(0,1)\) through an
inverse CDF to produce a target law.  Identities for the geometric uniform CDF
therefore induce identities for its inverse after monotonicity and range have
been proved.  Error transfer requires lower density or gap information,
because small CDF error can be magnified near flat endpoints.}

\glossaryentry{quantum-binomial-theorem}{Quantum binomial theorem}{\statuslean}{The quantum binomial theorem is the expansion of a power of a sum of
\(q\)-commuting variables by Gaussian coefficients.  Unlike the finite
commutative \(q\)-binomial product identity, it is an identity in a
noncommutative algebra.  The project formalizes the commutation relation and
then proves the expansion by induction using a \(q\)-Pascal recurrence.}

\glossaryentry{quantum-multinomial-theorem}{Quantum multinomial theorem}{\statuslean}{For a consistently ordered family of pairwise \(q\)-commuting variables, a
power of their sum expands through \(q\)-multinomial coefficients.  The precise
power of \(q\) depends on the chosen ordering and commutation convention.
Formal statements fix both, preventing a classical multinomial formula from
being imported into a noncommutative algebra without the required weights.}

\glossaryentry{quantum-plane}{Quantum plane}{\statuslean}{The quantum plane is the algebra generated by \(x,y\) with
\(yx=qxy\).  In this noncommutative setting,
\[
 (x+y)^n=\sum_{k=0}^{n}\qbinom{n}{k}{q} x^{\,n-k}y^k
\]
for the displayed commutation convention.  Reversing the relation or monomial
order replaces \(q\) by \(q^{-1}\) and introduces powers of \(q\).  The relation
therefore belongs to the statement, not merely to the proof.}

\glossaryentry{quotient-of-exponentials-approximation}{Quotient-of-exponentials approximation}{\statusclassical}{A quotient-of-exponentials approximation replaces a complicated product, recurrence, or sum by a ratio of exponential expressions suggested by leading logarithms. Such an ansatz can capture a coarse scale while missing a linear saddle residue, a periodic factor, or a normalization constant. Repository modules named for quotient-exponential mismatch formalize why certain attractive informal proxies cannot achieve the claimed error. Substitution into the exact stationary equation is used to expose the residual term.}

\section{R}\label{sec:letter-r}

\glossaryentry{r-stirling-number}{\(r\)-Stirling number}{\statusfrontier}{An \(r\)-Stirling number of the second kind counts set partitions in which
\(r\) distinguished elements lie in pairwise distinct blocks.  Boundary
conditions depend on whether the total set size includes those distinguished
elements.  The notation catalogue requires the convention to be stated
because several shifted definitions occur in the literature.}

\glossaryentry{rademacher-sign}{Rademacher sign}{\statusclassical}{The \(j\)-th Rademacher function alternates between \(+1\) and \(-1\) on
dyadic intervals of length \(2^{-j}\), equivalently reading one binary digit
of the argument.  Products of Rademacher functions form the Walsh system.
They are continuous analogues of digit characters and connect Thue--Morse
parity with dyadic harmonic analysis.}

\glossaryentry{rademacher-sine-product}{Rademacher sine-product identity}{\statusfrontier}{A Rademacher sine-product identity rewrites signs of dyadic sine or sinc
factors in terms of binary digit functions.  Magnitudes and signs are separated
so that the real Fourier lobes become a Walsh-like product.  Such identities
support exact sign classification and phase-aware extrapolation.}

\glossaryentry{radius-asymptotic}{Radius asymptotic}{\statusclassical}{A radius asymptotic describes the size of the neighborhood around a saddle on which a local expansion is used, often as a function of the large parameter. The radius must grow enough to capture Gaussian mass but remain small enough that higher phase terms are controlled. Optimizing this tradeoff is part of a quantitative saddle proof. The Fabius modules derive admissible central radii and show that tail and Taylor remainders meet the target inverse-\(\lambda\) order.}

\glossaryentry{radon-nikodym-derivative}{Radon--Nikodym derivative}{\statusclassical}{If a \(\sigma\)-finite measure \(\nu\) is absolutely continuous with
respect to \(\mu\), the Radon--Nikodym theorem gives a measurable density
\(\dd\nu/\dd\mu\), unique almost everywhere, such that
\(\nu(B)=\int_B(\dd\nu/\dd\mu)\,\dd\mu\).  Canonical continuous densities in
the Fabius development are selected representatives of this a.e. class.}

\glossaryentry{radon-transform}{Radon transform}{\statusfrontier}{The Radon transform integrates a multivariate function over affine
hyperplanes.  For a probability density it gives the densities of linear
projections.  Projecting a product or infinite convolution of segment laws
reduces multivariate Rvachev constructions to one-dimensional direction-profile
sinc products.}

\glossaryentry{ramification-index}{Ramification index}{\statusmixed}{If
\(f(x)-f(x_0)=c(x-x_0)^m+\cdots\) with \(c\ne0\), then \(m\) is the
ramification index of the local map.  The inverse has \(m\) Puiseux branches
with leading exponent \(1/m\).  The case \(m=1\) is an ordinary analytic
inverse; \(m=2\) gives a square-root branch point.}

\glossaryentry{random-series}{Random series}{\statusclassical}{A random series is a series of random variables. Its convergence can be studied almost surely, in probability, or in \(L^p\). The Fabius random series has bounded independent terms with geometrically summable weights, so it converges absolutely for every outcome. Its distributional self-similarity under removal of the first summand leads to the fixed-point and differential equations.}

\glossaryentry{rapid-decay}{Rapid decay}{\statusclassical}{A function \(g(t)\) has rapid, or Schwartz, decay if \(|t|^Ng(t)\to0\) for every \(N\), usually together with analogous bounds for derivatives. The Fourier transform of a compactly supported \(C^\infty\) function decays rapidly along the real axis. This property ensures absolute convergence of inversion integrals and Fourier coefficient series. The explicit sinc product supplies a direct route to rapid-decay estimates for the up-function transform.}

\glossaryentry{rarefaction}{Rarefaction of a sequence}{\statusfrontier}{Rarefaction extracts terms in one or more residue classes, such as
\(a(qn+r)\), or sums a sequence over those classes.  For automatic sequences
the rarefied subsequences belong to the kernel and satisfy a finite linear
system.  Thue--Morse rarefaction exposes spectral and arithmetic structure
hidden in the full sequence.}

\glossaryentry{rate-distortion-function}{Rate--distortion function}{\statusfrontier}{For a source \(X\) and distortion \(d\), the rate--distortion function is
\[
 R(D)=\inf I(X;\widehat X)
\]
over couplings with expected distortion at most \(D\).  It describes the
optimal information rate, not the error of one prescribed spline or prefix
quantizer.  Compact-support geometric sources yield structured candidate
couplings but still require an optimization theorem.}

\glossaryentry{rate-of-convergence}{Rate of convergence}{\statusclassical}{A rate of convergence quantifies how fast an error tends to zero, for example geometrically, polynomially, or as an inverse logarithm. A certified rate provides a function mapping tolerance to a sufficient index. Fabius fixed-point iterates converge geometrically, centered splines have a dyadic uniform rate, Fourier-Legendre coefficients decay rapidly, and the endpoint saddle expansion has inverse-Lambert orders. These rates serve different algorithms and should not be compared without matching cost models.}

\glossaryentry{ratio-set}{Ratio set}{\statusfrontier}{A ratio set is a finite family of small monomials used to witness that one
transseries monomial is asymptotically smaller than another.  Products of the
ratios generate a grid containing all relevant quotients.  It gives an
effective certificate for dominance and remainder control in grid-based
transseries calculations.}

\glossaryentry{rational-legendre-coefficient}{Rational Legendre coefficient}{\statuslean}{A rational Legendre coefficient is obtained by applying the exact rational
Rvachev moment functional to the rational polynomial \(P_n\):
\[
 c_n^{\Q}=\frac{2n+1}{2}\mathcal M_{\rm Rv}(P_n).
\]
After casting to \(\R\), it agrees with the analytic integral coefficient under
the established moment bridge.  This two-stage construction keeps exact
arithmetic independent of integration theory.}

\glossaryentry{rational-recurrence}{Rational recurrence}{\statusclassical}{A rational recurrence determines each term using rational operations on earlier terms. Triangular moment recurrences are typical. They prove rationality inductively but may obscure integrality and denominator growth. The Fabius arithmetic converts selected rational recurrences into division-free integer recurrences by multiplying through canonical factors, allowing parity and valuation proofs unavailable from the raw recurrence alone.}

\glossaryentry{rationality}{Rationality}{\statusclassical}{A quantity is rational when it belongs to \(\Q\). Rationality at all dyadic arguments is a remarkable property of the Fabius function, despite its definition through an infinite random series or transform product. It follows from terminating Taylor data and explicit finite Thue--Morse/moment formulas. Rationality alone does not identify the reduced denominator; the valuation and common-denominator theorems refine it.}

\glossaryentry{real-analytic-function}{Real-analytic function}{\statusclassical}{A real function is analytic at \(x_0\) if it agrees with a convergent power series on some neighborhood of \(x_0\). Real analyticity is stronger than \(C^\infty\) smoothness and entails local determination by derivatives. The Fabius function is analytic on the interiors of its constant tails but nowhere analytic on its transition interval; the up-function is analytic where it is identically zero outside its support and nowhere analytic inside. Endpoint flatness does not extend analyticity across the boundary.}

\glossaryentry{q-reciprocal-base-duality}{Reciprocal-base \(q\)-duality}{\statusmixed}{Reciprocal-base duality relates a formula at \(q\) to one at \(q^{-1}\) after
extracting a monomial power and sometimes a sign.  Gaussian palindromy and
\(q\)-Pochhammer reversal are canonical examples.  For probability laws, an
additional affine normalization may be needed.  The duality is algebraic for
finite products and analytic only on domains where the infinite objects are
defined.}

\glossaryentry{reciprocal-gamma-function}{Reciprocal gamma function}{\statusmixed}{The reciprocal gamma function \(1/\Gamma(z)\) is entire and has simple zeros
at \(0,-1,-2,\ldots\).  Its Weierstrass product and Taylor coefficients encode
Euler's constant and zeta values.  Entire-function treatment avoids carrying a
meromorphic denominator in products and derivative towers.}

\glossaryentry{reciprocal-power-table}{Reciprocal-power table}{\statusmixed}{A reciprocal-power table is a list of exact values, normalized numerators, or asymptotic data for \(F(2^{-n})\). These samples sit at the natural dyadic endpoint scales and are linked directly to half moments. The table is useful for discovering parity patterns, checking denominator formulas, and validating the periodic saddle asymptotic. Exact rational entries avoid the catastrophic underflow that affects direct floating-point evaluation for large \(n\).}

\glossaryentry{recurrence-relation}{Recurrence relation}{\statusclassical}{A recurrence relation defines terms of a sequence or family from earlier terms or smaller parameters. It is triangular when the new term appears with an invertible coefficient and only lower indices occur elsewhere. Recurrences can be linear or nonlinear, homogeneous or inhomogeneous, scalar or polynomial-valued. The Fabius directory uses them for moments, numerators, Bernoulli coefficients, saddle jets, Lambert coefficients, and exact dyadic evaluation.}

\glossaryentry{recursive-fast-name}{Recursive fast name}{\statusmixed}{A recursive fast name is a computable sequence of rational approximations whose error decreases at a prescribed exponential rate. Recursive means there is an effective algorithm for every precision, while fast refers to the convergence guarantee. The signed-natural dyadic representation in the repository is a concrete fast name. Composition with effective arithmetic and the finite-spline evaluator yields names for Fabius values.}

\glossaryentry{recursive-modulus}{Recursive modulus}{\statusclassical}{A recursive modulus is a computable rate function for continuity, convergence, or approximation. Given a desired output accuracy, it returns a sufficient input accuracy or truncation level. The Fabius computability theorem provides both a uniform-continuity modulus from the Lipschitz constant and an approximation modulus from the centered spline error. Together they make the qualitative algorithm uniform over real inputs.}

\glossaryentry{reduced-dyadic-level}{Reduced dyadic level}{\statuslean}{An interior dyadic rational is reduced at positive level \(n\) when it is
written \(m/2^n\) with \(0<m<2^n\) and \(m\) odd.  Oddness prevents the same
point from belonging to a coarser positive denominator level.  Exact
derivative-order, sign, denominator, and recursion statements use the reduced
level rather than an arbitrary presentation.}

\glossaryentry{reference-tail}{Reference tail}{\statusfrontier}{A reference tail is a simpler majorant or comparison integral used to bound the remote part of a contour integral. Its decay is chosen to dominate the true integrand after the central saddle region, while remaining explicitly integrable. Fabius saddle modules separate reference-tail estimates from local phase algebra, allowing the proof to show that omitted contour portions are below every retained inverse-saddle order.}

\glossaryentry{reference-weight}{Reference weight}{\statusfrontier}{A reference weight is a normalized model density, usually Gaussian, against which the local saddle integrand is compared. Moments of the reference weight generate universal double-factorial coefficients. The true normalized integrand is expanded as the reference weight times phase and amplitude corrections. Controlling the ratio on the central region and the tails produces both the leading mass and all higher saddle coefficients.}

\glossaryentry{refinement-equation}{Refinement equation}{\statusclassical}{A refinement equation expresses a function in terms of scaled and translated copies of itself, typically \(f(x)=\sum_kc_kf(2x-k)\). It is fundamental in wavelets, splines, and subdivision schemes. The up-function satisfies a differential/refinement identity derived from the Fabius equation; its Fourier transform then satisfies an infinite product relation. Support and normalization select the canonical compactly supported solution.}

\glossaryentry{refinement-operator}{Refinement operator}{\statusclassical}{A refinement operator maps a function to a combination or integral of its dilates and translates. Fixed points solve the associated refinement equation. Depending on the norm and normalization, the operator may be a contraction or preserve mass and positivity. Iterating the Fabius refinement operator produces increasingly fine probability or spline approximants and gives a constructive existence proof.}

\glossaryentry{reflection-symmetry}{Reflection symmetry}{\statusclassical}{Reflection symmetry about \(c\) means \(f(c+x)=f(c-x)\) for an even-type function, or \(F(c+x)=1-F(c-x)\) for a symmetric CDF. The bounded Fabius function satisfies \(F(1-x)=1-F(x)\), while the up-function is even about zero. Symmetry halves many integrals, kills odd moments or Legendre coefficients, and transfers endpoint and convexity results from one side to the other.}

\glossaryentry{regular-sequence}{Regular sequence}{\statusclassical}{A \(q\)-regular sequence has a \(q\)-kernel whose elements generate a
finitely generated module over the coefficient ring.  It need not take
finitely many values.  Repeated sums, polynomially weighted Thue--Morse
sequences, and many digit-counting functions naturally fall into this class,
which is stable under several algebraic operations.}

\glossaryentry{relative-entropy}{Relative entropy}{\statusfrontier}{For probability densities \(f,g\),
\[
 D(f\|g)=\int f\log(f/g)
\]
with the usual absolute-continuity convention.  It is nonnegative and vanishes
only for equality almost everywhere.  Comparisons between finite sinc-prefix
splines and the limiting up-density must control zeros before this integral is
asserted finite.}

\glossaryentry{relative-error}{Relative error}{\statusclassical}{The relative error of an approximation \(g\) to a nonzero quantity \(f\) is \(|g-f|/|f|\). For extremely small Fabius endpoint values, a tiny absolute error can still be a huge relative error. Logarithmic asymptotics are often analyzed first because direct values underflow, but exponentiating a bounded missing term causes an order-one relative oscillation. Exact rational benchmarks make relative-error assessment reliable.}

\glossaryentry{relative-saddle-mass}{Relative saddle mass}{\statusmixed}{Relative saddle mass is the ratio of the exact localized saddle integral to its leading Gaussian approximation. It tends to one and admits an inverse-saddle expansion. Expressing corrections relatively avoids repeatedly carrying the dominant exponential and square-root prefactors. The repository's mass-all-orders modules compute this ratio first and then take its formal logarithm to obtain the endpoint expansion for \(\log F\).}

\glossaryentry{relative-tail-logarithm}{Relative tail logarithm}{\statusfrontier}{For an infinite product \(\Phi=\Phi_NT_N\), the relative-tail logarithm is
\(\log T_N\) on a zero-free region with a specified logarithm branch.  Expanding
it in powers of the small tail scales yields multiplicative extrapolation
corrections.  Factoring out \(\Phi_N\) preserves the zeros already present in
the finite product.}

\glossaryentry{remainder-estimate}{Remainder estimate}{\statusclassical}{A remainder estimate bounds the error after truncating a series, Taylor expansion, contour localization, or asymptotic formula. A formal coefficient identity becomes an analytic theorem only when its remainder is controlled uniformly in the claimed regime. The Fabius development has separate modules for Taylor remainders, saddle central errors, minor arcs, tail flatness, and inverse-Lambert transfer. Their combination certifies each stated \(\bigO\)-term.}

\glossaryentry{removable-singularity}{Removable singularity}{\statusclassical}{A removable singularity is an isolated point where a function is initially undefined or apparently singular but has a finite limit and extends holomorphically or continuously. The sinc function \(\sin z/z\) has the removable value \(1\) at \(z=0\). Regularized Bose kernels likewise subtract a principal singular term to produce a smooth remainder, though the original integral may still require a finite part. Explicitly assigning removable values avoids domain defects in formal statements.}

\glossaryentry{removable-value-convention}{Removable-value convention}{\statusconvention}{A removable-value convention fills a point where a displayed quotient has
both numerator and denominator zero by its limiting analytic value.  Examples
include \(\sinc(0)=1\), Jackson quotients at special parameters, and polynomial
quotients constructed after cyclotomic cancellation.  The value must be
declared or proved removable before substitution; writing \(0/0\) and
informally cancelling at the evaluation point is not a definition.}

\glossaryentry{report-finite-fabius-approximant}{Report finite Fabius approximant}{\statuslean}{The report finite Fabius approximant is the explicitly named finite spline
used to mirror formulas in the exposition.  It is constructed from finitely
many uniform or dyadic convolution factors and has exact piecewise-polynomial
behavior.  Quarter-cell and related finite theorems concern this object; they
do not construct a finite inverse approximant unless a separate definition
does so.}

\glossaryentry{representation-independence}{Representation independence}{\statusclassical}{Representation independence means that a mathematical output depends only on the represented object, not on a nonunique encoding. Dyadic rationals can be written at multiple levels, real numbers can have two binary expansions, and signed dyadic numerators can have redundant positive-negative pairs. Exact Fabius evaluators prove compatibility under refinement and coding equivalence. This is necessary before an implementation can be regarded as a function on rationals or reals rather than on syntax.}

\glossaryentry{representation-operator}{Representation operator}{\statusmixed}{A representation operator maps coefficient data to a function, for example
\((Tc)(x)=\sum_k c_k\phi_k(x)\).  Boundedness, injectivity, surjectivity, and
stability depend on the coefficient and function spaces.  In the project,
Legendre synthesis, shifted-up synthesis, and sampling reconstruction are
different operators even when they represent the same polynomial.}

\glossaryentry{reproducing-kernel-hilbert-space}{Reproducing-kernel Hilbert space}{\statusfrontier}{An RKHS is a Hilbert space of functions in which each point-evaluation map is
continuous.  There is a unique kernel \(K\) with
\(f(x)=\langle f,K(\cdot,x)\rangle\).  Sampling stability, interpolation
matrices, and minimum-norm reconstruction can be phrased in this framework,
but a kernel's smoothness alone does not establish completeness or universality.}

\glossaryentry{reshetnikov-integers}{Reshetnikov integers}{\statusmixed}{Reshetnikov's integers are normalized quantities \(R_n\) formed from Fabius endpoint values or half moments by multiplying by factorial, power-of-two, double-factorial, and odd Mersenne factors according to the repository's displayed definition. The main arithmetic theorem proves that every \(R_n\) is a positive odd integer. Their integrality refines earlier conjectural observations and packages exact \(2\)-adic cancellation into a natural integer sequence.}

\glossaryentry{residue}{Residue}{\statusclassical}{The residue of a meromorphic function at \(z_0\) is the coefficient of \((z-z_0)^{-1}\) in its Laurent expansion. The residue theorem converts a contour shift into the sum of residues of crossed poles. In Mellin analysis, residues produce asymptotic terms. For the Fabius dyadic harmonic sum, residues at the imaginary lattice yield the Fourier coefficients of the periodic correction.}

\glossaryentry{resolvent-order-hierarchy}{Resolvent-order hierarchy}{\statuslean}{A resolvent-order hierarchy studies
\(\int(z-x)^{-\alpha}\,\dd\mu(x)\) as the positive order \(\alpha\) varies.
Differentiation in \(z\) raises integer order, while integral identities lower
general real order under explicit hypotheses.  The current formal results are
real-variable statements for \(z>1\), not a complex-order branch theory.}

\glossaryentry{retired-alias}{Retired alias}{\statushistorical}{A retired alias is a former command or spelling that may still occur in
archives or source-faithful quotations but is no longer emitted in normative
project prose.  The notation catalogue records a lossless mapping from retired
forms to canonical symbols when one exists.  An alias that merged genuinely
different normalizations is not mechanically replaceable and requires a
local semantic check.}

\glossaryentry{riemann-liouville-fractional-integral}{Riemann--Liouville fractional integral}{\statusclassical}{For \(\alpha>0\), the left-sided Riemann--Liouville integral is
\[
 (I_{a+}^{\alpha}f)(x)=\frac1{\Gamma(\alpha)}
 \int_a^x(x-t)^{\alpha-1}f(t)\,\dd t .
\]
A right-sided version reverses the interval.  Integer orders recover repeated
integration, and suitable function classes satisfy
\(I^\alpha I^\beta=I^{\alpha+\beta}\).}

\glossaryentry{riemann-zeta-function}{Riemann zeta function}{\statusclassical}{The Riemann zeta function is \(\zeta(s)=\sum_{n\ge1}n^{-s}\) for \(\Re s>1\), continued meromorphically to \(\C\) with a simple pole at \(s=1\). It appears in Mellin transforms of Bose-type kernels and in Euler--Maclaurin constants. The Fabius periodic Fourier coefficients use \(\zeta(1-\chi_k)\) at nonreal points. Nonvanishing on the line \(\Re s=1\) at these frequencies helps establish nonconstancy.}

\glossaryentry{riesz-basis}{Riesz basis}{\statusfrontier}{A Riesz basis is the image of an orthonormal basis under a bounded invertible
operator.  Equivalently, finite coefficient combinations satisfy two-sided
\(\ell^2\) norm bounds and the span is complete.  Exact reproduction by
translated up-functions does not by itself establish a Riesz basis; global
stability and completeness require separate proofs.}

\glossaryentry{rising-factorial}{Rising factorial}{\statusclassical}{The rising factorial is
\(x^{\overline n}=x(x+1)\cdots(x+n-1)\), with
\(x^{\overline0}=1\).  It is also the ordinary Pochhammer symbol
\((x)_n\) in hypergeometric notation.  The project keeps it distinct from the
\(q\)-Pochhammer symbol \((a;q)_n\) and from falling factorials.}

\glossaryentry{rodrigues-formula}{Rodrigues formula}{\statusclassical}{Rodrigues' formula constructs Legendre polynomials by \(P_n(x)=\frac1{2^nn!}\frac{\dd^n}{\dd x^n}(x^2-1)^n\). It proves degree, parity, endpoint normalization, and orthogonality after integration by parts. In the up-function Legendre analysis, it converts coefficients into integrals involving derivatives or moments, which can then be evaluated from the Fabius recurrence data.}

\glossaryentry{rogers-szego-polynomial}{Rogers--Szegő polynomial}{\statuslean}{The Rogers--Szegő polynomial is
\[
 H_n(z;q)=\sum_{k=0}^{n}\qbinom{n}{k}{q} z^k.
\]
It is a generating polynomial for one Gaussian row.  Its recurrence, symmetry,
and special values package several \(q\)-binomial identities.  It should not be
confused with the Hermite polynomial despite the common letter \(H\).}

\glossaryentry{root-of-unity-collision}{Root-of-unity collision}{\statusmixed}{A root-of-unity collision occurs when several factors \(1-q^m\) vanish at the
same base while numerator and denominator zeros interact.  The correct local
object is obtained by cyclotomic valuation or a Laurent expansion, not by
substituting into individual quotient factors.  Collisions govern zeros,
finite limits, derivative values, and ramification of inverse \(q\)-functions.}

\glossaryentry{rounding-to-nearest-dyadic}{Rounding to nearest dyadic}{\statusmixed}{Rounding to the nearest dyadic at precision \(p\) selects an integer \(a\) so that \(|x-a/2^p|\le2^{-p-1}\), with a specified tie convention. Combined with a Lipschitz bound, it converts input approximation error into output error. The computable Fabius algorithm can round a real name to a dyadic grid, evaluate a finite spline exactly there, and add the certified spline and input-propagation errors.}

\glossaryentry{rvachev-appell-polynomial}{Rvachev Appell polynomial}{\statuslean}{The Rvachev Appell polynomial is the centered Appell sequence associated with
the up-function probability law.  Its generating function is the exponential
kernel divided by the Rvachev moment-generating function.  Exact rational
moments determine all coefficients, and averaging a translated polynomial
against the Rvachev density recovers the corresponding monomial.}

\glossaryentry{rvachev-atom}{Rvachev atom}{\statuslean}{A Rvachev atom is a translate and scale of \(\Up\), multiplied by a
coefficient supplied by polynomial deconvolution or interpolation.  Compact
support makes only finitely many atoms active on a bounded interval.  Atom
families provide exact polynomial, Lagrange, and Legendre synthesis rather
than merely approximate basis expansions.}

\glossaryentry{rvachev-decoder}{Rvachev decoder}{\statuslean}{A Rvachev decoder is the coefficient functional or finite sum that combines
up-atoms to recover a target value or interpolating polynomial.  Its
normalization is chosen so that the sum of cardinal basis functions is one.
Decoder theorems state degree bounds, mesh divisibility, and exact finite index
ranges.}

\glossaryentry{rvachev-differential-difference-law}{Rvachev differential--difference law}{\statuslean}{Rvachev's up-function satisfies
\[
 \Up'(x)=2\bigl(\Up(2x+1)-\Up(2x-1)\bigr)
\]
for all real \(x\).  The equation combines differentiation, dilation, and two
translations.  Together with compact support, smoothness, evenness, and
normalization it forms the original characterization and generates higher
derivative and refinement identities.}

\glossaryentry{rvachev-polynomial-deconvolution}{Rvachev polynomial deconvolution}{\statuslean}{Rvachev polynomial deconvolution expresses a target polynomial as an average
or finite combination of translated Rvachev Appell polynomials.  It is exact
on each finite polynomial degree and should not be confused with unstable
deconvolution of arbitrary noisy functions.  The compact-support law supplies
the convolution operator; the Appell sequence supplies its polynomial inverse.}

\glossaryentry{rvachev-up-function}{Rvachev up-function}{\statusmixed}{Rvachev's up-function is the even compactly supported bump \(\Up(x)=F(1-|x|)\) obtained from the bounded Fabius function. It is nonnegative, supported exactly on \([-1,1]\), normalized by \(\Up(0)=1\), strictly increasing on \([-1,0]\), and strictly decreasing on \([0,1]\). It is the centered density form of the same probability law, has an entire sinc-product Fourier transform, satisfies refinement identities, and is a canonical atomic function.}

\section{S}\label{sec:letter-s}

\glossaryentry{saddle-central-arc}{Saddle central arc}{\statuslean}{The central arc is the portion of a contour within a shrinking neighborhood
of the saddle where a finite Taylor expansion of the phase and amplitude is
uniformly accurate.  Its width balances Taylor error against Gaussian tail
loss.  The central contribution yields the coefficient expansion; it is only
one part of the full inversion proof.}

\glossaryentry{saddle-coefficient-recurrence}{Saddle coefficient recurrence}{\statusmixed}{A saddle coefficient recurrence computes the coefficients of a local integral expansion from lower-order phase and amplitude data. Gaussian parity eliminates odd total powers, while each surviving coefficient is a finite combination of jets and double-factorial moments. The repository supplies recurrences suitable for Lean induction as well as explicit closed forms indexed by ordered compositions and subsets. The logarithmic coefficients are then obtained from the mass coefficients by a formal-log recurrence.}

\glossaryentry{saddle-equation}{Saddle equation}{\statuslean}{For an integral with exponent \(\Phi(z;\lambda)\), the saddle equation is
\(\partial_z\Phi=0\).  Solving it exactly or asymptotically selects the local
center of mass for steepest descent.  In the Fabius Bromwich problem it reduces
to a Lambert equation after isolating the dominant negative-Laplace terms.}

\glossaryentry{saddle-hessian}{Saddle Hessian}{\statusclassical}{In several variables the Hessian is the matrix of second derivatives of the
phase at the saddle.  A nonsingular Hessian gives a multidimensional Gaussian
main term with determinant factor \((\det(-\nabla^2\Phi))^{-1/2}\), after a
branch and orientation are chosen.  In one dimension it reduces to the second
derivative.}

\glossaryentry{saddle-jet}{Saddle jet}{\statusmixed}{A saddle jet is the finite list of derivatives of the phase or amplitude at the saddle, usually normalized by powers of the large parameter so that each component has a finite limit or periodic dependence. These jets determine every fixed order of the steepest-descent expansion. Repository modules derive closed forms for the Fabius saddle jets using Stirling-number algebra and the periodic Mellin correction, then feed them into generic coefficient machinery.}

\glossaryentry{saddle-mass}{Saddle mass}{\statusmixed}{The saddle mass is the integral contribution from the central neighborhood after the dominant exponential has been extracted. Its leading term is Gaussian, proportional to the inverse square root of the phase curvature. Higher corrections account for nonquadratic phase terms and amplitude variation. In the Fabius problem the normalized saddle mass has an all-orders inverse-\(\lambda\) expansion with periodic coefficients inherited from the negative-Laplace decomposition.}

\glossaryentry{saddle-point}{Saddle point}{\statusmixed}{A saddle point of a complex phase \(\Phi\) is a critical point where \(\Phi'(z_0)=0\). Along suitable directions the real part decreases away from the point, concentrating a contour integral. The Fabius saddle is a real negative point in the Bromwich variable, or equivalently a positive large parameter \(\lambda\) satisfying \(\lambda2^{-\lambda}=x\). Its exact location is expressed through \(W_{-1}\).}

\glossaryentry{saddle-point-method}{Saddle-point method}{\statusmixed}{The saddle-point method estimates integrals or coefficients with a large parameter by expanding the phase around a critical point. A complete proof identifies the saddle, controls contour deformation, evaluates the Gaussian central contribution, expands corrections, and bounds tails uniformly. The repository implements each of these steps for the Fabius endpoint and extends the calculation to all orders. The method is closely related to steepest descent and Laplace's method.}

\glossaryentry{saddle-radius}{Saddle radius}{\statusmixed}{The saddle radius is the cutoff separating the central region, where a Taylor expansion of the phase is accurate, from the tail or minor arc. It is typically a slowly growing function of the saddle parameter. Choosing it too small discards Gaussian mass; choosing it too large invalidates local remainder bounds. The Fabius radius modules prove explicit inequalities that balance these effects at each desired expansion order.}

\glossaryentry{saddle-tail-bound}{Saddle tail bound}{\statuslean}{A saddle tail bound controls the contour outside the central arc by proving a
strict drop in the real part of the phase or a transform decay estimate.
Without it, a local Gaussian calculation does not establish an asymptotic for
the original integral.  The current Fabius proof separates central, intermediate,
and remote arcs.}

\glossaryentry{saddlepoint-density-approximation}{Saddlepoint density approximation}{\statusfrontier}{A saddlepoint approximation evaluates a density through the solution
\(\Lambda'(t_x)=x\) of the tilted log-MGF equation:
\[
 f(x)\approx
 \frac{\exp(\Lambda(t_x)-t_xx)}
 {\sqrt{2\pi\Lambda''(t_x)}}.
\]
Higher corrections use tilted cumulants.  Endpoint degeneracy and lattice
effects require modifications.}

\glossaryentry{sampling-comb}{Sampling comb}{\statusmixed}{A sampling comb is a discrete family of nodes, often an affine lattice
\((k+\theta)/M\), together with weights.  In the Rvachev reconstruction
theorems, up-function atoms centered at the comb nodes reproduce polynomials
after deconvolution.  Mesh, phase, endpoint truncation, and normalization are
all part of the object.}

\glossaryentry{sampling-identity}{Sampling identity}{\statusclassical}{A sampling identity relates sums of function values on a lattice to sums of transform values on the reciprocal lattice. Poisson summation gives \(\sum_{m\in\Z}\Up(am)=a^{-1}\sum_{m\in\Z}\widehat\Up(m/a)\) under the repository's convention. Compact support can reduce the spatial sum to only a few terms for large spacing, while the transform side remains an infinite rapidly convergent series. This yields exact identities for selected up-values.}

\glossaryentry{scalar-extension}{Scalar extension}{\statusconvention}{Scalar extension maps an algebraic object over a field or ring into one over a larger scalar system, such as embedding \(\Q[x]\) into \(\C[x]\). Polynomial identities, finite sums, and recurrences are preserved by ring homomorphisms. The Fabius complex-shift modules prove rational identities once and extend them to complex parameters, keeping analytic convergence arguments separate from algebraic coercions.}

\glossaryentry{scaling-law}{Scaling law}{\statusclassical}{A scaling law relates values of a function or distribution under dilation. The global Fabius law \(\mathcal F'(x)=2\mathcal F(2x)\) and its derivative iterates are exact scaling laws. In transforms, dilation changes frequency inversely and contributes a Jacobian factor. Discrete-scale invariance at ratio two is the common source of dyadic arithmetic, sinc products, and log-periodic endpoint corrections.}

\glossaryentry{schroeder-equation}{Schröder equation}{\statusfrontier}{Schröder's equation
\[
 \phi(f(x))=\lambda\phi(x)
\]
conjugates a nonlinear map near a fixed point to multiplication by
\(\lambda=f'(x_0)\).  When \(\phi\) is locally invertible,
\(f^{\circ t}(x)=\phi^{-1}(\lambda^t\phi(x))\) suggests fractional iterates.
Parabolic or flat fixed points require different coordinates.}

\glossaryentry{schur-complement}{Schur complement}{\statusmixed}{For a block matrix
\(\begin{psmallmatrix}A&B\\C&D\end{psmallmatrix}\) with \(A\) invertible,
the Schur complement is \(D-CA^{-1}B\).  It controls determinants,
positive definiteness, conditional covariance, and least-squares residuals.
Exact product formulas for geometric-family prediction errors arise from
successive Schur complements of covariance matrices.}

\glossaryentry{schwartz-function}{Schwartz function}{\statusclassical}{A Schwartz function is \(C^\infty\) and, together with every derivative, decays faster than every polynomial: \(\sup_x|x^mf^{(n)}(x)|<\infty\) for all \(m,n\). Every compactly supported smooth function is Schwartz. Thus the up-function belongs to the Schwartz space, and its Fourier transform is also Schwartz on the real axis. This guarantees strong inversion and Poisson-summation properties.}

\glossaryentry{second-order-eulerian-number}{Second-order Eulerian number}{\statusfrontier}{Second-order Eulerian numbers count Stirling permutations by descents and
occur in expansions related to increasing trees and repeated derivative
operators.  They are not squares or iterates of ordinary Eulerian numbers.
The indexing and boundary convention are declared with the family.}

\glossaryentry{second-saddle-correction}{Second saddle correction}{\statusmixed}{The second saddle correction is the coefficient of \(1/\lambda^2\) in the normalized or logarithmic saddle expansion. It combines phase and amplitude derivatives through higher order and includes cross terms involving the first correction. The repository isolates this coefficient in dedicated modules both as a check on the generic all-orders machinery and as a sharper practical endpoint approximation.}

\glossaryentry{sectorial-asymptotic-expansion}{Sectorial asymptotic expansion}{\statusclassical}{A sectorial expansion is valid as a complex variable tends to a limit inside
a specified sector.  Powers, logarithms, and exponentials use branches
compatible with that sector.  Crossing a boundary ray can change exponentially
small terms or activate a Stokes phenomenon, even when the Poincaré power
series remains unchanged.}

\glossaryentry{self-convolution}{Self-convolution}{\statusclassical}{Self-convolution is \(f*f\), representing the sum of two independent variables with density \(f\). A self-similar density may satisfy an exact relation between its self-convolution and a scaled or shifted version of itself. The up-function has such identities inherited from splitting its infinite random series into interlaced dyadic scales. Fourier transforms turn them into elementary product identities.}

\glossaryentry{self-sampling-identity}{Self-sampling identity}{\statusfrontier}{A self-sampling identity reconstructs a function, transform, or polynomial
using samples of the same distinguished kernel rather than an unrelated
basis.  For the up-function, partition-of-unity and translate-synthesis
formulas supply exact self-referential sampling patterns.  Convergence and
finite support depend on the precise identity.}

\glossaryentry{self-similarity}{Self-similarity}{\statusclassical}{Self-similarity means that an object is reproduced, exactly or statistically, after rescaling and possibly translation or averaging. The Fabius random series splits into a first uniform component plus a scaled independent copy of itself. This distributional identity yields the contraction equation, differential scaling, and refinement structure. Because the scale ratio is discrete, the endpoint asymptotic retains periodic information in a logarithmic phase.}

\glossaryentry{semantic-api}{Semantic API}{\statusconvention}{A semantic API is a collection of notation commands named by mathematical
meaning rather than visual appearance.  Commands such as a specified Fourier
transform, binary digit sum, or generalized \(q\)-Pochhammer symbol encode
arity and normalization in their names.  This makes large-scale migration and
collision auditing possible and reduces the risk that identical glyphs conceal
different objects.}

\glossaryentry{sequential-computability}{Sequential computability}{\statusclassical}{A real function is sequentially computable if it maps every computable input sequence to a computable output sequence by one uniform algorithm. This formulation avoids choosing a particular representation for a single arbitrary real while still enforcing effective dependence on input data. In Grzegorczyk's framework, sequential computability plus effective uniform continuity defines computability of a continuous real function. The finite-spline Fabius algorithm supplies the sequential component.}

\glossaryentry{series-reversion-coefficient}{Series-reversion coefficient}{\statusmixed}{For
\(f(w)=\sum_{j\ge1}f_jw^j/j!\) with \(f_1\ne0\) and
\(f^{\langle-1\rangle}(z)=\sum g_nz^n/n!\), each \(g_n\) is a rational
polynomial in \(f_1^{-1},f_2,\ldots,f_n\).  Bell-polynomial and partition
formulas encode these coefficients.  Factorial normalization must be fixed
before comparing formulas.}

\glossaryentry{sharp-asymptotic}{Sharp asymptotic}{\statusclassical}{A sharp asymptotic includes every contribution needed to attain a specified small remainder, not merely the leading growth order. For \(F(x)\) near zero, sharpness requires the exact lower-Lambert phase, the nonconstant periodic correction, the Gaussian prefactor, and stated inverse-saddle errors. A coarser quadratic logarithmic law remains true but is not sharp enough for asymptotic equivalence or first correction terms.}

\glossaryentry{sharp-constant}{Sharp constant}{\statusclassical}{A constant in an inequality is sharp if it cannot be decreased while preserving the statement. Sharpness is usually proved by exhibiting points or sequences where equality holds or the ratio approaches the constant. The derivative bounds \(2^{\binom{m+1}{2}}\) for the Fabius-related functions are sharp because exact dyadic points attain them. The global Lipschitz constant \(2\) is similarly attained at the midpoint.}

\glossaryentry{sharp-derivative-bound}{Sharp derivative bound}{\statusmixed}{The repository proves \(\|F^{(m)}\|_\infty,\|\Up^{(m)}\|_\infty,\|\mathcal F^{(m)}\|_\infty\le2^{\binom{m+1}{2}}\), with equality in each case. The bound follows from derivative scaling and the unit bound on the signed extension; equality follows at carefully chosen dyadic arguments. These estimates imply effective flatness, coefficient decay, and stable uniform approximation. They are exact rather than merely exponential-order bounds.}

\glossaryentry{sharp-transfer}{Sharp transfer}{\statusclassical}{A sharp transfer theorem converts asymptotic information for one analytic object into equally precise information for another while preserving constants, periodic phases, and error orders. In the Fabius directory, transfer modules pass from negative-Laplace and saddle estimates to endpoint values, from the continuous endpoint formula to inverse-power sequences, and from the Lambert coordinate to elementary logarithms. The adjective sharp means no term at the target error scale is lost.}

\glossaryentry{sheffer-sequence}{Sheffer sequence}{\statuslean}{A Sheffer sequence has exponential generating function
\(A(z)e^{xB(z)}\), where \(A(0)\ne0\), \(B(0)=0\), and \(B'(0)\ne0\).
Appell sequences are the case \(B(z)=z\).  Sheffer addition, lowering
operators, and exponential Riordan arrays organize changes among polynomial
bases used in coefficient and reversion calculus.}

\glossaryentry{shifted-spline}{Shifted spline}{\statusclassical}{A shifted spline is a spline evaluated after translation, possibly by a complex parameter in a polynomial extension. Real shifts move support and knots; complex shifts are algebraic evaluations of the polynomial pieces rather than probability densities. The Fabius complex-shift spline modules prove convergence for fixed shifts and then combine the approximants with Gaussian \(q\)-binomial weights. Common-shift invariance belongs to the finite Taylor functional, not necessarily to each finite spline limit row.}

\glossaryentry{signed-global-series}{Signed global series}{\statusmixed}{The signed global Fabius representation writes the global extension
as a Thue--Morse-signed sum of integer translates of the bounded function, with
the normalization fixed in the defining module.  Because the bounded Fabius
function has constant tails and the translate coefficients cancel in complete
binary blocks, the formal expression must be read through the proved
locally-finite or finite-reduction identity appropriate to the theorem.  It
agrees with the bounded function on \([0,1]\), satisfies the pure global
differential-dilation law, and is not itself a cumulative distribution
function.}

\glossaryentry{signed-measure}{Signed measure}{\statusclassical}{A signed measure is a countably additive set function that may take positive and negative values but has finite total variation on the relevant space. It decomposes into positive and negative parts by the Jordan decomposition. Thue--Morse translate sums may be viewed as signed combinations of compactly supported pieces. The signed global Fabius function is not itself a probability CDF, even though it is built from the same local density.}

\glossaryentry{signed-natural-numerator}{Signed-natural numerator}{\statusclassical}{A signed-natural numerator is a pair \((c^+,c^-)\in\N^2\) representing the integer \(c^+-c^-\). It is redundant because many pairs represent the same integer, but addition, subtraction, and primitive recursion can be implemented using only natural numbers. At dyadic precision \(p\), the represented rational is \((c^+-c^-)/2^p\). This coding is used in the repository's fast names for computable real sequences.}

\glossaryentry{signed-stirling-first-kind}{Signed Stirling numbers of the first kind}{\statusclassical}{The signed Stirling number \(s(n,k)\) is the coefficient of \(x^k\) in
\(x^{\underline n}\).  It obeys
\(s(n+1,k)=s(n,k-1)-n\,s(n,k)\) and
\(s(0,0)=1\).  The sign convention is part of the name; unsigned first-kind
numbers are absolute values and count permutations by cycles.}

\glossaryentry{signed-thue-morse-prefix}{Signed Thue--Morse prefix}{\statuslean}{A signed Thue--Morse prefix is the finite list
\((\varepsilon_0,\ldots,\varepsilon_{N-1})\) or its partial sum.  Complete
powers-of-two prefixes sum to zero after the first level, while incomplete
prefixes obey a reflection recursion.  Repeated summation of these prefixes
produces polynomial arrays connected to Fabius dyadic values.}

\glossaryentry{signed-weight-random-series}{Signed-weight random series}{\statuslean}{A signed-weight random series permits positive and negative deterministic
coefficients in \(\sum a_jU_j\).  Absolute summability of the weights gives
almost-sure and \(L^p\) convergence for bounded summands and determines a
compact support interval by extremal choices.  Sign changes affect the affine
center and support but not the product structure of centered characteristic
functions.}

\glossaryentry{sinc-function}{Sinc function}{\statusclassical}{The sinc function is \(\sinc z=\sin z/z\) with removable value \(1\) at \(z=0\). It is entire and has simple zeros at nonzero integer multiples of \(\pi\) before normalization. The Fourier transform of a centered uniform density is a sinc factor. Multiplying the factors from all dyadic-scale uniforms gives the up-function's entire Fourier product.}

\glossaryentry{sinc-lobe-sign}{Sinc-lobe sign}{\statusmixed}{Between consecutive real zeros, a sinc factor has constant sign.  The sign
of a product is therefore determined by how many factors lie in negative
lobes.  For the dyadic up product this parity can be expressed through binary
or Thue--Morse data on integer-separated intervals.  Lobe sign is pointwise
real information, not a complex phase convention.}

\glossaryentry{sinc-product}{Sinc product}{\statusclassical}{The up-function Fourier transform is \(\widehat\Up(z)=\prod_{n=0}^{\infty}\sinc(\pi z/2^n)\). The product converges locally uniformly and defines an entire function. Its zeros at nonzero integers imply exact periodization and partition-of-unity identities; repeated factors and scale splitting encode convolution laws. Along the real axis, the accumulation of factors gives rapid decay.}

\glossaryentry{small-argument-asymptotic}{Small-argument asymptotic}{\statusclassical}{A small-argument asymptotic describes a function as \(x\to0\). For the Fabius function this regime is singular because all endpoint derivatives vanish, so no nontrivial Taylor series exists. The correct expansion is obtained indirectly through a negative-Laplace transform and a moving saddle. It includes a quadratic logarithmic main scale, nested-log corrections, a lower-Lambert phase, and smooth periodic coefficients.}

\glossaryentry{smooth-compact-support}{Smooth compact support}{\statusclassical}{A function has smooth compact support when it is \(C^\infty\) and vanishes outside a compact set. Such functions form the test-function space \(\mathcal D(\R)\) and are automatically Schwartz. Their Fourier transforms are entire of exponential type and rapidly decaying on the real axis. The up-function is a highly structured example, distinguished by its refinement equation and exact arithmetic.}

\glossaryentry{source-checkpoint}{Source checkpoint}{\statusconvention}{A source checkpoint is a recorded commit identifier at which a module
inventory or theorem list was compiled.  It makes a statement reproducible
even when the main branch continues to evolve.  A short hash is adequate as a
locator only when its repository is clear.  Checkpoints attached to individual
tranches do not imply that every current file was frozen at the same commit.}

\glossaryentry{source-module}{Source module}{\statusconvention}{A source module is a Lean file participating in the imported development.
Its mathematical subject may be narrow--for example one \(q\)-identity, one
moment recurrence, or one approximation bridge--even when it exports several
public declarations.  Module counts therefore measure formal organization,
not the number of independent mathematical theories or glossary concepts.}

\glossaryentry{source-pdf-parity}{Source--PDF parity}{\statusconvention}{Source--PDF parity means that a rendered PDF incorporates the mathematical
content currently present in its LaTeX source.  A live research repository can
contain post-render source additions, so a PDF may be valid yet incomplete
relative to the head source.  Audits record this condition explicitly; a
glossary based on current state must inspect source as well as retained PDFs.}

\glossaryentry{spectral-measure}{Spectral measure}{\statusclassical}{A spectral measure represents a self-adjoint operator or a moment
functional by integration against its spectrum.  For a cyclic vector,
moments are matrix elements of powers of the operator.  Orthogonal
polynomials, Jacobi matrices, Stieltjes transforms, and quadrature nodes are
different views of the same measure.}

\glossaryentry{spectral-profile-inversion}{Spectral profile inversion}{\statusfrontier}{If integer-zero multiplicities satisfy
\(m_v=\sum_{h\le v}c_h(v-h+1)\), then the profile is recovered by the second
difference
\(c_v=m_v-2m_{v-1}+m_{v-2}\).
Thus admissible zero ledgers are exactly the nonnegative discrete-convex
multiplicity sequences with the stated boundary conventions.}

\glossaryentry{spectral-zeta-function}{Spectral zeta function}{\statusfrontier}{A spectral zeta function is
\(\zeta_L(s)=\sum_{\lambda_j>0}\lambda_j^{-s}\) in its convergence half-plane,
continued when possible.  Poles encode eigenvalue asymptotics, and
\(-\zeta_L'(0)\) defines a regularized determinant when continuation is
regular at zero.  Any Fabius-transform version must state the chosen
multiplicity and zero enumeration.}

\glossaryentry{spline}{Spline}{\statusclassical}{A spline is a piecewise-polynomial function with prescribed continuity of derivatives across knots. Repeated convolution of box functions produces B-splines, whose degree and smoothness increase together. The finite Fabius approximants are explicit centered splines on dyadic knots. They converge uniformly to a \(C^\infty\) limit even though each finite approximant has only finite smoothness.}

\glossaryentry{square-summable-activation}{Square-summable activation profile}{\statuslean}{A square-summable activation statement controls a sequence of scale
differences, derivatives, or deficits in \(\ell^2\).  Such estimates are
stronger than pointwise convergence and support energy identities, Hilbert
space arguments, and stable reconstruction.  In the Fabius setting the scale
recursion and compact transition make the relevant activation increments
explicitly summable.}

\glossaryentry{stable-endpoint-evaluation}{Stable endpoint evaluation}{\statusmixed}{Stable endpoint evaluation computes extremely small Fabius values without losing all relative accuracy. Direct subtraction, naive transform inversion, or standard floating-point arithmetic may underflow or suffer cancellation. Reliable methods use exact dyadic formulas, high-precision logarithmic asymptotics with the periodic phase, or monotone recurrences. The repository discusses operation choice, residual checks, and Newton refinement in a regime where the value itself may be far below machine range.}

\glossaryentry{standardized-cumulant}{Standardized cumulant}{\statusmixed}{For variance \(\sigma^2>0\), the standardized cumulant is
\(\lambda_r=\kappa_r/\sigma^r\), \(r\ge3\).  It is invariant under positive
rescaling and measures non-Gaussian shape.  In geometric-uniform limits these
cumulants decay with explicit powers of the small mesh parameter and generate
Edgeworth corrections.}

\glossaryentry{standardized-geometric-uniform-law}{Standardized geometric uniform law}{\statusfrontier}{A standardized geometric uniform variable subtracts its mean and divides by
its standard deviation.  As the ratio approaches the Gaussian boundary
\(q\uparrow1\), standardized cumulants organize Edgeworth and
Cornish--Fisher expansions.  The standardization constant must be stated,
because source drafts also use centered variables on fixed support without
variance normalization.}

\glossaryentry{stars-and-bars}{Stars-and-bars principle}{\statusclassical}{Stars and bars is the bijection proving that weak compositions of \(n\)
into \(r\) parts are counted by \(\binom{n+r-1}{r-1}\).  Positive
compositions are obtained by subtracting one from every part.  The principle
explains binomial kernels in repeated summation and simplex-volume factors in
iterated integration.}

\glossaryentry{stationary-point}{Stationary point}{\statusclassical}{A stationary point of a differentiable real or complex phase is a point where the first derivative vanishes. It becomes a saddle suitable for asymptotics only when curvature and contour geometry make the contribution dominant. The Fabius audit examines an informal stationary substitution into a Stirling proxy and shows that the exact stationarity equation leaves a linear residual. Thus locating a stationary point does not by itself validate a claimed asymptotic simplification.}

\glossaryentry{steepest-descent}{Steepest descent}{\statusclassical}{Steepest descent is a contour method in which the path is chosen through a saddle along directions where the imaginary part of the phase is constant and the real part decreases most rapidly. It turns the local integral into a Gaussian model with controlled corrections. The Bromwich saddle method used for Fabius is a specialized steepest-descent analysis adapted to inverse Laplace integration and a dyadic transform product.}

\glossaryentry{steepest-descent-contour}{Steepest-descent contour}{\statusmixed}{A steepest-descent contour passes through a saddle along directions where the
imaginary part of the phase is constant and the real part decreases most
rapidly.  Deforming to it requires holomorphy and control of connecting arcs.
In a vertical Bromwich problem one may instead prove equivalent central and
tail estimates without naming a global contour.}

\glossaryentry{stein-equation}{Stein equation}{\statusfrontier}{For a test function \(h\), the Stein equation is
\[
 \mathcal Af_h=h-\int h\,\mathrm d\mu.
\]
Bounds on the solution \(f_h\) convert expectation identities into probability
metrics.  The norm, derivative order, and test class determine which distance
is controlled.}

\glossaryentry{stein-operator}{Stein operator}{\statusfrontier}{A Stein operator \(\mathcal A\) characterizes a probability law \(\mu\) on a
test class \(\mathcal T\) when
\(\int\mathcal Af\,\mathrm d\mu=0\) for every \(f\in\mathcal T\), together
with a uniqueness theorem.  An integration-by-parts identity alone is not a
complete characterization.  Proposed Fabius/Rvachev Stein operators retain
frontier status until the class and uniqueness are established.}

\glossaryentry{stieltjes-constant}{Stieltjes constant}{\statusclassical}{The Stieltjes constants \(\gamma_n\) are defined by the Laurent expansion \(\zeta(s)=1/(s-1)+\sum_{n\ge0}(-1)^n\gamma_n(s-1)^n/n!\), with \(\gamma_0=\gamma\). They appear when Mellin-zeta expressions are expanded near a pole to higher order. Fabius finite-part and all-orders calculations may involve them in constant or derivative coefficients of Gamma--zeta products.}

\glossaryentry{stieltjes-inversion-formula}{Stieltjes inversion formula}{\statuslean}{Stieltjes inversion reconstructs a measure from boundary values of its
Cauchy transform.  In the absolutely continuous case, the density is a
constant-sign multiple of
\(\lim_{\eta\downarrow0}\operatorname{Im}G(x+\ii\eta)\), determined by the
denominator convention.  Interval versions recover mass even when a density
is absent.}

\glossaryentry{stieltjes-moment-sequence}{Stieltjes moment sequence}{\statusmixed}{A sequence \((m_n)\) is a Stieltjes moment sequence when
\(m_n=\int_0^\infty x^n\,\mathrm d\mu(x)\) for a positive measure \(\mu\).
Equivalently, both the Hankel matrices \((m_{i+j})\) and shifted Hankel matrices
\((m_{i+j+1})\) are positive semidefinite.  Compactly supported Fabius/Rvachev
moments are Hausdorff moments after an affine change, hence automatically
determinate.}

\glossaryentry{stieltjes-pade-approximation}{Stieltjes--Padé approximation}{\statusfrontier}{For a Stieltjes function, continued-fraction convergents and diagonal Padé
approximants give alternating or monotone bounds on positive real rays under
standard hypotheses.  Moment positivity prevents spurious nonreal poles.
Applying this framework to transforms built from Rvachev moments yields
certified approximants distinct from direct spline truncations.}

\glossaryentry{stieltjes-perron-inversion}{Stieltjes--Perron inversion}{\statusclassical}{The Stieltjes--Perron formula recovers interval masses from the jump or
imaginary boundary values of a Cauchy transform.  Endpoint atoms receive
convention-dependent half weights in symmetric limits.  It connects analytic
resolvents to CDFs and spectral distributions.}

\glossaryentry{stieltjes-transform}{Stieltjes transform}{\statusmixed}{A Stieltjes transform commonly means the Cauchy transform of a measure on a
half-line, or the variant \(\int(x+z)^{-1}\,\dd\mu(x)\).  Sign and domain
conventions vary.  The repository prints the denominator and develops
fixed-point, logarithmic, and generalized-order resolvent hierarchies for
Fabius-related measures.}

\glossaryentry{stirling-formula}{Stirling formula}{\statusclassical}{Stirling's formula approximates the factorial: \(n!=\sqrt{2\pi n}(n/\e)^n(1+O(1/n))\), or logarithmically \(\log n!=n\log n-n+\frac12\log(2\pi n)+O(1/n)\). Explicit Robbins bounds provide a rigorous two-sided error. The Fabius audit uses Stirling's formula both in endpoint sequence asymptotics and to expose a residual linear term missed by an informal stationary substitution.}

\glossaryentry{stirling-numbers}{Stirling numbers}{\statusclassical}{Stirling numbers of the first kind convert falling factorials to ordinary powers, while those of the second kind convert ordinary powers to falling factorials. They also encode derivatives of logarithmic and exponential composites. The Fabius saddle-jet modules use Stirling-number identities to obtain closed forms for repeated derivatives of dyadic kernel expressions and Lambert-dependent phases. Sign conventions for first-kind numbers must be specified.}

\glossaryentry{stirling-second-kind}{Stirling numbers of the second kind}{\statusclassical}{The Stirling number
\(\genfrac{\{}{\}}{0pt}{}{n}{k}\) counts partitions of an \(n\)-element set into
\(k\) nonempty blocks.  It is the coefficient in
\(x^n=\sum_k\genfrac{\{}{\}}{0pt}{}{n}{k}x^{\underline k}\).
Its triangular recurrence and zero extensions are fixed at the boundaries.}

\glossaryentry{stokes-line}{Stokes line}{\statusfrontier}{A Stokes line is a ray or curve across which an exponentially small
contribution changes its asymptotic multiplier while dominant exponential
orders exchange activity.  The underlying analytic function remains
continuous after correct continuation.  A Stokes statement requires a chosen
phase, branch, and sector; it cannot be inferred from a real-axis expansion
alone.}

\glossaryentry{stopped-primitive}{Stopped primitive}{\statuslean}{A stopped primitive integrates only after a threshold, for example
\[
 (Pf)(x)=\int_a^x f(t)\,\dd t
\]
on \(x>a\) and assigns a fixed exterior value on the other side.  Stopping
preserves support or tail conventions while solving a derivative equation.
Iterated stopped primitives appear in CDF construction and compact-support
integration ladders.}

\glossaryentry{strang-fix-condition}{Strang--Fix condition}{\statusclassical}{A compactly supported generator reproduces polynomials through degree
\(d\) on a lattice when its Fourier transform has zeros of order at least
\(d+1\) at all nonzero dual-lattice points, together with a normalization at
zero.  These Strang--Fix conditions translate transform zero multiplicities
into approximation order and explain the role of \(v_2(M)\) in Rvachev comb
exactness.}

\glossaryentry{strict-interior-positivity}{Strict interior positivity}{\statuslean}{Strict interior positivity means \(\Up(x)>0\) exactly for
\(-1<x<1\), with zero values outside and at the endpoints.  It is stronger
than nonnegativity almost everywhere.  The result identifies the exact
pointwise support interior, proves strict increase of \(F\), and ensures that
the density relation has no hidden zero plateaus.}

\glossaryentry{strict-interior-range}{Strict interior range}{\statuslean}{The strict interior range statement says \(0<F(x)<1\) for
\(0<x<1\).  It follows from the strict positivity of the up-density and
monotonicity of the CDF.  This excludes flat subintervals, supplies injectivity
on the open unit interval, and makes the canonical inverse a genuine order
isomorphism there.}

\glossaryentry{strict-unimodality}{Strict unimodality}{\statusclassical}{A function is strictly unimodal if it increases strictly up to a unique mode and decreases strictly afterward. Rvachev's up-function is strictly increasing on \([-1,0]\), attains its unique maximum \(1\) at zero, and is strictly decreasing on \([0,1]\). Its derivative signs follow from the positive bounded Fabius function under affine scaling. Strictness strengthens mere symmetry and nonincreasing shape.}

\glossaryentry{strip-of-convergence}{Strip of convergence}{\statusclassical}{For a Mellin transform, the strip of convergence is the vertical region \(a<\Re s<b\) where the defining integral converges absolutely. Endpoint behavior at zero controls one boundary and behavior at infinity controls the other. Mellin inversion begins on a vertical line inside this strip; asymptotic analysis then uses meromorphic continuation and contour shifting. The regularized Fabius kernels enlarge the workable region by subtracting singular endpoint terms.}

\glossaryentry{sturm-liouville-equation}{Sturm-Liouville equation}{\statusclassical}{A Sturm-Liouville equation has the form \(-(p(x)y')'+q(x)y=\lambda w(x)y\) with boundary conditions. Its eigenfunctions are orthogonal with respect to the weight \(w\). Legendre polynomials solve the singular equation \(-((1-x^2)y')'=n(n+1)y\) on \((-1,1)\). This differential characterization yields recurrences and integration-by-parts identities used in the up-function's Legendre expansion.}

\glossaryentry{subgraph-fubini-identity}{Subgraph Fubini identity}{\statusclassical}{The measure of the subgraph
\(\{(x,t):0<t<f(x)\}\) can be computed by integrating first in \(t\) or first
in \(x\).  The two orders give \(\int f\) and an integral of level-set
measures.  Applied to a monotone CDF or inverse, this produces geometric
area identities and moment relations with endpoint terms made explicit.}

\glossaryentry{subspace-lattice}{Subspace lattice}{\statusclassical}{The subspaces of a finite-dimensional vector space over
\(\mathbb F_q\), ordered by inclusion, form a finite ranked lattice.
The number of \(k\)-dimensional subspaces is
\(\qbinom{n}{k}{q}\).  Intervals are themselves subspace lattices of quotient
spaces, which yields a closed M\"obius function.}

\glossaryentry{summability}{Summability}{\statusclassical}{Summability concerns assigning or proving limits for series, sometimes by methods extending ordinary convergence. In this directory most series are shown to converge in the ordinary sense, often absolutely or uniformly; Toeplitz rows are regular summability transformations carrying convergent sequences to the same limit. The term should not be read as permission to manipulate divergent series without a named regularization such as a finite part.}

\glossaryentry{superasymptotic-approximation}{Superasymptotic approximation}{\statusfrontier}{A superasymptotic approximation is an optimally or near-optimally truncated
asymptotic expansion whose error reaches the scale of the least term.  It
improves fixed-order asymptotics without introducing the next exponential
sector.  Hyperasymptotic methods go further by expanding that exponentially
small remainder.}

\glossaryentry{supergraph-fubini-identity}{Supergraph Fubini identity}{\statusclassical}{A supergraph identity integrates the set
\(\{(x,t):f(x)<t\}\) inside a bounded rectangle.  It is complementary to the
subgraph formula and naturally expresses deficits such as \(1-F(x)\).
Reflection of the Fabius CDF makes the subgraph and supergraph descriptions
equivalent after an affine change of variables.}

\glossaryentry{support}{Support}{\statusclassical}{The support of a continuous function is the closure of the set where it is nonzero. For a measure it is the smallest closed set of full measure, equivalently the points every neighborhood of which has positive mass. The up-function has support exactly \([-1,1]\), while the bounded Fabius probability law has support \([0,1]\). Endpoints belong to the support despite zero function values there.}

\glossaryentry{support-as-range}{Support-as-range principle}{\statuslean}{For a continuous nonnegative density, the topological support is the closure
of the set where the density is positive.  In the Fabius system the exact
interior positivity and endpoint zeros identify the support of \(\Up\) as
\([-1,1]\), while strict monotonicity identifies the range of \(F\) on
\([0,1]\) as \([0,1]\).  These facts justify canonical inverse domains rather
than merely almost-everywhere statements.}

\glossaryentry{family-support}{Support of an indexed family}{\statusclassical}{For an indexed family \((a_i)_{i\in I}\), its support is
\(\{i\in I:a_i\ne0\}\).  Finite support turns an apparently infinite sum into
a finite one, while local finiteness says only finitely many terms affect each
evaluation point.  This discrete notion is distinct from topological support
of a function or measure, though translate constructions use both
simultaneously.}

\glossaryentry{supremum-norm}{Supremum norm}{\statusclassical}{The supremum norm is \(\|f\|_\infty=\sup_x|f(x)|\) on a specified domain. Spaces of bounded continuous functions are complete in this norm, enabling the Banach fixed-point construction. Uniform convergence is convergence in supremum norm. Sharp derivative bounds and finite-spline approximation errors in the Fabius development are stated in this norm.}

\glossaryentry{survival-function}{Survival function}{\statusclassical}{The survival function of a real random variable is
\(\overline F(x)=\Prob(X>x)=1-F(x)\) under the stated endpoint convention.
Reflection can identify a survival tail with a CDF tail at a mirrored point.
Layer-cake formulas express positive moments as integrals of survival
probabilities.}

\glossaryentry{symmetric-mixed-difference}{Symmetric mixed difference}{\statuslean}{A symmetric mixed difference uses paired shifts
\(x\pm h_j/2\) or \(x\pm h_j\), producing parity-adapted cancellation.
Odd or even polynomial components then vanish automatically.  Such operators
connect centered moments, even up-function symmetry, and factorized sinc
products.}

\glossaryentry{symmetry}{Symmetry}{\statusclassical}{Symmetry is invariance or a simple covariance under a transformation. The bounded Fabius CDF has complement reflection symmetry about \(1/2\), and the up-function has even symmetry about zero. Symmetry determines means, cancels odd moments and coefficients, transfers endpoint formulas, and locates the unique inflection point. It also reduces computations to half of the support.}

\section{T}\label{sec:letter-t}

\glossaryentry{tail-estimate}{Tail estimate}{\statusclassical}{A tail estimate bounds the contribution from indices or integration regions beyond a cutoff. For infinite random sums it controls the omitted variables; for series it controls the remainder; for contour integrals it controls remote arcs. The dyadic random tail has an explicit deterministic bound, yielding the centered spline sandwich. Saddle tail estimates show that regions outside the central radius are negligible at the desired inverse-\(\lambda\) order.}

\glossaryentry{tail-flatness}{Tail flatness}{\statusclassical}{Tail flatness refers to rapid vanishing or derivative control in the remote part of a scaled kernel or contour integrand. In the Fabius saddle proof, it supplies uniform bounds when the local Taylor approximation is no longer used. The phrase is repository-specific: it packages estimates showing that the reference tail and true tail differ by less than the target asymptotic remainder.}

\glossaryentry{tanh-divided-slope}{Tanh divided slope}{\statuslean}{The tanh divided slope is the regularized quotient
\(\tanh x/x\), assigned the removable value \(1\) at \(x=0\), or an
affinely scaled analogue used to compare activation derivatives.  Positivity,
evenness, and monotonicity properties turn pointwise slope estimates into
summable scale bounds.  The current primary exposition includes this calculus
as a bridge between Fabius activation profiles and classical hyperbolic
functions.}

\glossaryentry{taylor-block}{Taylor block}{\statusclassical}{A Taylor block is a finite Taylor polynomial contribution associated with one dyadic scale in the binary reduction. Because derivatives at a dyadic center vanish above a known order, the block is exact rather than an approximation. Its coefficients are expressed through scaled values of the signed global extension. Gaussian \(q\)-binomial identities provide an alternative complex-shift representation of the same block.}

\glossaryentry{taylor-polynomial}{Taylor polynomial}{\statusclassical}{The degree-\(n\) Taylor polynomial of \(f\) at \(x_0\) is \(T_n(x)=\sum_{k=0}^{n}f^{(k)}(x_0)(x-x_0)^k/k!\). It approximates \(f\) locally when a remainder estimate is available. At reduced dyadic Fabius points, the formal Taylor series terminates exactly at the denominator exponent, yet the polynomial does not equal the function on any neighborhood. This is a striking separation between jet data and analyticity.}

\glossaryentry{taylor-series}{Taylor series}{\statusclassical}{The Taylor series of a smooth function at \(x_0\) is the formal series of its derivatives. Analyticity requires that this series converge to the function locally; smoothness alone does not. Fabius endpoint Taylor series are constant because of flatness, and interior dyadic Taylor series are finite polynomials, but neither represents the nonconstant function nearby. Transform Taylor series, by contrast, are entire and valid globally.}

\glossaryentry{terminating-basic-hypergeometric-series}{Terminating basic hypergeometric series}{\statuslean}{A basic hypergeometric series terminates when an upper parameter equals
\(q^{-N}\) for some \(N\in\N\), because \((q^{-N};q)_n=0\) for \(n>N\).
It is then a finite sum, so analytic convergence is irrelevant, though
denominator nonvanishing remains necessary.  Many Gaussian and
Fabius--Thue--Morse identities are best proved in this finite setting.}

\glossaryentry{terminating-taylor-series}{Terminating Taylor series}{\statusclassical}{A Taylor series terminates when all derivatives above some order vanish at the center, so the formal series is a polynomial. At \(x_0=a/2^n\) with odd \(a\), the Fabius derivatives vanish above order \(n\) because repeated dyadic scaling lands outside the active integer pattern. The order-\(n\) derivative is nonzero, so the degree is exactly \(n\). Termination does not imply local polynomiality for a merely smooth function.}

\glossaryentry{q-series-termination-index}{Termination index}{\statusconvention}{The termination index is the largest \(n\) for which a terminating
\(q\)-hypergeometric summand may be nonzero.  An upper parameter \(q^{-N}\)
usually gives index \(N\), but denominator zeros can shorten the meaningful
range or make a term undefined.  Explicitly recording the index keeps finite
proofs independent of convergence arguments.}

\glossaryentry{termwise-differentiation}{Termwise differentiation}{\statusclassical}{Termwise differentiation replaces the derivative of a series by the series of derivatives. It is valid under conditions such as locally uniform convergence of the derivative series together with convergence at one point, or normal convergence for holomorphic series. The rapid Fourier and periodic-coefficient convergence in the repository permits repeated differentiation. Formal power-series differentiation is algebraic and requires no analytic hypothesis, so the two settings must be distinguished.}

\glossaryentry{termwise-integration}{Termwise integration}{\statusclassical}{Termwise integration interchanges a sum and an integral. Tonelli's theorem covers nonnegative terms, Fubini covers absolute integrability, and dominated convergence covers limits of partial sums under a common majorant. Fabius moment, Fourier, and Laplace product derivations use these principles. A merely conditionally convergent signed series requires additional care and is not automatically integrable term by term.}

\glossaryentry{test-function-space}{Test-function space}{\statusclassical}{The test-function space \(\mathcal D(\R)\) consists of infinitely differentiable compactly supported functions, equipped with a standard locally convex topology. Its elements are used to define distributions and weak derivatives. The up-function is an explicit member of \(\mathcal D(\R)\). In the original literature it is notable as a highly structured test function satisfying exact refinement and partition identities.}

\glossaryentry{theta-quasi-periodicity}{Theta quasi-periodicity}{\statuslean}{For the multiplicative theta normalization,
\[
 \theta(qz;q)=-z^{-1}\theta(z;q).
\]
Iteration gives a monomial factor times \(\theta(z;q)\).  This is
quasi-periodicity rather than ordinary periodicity: the function changes by an
explicit multiplier.  The multiplier controls zero propagation and
root-of-unity scaling.}

\glossaryentry{three-term-recurrence}{Three-term recurrence for orthogonal polynomials}{\statusmixed}{Monic orthogonal polynomials obey
\[
 \pi_{n+1}(x)=(x-\alpha_n)\pi_n(x)-\beta_n\pi_{n-1}(x),
 \qquad \beta_n>0,
\]
for a positive measure with infinite support.  The coefficients are obtained
from inner products or Hankel determinants.  Symmetry of the measure forces
\(\alpha_n=0\).  This recurrence is distinct from the functional-differential
recurrences defining Fabius objects.}

\glossaryentry{thue-morse-bit}{Thue--Morse bit}{\statuslean}{The Thue--Morse bit is \(t_n=s_2(n)\bmod2\), where \(s_2(n)\) is the
binary digit sum.  The sign convention is
\(\varepsilon_n=(-1)^{t_n}\).  Bits satisfy
\(t_{2n}=t_n\) and \(t_{2n+1}=1-t_n\), while signs satisfy the multiplicative
versions.  The project uses both forms and never identifies a bit in
\(\{0,1\}\) with a sign in \(\{\pm1\}\).}

\glossaryentry{thue-morse-block-product}{Thue--Morse block product}{\statusclassical}{The generating polynomial of a complete Thue--Morse block factors as
\[
 \sum_{n=0}^{2^m-1}\varepsilon_n z^n
 =\prod_{j=0}^{m-1}(1-z^{2^j}).
\]
The factorization follows by binary digit independence.  Its zero of order
\(m\) at \(z=1\) is equivalent to vanishing signed power sums through degree
\(m-1\).}

\glossaryentry{thue-morse-gamma-product}{Thue--Morse gamma product}{\statuslean}{A Thue--Morse gamma product is a finite or convergent product of shifted gamma
factors whose exponents are Thue--Morse signs.  Prouhet cancellation removes
low-order terms in its logarithmic expansion.  Reciprocal-gamma normalization
turns pole bookkeeping into an entire zero divisor when possible.}

\glossaryentry{thue-morse-sequence}{Thue-Morse sequence}{\statusmixed}{The Thue--Morse sequence is the binary automatic sequence \(t_n=\wt(n)\bmod2\), satisfying \(t_{2n}=t_n\) and \(t_{2n+1}=1-t_n\). Its sign version is \(\varepsilon_n=(-1)^{t_n}\). Finite generating polynomials factor digitwise and have high-order zeros at one, producing Prouhet cancellation. The sequence governs the translate signs of the global Fabius extension and exact dyadic formulas.}

\glossaryentry{thue-morse-sign}{Thue-Morse sign}{\statusmixed}{The Thue--Morse sign is \(\varepsilon_n=(-1)^{\wt(n)}\). It obeys \(\varepsilon_{2n}=\varepsilon_n\), \(\varepsilon_{2n+1}=-\varepsilon_n\), and multiplicative block identities when binary digits do not overlap. Integer values of the signed global Fabius extension are tied to these signs. They also determine the signs of exact highest dyadic derivatives.}

\glossaryentry{tightness}{Tightness}{\statusclassical}{A family of probability measures on \(\R\) is tight if for every \(\varepsilon>0\) some compact set carries mass at least \(1-\varepsilon\) for every measure in the family. Tightness prevents mass from escaping to infinity and is central to weak compactness. Finite Fabius convolution laws have uniformly bounded compact supports, so tightness is immediate. This supports passage to the infinite convolution law.}

\glossaryentry{toeplitz-lemma}{Toeplitz lemma}{\statusclassical}{A Toeplitz lemma states that a triangular array of weights preserves the limit of a convergent sequence when row sums tend to one, individual fixed-column weights vanish or rows sample sufficiently late terms, and row total variations are uniformly bounded. The repository uses a finite-row variant: after reindexing, all sampled spline degrees are at least the row index, the mass is exactly one, and variation is uniformly bounded by a \(q\)-product ratio. Hence the weighted complex approximants converge to \(\mathcal F\).}

\glossaryentry{toeplitz-row}{Toeplitz row}{\statusclassical}{A Toeplitz row is one row \((\omega_{n,j})_j\) of a triangular summability matrix applied to a sequence. Regular rows have mass tending to one, vanishing influence from fixed early indices, and uniformly bounded total variation. The Fabius complex discrete limit obtains exact row mass and variation from two evaluations of the finite \(q\)-binomial theorem. Reindexing also guarantees that every row samples only spline degrees at least \(n\).}

\glossaryentry{tonelli-theorem}{Tonelli theorem}{\statusclassical}{Tonelli's theorem allows the order of integration or summation to be exchanged for nonnegative measurable functions, even before finiteness is known. It proves that the iterated integrals and product integral have the same extended value. In probability constructions with nonnegative densities and moments, it is often the cleanest justification. Alternating Thue--Morse sums require Fubini or a separate absolute-convergence argument instead.}

\glossaryentry{topological-support}{Topological support}{\statusclassical}{Topological support is the closed set of points at which an object cannot be locally discarded. For a continuous function it is the closure of its nonzero set; for a measure it consists of points whose every neighborhood has positive measure. The distinction from the set of nonzero values matters at flat endpoints. Thus \(\Up(\pm1)=0\) but \(\pm1\in\supp\Up\).}

\glossaryentry{total-variation}{Total variation}{\statusclassical}{For a signed or complex measure, total variation is the mass of its variation measure. For a finite weight row it is \(\sum_j|\omega_j|\), and for a function it is the supremum over partition increments. Uniform total-variation bounds make weighted averaging stable even with alternating coefficients. The \(q\)-binomial Toeplitz rows have an explicit uniform bound obtained by evaluating the finite \(q\)-binomial theorem at a negative argument.}

\glossaryentry{total-variation-distance}{Total-variation distance}{\statusclassical}{For probability measures \(\mu,\nu\),
\[
 d_{\mathrm{TV}}(\mu,\nu)=\sup_A|\mu(A)-\nu(A)|.
\]
If both have densities \(f,g\) with respect to a common measure, it equals
\(\frac12\int|f-g|\).  This is stronger than weak or Kolmogorov convergence in
many settings; the factor \(1/2\) is part of the project's convention.}

\glossaryentry{totalized-definition}{Totalized definition}{\statusconvention}{A totalized definition assigns a value at every input, including points where
the motivating partial formula would be undefined.  Lean often uses totalized
division, inverse, infinite sums, or branch placeholders; useful theorems then
state hypotheses under which the value matches the intended analytic object.
The repository pairs total definitions with explicit zero, endpoint, and
denominator lemmas so that totality is never mistaken for an unconditional
identity.}

\glossaryentry{touchard-polynomial}{Touchard polynomial}{\statusclassical}{The Touchard polynomial is
\(T_n(x)=\sum_k\genfrac{\{}{\}}{0pt}{}{n}{k}x^k\).
It is the \(n\)-th moment polynomial of a Poisson law with mean \(x\) and has
exponential generating function \(\exp(x(e^z-1))\).  At \(x=1\) it gives the
Bell number.}

\glossaryentry{transfer-theorem}{Transfer theorem}{\statusclassical}{A transfer theorem carries a property or asymptotic from one representation to another. Examples include transferring transform decay to smoothness, saddle estimates to coefficient asymptotics, or endpoint formulas to moment sequences. The Fabius directory contains explicit transfer modules because constants, normalizations, and periodic phases can be lost in an informal passage. A sharp transfer preserves the stated remainder order.}

\glossaryentry{translate-block}{Translate block}{\statuslean}{A translate block is a finite linear combination of consecutive translates
of one kernel, usually at a fixed scale.  Polynomial reproduction,
finite-difference annihilation, and local support are analyzed blockwise before
blocks are summed across scales or modes.  The term records geometry; it does
not imply orthogonality.}

\glossaryentry{translate-sum}{Translate sum}{\statusclassical}{A translate sum has the form \(\sum_kc_kf(x-k)\) or a scaled variant. With compact support it may be locally finite. Constant positive coefficients can form partitions of unity; Thue--Morse signs create a signed global extension. Fourier transformation converts translation into multiplication by phases, making periodicity and cancellation visible.}

\glossaryentry{translated-legendre-series}{Translated Legendre series}{\statusclassical}{A translated Legendre series expands a function on an interval other than \([-1,1]\) by composing Legendre polynomials with the affine map to the standard interval. The bounded Fabius function on \([0,1]\) can be related to the up-function series on \([-1,1]\) through translation and reflection. Coefficients and convergence must include the Jacobian and normalization of this change of variables.}

\glossaryentry{transmonomial}{Transmonomial}{\statusclassical}{A transmonomial is one multiplicative scale in a transseries, such as
\(x^\alpha(\log x)^m e^{-Ax}\).  Transmonomials form an ordered multiplicative
group.  The ordering records asymptotic dominance, not numerical size at one
finite point.}

\glossaryentry{transseries}{Transseries}{\statusmixed}{A transseries extends ordinary series by allowing ordered monomials involving
powers, logarithms, exponentials, and iterated combinations, for example
\(x^{-a}(\log x)^m e^{-bx}\).  Its support must be well based so that every
coefficient operation is locally finite.  Inversion and composition require
both formal support conditions and, for analytic use, branch and sector
conditions.}

\glossaryentry{transseries-composition}{Transseries composition}{\statusmixed}{Composition substitutes one large or small transseries into another.
Formally it requires support conditions making each output coefficient finite;
analytically the inner map must remain in the outer expansion's sector and
avoid branch cuts.  For \(t=1/X\) and \(L=\log X\), substituting
\(G=X+H\) changes both generators:
\(t\mapsto t/(1+tH)\) and \(L\mapsto L+\log(1+tH)\).}

\glossaryentry{transseries-reversion}{Transseries reversion}{\statusmixed}{Reversion constructs a compositional inverse in an ordered transseries class.
The leading monomial must be invertible and the correction must be smaller in
a way that makes coefficient determination triangular.  Polynomial--logarithmic
reversion solves successively for logarithmic coefficient blocks; multiple
critical points require Puiseux rather than ordinary Laurent powers.}

\glossaryentry{triangle-admissibility}{Triangle admissibility}{\statuslean}{A triple \((\ell,m,n)\) is triangle-admissible when each index is at most
the sum of the other two.  For Legendre products one also requires even total
degree.  These finite arithmetic conditions decide whether a Gaunt coefficient
can be nonzero before any factorial expression is evaluated.}

\glossaryentry{triangular-cancellation}{Triangular cancellation}{\statusclassical}{Triangular cancellation is a layered vanishing mechanism in which lower-degree contributions disappear successively, leaving a top-degree term. In the Fabius dyadic formula, Gaussian interpolation weights annihilate node polynomials below degree \(n\), while Thue--Morse signs independently annihilate lower power sums inside dyadic blocks. The surviving term has a triangular power-of-two normalization. Recognizing both layers explains why apparently complicated finite sums collapse exactly.}

\glossaryentry{triangular-recurrence}{Triangular recurrence}{\statusclassical}{A recurrence is triangular when the equation for index \(n\) involves the new term with a nonzero coefficient and only terms of smaller index otherwise. It can be solved inductively and often determines uniqueness. Moment, Bernoulli, numerator, and saddle coefficient recurrences in the repository have this form. Clearing denominators may convert a rational triangular recurrence into an integer division-free one.}

\glossaryentry{triangular-reduction}{Triangular reduction}{\statusclassical}{Triangular reduction is the exact decomposition of a dyadic numerator into a large block, a reflected remainder, and lower-level data so that evaluation proceeds inductively. The word triangular reflects both decreasing level and the triangular exponents \(\binom n2\) appearing in normalization. This reduction proves correctness of fast dyadic recursion and exposes binary reflection symmetries.}

\glossaryentry{truncated-power-basis}{Truncated-power basis}{\statusmixed}{The truncated-power function is
\((x-a)_+^m=\max(x-a,0)^m\).  Finite linear combinations of such functions,
plus a polynomial, represent univariate splines.  Inverse Fourier transforms of
finite sinc products admit exact alternating truncated-power formulas indexed
by signed subset sums, making knots and local polynomials explicit.}

\glossaryentry{truncation-error}{Truncation error}{\statusclassical}{Truncation error is the difference between an infinite object and a finite partial approximation. It may arise from omitting random summands, series terms, spline levels, or asymptotic corrections. The Fabius finite random tail gives a deterministic support bound and hence a uniform CDF error; Fourier and saddle truncations use summable or exponential estimates. Explicit truncation errors turn formulas into certified algorithms.}

\glossaryentry{two-adic-unit}{\(2\)-adic unit}{\statusclassical}{A rational or \(2\)-adic integer is a \(2\)-adic unit when its
\(2\)-adic valuation is zero.  For ordinary integers this means oddness.
Separating powers of two from unit factors clarifies exact dyadic
denominators, Reshetnikov odd numerators, and central-binomial valuations.}

\glossaryentry{two-kernel}{\(2\)-kernel of a sequence}{\statusclassical}{The \(2\)-kernel of a sequence \(a(n)\) is the family
\[
 \{\,n\mapsto a(2^en+r): e\ge0,\ 0\le r<2^e\,\}.
\]
A finite-valued sequence is \(2\)-automatic exactly when this kernel is finite.
The Thue--Morse sequence has a two-state kernel, while iterated prefix sums
typically produce finite-dimensional \(2\)-regular rather than finite-valued
behavior.}

\glossaryentry{two-pi-fourier-transform}{\(2\pi\)-frequency Fourier transform}{\statusconvention}{The \(2\pi\)-frequency convention is
\[
 \widehat f(\xi)=\int_{\R}f(x)e^{-2\pi\ii x\xi}\,\dd x .
\]
Under it, translation contributes \(e^{-2\pi\ii a\xi}\), differentiation
contributes \(2\pi\ii\xi\), and integer sampling is naturally paired with
integer frequencies.  The up-function product then uses
\(\sinc(\pi\xi/2^n)\) in the radian-normalized sinc.}

\glossaryentry{type-b-eulerian-number}{Type-B Eulerian number}{\statusfrontier}{Type-B Eulerian numbers count signed permutations by a type-B descent
statistic.  Their polynomials have different degree, symmetry, and recurrence
from the ordinary type-A Eulerian family.  The semantic family label prevents
a bare bracket notation from merging them.}

\glossaryentry{typed-map}{Typed map}{\statusclassical}{A typed map declaration \(f:A\to B\) includes both domain and codomain.
Restrictions, inverses, compositions, and supports depend on these types.
Many apparent formula conflicts in the Fabius corpus disappear once one
distinguishes a bounded real function, a function on a compact interval, a
formal power series, and a polynomial evaluated in a ring.  The type is part
of the mathematical statement, not implementation decoration.}

\section{U}\label{sec:letter-u}

\glossaryentry{unconditional-convergence}{Unconditional convergence}{\statusclassical}{A series converges unconditionally when every permutation of its
terms converges to the same value.  In finite-dimensional real or complex
spaces this is equivalent to absolute convergence, while in general Banach
spaces the notions differ.  Unconditional convergence licenses regrouping
without choosing a privileged order.  The repository invokes it for genuinely
infinite analytic and Hilbert-space expansions; locally finite Fabius translate
identities are order-independent for the simpler reason that only finitely
many summands are nonzero at each point.}

\glossaryentry{uniform-approximation}{Uniform approximation}{\statusclassical}{Uniform approximation means convergence in supremum norm: \(\sup_x|f_n(x)-f(x)|\to0\). It preserves continuity and provides one error bound for all inputs. The centered finite splines approximate the bounded Fabius function uniformly with an explicit dyadic rate, and the Fourier-Legendre series of the up-function converges uniformly as well. The proofs use different mechanisms: probabilistic coupling for splines and coefficient bounds for polynomials.}

\glossaryentry{uniform-asymptotic-expansion}{Uniform asymptotic expansion}{\statusmixed}{An expansion is uniform in a parameter \(u\in U\) when its remainder bound
holds with constants independent of \(u\).  Uniformity is indispensable near
moving saddles, for quantiles away from the endpoints, and for \(q\)-limits on
compact parameter sets.  A pointwise expansion for every fixed \(u\) does not
imply a uniform expansion.}

\glossaryentry{uniform-convergence}{Uniform convergence}{\statusclassical}{A sequence \(f_n\) converges uniformly to \(f\) if for every \(\varepsilon>0\) one index works for all points in the domain. Uniform limits of continuous functions are continuous, and uniform convergence permits integration interchange on finite measure domains. It is stronger than pointwise convergence. The Fabius development repeatedly upgrades structural approximations to uniform ones so they can support computability and exact shape arguments.}

\glossaryentry{uniform-distribution}{Uniform distribution}{\statusclassical}{The uniform distribution on \([a,b]\) has constant density \(1/(b-a)\) on that interval. Its characteristic function is a phase times a sinc factor, and its moments are elementary rational expressions. Independent dyadically scaled uniforms are the building blocks of the Fabius random series. Centering removes the phases in the product transform and yields the even up-density.}

\glossaryentry{uniform-legendre-convergence}{Uniform Legendre convergence}{\statuslean}{Uniform Legendre convergence means
\[
 \sup_{x\in[-1,1]}\left|f(x)-\sum_{n=0}^{N}c_nP_n(x)\right|\to0.
\]
It permits termwise evaluation and preserves continuity.  It is stronger than
\(L^2\) convergence; in the repository it follows from an explicit summable
majorant for the coefficients, not merely from completeness.}

\glossaryentry{uniform-norm}{Uniform norm}{\statusclassical}{Uniform norm is another name for the supremum norm on a function space. Convergence in this norm is uniform convergence. It is the metric used for the contraction fixed-point theorem, finite-spline error, and sharp derivative bounds in the Fabius development. On unbounded domains, boundedness or constant tails ensure the norm is finite.}

\glossaryentry{uniform-random-variable}{Uniform random variable}{\statusclassical}{A uniform random variable on \([a,b]\) assigns probability proportional to interval length. It can be realized as \(a+(b-a)U\) with \(U\sim\mathrm{Unif}[0,1]\). Boundedness gives all moments and exponential moments. The Fabius law is assembled from independent uniform variables with geometrically shrinking coefficients, so finite partial sums have spline densities.}

\glossaryentry{uniform-spline}{Uniform spline}{\statusclassical}{A uniform spline has knots on an equally spaced grid. Cardinal B-splines are the standard examples. Although the original Fabius summands occur at different dyadic scales, a common refinement converts finite convolutions into uniform-grid splines with explicit coefficients. Centered uniform splines form the constructive approximation used in the computability theorem.}

\glossaryentry{unimodality}{Unimodality}{\statusclassical}{A function or density is unimodal if it increases up to a mode and decreases afterward, allowing nonstrict plateaus in the basic definition. Strict unimodality requires strict inequalities away from the unique mode. The up-function is strictly unimodal with mode zero. Consequently the Fabius density rises then falls, and the CDF changes from convexity to concavity at its unique midpoint inflection.}

\glossaryentry{unique-midpoint-inflection}{Unique midpoint inflection}{\statuslean}{The bounded Fabius CDF changes from strict convexity to strict concavity at
\(x=1/2\), its unique inflection point.  Through
\(F''(x)=4\Up'(2x-1)\), this is equivalent to strict unimodality of the
up-density.  Reflection forces the midpoint location, while strict shape
theorems exclude additional inflections.}

\glossaryentry{uniqueness-of-periodic-correction}{Uniqueness of the periodic correction}{\statusmixed}{Suppose two continuous one-periodic functions inserted into the same Fabius endpoint main term both leave an error tending to zero. Their difference, evaluated along the unbounded lower-Lambert coordinate, tends to zero. Because every phase modulo one is approached by suitable endpoint sequences, periodicity and continuity force the difference to vanish identically. Thus the correction is uniquely determined, not merely one convenient Fourier representation.}

\glossaryentry{fabius-translate-on-unit-interval}{Unit-interval Fabius--up translate}{\statuslean}{On the unit interval, the bounded Fabius function is an affine translate of
the up-function, with the exact formula fixed by the chosen normalization.
The identity is not global because \(F\) has constant tails whereas
\(\Up\) is compactly supported and even.  Every use outside the stated
interval must instead employ the fold or density relation.}

\glossaryentry{unsigned-stirling-first-kind}{Unsigned Stirling numbers of the first kind}{\statusclassical}{The unsigned Stirling number
\(\genfrac{[}{]}{0pt}{}{n}{k}\) counts permutations of \(n\) elements having
exactly \(k\) cycles.  It equals \((-1)^{n-k}s(n,k)\).  These numbers appear
in logarithmic derivatives, rising-factorial expansions, and analytic
continuations of Bell-related sequences.}

\glossaryentry{up-synthesis-amplitude}{Up-synthesis amplitude}{\statusconvention}{An up-synthesis amplitude is the scalar multiplying one translated/dilated
up-atom.  It depends on the target polynomial, scale, centers, and chosen
normalization of the atom.  Amplitudes are not probability weights in general:
they may be signed and need not sum to one unless constant reproduction forces
that condition.}

\section{V}\label{sec:letter-v}

\glossaryentry{valuation}{Valuation}{\statusclassical}{A discrete prime valuation \(v_p(r)\) records the exponent of a prime \(p\) in a nonzero rational number. It is additive on products and satisfies a non-Archimedean inequality on sums. Valuations translate divisibility into arithmetic inequalities and identify reduced denominators. The Fabius arithmetic emphasizes \(v_2\), but odd-prime valuations also underlie least-common-multiple denominator bounds.}

\glossaryentry{valuation-histogram}{Valuation histogram}{\statusfrontier}{A valuation histogram counts how many entries in a finite family have each
\(p\)-adic valuation.  For a Pascal row it refines the count of coefficients
nonzero modulo \(p\); for rational Fabius data it records denominator depth.
Generating functions of these histograms reveal digit recursions and limiting
distributions.}

\glossaryentry{vandermonde-determinant}{Vandermonde determinant}{\statusclassical}{For distinct nodes \(x_0,\ldots,x_n\), the Vandermonde determinant is \(\det(x_i^j)_{0\le i,j\le n}=\prod_{i<j}(x_j-x_i)\), which is nonzero. It proves uniqueness of polynomial interpolation and invertibility of the evaluation map. On geometric nodes, its factors split into powers of two and \(q\)-Pochhammer products, explaining Gaussian interpolation weights in the dyadic formulas.}

\glossaryentry{vanishing-moment}{Vanishing moment}{\statusclassical}{A signed measure, sequence, or function has vanishing moments through order \(m-1\) if its integrals or sums against \(x^r\) vanish for \(r<m\). Vanishing moments annihilate low-degree polynomials and improve approximation or decay. Complete Thue--Morse blocks have vanishing discrete moments up to an order determined by block length. This is the Prouhet cancellation used in Fabius finite formulas.}

\glossaryentry{varentropy}{Varentropy}{\statusfrontier}{Varentropy is the variance of the information content
\(-\log f(X)\) for a random variable with density \(f\):
\[
 V(f)=\Var[-\log f(X)].
\]
It refines differential entropy and appears as a second-order fluctuation
quantity.  Compact support and smoothness do not automatically guarantee
finiteness because \(f\) vanishes at the boundary.}

\glossaryentry{variable-substitution}{Variable substitution}{\statusclassical}{Variable substitution transforms an integral or equation by a change of coordinates and includes the appropriate Jacobian. Affine substitutions connect \(F\) and \(\Up\); logarithmic substitution converts dilation to translation; saddle substitution centers and scales the Bromwich contour. A flawed stationary substitution can leave an unaccounted linear term, so exact algebra and error propagation are essential.}

\glossaryentry{variance}{Variance}{\statusclassical}{The variance of a square-integrable random variable is \(\Var(X)=\E[(X-\E X)^2]=\E[X^2]-(\E X)^2\). Variances add for independent variables. It is the second cumulant and sets the Gaussian scale in central approximations. For the dyadic uniform random series, the variance is an explicit geometric sum and is encoded among the centered moments.}

\glossaryentry{vertical-bromwich-line}{Vertical Bromwich line}{\statusclassical}{A vertical Bromwich line is the contour \(\Re s=c\) used in inverse Laplace integration, oriented from \(c-\ii\infty\) to \(c+\ii\infty\). The constant \(c\) is chosen inside the transform's inversion region and to the right of relevant singularities. Saddle analysis may shift or deform this line after proving analyticity and decay. Orientation and the factor \(1/(2\pi\ii)\) determine the sign and normalization.}

\glossaryentry{volterra-operator}{Volterra integral operator}{\statusclassical}{A Volterra operator has triangular kernel
\[
 (V_Kf)(x)=\int_a^xK(x,t)f(t)\,\dd t .
\]
Its iterates integrate over ordered simplices, and under standard bounds its
resolvent is represented by a Neumann series.  Fabius integral recursions,
stopped primitives, and fractional integrals are special Volterra
constructions.}

\glossaryentry{volterra-resolvent}{Volterra resolvent}{\statusclassical}{For \(I-\lambda V_K\), the Volterra resolvent is the inverse operator when
it exists, often written as \(I+\sum_{n\ge1}\lambda^nV_K^n\).
Triangular integration frequently gives convergence without a smallness
condition on \(\lambda\) over a compact interval because iterates acquire
factorial denominators.  Resolvents organize integral-equation fixed points
and perturbations.}

\section{W}\label{sec:letter-w}

\glossaryentry{walsh-hadamard-transform}{Walsh--Hadamard transform}{\statusclassical}{The Walsh--Hadamard transform replaces complex exponential characters of a
cyclic group by \(\{\pm1\}\)-valued characters of the binary cube.  It is
computed by repeated butterfly additions and subtractions.  Thue--Morse signs
are a highest-order Walsh character, so Prouhet cancellation and Boolean mixed
differences have a natural Walsh interpretation.}

\glossaryentry{wasserstein-distance}{Wasserstein distance}{\statusclassical}{The \(p\)-Wasserstein distance minimizes
\((\int|x-y|^p\,\dd\pi)^{1/p}\) over couplings \(\pi\) of two laws.  In one
dimension it can be computed from quantile functions.  Coupling a full
weighted random series with its finite truncation on the same probability
space gives an immediate Wasserstein bound by the deterministic or \(L^p\)
tail size.}

\glossaryentry{weak-composition}{Weak composition}{\statusclassical}{A weak composition of \(n\) into \(r\) parts is an \(r\)-tuple of
nonnegative integers summing to \(n\).  There are
\(\binom{n+r-1}{r-1}\) such tuples.  Weak compositions index repeated
convolutions, coefficient extractions, Bell-polynomial expansions, and
ordered distributions of derivative order.}

\glossaryentry{weak-convergence}{Weak convergence}{\statusclassical}{Probability measures \(\mu_n\) converge weakly to \(\mu\) if \(\int f\,\dd\mu_n\to\int f\,\dd\mu\) for every bounded continuous \(f\). Characteristic-function convergence often proves weak convergence. Finite convolution laws of the truncated Fabius random series converge weakly to the infinite law because the random variables converge almost surely and remain uniformly bounded. Additional regularity yields convergence of densities or CDFs.}

\glossaryentry{weak-star-convergence}{Weak-star convergence}{\statusclassical}{Weak-star convergence of finite measures tests convergence against continuous compactly supported functions, viewing measures as dual objects. For probability measures on a compact set it closely matches weak convergence. The repository's early measure bridge uses weak-star convergence of convolution measures before identifying the smooth density and uniform CDF limit. The star distinguishes dual-space convergence from weak convergence inside the primal function space.}

\glossaryentry{weierstrass-m-test}{Weierstrass M-test}{\statusclassical}{The Weierstrass M-test states that if \(|f_n(x)|\le M_n\) for every \(x\) and \(\sum M_n<\infty\), then \(\sum f_n\) converges absolutely and uniformly. On compact subsets it proves normal convergence of holomorphic series. The repository uses summable majorants for Fourier modes, global binary summands, and transform logarithms. Uniform convergence then licenses continuity and termwise operations.}

\glossaryentry{weierstrass-product}{Weierstrass product}{\statusclassical}{A Weierstrass product represents an entire function by factors prescribed by
its zeros, with canonical exponential corrections ensuring convergence.
Finite-genus products generalize polynomial factorization to infinite zero
sets.  The Rvachev sinc product is an especially explicit normally convergent
product whose dyadic factors expose zeros and multiplicities directly;
reciprocal-Gamma and theta products in the enlarged repository provide further
examples.}

\glossaryentry{weighted-path-expansion}{Weighted path expansion}{\statusclassical}{A weighted path expansion unfolds a recurrence as a sum over directed paths through its dependency graph, multiplying transition weights along each path. It converts recursive definitions into finite closed sums indexed by compositions or subsets. Saddle coefficient and moment recurrences can be represented this way. In Lean, path expansions can make positivity, denominator structure, and coefficient dependence easier to prove than repeated symbolic substitution.}

\glossaryentry{weighted-random-sum}{Weighted random sum}{\statusclassical}{A weighted random sum is \(\sum a_nX_n\). Its mean, variance, support, transforms, and convergence depend on both weights and component laws. Geometric weights give self-similarity and compact support when the variables are bounded. The Fabius law is the prototypical infinite weighted sum of independent uniforms, with base-two weights producing exact dyadic structure.}

\glossaryentry{weighted-uniform-law}{Weighted-uniform law}{\statuslean}{A weighted-uniform law is the distribution of a convergent sum of
independent uniforms with a general summable weight sequence.  It extends the
geometric family and admits transform, support, expectation, variance, and
continuous-linear mapping theorems under summability assumptions.  Special
self-similarity and exact \(q\)-identities arise only when the weights are
geometric.}

\glossaryentry{well-based-support}{Well-based support}{\statusmixed}{A support is well based when it contains no infinite sequence proceeding
indefinitely toward larger monomials in the chosen dominance order.  This
condition guarantees that the coefficient of a fixed monomial in a product is
a finite sum.  It is the transseries analogue of the lower-bounded support of a
formal Laurent series.}

\glossaryentry{well-founded-recursion}{Well-founded recursion}{\statusclassical}{Well-founded recursion defines an object at \(x\) from values at arguments
strictly smaller than \(x\) in a relation having no infinite descending chain.
It generalizes ordinary induction on natural numbers.  Fast dyadic evaluators,
binary reductions, and recursive numerator algorithms use measures such as
bit length, denominator level, or a positive grid numerator to justify
termination.  In formal proofs, supplying the decreasing measure is part of
the algorithm rather than an informal assurance.}

\glossaryentry{well-poised-basic-hypergeometric-series}{Well-poised basic hypergeometric series}{\statusclassical}{A basic hypergeometric series is well poised when its parameters satisfy a
symmetric set of multiplicative relations that pair numerator and denominator
factors.  Very-well-poised series impose an additional square-root pattern.
These conditions create cancellation and transformation formulas; they must
be printed because the adjectives alone do not determine a convention.}

\glossaryentry{wiener-chaos}{Wiener or orthogonal chaos}{\statusfrontier}{An orthogonal chaos decomposition writes
\(L^2(\Omega)=\bigoplus_{n\ge0}\mathcal H_n\) relative to a specified
independent-coordinate or Gaussian structure.  The projection \(J_n\) selects
the degree-\(n\) component.  For dyadic uniform coordinates the chaos basis is
built from the orthogonal polynomials of the uniform law, not Gaussian Hermite
polynomials.}

\glossaryentry{wigner-three-j-symbol}{Wigner \(3j\)-symbol}{\statusmixed}{A Wigner \(3j\)-symbol packages angular-momentum coupling coefficients.
For zero magnetic indices,
\[
 \int_{-1}^{1}P_\ell P_mP_n\,\dd x
 =2\begin{pmatrix}\ell&m&n\\0&0&0\end{pmatrix}^{\!2}.
\]
The symbol vanishes under the same parity and triangle obstructions as the
Gaunt coefficient.  The square removes its sign in this specialization.}

\glossaryentry{wright-omega-function}{Wright omega function}{\statusmixed}{Wright's omega function solves
\[
 \Omega(z)+\log\Omega(z)=z
\]
with a chosen logarithm branch, equivalently
\(\Omega(z)=W(e^z)\) under compatible continuation.  At large \(z\) it has a
polynomial--logarithmic expansion
\(z-\log z+\cdots\).  It removes one outer exponential from many Lambert
inversion problems.}

\section{Z}\label{sec:letter-z}

\glossaryentry{zero-bias-transform}{Zero-bias transform}{\statusfrontier}{For a mean-zero random variable \(X\) with variance \(\sigma^2>0\), a
zero-biased variable \(X^\ast\) satisfies
\[
 \mathbb E[Xf(X)]=\sigma^2\mathbb E[f'(X^\ast)]
\]
for an appropriate class of absolutely continuous \(f\).  The star denotes a
distributional transform, not complex conjugation.  Iterated zero-bias
constructions require recentering and variance bookkeeping.}

\glossaryentry{zero-of-entire-function}{Zero of an entire function}{\statusclassical}{A zero of order \(m\) of an entire function \(F\) at \(z_0\) satisfies \(F(z)=(z-z_0)^mG(z)\) with \(G(z_0)\ne0\). Zeros are isolated unless the function is identically zero. The sinc product locates transform zeros at dyadically structured frequencies, with multiplicities determined by how many factors vanish. These zeros imply exact Fourier coefficient vanishings, partitions of unity, and odd-harmonic cosine expansions.}

\glossaryentry{zero-run}{Zero run}{\statusclassical}{A zero run is a consecutive block of zero digits or zero values. In binary recursions, long zero runs correspond to inactive dyadic scales and can be skipped by highest-set-bit algorithms. In coefficient sequences, a zero run may reflect parity, support, or finite-degree annihilation. The Fabius binary evaluator exploits bit structure rather than iterating through every zero position.}

\glossaryentry{zeta-regularization}{Zeta regularization}{\statusclassical}{Zeta regularization assigns finite values to certain divergent products or sums by analytic continuation of a zeta-type function. It is conceptually related to, but distinct from, the Hadamard finite part used in the Fabius Mellin analysis. The repository's Gamma--zeta coefficients arise from ordinary meromorphic continuation and residue calculus; no arbitrary assignment of a divergent original series is intended. The term is included because zeta-based regularization language appears in discussions of finite parts and must be distinguished carefully.}

\glossaryentry{zeta-regularized-determinant}{Zeta-regularized determinant}{\statusfrontier}{When a spectral zeta function is regular at \(s=0\), the
zeta-regularized determinant is
\(\det_\zeta L=\exp(-\zeta_L'(0))\).  It extends finite products of positive
eigenvalues but is not an ordinary convergent product in general.  Scale
changes contribute explicit powers governed by \(\zeta_L(0)\).}

\glossaryentry{zonoid}{Zonoid}{\statusfrontier}{A zonoid is a Hausdorff limit of zonotopes, equivalently a convex body
represented by an even measure on directions through its support function.
Infinite sums of shrinking random segments suggest zonoid support sets in
multivariate atomic-function generalizations.}

\glossaryentry{zonotope}{Zonotope}{\statusfrontier}{A zonotope is a Minkowski sum of finitely many line segments.  It is the image
of a cube under a linear map.  Finite sums of independent uniform segment
variables have box-spline densities supported on zonotopes, linking
multivariate Rvachev-type laws to spline geometry.}

\appendix
\section{Normalization and notation guardrails}
The glossary is descriptive, while the live notation catalogue is normative.
The following table records the distinctions most likely to corrupt a theorem
when a formula is copied between documents.
\begingroup\small\renewcommand{\arraystretch}{1.15}
\begin{longtable}{@{}L{0.22\textwidth}L{0.29\textwidth}L{0.43\textwidth}@{}}
\toprule Object & Canonical form & Guardrail\\ \midrule\endfirsthead
\toprule Object & Canonical form & Guardrail\\ \midrule\endhead
Bounded Fabius function & \(F:\R\to[0,1]\) & Clamped CDF: \(F=0\) on \((-\infty,0]\), \(F=1\) on \([1,\infty)\). \\
Global signed extension & \(\mathcal F:\R\to\R\) & Agrees with \(F\) on \([0,1]\); satisfies the global dilation--differential law. \\
Rvachev up-function & \(\Up(x)=F(1-|x|)\) & Even bump, support \([-1,1]\), \(\Up(0)=1\); not the bounded CDF. \\
Ordinary Fabius inverse & \(F^{-1}:[0,1]\to[0,1]\) & Inverse of the strictly increasing restriction \(F|_{[0,1]}\). \\
Clamped Fabius inverse & \(F^{-1}(\min\{1,\max\{0,y\}\})\) & Global real input, constant outside the unit ordinate interval. \\
Thue--Morse bit/sign & \(t_n=s_2(n)\bmod2\), \(\varepsilon_n=(-1)^{t_n}\) & Bit values \(0,1\) and signs \(\pm1\) are different types. \\
Radian sinc & \(\sinc z=\sin z/z\) & The argument already contains any required factor of \(\pi\). \\
\(2\pi\)-Fourier transform & \(\widehat f(\xi)=\int f(x)e^{-2\pi i x\xi}\,dx\) & Integer sampling and frequency lattices use this convention. \\
Angular Fourier transform & \(\widehat f(\omega)=\int f(x)e^{-i\omega x}\,dx\) & Related by \(\widehat f_{2\pi}(\xi)=\widehat f_{\rm ang}(2\pi\xi)\). \\
Finite \(q\)-Pochhammer & \((a;q)_n=\prod_{j=0}^{n-1}(1-aq^j)\) & Length \(n\in\N\); the empty product at \(n=0\) is one. \\
Infinite \(q\)-Pochhammer & \((a;q)_\infty=\prod_{j\ge0}(1-aq^j)\) & Analytic product for \(|q|<1\); zeros and logarithm branches are tracked. \\
Gaussian coefficient & \(\qbinom{n}{k}{q}\) & Canonical polynomial continuation, not merely a quotient away from roots of unity. \\
Geometric-uniform law & \((1-q)\sum_{j\ge0}q^jU_j\) & Support-normalized on \([0,1]\); centered and standardized variants are separate. \\
Legendre polynomial & \(P_n(1)=1\) & Classical, neither monic nor orthonormal; norm squared \(2/(2n+1)\). \\
Lambert branches & \(W_0\) and \(W_{-1}\) & Principal and lower real branches; endpoint saddles require the branch stated by the theorem. \\
\bottomrule\end{longtable}\endgroup
\section{Representative Lean crosswalk}
This table is intentionally representative rather than a substitute for the
675-module generated inventory.  It gives high-value entry points from
mathematical terminology to active module families.  Exact declaration names
and imports remain controlled by the repository facade and generated census.
\begingroup\footnotesize\renewcommand{\arraystretch}{1.12}
\begin{longtable}{@{}>{\raggedright\arraybackslash}p{0.47\textwidth}>{\raggedright\arraybackslash}p{0.47\textwidth}@{}}
\toprule Module & Mathematical interface\\ \midrule\endfirsthead
\toprule Module & Mathematical interface\\ \midrule\endhead
\nolinkurl{Basic.lean} & \texttt{IsFabius}, foundational tails/reflection/differential hypotheses \\
\nolinkurl{Differential.lean} & iterated derivative scaling and dyadic derivative formulas \\
\nolinkurl{Arithmetic.lean} & exact Fabius arithmetic and recurrence infrastructure \\
\nolinkurl{Existence.lean} & existence of the canonical bounded Fabius object \\
\nolinkurl{OriginalUniqueness.lean} & original compact-support predicate and uniqueness bridges \\
\nolinkurl{FabiusInverse.lean} & \texttt{fabiusInv}, strict monotonicity, order isomorphism \\
\nolinkurl{GlobalExtension.lean} & global signed extension, translate cells, global calculus \\
\nolinkurl{FourierProduct.lean} & infinite sinc-product identity \\
\nolinkurl{FourierAnalytic.lean} & analytic Fourier consequences and zero structure \\
\nolinkurl{ProbabilityRepresentation.lean} & random-series law and distributional representation \\
\nolinkurl{StepApproximationLimit.lean} & finite spline/step approximants and convergence \\
\nolinkurl{FabiusComputable.lean} & effective evaluation and computability interface \\
\nolinkurl{FabiusLogScale.lean} & endpoint logarithmic scale \\
\nolinkurl{FabiusLogSquaredAsymptotic.lean} & quadratic logarithmic endpoint law \\
\nolinkurl{NegativeLaplace.lean} & negative-Laplace product and decomposition \\
\nolinkurl{MellinBose.lean} & Mellin/Bose kernel and finite-part infrastructure \\
\nolinkurl{PeriodicCorrection.lean} & period-one endpoint correction \\
\nolinkurl{FabiusLambertPhase.lean} & Lambert phase coordinates \\
\nolinkurl{BromwichSaddle.lean} & Bromwich inversion and saddle localization \\
\nolinkurl{FabiusSharpAsymptotic.lean} & sharp endpoint asymptotic \\
\nolinkurl{FabiusFullAsymptoticExpansion.lean} & all-orders endpoint expansion \\
\nolinkurl{ThueMorseAutocorrelation.lean} & finite-window and dyadic autocorrelations \\
\nolinkurl{ThueMorseAutocorrelationLimit.lean} & normalized limit recursion \\
\nolinkurl{QPochhammerEntire.lean} & finite/infinite \(q\)-product analytic structure \\
\nolinkurl{GeometricPochhammerNormalConvergence.lean} & normal convergence on compact sets \\
\nolinkurl{QPochhammerDissection.lean} & residue-class product dissections \\
\nolinkurl{QPochhammerInfinite.lean} & infinite \(q\)-Pochhammer identities \\
\nolinkurl{GaussianBinomialPolynomialStructure.lean} & polynomiality and coefficient structure \\
\nolinkurl{GaussianBinomialPalindromic.lean} & reciprocal-base palindromy \\
\nolinkurl{GaussianBinomialAtNegOneDerivative.lean} & value and derivative phenomena at \(q=-1\) \\
\nolinkurl{QBinomialTheoremInfinite.lean} & analytic infinite \(q\)-binomial theorem \\
\nolinkurl{BasicHypergeometricSeries.lean} & normalized \({}_r\phi_s\) series \\
\nolinkurl{HeineTransformation.lean} & Heine's transformation \\
\nolinkurl{QGaussSummation.lean} & \(q\)-Gauss summation \\
\nolinkurl{QPfaffSaalschutz.lean} & terminating balanced \({}_3\phi_2\) summation \\
\nolinkurl{JacksonIntegral.lean} & Jackson integral and \(q\)-calculus identities \\
\nolinkurl{ThetaQuasiPeriodicity.lean} & multiplicative theta transformation law \\
\nolinkurl{JacobiTripleProduct.lean} & bilateral series/product identity \\
\nolinkurl{QLucas.lean} & root-of-unity digit reduction \\
\nolinkurl{CyclotomicDivisibility.lean} & cyclotomic exponents and divisibility \\
\nolinkurl{QBetaIntegral.lean} & \(q\)-beta integral \\
\bottomrule\end{longtable}\endgroup
\section{Coverage map and edition delta}
The previous live glossary exposed 437 headwords.  This edition exposes
976, a net increase of 539.  The increase is not a
mechanical declaration dump: synonyms and theorem-name variants are merged,
while materially different normalizations, branches, equality modes, and
research statuses receive separate entries.
\begingroup\small\renewcommand{\arraystretch}{1.15}
\begin{longtable}{@{}L{0.70\textwidth}r@{}}
\toprule Broad domain & Headwords\\ \midrule\endfirsthead
\toprule Broad domain & Headwords\\ \midrule\endhead
Foundations, formalization, notation, and general analysis & 320 \\
Thue--Morse, dyadic arithmetic, and coefficient calculus & 134 \\
Probability, moments, operators, and information & 99 \\
Approximation, splines, orthogonal systems, and representations & 92 \\
Asymptotics, inversion, transseries, and dynamics & 87 \\
\(q\)-series, special functions, and geometric calculus & 87 \\
Fourier, transforms, sampling, and spectra & 82 \\
Fabius, inverse Fabius, and Rvachev core & 75 \\
\midrule Total & 976\\ \bottomrule\end{longtable}\endgroup
\subsection{Status distribution}
\begingroup\small\renewcommand{\arraystretch}{1.12}
\begin{tabular}{@{}lr@{}}\toprule Status & Headwords\\\midrule
Classical & 468 \\
Lean-proved & 179 \\
Project/mixed & 151 \\
Frontier & 120 \\
Convention & 56 \\
Historical & 2 \\
\bottomrule\end{tabular}\endgroup
\section{Audit method and maintenance contract}
The update was produced by comparing the prior glossary against the current
repository facade, generated declaration census, notation catalogue, primary
exposition, Lean walkthrough, active frontier synthesis, and recursive active
TeX inventory.  The audit normalized aliases before counting concepts.  It
then screened each retained or added entry for the following fields:
\begin{enumerate}
\item exact mathematical object and argument types;
\item domain, codomain, support, and endpoint conventions;
\item characteristic definition or formula;
\item normalization, transform convention, and branch data;
\item pointwise, almost-everywhere, formal, or asymptotic equality mode;
\item current proof status and the boundary between Lean and frontier prose;
\item collisions with the bounded/global/up triad, bit/sign conventions,
Fourier normalizations, and overloaded one-letter symbols.
\end{enumerate}
A future glossary update is incomplete unless it reruns the module/declaration
census, the strict notation audit, the active-TeX inventory, alias checks, and
rendered-PDF inspection.  Numerical evidence and formal transseries algebra
must never be silently promoted to an analytic or kernel-checked theorem.
\section{Repository sources inspected}
All repository links below were accessed on 3 September 2026 from the
\texttt{main} branch.  The formal module/declaration census was the repository's
1 September 2026 checkpoint; the active TeX and notation counts were re-derived
by the notation catalogue on 3 September 2026.
\begingroup\small\renewcommand{\arraystretch}{1.12}
\begin{longtable}{@{}L{0.27\textwidth}L{0.67\textwidth}@{}}
\toprule Source family & Location\\ \midrule\endfirsthead
\toprule Source family & Location\\ \midrule\endhead
Live project root & \url{https://github.com/VladimirReshetnikov/ProveIt/tree/main/Analysis/FabiusFunction} \\
Current documentation root & \url{https://github.com/VladimirReshetnikov/ProveIt/tree/main/Analysis/FabiusFunction/docs} \\
Prior glossary & \url{https://github.com/VladimirReshetnikov/ProveIt/tree/main/Analysis/FabiusFunction/docs/FabiusFunction_Mathematical_Glossary} \\
Mathematical notation catalogue & \url{https://github.com/VladimirReshetnikov/ProveIt/tree/main/Analysis/FabiusFunction/docs/FabiusFunction_Mathematical_Notation_Catalogue} \\
Primary exposition & \url{https://github.com/VladimirReshetnikov/ProveIt/tree/main/Analysis/FabiusFunction/docs/Fabius_Function_and_Rvachev_Up} \\
Lean walkthrough & \url{https://github.com/VladimirReshetnikov/ProveIt/tree/main/Analysis/FabiusFunction/docs/fabius_lean_walkthrough} \\
Semi-formalized frontier synthesis & \url{https://github.com/VladimirReshetnikov/ProveIt/tree/main/Analysis/FabiusFunction/docs/semi-formalized-research-frontiers} \\
\bottomrule\end{longtable}\endgroup
\section{Selected external reference works}
\begin{thebibliography}{99}
\bibitem{AriasUp}
J.~Arias de Reyna,
\emph{An infinitely differentiable function with compact support: definition and properties},
arXiv:1702.05442, 2017.

\bibitem{AriasArithmetic}
J.~Arias de Reyna,
\emph{Arithmetic of the Fabius function},
\emph{Integers} \textbf{18} (2018), Article A51; arXiv:1702.06487.

\bibitem{Haugland}
J.~K.~Haugland,
\emph{Evaluating the Fabius function},
arXiv:1609.07999, 2016.

\bibitem{AlloucheShallit}
J.-P.~Allouche and J.~Shallit,
\emph{Automatic Sequences: Theory, Applications, Generalizations},
Cambridge University Press, 2003.

\bibitem{GasperRahman}
G.~Gasper and M.~Rahman,
\emph{Basic Hypergeometric Series}, second edition,
Cambridge University Press, 2004.

\bibitem{Corless}
R.~M.~Corless, G.~H.~Gonnet, D.~E.~G.~Hare, D.~J.~Jeffrey, and D.~E.~Knuth,
On the Lambert \(W\) function,
\emph{Advances in Computational Mathematics} \textbf{5} (1996), 329--359.

\bibitem{deBoor}
C.~de Boor,
\emph{A Practical Guide to Splines}, revised edition,
Springer, 2001.

\bibitem{Szego}
G.~Szegő,
\emph{Orthogonal Polynomials}, fourth edition,
American Mathematical Society, 1975.
\end{thebibliography}

\begin{scopebox}
\textbf{Snapshot limitation.}
The repository is a living project.  This glossary is complete relative to the
active \texttt{main}-branch corpus and generated inventories retrieved on the
dates printed above.  Later modules or documentation changes require a fresh
audit; archived terminology remains historical unless promoted back into the
active tree.
\end{scopebox}

\end{document}